\documentclass[11pt]{amsart}
\usepackage[T1]{fontenc}
\usepackage{lmodern}
\input{glyphtounicode}
\pdfgentounicode=1
\usepackage{amsmath,amssymb,mathrsfs}
\usepackage[table]{xcolor}
\usepackage{array}
\usepackage{booktabs}
\usepackage{enumitem}
\usepackage[margin=1in]{geometry}
\usepackage{url}
\usepackage[pdfauthor={311859884+mt0-svg@users.noreply.github.com},
            pdftitle={Adversarial bandit convex optimisation: regret d\textasciicircum(3/2) sqrt(n) with polynomial-time steps on a real RAM, and a matching lower bound},
            pdfkeywords={bandit convex optimisation, adversarial bandits, adaptive adversary, minimax regret, lower bound, online Newton step, information ratio, exploration by optimisation, Gaussian exploration},
            bookmarksnumbered=true]{hyperref}
\hypersetup{colorlinks=true, linkcolor=blue, citecolor=blue, urlcolor=blue}
\emergencystretch=3em

\newtheorem{theorem}{Theorem}[section]
\newtheorem{lemma}[theorem]{Lemma}
\newtheorem{proposition}[theorem]{Proposition}
\newtheorem{corollary}[theorem]{Corollary}
\theoremstyle{definition}
\newtheorem{definition}[theorem]{Definition}
\newtheorem{remark}[theorem]{Remark}
\newtheorem{algorithm}[theorem]{Algorithm}
\numberwithin{equation}{section}

% The machine named in every sentence that cites Conjecture 1.
\newcommand{\costmodel}{a real RAM with exact square root, exponential and logarithm, Gaussian draws and spectral functions}
\newcommand{\repo}{https://github.com/mt0-svg/adversarial-bco}
% robust, so that amsart does not uppercase the address in the author line
\DeclareRobustCommand{\authorlink}{\href{\repo/issues/new/choose}{\nolinkurl{311859884+mt0-svg@users.noreply.github.com}}}

\newcommand{\R}{\mathbb{R}}
\newcommand{\E}{\mathbb{E}}
\newcommand{\Prob}{\mathbb{P}}
\newcommand{\cK}{\mathcal{K}}
\newcommand{\cJ}{\mathcal{J}}
\newcommand{\cC}{\mathcal{C}}
\newcommand{\cG}{\mathscr{G}}
\newcommand{\cF}{\mathscr{F}}
\newcommand{\Reg}{\mathfrak{R}}
\newcommand{\BReg}{\mathfrak{BR}}
\newcommand{\diam}{\operatorname{diam}}
\newcommand{\dist}{\operatorname{dist}}
\newcommand{\KL}{\operatorname{KL}}
\newcommand{\Var}{\operatorname{Var}}
\newcommand{\Mills}{\operatorname{M}}
\newcommand{\tr}{\operatorname{tr}}
\newcommand{\ip}[1]{\langle #1\rangle}
\newcommand{\lnp}{\ln_{+}}
\newcommand{\fbar}{\bar f}
\newcommand{\Fbar}{\bar F}

\begin{document}

\title[Adversarial bandit convex optimisation]{Adversarial bandit convex optimisation: regret \texorpdfstring{$d^{3/2}\sqrt{n}$}{d^{}(3/2) sqrt(n)} with polynomial-time steps on a real RAM, and a matching lower bound}
\author{\authorlink}
\thanks{2020 {\em Mathematics Subject Classification.} Primary: 68Q32; Secondary: 90C25, 62L05, 03B35.\\\indent
{\em Key words and phrases.} Bandit convex optimisation; adversarial bandits; adaptive adversary; minimax regret; lower bound; online Newton step; information ratio; exploration by optimisation.}

\begin{abstract}
In adversarial bandit convex optimisation a learner plays $n$ rounds in a convex body of $\mathbb R^d$ against convex losses in $[0,1]$, seeing only its loss. Bubeck, Eldan and Lee conjectured that some learner whose rounds take time polynomial in $d$ and $\ln n$ on polytopes with polynomially many rational constraints has regret $d^{3/2}\sqrt n$ up to logarithms, and that none does better. In their game, against an adaptive adversary, we prove the first on their polytopes, timed on a real RAM with square root, exponential, logarithm, Gaussian draws and spectral functions: a learner has upper expected regret at most a constant times $d^{3/2}\sqrt n\,\ln^5(2dn)$. We prove the second on the unit ball: for $d\ge4\cdot10^6$, $n\ge d^{10}$ every learner has expected regret at least $7\cdot10^{-9}d^{3/2}\sqrt n/\ln d$ against some adversary. For fixed losses on bodies of diameter $D$ containing a unit ball, an explicit learner has expected regret $O(d^{3/2}\sqrt n\,\ln^{3/2}(nD))$.
\end{abstract}

\maketitle

\definecolor{leangreen}{RGB}{0,122,61}
\newcommand{\badge}[2]{\colorbox{#1}{\strut\textcolor{white}{\sffamily\bfseries\small #2}}}
\begin{center}
\small
\renewcommand{\arraystretch}{1.5}
\begin{tabular}{>{\raggedright\arraybackslash}p{0.2\textwidth} >{\raggedright\arraybackslash}p{0.7\textwidth}}
\rowcolor{black!6}
\multicolumn{2}{>{\raggedright\arraybackslash}p{0.93\textwidth}}{\textit{This is AI-generated research: the results, proofs and code were found and written by AI. Credit goes to all the humans whose work it builds on.}}\\
\badge{leangreen}{Proved in Lean 4} & Theorems~\ref{thm:upper}, \ref{thm:polytope}, \ref{thm:lower} and~\ref{thm:explicit}, Remark~\ref{rem:upper-prob}, Corollaries~\ref{cor:poly} and~\ref{cor:minimax}, Proposition~\ref{prop:body}, Theorem~\ref{thm:dil} and Lemma~\ref{lem:marg} are proved in Lean~4 with Mathlib, from statements written out from the definitions of the paper, with only the standard axioms; Comparator checks the statements (Section~\ref{sec:formal}). Every other numbered statement is a step of these proofs, stated as the formal proof proves it.\\
\rowcolor{black!6}
\multicolumn{2}{>{\raggedright\arraybackslash}p{0.93\textwidth}}{Lean proofs, code and the \TeX{} source: \href{\repo}{\texttt{github.com/mt0-svg/adversarial-bco}}}\\
\end{tabular}
\end{center}

\clearpage

\section{Introduction}\label{sec:intro}

This paper proves the two sentences of Conjecture~1 of Bubeck, Eldan and Lee on adversarial bandit convex optimisation in the game in which they state it, the first on their polytopes with the time of a step counted on \costmodel, the second on the unit ball. On the polytopes of the conjecture, given by their constraints, a learner whose every round takes a number of steps of that machine polynomial in $d$ and $\ln n$ has upper expected regret at most a constant times $d^{3/2}\sqrt n\,\ln^5(2dn)$ (Theorem~\ref{thm:polytope}), and on the unit ball no learner, whatever its time, does better for large $d$ and $n$ (Theorem~\ref{thm:lower}). The bounds on the regret are properties of the learner as a map from its Gaussian draws and observations to its plays, and of every learner for the lower bound; the machine enters only the count of steps. For losses fixed in advance the paper also gives an explicit learner whose expected regret, on a body that contains a unit ball, is at most a constant times $d^{3/2}\sqrt n\,\ln^{3/2}(nD)$, where $D$ bounds the diameter of the body.

\subsection{The game and the conjecture}

Let $\cK\subset\R^d$ be a convex body, that is a compact convex set with nonempty interior, and let $\cG(\cK)$ be the set of convex functions $\cK\to[0,1]$. The game lasts $n$ rounds; in round $t=1,\dots,n$:
\begin{enumerate}[label=(\arabic*)]
\item the adversary chooses a loss $f_t\in\cG(\cK)$, as a function of the earlier plays $x_1,\dots,x_{t-1}$ and observations $o_1,\dots,o_{t-1}$, and reveals nothing;
\item the learner chooses $x_t\in\cK$, possibly at random, as a function of $x_1,o_1,\dots,x_{t-1},o_{t-1}$ and of its own randomisation, and observes $o_t=f_t(x_t)$ and nothing else.
\end{enumerate}
A \emph{strategy} of the adversary is the family of maps $\sigma_t$ that give $f_t$ from the history; the adversary is \emph{oblivious} when $\sigma_t$ reads nothing, that is when $f_1,\dots,f_n$ are fixed before the game. The regret of a run is
\begin{equation}\label{eq:regret-run}
\Reg_n=\sum_{t=1}^n f_t(x_t)-\inf_{x\in\cK}\sum_{t=1}^n f_t(x),
\end{equation}
against the losses chosen on that run. Bubeck, Eldan and Lee state the game with the adversary choosing $f_t$ ``simultaneously'' with $x_t$ \cite[Section~1.1]{BEL21}, and their proof of \cite[Theorem~1]{BEL21} lets $f_t$ depend on the earlier plays; the adversary cannot read $x_t$ itself, against which no learner has regret $o(n)$. In their notation, where $n$ is the dimension and $T$ the horizon, they conjecture \cite[Conjecture~1]{BEL21}:
\begin{quote}
\emph{There exists an algorithm such that each step takes $\mathrm{poly}(n\log(T))$-time (under the same assumption on $\cK$ as in Theorem~1) and that achieves $\E R_T=\tilde O(n^{1.5}\sqrt T)$. Furthermore, no algorithm can achieve a better regret bound for large $n$ and $T$.}
\end{quote}
The assumption of their Theorem~1 is that $\cK$ is a polytope described by $\mathrm{poly}(n)$ constraints with rational coefficients whose numerators and denominators are bounded by $\mathrm{poly}(T)$ in absolute value; they normalise $\cK$ to contain a unit ball and to have diameter at most $T$. Their time bounds count evaluations of functions and membership tests \cite[Section~6]{BEL21} and do not name a machine; on rational data the usual reading of time is the number of bit operations, and the results below count steps of a machine on real numbers instead.

\subsection{The results}

Time is counted on \costmodel{} \cite{BSS89}. A \emph{program} is a finite list of instructions with rational constants; one step is one of: an operation $+,-,\times,/$, a square root, an exponential, a logarithm or a comparison of real numbers; an operation on natural numbers used as addresses; one of two spectral functions of a symmetric $d\times d$ matrix, the clamp of its eigenvalues to an interval or the projection on its eigenvectors with eigenvalue above a threshold; one draw of a standard Gaussian number; one play; and, in Theorem~\ref{thm:upper}, one membership question ``$x\in\cK$?'' (Definition~\ref{def:machine}). The observation may carry a noise: $o_t=f_t(x_t)+\varepsilon_t$, where the sequence $(\varepsilon_t)$ has a law $\nu$, independent of the draws of the learner, under which, given $\varepsilon_1,\dots,\varepsilon_{t-1}$, the variable $\varepsilon_t$ has mean $0$ and $\E e^{\varepsilon_t^2}\le2$; we call this hypothesis (N). The game of \cite{BEL21} is the case $\nu=\delta_0$. A strategy here is any family of maps $\sigma_t$ from $(\R^d\times\R)^{t-1}$ to $\cG(\cK)$ such that, at each point $y\in\cK$, $\sigma_t(\cdot)(y)$ is a Borel function of the history. Strategies are deterministic. A randomised adversary, a seed drawn independently of the game and a strategy for each value of it, is a mixture of strategies, and the upper bounds below hold for each of them; their average over the seed would need the regret to be measurable jointly in the seed and the run, which is not shown.

\begin{theorem}\label{thm:upper}
There are a program $P$, real constants $C,C'$ and an integer $k\ge0$ with the following property. Let $d,n\ge1$ be integers, $0<r_0\le R_0$, $\cK\subset\R^d$ a compact convex set with $B(0,r_0)\subset\cK\subset B(0,R_0)$, $\sigma$ a strategy on $\cK$ for $n$ rounds and $\nu$ a law of the noise with (N), and put $L=\ln(2dn(n+1))$. Run $P$ on the inputs $d,n,r_0,R_0$ and $1/(n+1)$ against $\sigma$, its membership questions answered for $\cK$ and the noise drawn from $\nu$. Then on every run the points played in the rounds $1,\dots,n$ lie in $\cK$, the first play comes within $C'\bigl(dL(2+\ln(R_0/r_0))\bigr)^k$ steps and each play of the rounds $2,\dots,n$ within $C'(dL)^k$ steps of the one before; and the upper expectation of the positive part of the regret is at most $C\,d^{3/2}\sqrt n\,L^5+1$: there is a measurable function $R\ge0$ of the Gaussian sequence and the noise such that $\Reg_n\le R$ on every run and
\[
\E R\le C\,d^{3/2}\sqrt n\,L^5+1 .
\]
\end{theorem}

Every round after the first takes a number of machine steps polynomial in $d$ and $\ln n$; the first round also depends on $\ln(R_0/r_0)$, which under the normalisation of \cite{BEL21} is at most $\ln(2n)$. Nothing is amortised over the rounds. The regret is taken against the losses chosen on the run. The theorem does not say that it is a measurable function of the run: a loss of $\cG(\cK)$ may take any values on the boundary of $\cK$ that keep it convex and at most $1$, and the theorem says nothing of where in $\cK$ the points played fall. So it bounds the upper expectation of $\max(\Reg_n,0)$, the infimum of $\E R$ over the measurable $R\ge\max(\Reg_n,0)$; when the regret is measurable and integrable, $\E\Reg_n\le\E R$. The term $+1$ does not weaken the order of the bound: $L\ge\ln4>1$, so $d^{3/2}\sqrt n\,L^5\ge1$ and $C+1$ in place of $C$ absorbs it. The exponent $5$ of the logarithm is that of the proof; Conjecture~1 of \cite{BEL21} asks for $d^{3/2}\sqrt n$ up to logarithmic factors, with no exponent fixed. The program is the learner of Section~\ref{sec:poly}: an online Newton method in the form of \cite[Chapter~11]{Lat24}, with Gaussian exploration, rank-one updates, a restart test written as a quadratic, a positioning of $\cK$ from membership questions, and projections computed by an ellipsoid method. The learner reads a confidence $\delta$, and its regret satisfies the following bound in probability.

\begin{remark}\label{rem:upper-prob}
There are a program $P$, real constants $C,C'$ and an integer $k\ge0$ with the following property. Let $d,n,r_0,R_0,\cK,\sigma,\nu$ be as in Theorem~\ref{thm:upper}, $0<\delta<1$, and $L_\delta=\ln(2dn/\delta)$. Run $P$ on the inputs $d,n,r_0,R_0,\delta$ against $\sigma$. Then on every run the points played in the rounds $1,\dots,n$ lie in $\cK$, the first play comes within $C'\bigl(dL_\delta(2+\ln(R_0/r_0))\bigr)^k$ steps and each play of the rounds $2,\dots,n$ within $C'(dL_\delta)^k$ steps of the one before; and the runs with
\[
\Reg_n>C\,d^{3/2}\sqrt n\,L_\delta^5
\]
have outer probability at most $\delta$.
\end{remark}

Theorem~\ref{thm:upper} follows, with the same program, at $\delta=1/(n+1)$. Increasing $C$ keeps the remark true, so let $C\ge0$. Every point played lies in $\cK$ and every loss takes values in $[0,1]$ on $\cK$, so $\Reg_n\le n$ on every run. The runs above $C\,d^{3/2}\sqrt n\,L^5$ lie in a measurable set $A$ of probability at most $1/(n+1)$, and $R=C\,d^{3/2}\sqrt n\,L^5+n\mathbf 1_A$ is a measurable majorant of the regret with $\E R\le C\,d^{3/2}\sqrt n\,L^5+n/(n+1)$.

On the polytopes of \cite{BEL21} no membership oracle is needed: a membership question is $m$ inner products, and a centre is found from the data.

\begin{theorem}\label{thm:polytope}
\textup{(a)} There are a program $P_{\mathrm{pol}}$ with no membership question, real constants $C,C_1$ and an integer $k_1\ge0$ with the following property. Let $d,n,m,M\ge1$ be integers, $0<\delta<1$, and $\cK=\{x\in\R^d:\ip{a_i,x}\le b_i,\ i=1,\dots,m\}$ a polytope of diameter at most $n$ that contains a closed ball of radius $1$, whose coefficients are rational numbers with numerators and denominators at most $M$ in absolute value, in lowest terms; let $\sigma$ be a strategy on $\cK$ for $n$ rounds and $\nu$ a law of the noise with (N). Run $P_{\mathrm{pol}}$ on $d,n,m,\delta$ and the $m(d+1)$ numbers $a_i,b_i$ against $\sigma$, the noise drawn from $\nu$. Then on every run the points played in the rounds $1,\dots,n$ lie in $\cK$, the first play comes within $C_1\bigl(dmL_\delta\ln(2M)\bigr)^{k_1}$ steps and each play of the rounds $2,\dots,n$ within $C_1(dmL_\delta)^{k_1}$ steps of the one before; and the runs with $\Reg_n>C\,d^{3/2}\sqrt n\,L_\delta^5$ have outer probability at most $\delta$.

\textup{(b)} There are a program $P'_{\mathrm{pol}}$ with no membership question, real constants $C',C'_1$ and an integer $k'_1\ge0$ with the following property. Let $d,n,m,M,\cK,\sigma,\nu$ be as in \textup{(a)} and $L=\ln(2dn(n+1))$. Run $P'_{\mathrm{pol}}$ on $d,n,m$, $1/(n+1)$ and the numbers $a_i,b_i$ against $\sigma$, the noise drawn from $\nu$. Then on every run the points played in the rounds $1,\dots,n$ lie in $\cK$, the first play comes within $C'_1\bigl(dmL\ln(2M)\bigr)^{k'_1}$ steps and each play of the rounds $2,\dots,n$ within $C'_1(dmL)^{k'_1}$ steps of the one before; and the upper expectation of the positive part of the regret is at most $C'\,d^{3/2}\sqrt n\,L^5+1$: there is a measurable function $R\ge0$ of the Gaussian sequence and the noise such that $\Reg_n\le R$ on every run and $\E R\le C'\,d^{3/2}\sqrt n\,L^5+1$.
\end{theorem}

Part \textup{(b)} follows from part \textup{(a)}, with the same program, at $\delta=1/(n+1)$, as Theorem~\ref{thm:upper} follows from Remark~\ref{rem:upper-prob}. The program is not given $M$. The hypotheses have no instance with $n<2$, since a ball of radius $1$ has diameter $2$; this is the normalisation of \cite{BEL21}, whose horizon exceeds the diameter of the body. For their polytopes, with $m$ at most a power of $d$ and $M$ at most a power of $n$, every round, the first included, takes a number of steps of \costmodel{} polynomial in $d$ and $\ln(n/\delta)$, and in \textup{(b)} polynomial in $d$ and $\ln n$. The program runs the program of Remark~\ref{rem:upper-prob} on $\cK$ translated by a centre that it finds first, with every membership question replaced by the $m$ inequalities; part \textup{(a)} is Remark~\ref{rem:upper-prob} with an input that \cite{BEL21} give and nothing more (Section~\ref{sec:poly-polytope}).

The second half of the conjecture is a lower bound on every learner, and it holds whatever the machine. A \emph{learner} on $\cK$ is a probability space $(\Omega,\mu)$ and maps $\pi_t$, $t=1,\dots,n$, from $\Omega\times(\R^d\times\R)^{t-1}$ to $\cK$, jointly measurable, with $x_t=\pi_t(\omega,x_1,o_1,\dots,x_{t-1},o_{t-1})$; nothing is asked of its time.

\begin{theorem}\label{thm:lower}
Let $B$ be the closed unit ball of $\R^d$, $d\ge4\cdot10^6$ and $n\ge d^{10}$. For every learner on $B$ there is a strategy $\sigma$ on $B$ for $n$ rounds, whose loss of round $t$ depends only on $x_1,\dots,x_{t-1}$ and is $\frac14$-Lipschitz, such that in the game without noise the regret of the learner against $\sigma$ is integrable and
\[
\E\,\Reg_n\ge7\cdot10^{-9}\,\frac{d^{3/2}\sqrt n}{\ln d}.
\]
\end{theorem}

Theorem~\ref{thm:polytope}\textup{(b)} is the first sentence of Conjecture~1, on the polytopes of \cite{BEL21} and with time counted on \costmodel. Theorem~\ref{thm:lower} is the second, on the ball $B$, up to the factor $\ln d$ and the powers of $\ln(2dn)$ of the first. The ball is not a polytope; on it the upper bound is Theorem~\ref{thm:upper}, with its membership questions. The strategies of Theorem~\ref{thm:lower} are strategies of Theorem~\ref{thm:upper} for $\cK=B$, with the radii $r_0=R_0=1$ of that theorem, and $\nu=\delta_0$. Against each of them the program of Theorem~\ref{thm:upper} has upper expected regret at most $Cd^{3/2}\sqrt n\,\ln^5(2dn(n+1))+1$, and for $d\ge4\cdot10^6$ and $n\ge d^{10}$ every learner has expected regret at least $7\cdot10^{-9}d^{3/2}\sqrt n/\ln d$ against one of them. They are not shown to meet on one learner: the program of Theorem~\ref{thm:upper} is not shown to be a learner in the sense above, with jointly measurable maps, so Theorem~\ref{thm:lower} is not shown to apply to it, and Theorem~\ref{thm:upper} bounds an upper expectation of its regret. Theorem~\ref{thm:lower} needs the adversary to read the plays; for losses fixed in advance nothing is claimed, and the exponent of $d$ in the minimax regret in the game without noise lies between $1$ and $3/2$ (Section~\ref{sec:limits}).

For losses fixed in advance the paper gives a second learner, explicit and inefficient, whose bound is in expectation with a better power of the logarithm and explicit constants. For it the regret of a learner with actions $X_1,\dots,X_n$ against $f=(f_1,\dots,f_n)$ fixed in advance is
\begin{equation}\label{eq:regret}
\Reg_n(f)=\E\Bigl[\sum_{t=1}^n f_t(X_t)\Bigr]-\inf_{x\in\cK}\sum_{t=1}^n f_t(x),
\end{equation}
the expectation being over the randomisation of the learner only. The infimum need not be attained, since a convex function on $\cK$ need not be lower semicontinuous at the boundary. The learner keeps a finite net $\cC$ of a body slightly inside $\cK$ and, in each round, solves one convex program for the law of its action and for a table of loss estimates; it reads the observation only through a coin with success probability $f_t(x_t)$. It is exploration by optimisation \cite{LG21} with the coin of \cite{AG12,XZ25}, and it is stated as Algorithm~\ref{alg:exo} in Section~\ref{sec:explicit}. The theorem bounds its regret against every sequence fixed in advance, and against an adaptive adversary in the weaker form of the pseudo-regret, where the best point is chosen after the expectation.

\begin{theorem}\label{thm:explicit}
Let $d\ge1$ and $n\ge2$ be integers, $\cK\subset\R^d$ a convex body containing $B(x_0,1)$, $D\ge\diam\cK$, and
\[
\beta=1.2\cdot10^9\,d^2\ln(ndD)\ln(2n).
\]
The learner of Algorithm~\ref{alg:exo}, run with $\varepsilon=1/(n^2(2\beta)^{1/2})$, $\eta=(8\ln m/(n\beta))^{1/2}$ and $\mathrm{tol}=1/n$, satisfies for every $f\in\cG(\cK)^n$
\[
\Reg_n(f)\le6.71+\Bigl(\tfrac12\,\beta\,n\ln m\Bigr)^{1/2}\le6.71+\Bigl(\tfrac12\,\beta\,n\,d\ln\bigl(5n^2(2\beta)^{1/2}D\bigr)\Bigr)^{1/2}
\]
and $\Reg_n(f)\le C_X\,d^{3/2}\sqrt n\,\ln^{3/2}(nD)$ with $C_X=1.9\cdot10^5$; here $m\le(5n^2(2\beta)^{1/2}D)^d$ is the size of its net. For $d\ge n$ the last bound follows from $\Reg_n(f)\le n$ alone. The same bounds hold for $\sup_{x\in\cK}\E\sum_t(f_t(X_t)-f_t(x))$ when each $f_t$ is chosen as a function of the earlier actions and coins of the learner.
\end{theorem}

The learner needs $D$, $x_0$, $n$ and a membership test for the inner body $\cJ$ of Section~\ref{sec:proof}. The learner is explicit and not efficient: the bound on $m$ is exponential in $d$, and step~1 of each round minimises over a convex set the maximum of affine functions indexed by a convex set of loss vectors on the net. The constant $C_0=1.2\cdot10^9$ in $\beta$ is an upper bound for $4\cdot 1.6\cdot10^6\,C_T$, with $C_T=(2+4\ln C_1)/\ln2+48$ and $C_1=52(\pi/2)^{1/2}\frac{17}{16}(1+10^8)$, the constant of the one-round inequality (Section~\ref{sec:body}); and $C_X=1.9\cdot10^5=(27C_0)^{1/2}+10^4$, where $(27C_0)^{1/2}$ comes from the exploration term $(\frac12\beta n\ln m)^{1/2}$, with $C_0$ through $\beta$ and $27/2$ through the logarithm of the size of the net, and $10^4$ covers the additive constant $6.71$.

A policy is a family $\pi=(\pi_1,\dots,\pi_n)$ of maps, $\pi_t$ sending each history $(x_1,o_1,\dots,x_{t-1},o_{t-1})\in(\cK\times[0,1])^{t-1}$ to a probability law on $\cK$. When all its laws are supported by one finite set $\cC\subset\cK$, we call it a \emph{$\cC$-policy}; no measurability is asked of the maps $\pi_t$, and the law of $(x_1,\dots,x_n)$ against $f$ is the law on $\cC^n$ given by the finite product
\begin{equation}\label{eq:play}
\Prob_{\pi,f}(y_1,\dots,y_n)=\prod_{t=1}^n\pi_t\bigl(y_1,f_1(y_1),\dots,y_{t-1},f_{t-1}(y_{t-1})\bigr)(\{y_t\}).
\end{equation}
Its regret $\Reg_n(\pi,f)$ is \eqref{eq:regret} with this law, and the minimax regret $\Reg_n^\star(\cG(\cK))$ is the infimum, over finite sets $\cC\subset\cK$ and $\cC$-policies $\pi$, of $\sup_{f\in\cG(\cK)^n}\Reg_n(\pi,f)$. An upper bound on it bounds the infimum over any larger class of policies whose law of play is defined. Marginalising the coins of the learner gives a $\cC$-policy with the same law of play (Lemma~\ref{lem:marg}), hence the following bound.

\begin{corollary}\label{cor:minimax}
Let $\cK\subset\R^d$ be a convex body containing a Euclidean ball of radius $1$ and $D=\diam\cK$. For every $n\ge1$,
\[
\Reg_n^\star(\cG(\cK))\le C_X\,d^{3/2}\sqrt n\,\ln^{3/2}(nD).
\]
\end{corollary}

Both results rest on a one-round inequality, Proposition~\ref{prop:body}: for a convex body $\cJ$ containing a ball of radius $1/2$ and every finitely supported law $\mu$ on the convex $n$-Lipschitz functions $\cJ\to[0,1]$, some law $\rho$ on $\cJ$ has regret at most $1/n$ plus $(\beta V)^{1/2}$, where $V$ is the variance of the loss under $\rho\otimes\mu$ and $\beta$ is of order $d^2$ up to logarithms. The inequality of \cite[Theorem~3]{Lat20} has $\beta$ of order $d^4\ln(nd/m)$, for losses with strong convexity parameter $m$. The passage from this inequality to a learner is in print, as the next subsection explains.

\subsection{Earlier bounds}

The comparisons with print in this section rest on a search, up to 11 November 2026, of arXiv (listings by date and by author), of the works citing \cite{BEL21,Lat20,Lat24,FHLM24,RH26} in Semantic Scholar and OpenAlex, of Zenodo, and of the Lean libraries. For $d=1$ the minimax regret is of order $\sqrt n$ up to logarithms \cite{BDKP15}, and the first bound polynomial in $d$ times $\sqrt n$ is \cite{BE18}; both go through the minimax theorem and a Bayesian argument, which for \cite{BE18} gives an algorithm with running time exponential in a polynomial of $n$ \cite{BEL21}. The kernel method of \cite{BEL21} has regret $d^{9.5}\sqrt n$ up to logarithms with probability $1-1/n$, in the game above \cite[Theorem~1]{BEL21}; its version that takes time $\mathrm{poly}(d\ln n)$ per step pays a further factor $\mathrm{poly}(d)\,n^{o(1)}$, and \cite{Lat24} lists its polynomial-time version at $d^{10.5}\sqrt n$. For losses fixed in advance the best bound in print is $C\,d^{5/2}\sqrt n\ln(nD)$ \cite[Theorem~1]{Lat20}, by existence; explicit learners at $d^{5/2}\sqrt n$ up to logarithms, inefficient in $d$, are exploration by optimisation \cite[Section~10]{LG21}, its form in \cite[Section~8.5]{Lat24}, and the adaptive minimax sampling of \cite{XZ25}, which already replaces the observation by a coin and which its authors describe as the first algorithm with a finite running time at that order; against an adaptive adversary these bound the pseudo-regret. The best polynomial-time bound in print is that of an online Newton method in \cite[Chapter~11]{Lat24}, a refinement of the method of \cite{FHLM24}, whose own bound is $d^{3.5}\sqrt n$: regret of order $d^{2.5}\sqrt n$ times a power of $\ln(1/\delta)$, with probability $1-\delta$, for losses fixed in advance in $[0,1]$ on a body with $B(0,1)\subset\cK\subset B(0,2d)$, with rounds whose time grows with $t$. In the stochastic version, where one convex function is observed with noise, the regret is at most of order $d^{3/2}\sqrt n$ up to logarithms and lower-order terms, for Lipschitz losses on all of $\R^d$ whose minimiser lies in a known ball of radius $r$, with a further factor $1+r/d$ \cite{LG23}, and for bounded losses on a body in a suitable position \cite{FHLM24}.
In this stochastic version the preprint \cite{RH26} gives a lower bound of order $d^{4/3}\sqrt n$ for $n\ge d^{10/3}$ on the ball, with Gaussian noise.
Linear losses give the lower bound of order $d\sqrt n$ \cite{DHK08} (see \cite[Section~2.4]{Lat24}), and for linear losses this order is also an upper bound against an adaptive adversary, up to a factor $(\ln(n/\delta))^{1/2}$ with probability $1-\delta$ \cite[Theorem~5]{ZL22}. For bandit principal component analysis, whose gains $w^{\mathsf T}G_tw\in[0,1]$ are convex on the ball, so that the losses are concave, \cite[Theorem~2 and Section~4]{BOS26} give a lower bound of order $d^{3/2}\sqrt n$ up to logarithms on the pseudo-regret, with an adversary that removes from its instance the subspace the learner has explored and discloses that subspace. The monograph \cite{Lat24} surveys the field.

Theorem~\ref{thm:upper} lowers the exponent of $d$ for polynomial time to $3/2$: from the $5/2$ of \cite[Chapter~11]{Lat24} for losses fixed in advance, and against an adaptive adversary from the $9.5$ of \cite[Theorem~1]{BEL21}, whose polynomial-time version pays a further factor $\mathrm{poly}(d)\,n^{o(1)}$; it makes the time of every round after the first polynomial in $d$ and $\ln n$. Theorem~\ref{thm:lower} raises the exponent of $d$ in a lower bound for convex losses against an adaptive adversary, in the game without noise, from $1$ to $3/2$, and separates them there from linear losses by a factor $d^{1/2}/\ln d$. Theorem~\ref{thm:explicit} lowers the exponent for losses fixed in advance from $5/2$ to $3/2$ by a route of its own. What it adds is the one-round inequality, Proposition~\ref{prop:body}, with Theorem~\ref{thm:dil} and the dilution behind it; Theorem~\ref{thm:explicit} follows by a converter that is in print, exploration by optimisation \cite{LG21}, which we prove in full and credit component by component (Section~\ref{sec:explicit}), where \cite{Lat20} goes through minimax duality and a Bayesian argument instead. In the minimax step of a round the compactness is asked of the adversary's side only, a compact set of finite dimension, and no Tychonoff argument enters, where \cite{XZ25} applies Tychonoff's theorem to the learner's side; that finite-dimensional set is the parameterisation of the adversary with Bernoulli rewards of \cite[Appendix, Concave parameterization with Bernoulli reward]{XZ25}, and the closed form of the value of a round (Lemma~\ref{lem:dual}) is the formula given there for a fixed law of play, minimised over that law, with losses in place of rewards. We claim no gain in complexity over \cite{XZ25}. The bound carries a factor $\ln^{3/2}(nD)$ instead of the $\ln(nD)$ of \cite{Lat20}.

\subsection{The proofs}

The program of Theorem~\ref{thm:upper} is an online Newton method; its analysis follows \cite[Chapter~11]{Lat24} and changes it in four places, each removing a factor of $d$ or of $n$ (Section~\ref{sec:poly}).
\begin{enumerate}[label=(\arabic*)]
\item The body is first put in a Gaussian position from membership questions: an affine map after which $B(0,1)\subset\cK\subset B(0,R)$ with $R$ of order $d(\ln d)^{1/2}$ and the Gaussian mean width of $\cK$ of order $d^{-1/2}$, found with probability $1-\delta$ by a stochastic descent on balls of growing radius (Theorem~\ref{thm:pg}). This is the only step whose time involves $\ln(R_0/r_0)$.
\item In each round the learner plays a Gaussian point $X_t$ of law $N(\mu_t,(1-\lambda)^2\Sigma_t)$, and one observation gives unbiased estimates of the gradient and the Hessian, at $\mu_t$, of the optimistic surrogate of the loss \cite{LG21a,Lat24}, with no density ratio. The precision $\Sigma_t^{-1}$ grows by a fixed fraction of the estimated curvature, a rank-one update along the direction played, and is never decayed; the curvature it leaves unused is the bonus that makes the restart test fire when the best point leaves the region, which in \cite[Chapter~11]{Lat24} a decay of the covariance supplies at the cost of a factor $d$. The covariance, its inverse and a factor of it are updated in closed form.
\item The region of the next iterate is one sublevel set of a quadratic, and the restart test is the maximum of a concave quadratic over it; a restart costs nothing to the regret against the best point of the whole game. The regret of the learner is bounded on the intersection of a few events, each controlled by a martingale inequality with a predictable variance, with probabilities summing to $\delta$.
\item The gauge of the body is computed by bisection with a random dither, which keeps the probability of the events of the analysis close to their values with the exact gauge, by a comparison of one-dimensional Gaussian laws, and each projection by an ellipsoid method with a separation obtained from membership questions; no step reads a list that grows with $t$.
\item Against a strategy the learner, on each run, plays as it does against the losses that the strategy chose on that run; the statements in probability are proved for losses that depend on the past, by conditioning each round on what precedes its Gaussian draw.
\end{enumerate}

The adversary of Theorem~\ref{thm:lower} hides a minimiser $x^\circ$ drawn uniformly from $\frac12B$ and keeps $\Pi_t$, the law of $x^\circ$ given the history and the answers of a genie described below (Section~\ref{sec:lower}).
\begin{enumerate}[label=(\arabic*)]
\item The loss of round $t$ is $\frac14$ plus a \emph{dent} plus a constant of size at most $\frac18$, a clamped Gaussian draw of the adversary. The dent is the lower convex hull of a bowl centred at the mean of $\Pi_t$, in a metric given by its covariance clipped from below, and of the point $(x^\circ,0)$. The loss is convex, $\frac14$-Lipschitz, with values in $[\frac18,\frac78]$, and minimal at $x^\circ$; the regret is the sum of the dents at the plays.
\item After each play the genie tells the adversary in which of four cells, cut by the play, $x^\circ$ lies. On the regular cell the dent is at least half the bowl at every play, so each round costs the learner a fixed depth $g$, while the observation carries information of order $g^2/d^2$.
\item The other cells are paid by a ledger, half the logarithm of the determinant of the posterior covariance plus a floor: an exact identity gives its expected decrease in a round, and each irregular cell forces a decrease proportional to its mass. A posterior that loses half its directions has lost entropy of order $d\ln n$.
\item Either the posterior keeps its directions and the learner pays $g$ on most rounds, or it loses them and the learner has gathered order $d\ln n$ nats at $g^2/d^2$ per round; the depth $g$ of order $d^{3/2}/\sqrt n$ balances the two. The genie speaks only on cells heavy enough to matter, which removes a cap on the horizon. A bound in expectation over the adversary's randomness holds for one value of it.
\end{enumerate}

The proof of Theorem~\ref{thm:explicit} changes the one-round step of \cite{Lat20} and replaces its converter by exploration by optimisation.
\begin{enumerate}[label=(\arabic*)]
\item The learner plays in the inner body $\cJ$, on which every loss is $n$-Lipschitz, at a cost of $1$ (Lemma~\ref{lem:inner}, Section~\ref{sec:reduce}).
\item The class of convex $n$-Lipschitz functions on $\cJ$ satisfies a one-round inequality $\mathrm{(H)}[\alpha,\beta]$, between the regret of an exploratory law and the variance of the losses under it, with $\alpha=1/n$ and $\beta$ of order $d^2\ln(ndD)\ln n$ (Proposition~\ref{prop:body}, Section~\ref{sec:body}), against $d^4\ln(nd/m)$ in \cite[Theorem~3]{Lat20} for a class with strong convexity parameter $m$. It follows from an inequality on $\R^d$ (Theorem~\ref{thm:dil}, Section~\ref{sec:core}) by an extension of the losses with a penalty for leaving $\cJ$, a projection and Gaussian concentration (Appendix~\ref{app:gauss}).
\item Exploration by optimisation on a net of $\cJ$ turns $\mathrm{(H)}[\alpha,\beta]$ into regret at most $n\alpha$ plus a term of order $(\beta nd\ln(nD))^{1/2}$ and the cost of the net (Theorem~\ref{thm:exo}, Section~\ref{sec:explicit}). The coin makes each round a convex program in finite dimension, a minimax inequality with compactness on one side exchanges the learner and the adversary in that round, and the value of the round is bounded through $\mathrm{(H)}$ and Pinsker's inequality, as in \cite[proof of Theorem~2]{Lat20}.
\end{enumerate}

The inequality on $\R^d$ is proved with Gaussian exploration. For a temperature $\lambda>0$, the Gaussian $N(b,A)$ that minimises the expected mean loss minus $\lambda\ln\det A^{1/2}$ is the Gaussian variational approximation of the law with density proportional to $e^{-\Fbar/\lambda}$ \cite{OA09,CB13}; its stationarity conditions say that the mean loss is flat on average and has curvature $\lambda A^{-1}$ on average (Lemma~\ref{lem:ec}). For each loss of the family, the decrease from the centre $b$ toward the loss's own minimiser is turned, by Gaussian integration by parts on a half space, into a covariance between the loss and a score whose second moment is bounded independently of the direction and the distance (Lemma~\ref{lem:inf}). The Cauchy and Schwarz inequality then bounds the gap of each loss at $b$ by $d$ times the root mean square of its deviation from the mean loss, up to a logarithm of the Lipschitz constant that enters through the Mills ratio (Lemma~\ref{lem:half}, Proposition~\ref{prop:core}). The integration by parts is the infinitesimal form of the change of measure behind the optimistic surrogate of \cite{BEL21,LG21a}. Finally the regret may be carried by losses of small total weight whose gains are much larger than the floor; the exploratory law mixes a Gaussian at the scale of those gains with a small weight, chosen among the scales by a layer-cake argument that costs a logarithm (Theorem~\ref{thm:dil}).

The exponent $3/2$ is the limit of this route: for the unit ball and $d>256$, \cite[Theorem~12]{BLS25} gives a law on convex losses under which every exploratory law has regret at least $2^{-19}(d/\ln d)$ times the square root of its variance; its authors conclude: ``At least naively, this means that the information-theoretic machinery will not yield a bound on the regret for general bandit convex optimisation that is better than $\tilde O(d^{1.5}\sqrt n)$.'' This is not a lower bound on the regret.

\subsection{Conventions}

Throughout, $|\cdot|$ is the Euclidean norm, $\ip{\cdot,\cdot}$ the inner product, $B(x,r)$ the closed ball, $\ln$ the natural logarithm and $\lnp=\max(\ln,0)$. A law is a Borel probability measure; $N(b,A)$ is the Gaussian law of mean $b$ and covariance $A$, $\varphi$ and $Q$ are the standard Gaussian density and upper tail on $\R$, and $\Mills=Q/\varphi$ is the Mills ratio. Constants written $C$, $C_0$, $C'$ are absolute.


\section{The inner body and the one-round inequality}\label{sec:proof}

This section states the two results that the proof of Theorem~\ref{thm:explicit} combines with the converter of Section~\ref{sec:explicit}. For $f\colon\cJ\to\R$ continuous on a compact set $\cJ$ we write $f_\star=\min_{\cJ}f$.

The learner plays on a smaller body, on which every loss is Lipschitz. For a convex body $\cK$ and $n\ge1$ let
\[
\cJ=\{x\in\cK: B(x,1/n)\subset\cK\},
\]
and for a compact set $\cJ$ let $\cF(\cJ,n)$ be the set of convex functions $f\colon\cJ\to[0,1]$ with $|f(x)-f(y)|\le n|x-y|$ for $x,y\in\cJ$. The next lemma passes from $\cK$ to $\cJ$; it is proved in Section~\ref{sec:reduce}.

\begin{lemma}\label{lem:inner}
Let $\cK\subset\R^d$ be a convex body containing $B(x_0,1)$, $D=\diam\cK$ and $n\ge2$.
\begin{enumerate}[label=(\alph*)]
\item $\cJ$ is a convex body containing $B(x_0,1-1/n)$, and $1\le\diam\cJ\le D$.
\item For every $f\in\cG(\cK)$ the restriction of $f$ to $\cJ$ lies in $\cF(\cJ,n)$.
\item For $f_1,\dots,f_n\in\cG(\cK)$, $\inf_{\cJ}\sum_t f_t\le\inf_{\cK}\sum_t f_t+1$.
\end{enumerate}
\end{lemma}

On $\cJ$ the regret is controlled by a one-round inequality between the regret of an exploratory law and the variance of the losses under it. For a class $\cF$ of continuous functions on a compact set $\cJ$ and $\alpha,\beta\ge0$, we say that $\cF$ satisfies $\mathrm{(H)}[\alpha,\beta]$ if for every finitely supported probability law $\mu$ on $\cF$ whose mean $\fbar=\int f\,d\mu(f)$ lies in $\cF$ there is a Borel probability law $\rho$ on $\cJ$ with
\begin{equation}\label{eq:H}
\int\fbar\,d\rho-\int f_\star\,d\mu(f)\le\alpha+\Bigl(\beta\int\!\!\int\bigl(\fbar(x)-f(x)\bigr)^2\,d\rho(x)\,d\mu(f)\Bigr)^{1/2}.
\end{equation}
The left side is the regret of a play drawn from $\rho$ against a loss drawn from $\mu$, measured against the minimum of that loss; the integral on the right is the variance of the loss under $\rho\otimes\mu$, which is what one observation reveals. Section~\ref{sec:explicit} turns \eqref{eq:H} into a regret bound. The hypothesis is supplied by the following proposition, the new step of the paper. It is proved in Section~\ref{sec:body} from a one-round inequality on $\R^d$, Theorem~\ref{thm:dil}.

\begin{proposition}\label{prop:body}
Let $d\ge1$ and $n\ge2$ be integers, and let $\cJ\subset\R^d$ be a convex body containing a ball of radius $1/2$, of diameter $D_\cJ$. Then $\cF(\cJ,n)$ satisfies $\mathrm{(H)}[1/n,\beta]$ with $\beta=C_0\,d^2\ln(ndD_\cJ)\ln(2n)$ and $C_0=1.2\cdot10^9$.
\end{proposition}


The constant $\beta$ is of order $d^2$ up to logarithms; the printed value \cite[Theorem~3]{Lat20} is of order $d^4\ln(nd/m)$ for a class with strong convexity parameter $m$. Since $\cF(\cJ,n)$ is closed under finite convex combinations, the mean $\fbar$ in \eqref{eq:H} always lies in it. The three logarithms of the final bound have different sources: one comes from the dilution over the scales of the losses in Theorem~\ref{thm:dil}, one from the Gaussian tail in the second-order term of Lemma~\ref{lem:half}, and one from the entropy $d\ln(1/\varepsilon)$ of the net in Theorem~\ref{thm:exo}.

\section{A one-round inequality on \texorpdfstring{$\R^d$}{R\^{}d}}\label{sec:core}

This section proves the one-round inequality for a finite family of convex Lipschitz functions on all of $\R^d$; Section~\ref{sec:body} brings it back to a convex body.

\subsection{Setting and statement}\label{sec:core-setting}

A \emph{family} is a finite set of indices $c$ with weights $\mu_c>0$, $\sum_c\mu_c=1$, and convex $L$-Lipschitz functions $F_c\colon\R^d\to\R$, each attaining its minimum at some point $x_c$, whose mean $\Fbar=\sum_c\mu_cF_c$ is coercive: $\Fbar(x)\ge a_1|x|-a_2$ for some $a_1>0$ and $a_2$. Then $\Fbar$ is convex, $L$-Lipschitz and attains its minimum. The family is \emph{regular} if moreover every $F_c$ is of class $C^2$ with a Lipschitz gradient; then so is $\Fbar$. Write
\[
u_c=F_c-\Fbar,\qquad V(x)=\sum_c\mu_c\,u_c(x)^2,\qquad \Delta_\star=\min\Fbar-\sum_c\mu_c\min F_c .
\]
The \emph{floor} $\Delta_\star$ is the regret of the best fixed point: for a Borel probability law $p$ on $\R^d$ with $\int|x|\,dp<\infty$ put
\[
\Delta(p)=\int\Fbar\,dp-\sum_c\mu_c\min F_c\ \ (\ge\Delta_\star),\qquad \mathcal I(p)=\int V\,dp .
\]
The minima of the $F_c$ are not assumed equal: this is what a round of the learner produces (Lemma~\ref{lem:conv}), a family each of whose members is compared with its own minimum.

Gaussian laws enter through a temperature $\lambda>0$. For $b\in\R^d$ and $S$ symmetric positive definite let $Z$ be a standard Gaussian vector of $\R^d$, $w=SZ$, $x=b+w$, so that $x$ has law $N(b,A)$ with $A=S^2$, and put
\begin{equation}\label{eq:entropic}
\mathcal E_\lambda(b,S)=\E\,\Fbar(b+SZ)-\lambda\ln\det S .
\end{equation}
The minimiser $(b,S)$ of $\mathcal E_\lambda$, which exists and is unique by Lemma~\ref{lem:ec}, or the law $N(b,A)$, is the \emph{entropic Gaussian of $\Fbar$ at temperature $\lambda$}. It balances the expected value of the mean loss against the volume of the Gaussian: a cold Gaussian sits near the minimiser of $\Fbar$, a hot one spreads until its cost $\E\Fbar(x)-\min\Fbar$, at most $\lambda d$ by Lemma~\ref{lem:ec} when $\Fbar$ is of class $C^1$, matches the temperature.

\begin{theorem}\label{thm:dil}
Let $(\mu_c,F_c)_c$ be a family with $\Delta_\star>0$, let $b^\star$ be a minimiser of $\Fbar$, and let $g_c=F_c(b^\star)-\min F_c$ and $g_{\max}=\max_cg_c$. There is a scale $h\in[\Delta_\star/2,g_{\max}]$, one of the values $g_c$, with the following property. Let $(b_h,S_h)$ be the entropic Gaussian of $\Fbar$ at temperature $\lambda_h=h/(16d)$, $A_h=S_h^2$, $s_h$ the largest eigenvalue of $S_h$, $T_h=(2\lnp(26\,d\,L\,s_h/h))^{1/2}$ with $\lnp=\max(\ln,0)$, $q=\Delta_\star/(4h)$, and
\[
p_h=(1-q)\,\delta_{b^\star}+\tfrac q2\,\delta_{b_h}+\tfrac q2\,N(b_h,A_h).
\]
Then $\Delta(p_h)\le\frac{65}{64}\Delta_\star$,
\begin{equation}\label{eq:P3}
V(b^\star)+V(b_h)+\E\,V(b_h+S_hZ)\ge\frac{h\,\Delta_\star}{2K_S\,(1+T_h)^2\,d^2\,\ln(2g_{\max}/\Delta_\star)},
\end{equation}
where $K_S=2(32/\kappa)^2(13/4+2\kappa_2)$, $\kappa=\varphi(0)-\varphi(1)+Q(1)$ and $\kappa_2=1-Q(1)-\varphi(1)$, and
\[
\Delta(p_h)^2\le K_D\,(1+T_h)^2\,d^2\,\ln\bigl(2g_{\max}/\Delta_\star\bigr)\;\mathcal I(p_h),\qquad K_D=1.6\cdot10^6 .
\]
\end{theorem}

Since $\sum_c\mu_cg_c=\Fbar(b^\star)-\sum_c\mu_c\min F_c=\Delta_\star$, $g_{\max}\ge\Delta_\star$ and the logarithm is at least $\ln2$. The theorem is an information-ratio bound of order $d^2$ up to logarithms. The floor may be carried by components of small total weight whose gains $g_c$ are much larger than $\Delta_\star$; a Gaussian whose cost is of the order of the floor is then too narrow to read them. The law $p_h$ takes a hotter Gaussian, at the scale $h$ of the gains that carry the floor, and dilutes it with a mass $q$ small enough to keep its cost $q\lambda_hd=\Delta_\star/64$. The two point masses at $b^\star$ and at the centre $b_h$ read the components whose gain at $b_h$ is lost because the centre has moved toward them (Step~1 of the proof). The bound \eqref{eq:P3} is the lower bound on the variance that Steps~1 to~3 of the proof produce; Section~\ref{sec:body} uses it.

The proof rests on one inequality for one component, Proposition~\ref{prop:core} below, which comes from three lemmas. Lemma~\ref{lem:ec} gives the entropic Gaussian and its stationarity conditions; Lemma~\ref{lem:inf} is the gradient inequality of a component toward its minimiser, integrated by parts against the Gaussian on a half space; Lemma~\ref{lem:half} bounds below the second-order term that the integration by parts produces. The three lemmas and the proposition are statements about functions on $\R^d$ and the standard Gaussian; this section proves Theorem~\ref{thm:dil} from them for regular families, and Appendix~\ref{app:calc} extends it to every family by a limit.

\subsection{The three lemmas}

The stationarity conditions of the entropic Gaussian are used in Proposition~\ref{prop:core}, and its cost in Theorem~\ref{thm:dil}. Every law below of the form $N(b,A)$ with $A$ positive definite has a density.

\begin{lemma}[entropic Gaussian]\label{lem:ec}
Let $\Fbar\colon\R^d\to\R$ be convex, Lipschitz and coercive, and $\lambda>0$. Then $\mathcal E_\lambda$ has exactly one minimiser $(b,S)$, and it minimises also the extension $\mathcal E_\lambda(b,S)=\E\,\Fbar(b+SZ)-\lambda\ln|\det S|$ of \eqref{eq:entropic} to the invertible $S$. Every $b$ and every linear map $S$ have
\[
\mathrm{(S4)}\ \ \E\,\Fbar(b+SZ)\ge\Fbar(b).
\]
If $\Fbar$ is moreover of class $C^1$, every minimiser $(b,S)$ of $\mathcal E_\lambda$ over the invertible $S$ has, with $x=b+SZ$,
\begin{gather*}
\mathrm{(S1)}\ \ \E\,\nabla\Fbar(x)=0,\qquad\mathrm{(S2)}\ \ \E\bigl[\ip{v,Z}\ip{\nabla\Fbar(x),Sv'}\bigr]=\lambda\ip{v,v'}\ \ (v,v'\in\R^d),\\
\mathrm{(S3)}\ \ \E\,\Fbar(x)\le\Fbar(z)+\lambda d\ \ (z\in\R^d).
\end{gather*}
\end{lemma}

Conditions (S1) and (S2) are the stationarity conditions of the Gaussian variational approximation \cite{OA09}, and the objective \eqref{eq:entropic} is convex in $(b,S)$ \cite{CB13}; the proof below uses neither fact. With $z$ a minimiser of $\Fbar$, (S3) bounds the cost $\E\Fbar(x)-\min\Fbar$ by $\lambda d$.

The other lemmas are stated in the coordinates of the Gaussian. For $(b,S)$ with $S$ invertible and $F\colon\R^d\to\R$ write $F^{b,S}(g)=F(b+Sg)$, so that $\E F(b+SZ)=\E F^{b,S}(Z)$. If $F$ is convex and $L$-Lipschitz, $F^{b,S}$ is convex and $L\|S\|$-Lipschitz, $\|S\|$ the operator norm. If $\Fbar$ is differentiable, $\ip{\nabla\Fbar^{b,S}(g),v}=\ip{\nabla\Fbar(b+Sg),Sv}$, so that for $G=\Fbar^{b,S}$ conditions (S1), (S2) and (S4) read
\[
\E\,\nabla G(Z)=0,\qquad\E\bigl[\ip{v,Z}\ip{\nabla G(Z),v'}\bigr]=\lambda\ip{v,v'}\ \ (v,v'\in\R^d),\qquad\E\,G(Z)\ge G(0).
\]
A component $c$ with $x_c\ne b$ sits in these coordinates at $y_c=S^{-1}(x_c-b)$. Its \emph{depth} $\vartheta_c=|y_c|$ is the distance from $b$ to $x_c$ in the metric of the Gaussian, and its \emph{gap} at $b$ is $\gamma_c=F_c(b)-\min F_c=F_c^{b,S}(0)-F_c^{b,S}(y_c)$. Along the segment from $0$ to $y_c$ the convex function $F_c^{b,S}$ decreases at rate at least $\gamma_c/\vartheta_c$. Lemma~\ref{lem:inf} reads this decrease from the Gaussian: at a point $g$ of a half space it moves $g$ by an infinitesimal step toward a point of the segment, and the gradient inequality, integrated by parts, turns the decrease into a covariance between $u_c^{b,S}$ and an explicit score.

Let $\varphi$ and $Q$ be the standard Gaussian density and upper tail on $\R$, $Z_1$ a standard Gaussian variable, $\Mills(T)=Q(T)/\varphi(T)$ the Mills ratio, and
\begin{gather*}
\kappa=\E\min(Z_1^+,1)=\varphi(0)-\varphi(1)+Q(1),\\
\kappa_2=\E\bigl[(\min(Z_1,1)Z_1-\mathbf 1\{Z_1<1\})^2;Z_1>0\bigr]=1-Q(1)-\varphi(1).
\end{gather*}

\begin{lemma}[infinitesimal see-saw]\label{lem:inf}
Let $F,G\colon\R^d\to\R$ be differentiable with bounded gradients, $F$ convex, $e$ a unit vector, $R>0$ and $\Gamma$ a real number with $F(Re)\le F(0)-\Gamma$. Put $\zeta=\ip{e,Z}$, $\Omega=\{\zeta>0\}$, $k=\min(\zeta^+,1)$, $m=\mathbf 1_\Omega/R$ and
\[
\ell=\mathbf 1_\Omega\Bigl[\frac{|Z|^2-d}R-k\zeta+\mathbf 1\{\zeta<1\}\Bigr].
\]
Then
\begin{multline}\label{eq:inf}
\frac{\kappa\,\Gamma}R+\E\bigl[k\ip{\nabla G(Z),e}\bigr]-\E\bigl[m\bigl(\ip{\nabla G(Z),Z}-G(Z)+G(0)\bigr)\bigr]\\
\le\E\bigl[\bigl(F(Z)-G(Z)\bigr)(\ell-m)\bigr]+\frac{F(0)-G(0)}{2R}.
\end{multline}
Moreover $\E k=\kappa$, $\E m=1/(2R)$ and $\E(\ell-m)^2\le(2d+1)/R^2+2\kappa_2$.
\end{lemma}

The three terms on the left are the \emph{tilt}, the decrease of $F$ toward the point $Re$; the \emph{second-order term}, which reads the curvature of $G$ in the direction $e$; and the \emph{dilation term}, the price of the radial part of the step. With $F=F_c^{b,S}$ and $G=\Fbar^{b,S}$ the right side is a covariance of $u_c^{b,S}$ in the variance currency of \eqref{eq:H}: the score $\ell-m$ has a second moment bounded whatever its direction, and by $3+2\kappa_2$ for $R\ge d^{1/2}$. The proof of Proposition~\ref{prop:core} bounds the dilation term by (S2) and (S4), the tilt by the gap and the second-order term by the next lemma.

\begin{lemma}[second-order term]\label{lem:half}
Let $e$ be a unit vector, $\zeta=\ip{e,Z}$, $k=\min(\zeta^+,1)$, $\lambda$ real, $L\ge0$, and $H\colon\R^d\to\R$ differentiable with $|H|\le L$, a bounded gradient and $\ip{\nabla H,e}\ge0$, such that $\E H(Z)=0$ and $\E[\zeta H(Z)]=\lambda$. Then for every $T\ge0$
\[
\E\bigl[k\,H(Z)\bigr]\ge\kappa\,\Mills(T)\bigl(\lambda-2L\varphi(T)\bigr).
\]
\end{lemma}

The lemma is applied to $H=\ip{\nabla G,e}$, for which (S1) and (S2) give $\E H(Z)=0$ and $\E[\zeta H(Z)]=\lambda$: the full Gaussian sees the curvature of $G$ in the direction $e$ with weight $\lambda$, and the half space sees at least the fraction $\kappa\Mills(T)$ of the part of it that sits within $T$ of the centre, the Lipschitz constant bounding the rest. This is where the Lipschitz constant enters the final bound, through $\Mills(T)\ge1/(1+T)$ and a logarithm.

\subsection{The core inequality and the proof of Theorem~\ref{thm:dil}}

The three lemmas combine into one inequality for one component that is far above its minimum at the centre of a Gaussian that satisfies (S1), (S2) and (S4); Step~1 of the proof of Theorem~\ref{thm:dil} applies it, in the coordinates of the entropic Gaussian, to every component that keeps its gain when the centre moves.

\begin{proposition}[core inequality]\label{prop:core}
Let $d\ge1$, $L'\ge0$, and $F,G\colon\R^d\to\R$ convex and differentiable with $|\nabla F|,|\nabla G|\le L'$, the gradient of $G$ differentiable with a bounded derivative and $\ip{\nabla^2G(g)v,v}\ge0$ for all $g,v$. Let $h>0$, $\lambda=h/(16d)$, and assume
\[
\E\,\nabla G(Z)=0,\qquad\E\bigl[\ip{v,Z}\ip{\nabla G(Z),v'}\bigr]=\lambda\ip{v,v'}\ \ (v,v'\in\R^d),\qquad G(0)\le\E\,G(Z).
\]
Let $T\ge0$ with $2L'\varphi(T)\le h/(32d)$, and assume $F(y)\le F(0)-h/2$ for some $y$. Then
\[
\Bigl(\frac{\kappa\,h}{32(1+T)d}\Bigr)^2\le\Bigl(3+2\kappa_2+\frac1{4d}\Bigr)\Bigl(\E\bigl(F(Z)-G(Z)\bigr)^2+\bigl(F(0)-G(0)\bigr)^2\Bigr).
\]
\end{proposition}

\begin{proof}
Since $h>0$, $y\ne0$. Put $R=|y|$ and $e=y/R$, and apply Lemma~\ref{lem:inf} with $\Gamma=h/2$. We bound the three terms on the left of \eqref{eq:inf}.

Tilt. It is $\kappa h/(2R)$.

Dilation. By the gradient inequality for the convex $G$ at $g$ toward $0$, $G(0)\ge G(g)-\ip{\nabla G(g),g}$, so the integrand $\ip{\nabla G(Z),Z}-G(Z)+G(0)$ is nonnegative. Its integral over $\Omega$ is therefore at most its integral over all of $\R^d$, which is $\lambda d-(\E G(Z)-G(0))\le\lambda d$, by the second assumption with $v=v'$ running over an orthonormal basis and by the third. The dilation term is at most $\lambda d/R=h/(16R)$.

Second order. The function $H=\ip{\nabla G,e}$ has $|H|\le L'$, a bounded gradient and $\ip{\nabla H,e}=\ip{\nabla^2G\,e,e}\ge0$; $\E H(Z)=0$ by the first assumption and $\E[\zeta H(Z)]=\lambda$ by the second with $v=v'=e$. Lemma~\ref{lem:half} and $2L'\varphi(T)\le h/(32d)=\lambda/2$ give a second-order term at least $\kappa\Mills(T)\lambda/2=\kappa\Mills(T)h/(32d)$.

So the left side of \eqref{eq:inf} is at least $(\kappa/2-1/16)h/R+\kappa\Mills(T)h/(32d)$, with $\kappa/2-1/16\ge\kappa/32>0$, since $\kappa\ge0.312$ (checked in Lean). Put $\omega=\min(1,R/d^{1/2})$. If $R\le d^{1/2}$, then $\omega=R/d^{1/2}$ and $\omega$ times the left side is at least $(\kappa/2-1/16)h/d^{1/2}\ge(\kappa/32)h/d$. If $R\ge d^{1/2}$, then $\omega=1$ and the left side is at least $\kappa\Mills(T)h/(32d)$. Since $\Mills(T)\ge1/(1+T)$ (Lemma~\ref{lem:mills}), in both cases $\omega$ times the left side is at least $\kappa h/(32(1+T)d)$.

By the Cauchy and Schwarz inequality and Lemma~\ref{lem:inf}, the right side of \eqref{eq:inf} is at most $(\E(F(Z)-G(Z))^2)^{1/2}((2d+1)/R^2+2\kappa_2)^{1/2}+|F(0)-G(0)|/(2R)$. Since $\omega\le1$ and $\omega/R\le d^{-1/2}$, $\omega^2((2d+1)/R^2+2\kappa_2)\le(2d+1)/d+2\kappa_2\le3+2\kappa_2$ and $\omega/(2R)\le1/(2d^{1/2})$. Hence
\[
\frac{\kappa\,h}{32(1+T)d}\le(3+2\kappa_2)^{1/2}\bigl(\E(F(Z)-G(Z))^2\bigr)^{1/2}+\frac1{2d^{1/2}}\,|F(0)-G(0)|,
\]
and $(\alpha_1y_1+\alpha_2y_2)^2\le(\alpha_1^2+\alpha_2^2)(y_1^2+y_2^2)$ gives the claim.
\end{proof}

The proof above and the proof of Lemma~\ref{lem:half} use two properties of the Mills ratio.

\begin{lemma}\label{lem:mills}
For $T\ge0$, $1/(1+T)\le\Mills(T)\le\Mills(0)=(\pi/2)^{1/2}$, and $\Mills$ is nonincreasing.
\end{lemma}

\begin{proof}
Let $v(T)=Q(T)-\varphi(T)/(1+T)$. Since $\varphi'(T)=-T\varphi(T)$, $v'(T)=-\varphi(T)+\varphi(T)\bigl(T(1+T)+1\bigr)/(1+T)^2=-T\varphi(T)/(1+T)^2\le0$, and $v(T)\to0$ as $T\to\infty$, so $v\ge0$. For the monotonicity, $Q(T)=\int_T^\infty\varphi\le\int_T^\infty(s/T)\varphi(s)\,ds=\varphi(T)/T$ for $T>0$, and $\Mills'(T)=-1+T\Mills(T)\le0$.
\end{proof}

\begin{proof}[Proof of Theorem~\ref{thm:dil}]
Let the family be regular; Lemma~\ref{lem:approx} of Appendix~\ref{app:calc} gives the theorem for every family from this case. For $h>0$ let $\Lambda(h)=\sum_{c:g_c\ge h}\mu_c$, let $(b_h,S_h)$ be the entropic Gaussian at temperature $\lambda_h=h/(16d)$ and $x=b_h+S_hZ$, and put $V_{\mathrm G}(h)=\E\,V(x)$. The minimiser $b^\star$ is fixed first, before any scale; the proof uses of it only that it minimises $\Fbar$. By Lemma~\ref{lem:ec}, $(b_h,S_h)$ minimises $\mathcal E_{\lambda_h}$ also over the invertible $S$, so (S1) to (S4) hold for it.

Step 1, one scale. We show that for every $h>0$, with $K_S$ the constant of \eqref{eq:P3},
\begin{equation}\label{eq:scale}
\Lambda(h)\,h^2\le K_S\,(1+T_h)^2\,d^2\bigl(V_{\mathrm G}(h)+V(b_h)+V(b^\star)\bigr).
\end{equation}
Let $c$ have $g_c\ge h$. Then $F_c(b_h)-\min F_c=g_c+F_c(b_h)-F_c(b^\star)$. Either this is at least $h/2$, and we call $c$ \emph{good}, or $F_c(b_h)-F_c(b^\star)<h/2-g_c\le-h/2$, and we call $c$ \emph{moved}. For a moved $c$, since $\Fbar(b_h)\ge\Fbar(b^\star)$,
\[
u_c(b_h)-u_c(b^\star)=\bigl(F_c(b_h)-F_c(b^\star)\bigr)-\bigl(\Fbar(b_h)-\Fbar(b^\star)\bigr)<-h/2,
\]
so $u_c(b_h)^2+u_c(b^\star)^2\ge h^2/8$. If the moved components have total weight at least $\Lambda(h)/2$, then $V(b_h)+V(b^\star)\ge\Lambda(h)h^2/16$, and $16\le K_S$. Otherwise the good components have weight at least $\Lambda(h)/2$. For a good $c$ apply Proposition~\ref{prop:core} to $F=F_c^{b_h,S_h}$ and $G=\Fbar^{b_h,S_h}$, with $L'=Ls_h$ and $T=T_h$. Both are convex and differentiable with gradients bounded by $Ls_h$; the gradient of $G$ is differentiable with a bounded derivative since the family is regular, and $\ip{\nabla^2G\,v,v}\ge0$ since $G$ is convex; the three assumptions on $G$ are (S1), (S2) and (S4) in the coordinates of the Gaussian; and $F(y)\le F(0)-h/2$ at $y=S_h^{-1}(x_c-b_h)$, since $F_c(x_c)=\min F_c\le F_c(b_h)-h/2$. For $T_h$: if $26dLs_h/h\ge1$, then $\varphi(T_h)=(2\pi)^{-1/2}h/(26dLs_h)\le h/(64dLs_h)$ because $26(2\pi)^{1/2}\ge64$ (checked in Lean); if $26dLs_h/h<1$, then $T_h=0$ and $64dLs_h/h<64/26\le(2\pi)^{1/2}$, so $2Ls_h\varphi(0)<h/(32d)$. In both cases $2Ls_h\varphi(T_h)\le h/(32d)$. Since $\E(F(Z)-G(Z))^2=\E u_c(x)^2$ and $F(0)-G(0)=u_c(b_h)$, Proposition~\ref{prop:core} gives
\[
\Bigl(\frac{\kappa\,h}{32(1+T_h)d}\Bigr)^2\le\Bigl(3+2\kappa_2+\frac1{4d}\Bigr)\bigl(\E u_c(x)^2+u_c(b_h)^2\bigr).
\]
Multiplying by $\mu_c$ and summing over the good components,
\[
\tfrac12\Lambda(h)\,h^2\le(32/\kappa)^2(1+T_h)^2d^2\bigl(3+2\kappa_2+1/(4d)\bigr)\bigl(V_{\mathrm G}(h)+V(b_h)\bigr),
\]
which is \eqref{eq:scale}, since $1/(4d)\le1/4$.

Step 2, the law at a scale $h\ge\Delta_\star/2$. Then $q=\Delta_\star/(4h)\le1/2$. By (S4) and (S3) with $z=b^\star$,
\[
\Delta(p_h)=\Delta_\star+\tfrac q2\bigl(\Fbar(b_h)-\min\Fbar\bigr)+\tfrac q2\bigl(\E\Fbar(x)-\min\Fbar\bigr)\le\Delta_\star+q\lambda_hd=\Delta_\star+\frac{\Delta_\star}{4h}\cdot\frac h{16}=\frac{65}{64}\Delta_\star .
\]
Since $1-q\ge1/2\ge q/2$, \eqref{eq:scale} gives
\[
\mathcal I(p_h)\ge\tfrac q2\bigl(V(b^\star)+V(b_h)+V_{\mathrm G}(h)\bigr)\ge\frac{\Delta_\star\,h\,\Lambda(h)}{8K_S(1+T_h)^2d^2}.
\]

Step 3, the scale. By the layer-cake formula $\Delta_\star=\sum_c\mu_cg_c=\int_0^{g_{\max}}\Lambda(t)\,dt$, and at most $\Delta_\star/2$ of it comes from $t<\Delta_\star/2$ since $\Lambda\le1$. So
\[
\frac{\Delta_\star}2\le\int_{\Delta_\star/2}^{g_{\max}}t\Lambda(t)\,\frac{dt}t\le\Bigl(\sup_{t\in[\Delta_\star/2,g_{\max}]}t\Lambda(t)\Bigr)\ln\frac{2g_{\max}}{\Delta_\star}.
\]
The function $\Lambda$ is constant on each interval between consecutive values of the $g_c$, closed on the right, so $t\Lambda(t)$ increases there, and the supremum is attained at a value $h=g_c\in[\Delta_\star/2,g_{\max}]$. For that $h$, \eqref{eq:scale} and $h\Lambda(h)\ge\Delta_\star/(2\ln(2g_{\max}/\Delta_\star))$ give \eqref{eq:P3}, Step 2 gives $\mathcal I(p_h)\ge\Delta_\star^2/(16K_S(1+T_h)^2d^2\ln(2g_{\max}/\Delta_\star))$, and
\[
\Delta(p_h)^2\le\Bigl(\frac{65}{64}\Bigr)^2\Delta_\star^2\le\Bigl(\frac{65}{64}\Bigr)^2 16K_S\,(1+T_h)^2d^2\ln\frac{2g_{\max}}{\Delta_\star}\,\mathcal I(p_h),
\]
with $(65/64)^2\cdot16K_S\le K_D$ (checked in Lean).
\end{proof}

\subsection{Proofs of the three lemmas}\label{sec:core-proofs}

The integrations by parts of Lemma~\ref{lem:inf} take place on a half space. The next lemma does them by two computations in one variable, a dilation and a translation, with no boundary term. Fix a unit vector $e$, $\zeta=\ip{e,Z}$ and $\Omega=\{\zeta>0\}$, and let $\gamma$ be the standard Gaussian density of $\R^d$.

\begin{lemma}[integration by parts on a half space]\label{lem:ibp}
Let $G\colon\R^d\to\R$ be differentiable, with $|G(g)|\le a\,e^{c|g|}$ and $|\nabla G(g)|\le a\,e^{c|g|}$ for all $g$, for some constants $a$ and $c$. Then
\begin{align*}
&\mathrm{(a)}\ \ \E\bigl[\mathbf 1_\Omega\ip{\nabla G(Z),Z}\bigr]=\E\bigl[\mathbf 1_\Omega G(Z)\bigl(|Z|^2-d\bigr)\bigr],\\
&\mathrm{(b)}\ \ \E\bigl[\min(\zeta^+,\beta)\ip{\nabla G(Z),e}\bigr]=\E\bigl[G(Z)\bigl(\min(\zeta^+,\beta)\,\zeta-\mathbf 1\{0<\zeta<\beta\}\bigr)\bigr]\ \text{ for every }\beta\ge0.
\end{align*}
\end{lemma}

\begin{proof}
All the expectations are finite, since $\E\,e^{c'|Z|}|Z|^j<\infty$ for all $c'$ and $j$.

(a) For $t>0$ let $\Phi(t)=\int_\Omega G(ty)\gamma(y)\,dy$. For fixed $y$ the derivative of $t\mapsto G(ty)$ is $\ip{\nabla G(ty),y}$, bounded for $t\in[1/2,2]$ by $a\,e^{2c|y|}|y|$, which is integrable; by the mean value theorem and dominated convergence $\Phi'(1)=\E[\mathbf 1_\Omega\ip{\nabla G(Z),Z}]$. On the other hand $\Omega$ is a cone, so the change of variables $z=ty$ gives $\Phi(t)=\int_\Omega G(z)\,t^{-d}\gamma(z/t)\,dz$. With $\gamma(z)=(2\pi)^{-d/2}e^{-|z|^2/2}$, the derivative of $t\mapsto t^{-d}\gamma(z/t)$ is $t^{-d}\gamma(z/t)(|z|^2/t^3-d/t)$, which for $t\in[1/2,2]$ is at most $2^d(2\pi)^{-d/2}e^{-|z|^2/8}(8|z|^2+2d)$ in absolute value; times $a\,e^{c|z|}$ this bound is integrable. Dominated convergence gives $\Phi'(1)=\int_\Omega G(z)\gamma(z)(|z|^2-d)\,dz$, the right side of (a).

(b) Put $W'=Z-\zeta e$; it is independent of $\zeta$, which is a standard Gaussian variable. Fix $W'$ and let $f(t)=G(te+W')$, so that $f'(t)=\ip{\nabla G(te+W'),e}$ and $|f(t)|,|f'(t)|\le a\,e^{c|W'|}e^{c|t|}$. The function $g(t)=\min(t^+,\beta)\varphi(t)$ is continuous and vanishes on $(-\infty,0]$; it is of class $C^1$ on $[0,\beta]$ and on $[\beta,\infty)$ with $g'(t)=(\mathbf 1\{t<\beta\}-\min(t,\beta)\,t)\varphi(t)$. Integrating $g f'$ by parts on $[0,\beta]$ and on $[\beta,r]$, the terms at $\beta$ cancel and the term at $0$ vanishes; the term $g(r)f(r)$ tends to $0$ as $r\to\infty$, and both integrals converge absolutely, so
\[
\int\min(t^+,\beta)f'(t)\varphi(t)\,dt=\int f(t)\bigl(\min(t^+,\beta)\,t-\mathbf 1\{0<t<\beta\}\bigr)\varphi(t)\,dt.
\]
Taking the expectation in $W'$ gives (b); Fubini's theorem applies, both integrands being bounded by $a\,e^{c|W'|}e^{c|t|}(\beta|t|+\beta+1)\varphi(t)$.
\end{proof}

\begin{proof}[Proof of Lemma~\ref{lem:ec}]
Existence. Let $S$ be positive definite, and $s$ and $s_{\min}$ its largest and smallest eigenvalues. For a convex $G$ and a centred Gaussian vector $w$, $w$ and $-w$ have the same law, so $\E G(b+w)=\E[(G(b+w)+G(b-w))/2]\ge G(b)$; we call this symmetrisation. It gives $\E|b+SZ|\ge|b|$ and, from $|b+SZ|+|b-SZ|\ge2|SZ|$, $\E|b+SZ|\ge\E|SZ|\ge\E|\ip{v,SZ}|=s\,\E|Z_1|=s(2/\pi)^{1/2}$ for a unit eigenvector $v$ of $S$ with eigenvalue $s$. With $\Fbar(x)\ge a_1|x|-a_2$,
\[
\mathcal E_\lambda(b,S)\ge a_1\max\bigl(|b|,(2/\pi)^{1/2}s\bigr)-a_2-\lambda(d-1)\ln s-\lambda\ln s_{\min}.
\]
So $\mathcal E_\lambda$ tends to $+\infty$ when $|b|+s\to\infty$ or, with $s$ bounded, when $s_{\min}\to0$. It is continuous (dominated convergence, $\Fbar$ being Lipschitz), so it attains its minimum on a compact set $\{|b|\le r_1,\ r_2\le s_{\min},\ s\le r_1\}$ that contains a point where its value is below its infimum outside. The minimiser is unique by Lemma~\ref{lem:unique}, and (S4) is symmetrisation, $SZ$ being a centred Gaussian vector for every linear $S$.

Invertible $S$. For $S$ invertible let $P=(SS^{\mathsf T})^{1/2}$, which is positive definite. The centred Gaussian vectors $SZ$ and $PZ$ have the same covariance $SS^{\mathsf T}$, hence the same law, and $|\det S|=\det P$. So $\mathcal E_\lambda(b,S)=\mathcal E_\lambda(b,P)$: the infimum over the invertible $S$ is the minimum over the positive definite $S$, and the minimiser over the positive definite $S$ minimises over the invertible $S$.

Stationarity. Let $\Fbar$ be of class $C^1$ and $L$-Lipschitz, and $(b,S)$ minimise $\mathcal E_\lambda$ over the invertible $S$. For $db\in\R^d$ and a $d\times d$ matrix $dS$, the derivative of $t\mapsto\Fbar(b+t\,db+(S+t\,dS)Z)$ is bounded by $L(|db|+\|dS\||Z|)$, which is integrable, and $S+t\,dS$ is invertible for $t$ small; so $t\mapsto\mathcal E_\lambda(b+t\,db,S+t\,dS)$ is differentiable at $0$ with derivative $\E\ip{\nabla\Fbar(x),db+dS\,Z}-\lambda\tr(S^{-1}dS)$, and it vanishes there. With $dS=0$ this is (S1). With $db=0$ and $dS=Sv'v^{\mathsf T}$, for which $dS\,Z=\ip{v,Z}Sv'$ and $\tr(S^{-1}dS)=\ip{v,v'}$, it is (S2).

(S3). Write $w=SZ=\sum_i\ip{e_i,Z}Se_i$ for an orthonormal basis $(e_i)$; (S2) with $v=v'=e_i$, summed over $i$, gives $\E\ip{\nabla\Fbar(x),w}=\lambda d$. Let $z\in\R^d$ and $j(r)=\E\Fbar(z+r(x-z))$ for $r\in[0,1]$. It is convex and differentiable, with $j'(r)=\E\ip{\nabla\Fbar(z+r(x-z)),x-z}$ (the integrand is bounded by $L|x-z|$), $j(0)=\Fbar(z)$ and $j(1)=\E\Fbar(x)$. By (S1), $j'(1)=\ip{\E\nabla\Fbar(x),b-z}+\E\ip{\nabla\Fbar(x),w}=\lambda d$, and convexity gives $j(1)-j(0)\le j'(1)$.
\end{proof}

\begin{proof}[Proof of Lemma~\ref{lem:inf}]
Put $v=\mathbf 1_\Omega Z/R-k\,e$.

(i) The gradient inequality. Let $g\in\R^d$ with $\ip{e,g}>0$, and $k=\min(\ip{e,g},1)\in(0,1]$. The point $y=kRe$ lies on the segment from $0$ to $Re$, so by convexity $F(y)\le(1-k)F(0)+kF(Re)\le F(0)-k\Gamma$. The gradient inequality $F(y)\ge F(g)+\ip{\nabla F(g),y-g}$ of the convex differentiable $F$, divided by $R$, gives
\[
\ip{\nabla F(g),g/R-ke}\ge\frac{F(g)-F(y)}R\ge\frac{F(g)-F(0)}R+\frac{k\,\Gamma}R .
\]
So $\ip{\nabla F(Z),v}\ge m(F(Z)-F(0))+k\Gamma/R$ on $\Omega$, and outside $\Omega$ both sides vanish.

(ii) Integration by parts. Both $F$ and $G$ are Lipschitz, so they and their gradients are bounded by $a\,e^{|g|}$ for some $a$. For $H\in\{F,G\}$, Lemma~\ref{lem:ibp}(a) divided by $R$ minus Lemma~\ref{lem:ibp}(b) with $\beta=1$ gives $\E\ip{\nabla H(Z),v}=\E[H(Z)\ell]$, since $k\zeta=\mathbf 1_\Omega k\zeta$ and $\mathbf 1\{0<\zeta<1\}=\mathbf 1_\Omega\mathbf 1\{\zeta<1\}$. With (i) and $\E k=\kappa$,
\[
\E[F(Z)\ell]\ge\E\bigl[m\bigl(F(Z)-F(0)\bigr)\bigr]+\frac{\kappa\,\Gamma}R .
\]

(iii) The difference. Subtract $\E[G(Z)\ell]=\E[m\ip{\nabla G(Z),Z}]-\E[k\ip{\nabla G(Z),e}]$ and write $F(Z)-F(0)=(F(Z)-G(Z))-(F(0)-G(0))+(G(Z)-G(0))$:
\begin{multline*}
\E\bigl[(F(Z)-G(Z))(\ell-m)\bigr]+\bigl(F(0)-G(0)\bigr)\E m\\
\ge\frac{\kappa\,\Gamma}R+\E\bigl[k\ip{\nabla G(Z),e}\bigr]-\E\bigl[m\bigl(\ip{\nabla G(Z),Z}-G(Z)+G(0)\bigr)\bigr],
\end{multline*}
which is \eqref{eq:inf} once $\E m=1/(2R)$.

(iv) The moments. Since $\zeta$ is a standard Gaussian variable, $\E k=\E\min(Z_1^+,1)=\kappa$ and $\E m=\Prob(\zeta>0)/R=1/(2R)$. Write $|Z|^2=\zeta^2+\Xi$, with $\Xi$ independent of $\zeta$ and of chi-square law with $d-1$ degrees of freedom. On $\Omega$, $\ell-m=(|Z|^2-d-1)/R-(k\zeta-\mathbf 1\{\zeta<1\})$. The first term depends on $\zeta$ through its square, so $\E[(|Z|^2-d-1)^2;\Omega]=\frac12\E(|Z|^2-d-1)^2=\frac12(\Var|Z|^2+1)=(2d+1)/2$. The second term is $\min(\zeta,1)\zeta-\mathbf 1\{\zeta<1\}$ on $\Omega$, of second moment $\kappa_2$ there. The inequality $(y_1+y_2)^2\le2y_1^2+2y_2^2$ gives $\E(\ell-m)^2\le(2d+1)/R^2+2\kappa_2$.
\end{proof}

\begin{proof}[Proof of Lemma~\ref{lem:half}]
Let $W'=Z-\zeta e$, independent of $\zeta$, and $\psi(\theta)=\E H(\theta e+W')$. Differentiating under the expectation, the gradient of $H$ being bounded, $\psi$ is differentiable with $\psi'(\theta)=\E\ip{\nabla H(\theta e+W'),e}\ge0$, and $\psi'$ is bounded. So $|\psi|\le L$, $\psi$ is nondecreasing and $\int_\R\psi'\le2L$. By independence and Fubini's theorem, $\E[f(\zeta)H(Z)]=\E[f(\zeta)\psi(\zeta)]$ for every Borel $f$ with $\E|f(\zeta)|<\infty$. With $f=1$, $\E\psi(\zeta)=0$; with $f(t)=t$, $\E[\zeta\psi(\zeta)]=\lambda$; with $f=k$, $\E[kH(Z)]=\E[k\psi(\zeta)]$.

The covariance formula. For every real $t$, $\psi(t)=\psi(0)+\int\psi'(\theta)(\mathbf 1\{\theta<t\}-\mathbf 1\{\theta<0\})\,d\theta$. Multiply by $f(\zeta)-\E f(\zeta)$ at $t=\zeta$ and take the expectation; Fubini's theorem applies, the double integral of the absolute values being at most $2L\,\E|f(\zeta)-\E f(\zeta)|$. So $\operatorname{Cov}(f(\zeta),\psi(\zeta))=\int\psi'(\theta)\operatorname{Cov}(f(\zeta),\mathbf 1\{\zeta>\theta\})\,d\theta$, and, since $\E\psi(\zeta)=0$,
\[
\E[k\,\psi(\zeta)]=\int K(\theta)\,\psi'(\theta)\,d\theta,\qquad\lambda=\operatorname{Cov}\bigl(\zeta,\psi(\zeta)\bigr)=\int\varphi(\theta)\,\psi'(\theta)\,d\theta,
\]
with $K(\theta)=\operatorname{Cov}(k,\mathbf 1\{\zeta>\theta\})$ and $\operatorname{Cov}(\zeta,\mathbf 1\{\zeta>\theta\})=\E[\zeta;\zeta>\theta]=\varphi(\theta)$.

The kernel. We claim $K(\theta)\ge\kappa Q(|\theta|)$ for every $\theta$. For $\theta\le0$, $k$ vanishes where $\zeta<\theta$, so $K(\theta)=\E k-\E k\,Q(\theta)=\kappa Q(|\theta|)$. For $\theta\ge1$, $K(\theta)=Q(\theta)-\kappa Q(\theta)\ge\kappa Q(\theta)$ since $\kappa<1/2$. For $0\le\theta\le1$, $K(\theta)=\varphi(\theta)-\varphi(1)+Q(1)-\kappa Q(\theta)$, and the claim reads $g(\theta)=\varphi(\theta)-\varphi(1)+Q(1)-2\kappa Q(\theta)\ge0$. Here $g(0)=\kappa-\kappa=0$ by the value of $\kappa$, $g'(\theta)=\varphi(\theta)(2\kappa-\theta)$ is positive and then negative on $[0,1]$, and $g(1)=Q(1)(1-2\kappa)>0$; so $g\ge0$ on $[0,1]$.

Since $Q(|\theta|)=\Mills(|\theta|)\varphi(\theta)$ and $\Mills$ is nonincreasing (Lemma~\ref{lem:mills}), $K(\theta)\ge\kappa\Mills(T)\varphi(\theta)$ for $|\theta|\le T$, and $K\ge0$ everywhere. Hence
\[
\E[k\,\psi(\zeta)]\ge\kappa\Mills(T)\int_{|\theta|\le T}\varphi\,\psi'\ge\kappa\Mills(T)\Bigl(\lambda-\varphi(T)\int_\R\psi'\Bigr)\ge\kappa\Mills(T)\bigl(\lambda-2L\varphi(T)\bigr).
\]
\end{proof}

\section{The one-round inequality on a convex body}\label{sec:body}

This section proves Proposition~\ref{prop:body}. The losses of a law $\mu$ on $\cF(\cJ,n)$ are extended to $\R^d$, a penalty for leaving $\cJ$ is added so that Theorem~\ref{thm:dil} applies, and the law it gives is projected back onto $\cJ$. The projection changes the variance by the mass that the Gaussian puts outside $\cJ$, which the penalty and Gaussian concentration make small.

\begin{proof}[Proof of Proposition~\ref{prop:body}]
Since $\cJ$ contains a ball of radius $1/2$, $D_\cJ\ge1$; put $z=ndD_\cJ\ge2$. Fix
\[
\Gamma=C'z^{10},
\]
with $C'=10^8$, the constant of Lemma~\ref{lem:penalty} below; in particular $\Gamma\ge1$. Let $\mu$ be a finitely supported law on $\cF(\cJ,n)$, written as weights $\mu_c>0$ on functions $f_c$, with mean $\fbar$, and let $\Delta_\star=\min_\cJ\fbar-\sum_c\mu_c\min_\cJ f_c$. If $\Delta_\star\le1/n$, the point mass at a minimiser of $\fbar$ satisfies \eqref{eq:H} with $\alpha=1/n$. Assume $\Delta_\star>1/n$.

Extension and penalty. Let $\mathrm{pr}$ be the nearest-point projection onto $\cJ$ and $r(x)=\dist(x,\cJ)=|x-\mathrm{pr}(x)|$, a convex $1$-Lipschitz function. Put
\[
\tilde f_c(x)=\min_{y\in\cJ}\bigl(f_c(y)+n|x-y|\bigr),\qquad F_c=\tilde f_c+\Gamma r .
\]
The function $\tilde f_c$ is convex (an infimal convolution of convex functions) and $n$-Lipschitz, equals $f_c$ on $\cJ$ since $f_c$ is $n$-Lipschitz there, and satisfies $f_c(\mathrm{pr}\,x)\le\tilde f_c(x)\le f_c(\mathrm{pr}\,x)+nr(x)$: the upper bound is the choice $y=\mathrm{pr}\,x$, the lower bound holds because $|x-y|\ge|\mathrm{pr}\,x-y|$ for $y\in\cJ$ and $f_c$ is $n$-Lipschitz. So $F_c$ is convex and $(n+\Gamma)$-Lipschitz, $\min F_c=\min_\cJ f_c$, attained in $\cJ$, and the mean $\Fbar$ satisfies $\Fbar(x)\ge\fbar(\mathrm{pr}\,x)+\Gamma r(x)$. Hence $\Fbar$ is coercive, $\min\Fbar=\min_\cJ\fbar$ is attained at a point $b^\star\in\cJ$, the floor of the family $(\mu_c,F_c)$ is $\Delta_\star$, and its gains $g_c=f_c(b^\star)-\min_\cJ f_c$ lie in $[0,1]$. The differences $u_c=F_c-\Fbar$ do not depend on $\Gamma$, and with $u_c^\cJ=f_c-\fbar$ on $\cJ$, $|u_c(x)-u_c^\cJ(\mathrm{pr}\,x)|\le nr(x)$. With $V^\cJ=\sum_c\mu_c(u_c^\cJ)^2$ this gives
\begin{equation}\label{eq:proj-V}
V(x)\le2V^\cJ(\mathrm{pr}\,x)+2n^2r(x)^2 .
\end{equation}

The law. Theorem~\ref{thm:dil} applied to $(\mu_c,F_c)$ with $b^\star$ gives a scale $h\in[\Delta_\star/2,g_{\max}]\subset(1/(2n),1]$, the entropic Gaussian $(b_h,S_h)$ at temperature $\lambda_h=h/(16d)$, with $s_h$ the largest eigenvalue of $S_h$, and the law $p_h$. Let $\rho$ be the image of $p_h$ under $\mathrm{pr}$, a law on $\cJ$. Since $\fbar\circ\mathrm{pr}\le\Fbar$, the regret of $\rho$ is at most $\Delta(p_h)\le\frac{65}{64}\Delta_\star$. The law $p_h$ puts mass at least $q/2$ on each of $\delta_{b^\star}$, $\delta_{b_h}$ and $N(b_h,A_h)$, so \eqref{eq:P3}, with $h\ge\Delta_\star/2$, $K_S\le10^5$ (checked in Lean) and $g_{\max}\le1$, gives
\begin{equation}\label{eq:iota}
\mathcal I(p_h)\ge\frac q2\bigl(V(b^\star)+V(b_h)+\E V(x)\bigr)\ge\frac q2\cdot4\iota,\qquad \iota=\frac{\Delta_\star^2}{16\cdot10^5(1+T_h)^2d^2\ln(2/\Delta_\star)},
\end{equation}
where $x=b_h+S_hZ$ and $q=\Delta_\star/(4h)$.

The mass outside $\cJ$. Let $\bar m=\E r(x)$. Since $\Fbar(b_h)\ge\min\Fbar=\Fbar(b^\star)$,
\[
\Delta(p_h)=\Delta_\star+\tfrac q2\bigl(\Fbar(b_h)-\min\Fbar\bigr)+\tfrac q2\bigl(\E\Fbar(x)-\min\Fbar\bigr)\ge\Delta_\star+\tfrac q2\bigl(\E\Fbar(x)-\min\Fbar\bigr),
\]
and $\Delta(p_h)\le\frac{65}{64}\Delta_\star$ gives $\E\Fbar(x)-\min\Fbar\le(2/q)\Delta_\star/64=h/8$. Since $\Fbar\ge\min\Fbar+\Gamma r$, $\Gamma\bar m\le h/8\le1/8$, so $\bar m\le1/(8\Gamma)\le1/8$, and $r(b_h)\le\bar m$ by Jensen's inequality. The triangle inequality through $\mathrm{pr}\,x$ and $\mathrm{pr}\,b_h$ gives $|x-b_h|\le r(x)+D_\cJ+r(b_h)$, and $\E|x-b_h|=\E|S_hZ|\ge s_h(2/\pi)^{1/2}$, so
\begin{equation}\label{eq:sh}
s_h\le(\pi/2)^{1/2}(D_\cJ+2\bar m)\le(\pi/2)^{1/2}(D_\cJ+1/4).
\end{equation}
The function $Z\mapsto r(b_h+S_hZ)$ is $s_h$-Lipschitz, so Theorem~\ref{thm:gauss} gives $\Prob(r(x)\ge\bar m+t)\le\exp(-2t^2/(\pi^2s_h^2))=\exp(-t^2/(2s'^2))$ for $t\ge0$, with $s'=(\pi/2)s_h$. Fix $t_0\ge0$ and put $\tau=\bar m+t_0$. Pointwise $r^2\le2\tau r+(r-\tau)_+^2$, and
\[
\E(r(x)-\tau)_+^2=\int_0^\infty2t\,\Prob\bigl(r(x)>\tau+t\bigr)\,dt\le\int_0^\infty2t\,e^{-(t_0^2+t^2)/(2s'^2)}\,dt=2s'^2e^{-t_0^2/(2s'^2)},
\]
so $\E r(x)^2\le2(\bar m+t_0)\bar m+2s'^2e^{-t_0^2/(2s'^2)}$. Take
\[
t_0=s'\bigl(2\lnp(8n^2s'^2/\iota)\bigr)^{1/2}.
\]
If $8n^2s'^2\ge\iota$ then $2n^2s'^2e^{-t_0^2/(2s'^2)}=\iota/4$; otherwise $t_0=0$ and $2n^2s'^2<\iota/4$. We check that
\begin{equation}\label{eq:penalty}
\Gamma\ge\frac{n^2(3+2t_0)}{4\iota}.
\end{equation}
Let $A_0=10^5$, $c_1=(\pi/2)^{1/2}\frac{17}{16}$ and $T(z)$, $\iota_0(z)$, $s_0(z)$ be as in Lemma~\ref{lem:penalty}. By \eqref{eq:sh}, $s_h\le c_1z$ and $s'\le s_0(z)=(\pi/2)c_1z$. The bound $26d(n+\Gamma)s_h/h\le C_1z^{12}$ at the end of this proof, which does not use \eqref{eq:penalty}, gives $T_h\le T(z)$. Since $\Delta_\star^2/\ln(2/\Delta_\star)$ increases in $\Delta_\star\in(1/n,1]$, $nd\le z$ and $\ln(2n)\le\ln(2z)$, $\iota\ge\iota_0(z)$. Since $t_0$ is nondecreasing in $s'$ and in $1/\iota$, $t_0\le s_0(2\ln(8z^2s_0^2/\iota_0))^{1/2}$. So the right side of \eqref{eq:penalty} is at most $\mathrm{need}(z)$, which is at most $\Gamma=C'z^{10}$ by Lemma~\ref{lem:penalty}.
Since $\bar m\le1$, $n^2(\bar m^2+\E r(x)^2)\le n^2\bar m(3+2t_0)+\iota/4\le n^2(3+2t_0)/(8\Gamma)+\iota/4\le3\iota/4$.

The projection. Since $r(b^\star)=0$ and $r(b_h)\le\bar m$,
\[
n^2\int r^2\,dp_h\le\frac{q}{2}\,n^2\bigl(\bar m^2+\E r(x)^2\bigr)\le\frac{3q\iota}{8}\le\frac{3}{16}\,\mathcal I(p_h)
\]
by \eqref{eq:iota}, and \eqref{eq:proj-V} gives $\int V^\cJ\,d\rho\ge\frac12\mathcal I(p_h)-n^2\int r^2\,dp_h\ge\frac5{16}\mathcal I(p_h)$. With Theorem~\ref{thm:dil}, $\ln(2g_{\max}/\Delta_\star)\le\ln(2/\Delta_\star)<\ln(2n)$ and $16/5\le4$,
\[
\Bigl(\int\fbar\,d\rho-\sum_c\mu_c\min_\cJ f_c\Bigr)^2\le\Delta(p_h)^2\le4\cdot1.6\cdot10^6\,(1+T_h)^2d^2\ln(2n)\int V^\cJ\,d\rho .
\]
The integral on the right is the double integral of \eqref{eq:H}. It remains to bound $T_h$. With $n+\Gamma\le(1+C')z^{10}$, $h>1/(2n)$, $nd\le z$ and \eqref{eq:sh} with $D_\cJ+1/4\le\frac54D_\cJ\le\frac{17}{16}z$,
\[
\frac{26d(n+\Gamma)s_h}{h}\le C_1z^{12},\qquad C_1=52\,(\pi/2)^{1/2}\,\tfrac{17}{16}\,(1+C'),
\]
so $(1+T_h)^2\le2+2T_h^2\le2+4\ln C_1+48\ln z\le C_T\ln z$ with $C_T=(2+4\ln C_1)/\ln2+48$, since $\ln z\ge\ln2$. This is \eqref{eq:H} with $\alpha=1/n$ and $\beta=4\cdot1.6\cdot10^6\,C_Td^2\ln(ndD_\cJ)\ln(2n)$, and $4\cdot1.6\cdot10^6\,C_T\le C_0$ (checked in Lean); a larger $\beta$ weakens \eqref{eq:H}.
\end{proof}

The bounds on $s_h$, $h$, $T_h$, $\iota$ and $t_0$ in the proof are uniform over the law $\mu$ and the scale; this is what lets one penalty $\Gamma$, fixed before $\mu$, serve every law.

\begin{lemma}\label{lem:penalty}
Let $A_0=10^5$, $c_1=(\pi/2)^{1/2}\frac{17}{16}$, $C'=10^8$ and $C_1=52c_1(1+C')$, and for real $z\ge2$ put
\begin{gather*}
T(z)=\bigl(2\ln(C_1z^{12})\bigr)^{1/2},\qquad \iota_0(z)=\frac1{16A_0(1+T(z))^2z^2\ln(2z)},\qquad s_0(z)=\frac\pi2c_1z,\\
\mathrm{need}(z)=\frac{z^2}{4\iota_0}\Bigl(3+2s_0\bigl(2\ln(8z^2s_0^2/\iota_0)\bigr)^{1/2}\Bigr).
\end{gather*}
Then $\mathrm{need}(z)\le C'z^{10}$ for every real $z\ge2$.
\end{lemma}

\begin{proof}
Write $l=\ln z\ge\ln2$, $T=T(z)$ and $\Theta=\ln(8z^2s_0^2/\iota_0)=\ln(32\pi^2A_0c_1^2)+6l+2\ln(1+T)+\ln\ln(2z)$, so that $\mathrm{need}(z)=4A_0(1+T)^2z^4\ln(2z)(3+R)$ with $R=2s_0(2\Theta)^{1/2}=\pi c_1z(2\Theta)^{1/2}$, and
\[
\ln\frac{\mathrm{need}(z)}{C'z^{10}}=\ln\frac{4A_0}{C'}+2\ln(1+T)+\ln\ln(2z)+\ln(3+R)-6l,\qquad T^2=2\ln C_1+24l .
\]
All these are differentiable functions of $l$. The derivative of $\ln\ln(2z)$ is $1/\ln(2z)\le1/\ln4$. Since $dT/dl=12/T$, the derivative of $2\ln(1+T)$ is $24/(T(1+T))$, which decreases in $T$, and $T$ increases with $l$; so it is at most $q_2=24/(T_2(1+T_2))$ with $T_2=T(2)$. Each term of $\Theta$ increases with $l$, so $\Theta\ge\Theta_2$, its value at $z=2$, and $d\Theta/dl=6+24/(T(1+T))+1/\ln(2z)\le6+q_2+1/\ln4$. Since $dR/dl=R(1+(d\Theta/dl)/(2\Theta))$, the derivative of $\ln(3+R)$ is at most $1+(6+q_2+1/\ln4)/(2\Theta_2)$. Hence
\[
\frac{d}{dl}\ln\frac{\mathrm{need}(z)}{C'z^{10}}\le\frac1{\ln4}+q_2+1+\frac{6+q_2+1/\ln4}{2\Theta_2}-6 .
\]
Since $T_2\ge7$ and $\Theta_2\ge20$ (checked in Lean), $q_2\le1$, and $\ln4\ge1$, so the right side is at most $1+1+1+8/40-6<0$. So the ratio $\mathrm{need}(z)/(C'z^{10})$ decreases on $[2,\infty)$, and it is at most $1$ at $z=2$, since $\mathrm{need}(2)\le2^{10}C'$ (checked in Lean).
\end{proof}

\section{A cover of the inner body}\label{sec:ids}

The learner of Section~\ref{sec:explicit} plays on a finite subset of $\cJ$, so the comparator and the exploratory laws of $\mathrm{(H)}$ must be moved onto it. This section gives the cover and bounds the cost of the move; the argument is the one of \cite[proof of Theorem~2]{Lat20}.

Throughout, $\cJ\subset\R^d$ is a compact set of diameter $D_\cJ>0$, $n\ge1$, $L>0$, $\alpha,\beta\ge0$, and $\cF$ is a class of $L$-Lipschitz functions $\cJ\to[0,1]$ that satisfies $\mathrm{(H)}[\alpha,\beta]$. Each $f\in\cF$ is continuous on the compact set $\cJ$, so $f_\star=\min_\cJ f$ and the minimum over $\cJ$ of a finite sum of members of $\cF$ exist.

The next lemma gives a finite set and its size; it bounds the size of the net of the learner (Lemma~\ref{lem:net}).

\begin{lemma}\label{lem:cover}
Let $\varepsilon>0$.
\begin{enumerate}[label=(\alph*)]
\item Every finite set $Y\subset\cJ$ with $|y-y'|\ge\varepsilon$ for distinct $y,y'\in Y$ has at most $(1+2D_\cJ/\varepsilon)^d$ elements.
\item There is a nonempty finite set $\cC\subset\cJ$ of at most $(1+2D_\cJ/\varepsilon)^d$ elements, with $|y-y'|\ge\varepsilon$ for distinct $y,y'\in\cC$ and with every point of $\cJ$ at distance less than $\varepsilon$ from $\cC$.
\end{enumerate}
\end{lemma}

\begin{proof}
(a) We may assume $Y$ nonempty. The open balls of radius $\varepsilon/2$ about the points of $Y$ are disjoint, since a common point $z$ of two of them would give $|y-y'|\le|y-z|+|z-y'|<\varepsilon$. Fix $y_0\in Y$; each ball lies in the ball of radius $D_\cJ+\varepsilon/2$ about $y_0$. Comparing Lebesgue measures, $Y$ has at most $(D_\cJ+\varepsilon/2)^d/(\varepsilon/2)^d$ elements.

(b) By (a) the cardinalities of the finite separated subsets of $\cJ$ are bounded, and a one-point subset is one of them; let $\cC$ be one of largest cardinality. If a point of $\cJ$ were at distance at least $\varepsilon$ from all of $\cC$, it could be added to $\cC$.
\end{proof}

Fix $\varepsilon>0$, a finite set $\cC\subset\cJ$ such that every point of $\cJ$ is at distance less than $\varepsilon$ from $\cC$, such as the set of Lemma~\ref{lem:cover}(b), and an enumeration $c_1,\dots,c_N$ of $\cC$. Let $\Pi\colon\cJ\to\cC$ send $x$ to the point of least index among the points of $\cC$ nearest to $x$. Its fibres are defined by finitely many inequalities between continuous functions, so $\Pi$ is Borel, and $|x-\Pi(x)|<\varepsilon$. For a Borel probability law $\rho$ on $\cJ$, $\Pi_\#\rho$ is its image, a law on $\cC$.

The next lemma bounds the cost of moving the laws onto $\cC$, for Lemma~\ref{lem:conv} and Theorem~\ref{thm:exo}. For a finitely supported law $\mu$ on $\cF$ with mean $\fbar$ and a Borel probability law $\sigma$ on $\cJ$ put
\[
V_\mu(\sigma)=\int\!\!\int\bigl(\fbar(x)-f(x)\bigr)^2\,d\sigma(x)\,d\mu(f).
\]

\begin{lemma}\label{lem:moving}
\begin{enumerate}[label=(\alph*)]
\item For $f\in\cF^n$, $\min_\cC\sum_tf_t\le\min_\cJ\sum_tf_t+nL\varepsilon$.
\item Let $\mu$ be a finitely supported law on $\cF$ with mean $\fbar$, $\rho$ a Borel probability law on $\cJ$ and $\xi=\Pi_\#\rho$. Then $\int\fbar\,d\xi\le\int\fbar\,d\rho+L\varepsilon$ and $V_\mu(\rho)^{1/2}\le V_\mu(\xi)^{1/2}+L\varepsilon$. If $\rho$ satisfies \eqref{eq:H} for $\mu$, then
\begin{equation}\label{eq:moved}
\int\fbar\,d\xi-\int f_\star\,d\mu(f)\le\alpha+L\varepsilon\bigl(1+\beta^{1/2}\bigr)+\bigl(\beta V_\mu(\xi)\bigr)^{1/2}.
\end{equation}
\end{enumerate}
\end{lemma}

\begin{proof}
(a) Let $x^\star$ minimise $\sum_tf_t$ on $\cJ$; since each $f_t$ is $L$-Lipschitz, $\sum_tf_t(\Pi(x^\star))\le\sum_tf_t(x^\star)+nL\varepsilon$.

(b) The first inequality is $\int(\fbar(\Pi(x))-\fbar(x))\,d\rho(x)\le L\varepsilon$. For the second, work in $L^2(\rho\otimes\mu)$, a finite sum of copies of $L^2(\rho)$, with $u(x,f)=f(x)-\fbar(x)$. Then $\|u\|^2=V_\mu(\rho)$ and, by the definition of the image law, $\|u\circ(\Pi\times\mathrm{id})\|^2=V_\mu(\xi)$. Fix $x\in\cJ$ and put $W(f)=f(x)-f(\Pi(x))$, so that $|W|\le L\varepsilon$ on the support of $\mu$. Since $\fbar$ is the mean of $\mu$,
\[
u(x,f)-u(\Pi(x),f)=W(f)-\int W\,d\mu,\qquad\int\bigl(u(x,f)-u(\Pi(x),f)\bigr)^2\,d\mu(f)\le\int W^2\,d\mu\le L^2\varepsilon^2,
\]
the middle quantity being the variance of $W$ under $\mu$. Integrating in $x$ against $\rho$ and applying Minkowski's inequality gives $\|u\|\le\|u\circ(\Pi\times\mathrm{id})\|+L\varepsilon$. Finally \eqref{eq:moved} follows from the first inequality, \eqref{eq:H} and the second.
\end{proof}


\section{Lipschitz losses on an inner body}\label{sec:reduce}

This section proves Lemma~\ref{lem:inner}. A convex function with values in $[0,1]$ can be arbitrarily steep near the boundary of $\cK$, and on the boundary it need not even be Borel; on the inner body $\cJ$ it is $n$-Lipschitz, and playing in $\cJ$ costs at most $1$ against the best point of $\cK$.

\begin{proof}[Proof of Lemma~\ref{lem:inner}]
(a) The set $\cJ$ is convex: if $B(x,1/n)$ and $B(y,1/n)$ lie in $\cK$ and $\theta\in[0,1]$, a point $\theta x+(1-\theta)y+u$ with $|u|\le1/n$ is the convex combination $\theta(x+u)+(1-\theta)(y+u)$ of two points of $\cK$. It is closed, since $\cK$ is, and bounded, so it is compact. For $x\in\cK$ the point $y=x_0+(1-1/n)(x-x_0)$ lies in $\cJ$: for $|u|\le1/n$, $y+u=(1-1/n)x+\frac1n(x_0+nu)$ with $x_0+nu\in B(x_0,1)\subset\cK$. Taking $x\in B(x_0,1)$ shows $B(x_0,1-1/n)\subset\cJ$, so $\cJ$ has nonempty interior and diameter at least $2(1-1/n)\ge1$; it lies in $\cK$, so its diameter is at most $D$.

(b) Let $f\in\cG(\cK)$, $x\ne y$ in $\cJ$, $s=|y-x|$ and $z=y+\frac1n(y-x)/s$, a point of $B(y,1/n)\subset\cK$. With $\theta=(1/n)/(s+1/n)$, $y=\theta x+(1-\theta)z$, and convexity gives $f(y)-f(x)\le(1-\theta)(f(z)-f(x))\le1-\theta=s/(s+1/n)\le ns$. Exchanging $x$ and $y$ gives $|f(x)-f(y)|\le n|x-y|$; the restriction lies in $\cF(\cJ,n)$.

(c) Let $x\in\cK$ and $y=x_0+(1-1/n)(x-x_0)=(1-\frac1n)x+\frac1nx_0$, a point of $\cJ$ by (a). By convexity and $0\le f_t\le1$, $f_t(y)\le f_t(x)+\frac1n(f_t(x_0)-f_t(x))\le f_t(x)+\frac1n$. Summing over $t$ and taking the infimum over $x\in\cK$ gives the claim.
\end{proof}


\section{An explicit learner}\label{sec:explicit}

This section proves Theorem~\ref{thm:explicit} and Corollary~\ref{cor:minimax}. The learner runs exponential weights on a finite net $\cC$ of the inner body $\cJ$; in each round it solves a convex program for the law of its action and for a table of loss estimates, and it reads its observation only through a coin of that mean. None of the parts is new, and we credit each where it enters: the algorithm and its regret decomposition are exploration by optimisation \cite{LG21}, \cite[Section~8.5]{Lat24}; the coin is the device of \cite{AG12}, used as here, to make each round a program in finite dimension, by \cite{XZ25}, whose parameterisation of the adversary with Bernoulli rewards is the set $\mathcal A$ of Lemma~\ref{lem:advside} and whose formula for a fixed law of play, minimised over that law, gives the bound of Lemma~\ref{lem:dual}, with losses in place of rewards \cite[Appendix, Concave parameterization with Bernoulli reward]{XZ25}; the step from the one-round inequality to the information ratio of a round (Lemma~\ref{lem:conv}) is a step of the proof of \cite[Theorem~2]{Lat20}, used also in \cite{Lat24,XZ25}, and \cite{FKQR21} bounds its decision-estimation coefficient for convex bandits by a variant of \cite[Theorem~3]{Lat20}. The section is written so that its only input from the rest of the paper is the one-round inequality, Proposition~\ref{prop:body}; it uses besides Lemmas~\ref{lem:inner}, \ref{lem:cover} and~\ref{lem:moving}, Pinsker's inequality for two Bernoulli laws (Lemma~\ref{lem:pinsker}), the minimax theorem of Sion and the smoothing of the proof of Theorem~\ref{thm:gauss}; no compactness of a space of policies and no Tychonoff argument enters. Lemmas~\ref{lem:ew} and~\ref{lem:br} give the regret decomposition over the rounds; Lemmas~\ref{lem:mm1} and~\ref{lem:advside} exchange the learner and the adversary in one round; Lemmas~\ref{lem:gibbs} and~\ref{lem:dual} bound the value of the adversary, and Lemma~\ref{lem:conv} and Proposition~\ref{prop:value} bound it through $\mathrm{(H)}$; Lemmas~\ref{lem:class} and~\ref{lem:net} give the class on the net and the net, the latter through Lemma~\ref{lem:cover}. The proof of Theorem~\ref{thm:exo} assembles these and Theorem~\ref{thm:explicit} is Theorem~\ref{thm:exo} with the parameters chosen.

Throughout, $\cK$ is a convex body of $\R^d$ containing $B(x_0,1)$, $D\ge\diam\cK$ is a number known to the learner, $n\ge2$, $\cJ$ is the inner body of Section~\ref{sec:proof} and $L=n$. Every probability space below is finite, and no measurability is asked of anything.

\subsection{The algorithm}

For a finite set $\cC\subset\R^d$ let $\Delta(\cC)$ be the probability vectors on $\cC$, $\Delta^\circ(\cC)$ those with all coordinates positive, $\ip{u,v}=\sum_{x\in\cC}u(x)v(x)$ for $u,v\in\R^\cC$, and
let $\mathcal P_L$ be the set of the vectors $v\in[0,1]^\cC$ for which there are $s_x\in\R^d$ ($x\in\cC$) with $|s_x|\le L$ and
\[
v(z)\ge v(x)+\ip{s_x,z-x}\qquad\text{for all }x,z\in\cC,
\]
that is, the restrictions to $\cC$ of the convex $L$-Lipschitz functions on $\R^d$ with values in $[0,1]$ on $\cC$ (Lemma~\ref{lem:class}). For $q\in\Delta^\circ(\cC)$ and $\eta>0$, let $\mathcal S_q$ be the set of pairs $(p,g)$, $p\in\Delta(\cC)$ and $g=(g(x,b))_{x\in\cC,\,b\in\{0,1\}}$ with $g(x,b)\in\R^\cC$, such that for every $(x,b)$
\begin{equation}\label{eq:Sq}
\sum_{z\in\cC}q(z)\exp\bigl(-\eta\,g(x,b)(z)/p(x)\bigr)\le1\ \text{ if }p(x)>0,\qquad g(x,b)\ge0\ \text{ if }p(x)=0 .
\end{equation}
For $(p,g)\in\mathcal S_q$, $f\in[0,1]^\cC$ and $y\in\cC$ put
\begin{gather}
\mathcal R(p,g;f,y)=\sum_xp(x)f(x)-f(y)+\sum_x\bigl[(1-f(x))\,g(x,0)(y)+f(x)\,g(x,1)(y)\bigr],\label{eq:Rround}\\
\overline{\mathcal R}(p,g)=\max_{y\in\cC}\ \max_{f\in\mathcal P_L}\mathcal R(p,g;f,y).\notag
\end{gather}

\begin{algorithm}[exploration by optimisation with a coin]\label{alg:exo}
\leavevmode
\begin{center}
\fbox{\begin{minipage}{0.93\textwidth}
\textbf{Input:} $d$, $n\ge2$, a number $D\ge\diam\cK$, a point $x_0$ with $B(x_0,1)\subset\cK$, a membership test for $\cJ$; parameters $\varepsilon\in(0,1/4]$, $\eta>0$, $\mathrm{tol}>0$; $L=n$.

\textbf{Net.} Let $h=\varepsilon/(2d^{1/2}(D+1))$. Scan, in a fixed order, the points of $h\mathbb Z^d$ in the cube $x_0+[-D,D]^d$ that lie in $\cJ$, and keep a point if its distance to every point kept so far is at least $\varepsilon/2$. Let $\cC$ be the set kept and $m=|\cC|$.

\textbf{State:} $q_t\in\Delta^\circ(\cC)$, with $q_1$ uniform.

\textbf{Round} $t=1,\dots,n$:
\begin{enumerate}[label=\arabic*.,leftmargin=2em]
\item Compute from $q_t$, by a deterministic procedure, $(P_t,G_t)\in\mathcal S_{q_t}$ with $\overline{\mathcal R}(P_t,G_t)\le\inf_{\mathcal S_{q_t}}\overline{\mathcal R}+\mathrm{tol}$.
\item Draw $X_t$ from $P_t$, play it, and observe $o_t=f_t(X_t)$.
\item Draw $B_t\in\{0,1\}$ with $\Prob(B_t=1)=o_t$, with fresh randomness, and forget $o_t$.
\item Put $\hat\ell_t=G_t(X_t,B_t)/P_t(X_t)\in\R^\cC$.
\item Put $q_{t+1}(y)=q_t(y)e^{-\eta\hat\ell_t(y)}\big/\sum_{z}q_t(z)e^{-\eta\hat\ell_t(z)}$ for $y\in\cC$.
\end{enumerate}
\end{minipage}}
\end{center}
\end{algorithm}

The infimum in step~1 is finite (Proposition~\ref{prop:value}) and need not be attained, which is why the step asks for a tolerance. Since $X_t$ is drawn from $P_t$, $P_t(X_t)>0$ and $\hat\ell_t$ is a finite vector. Step~1 minimises over the convex set $\mathcal S_{q_t}$ the maximum of the affine functions $(p,g)\mapsto\mathcal R(p,g;f,y)$, $f\in\mathcal P_L$, $y\in\cC$ (Lemma~\ref{lem:advside}). The quantity $\mathcal R$ bounds, for $X$ drawn from $p$, $B$ a coin of mean $f(X)$ and $\hat\ell=g(X,B)/p(X)$, the expectation of the loss of the action against the comparator $y$, plus the estimated loss of $y$, plus the increment of the log-partition function of the weights, which \eqref{eq:Sq} makes nonpositive (Lemma~\ref{lem:br}); exponential weights add these bounds over the rounds (Lemma~\ref{lem:ew}). The regret of the learner against $f=(f_1,\dots,f_n)$ is $\Reg_n(f)=\E\sum_tf_t(X_t)-\inf_\cK\sum_tf_t$, the expectation being over its draws, actions and coins.

\begin{theorem}\label{thm:exo}
Let $\alpha,\beta\ge0$ be such that $\cF(\cJ,n)$ satisfies $\mathrm{(H)}[\alpha,\beta]$, and let $\varepsilon\in(0,1/4]$, $\eta>0$ and $\mathrm{tol}>0$. For every $f\in\cG(\cK)^n$ fixed before the game, the learner of Algorithm~\ref{alg:exo} has
\[
\Reg_n(f)\le1+nL\varepsilon+\frac{\ln m}\eta+n\Bigl[\kappa_\varepsilon\alpha+L\varepsilon\bigl(1+\beta^{1/2}\bigr)+\frac{\eta\beta}8+\mathrm{tol}\Bigr],\qquad\kappa_\varepsilon=1+2L\varepsilon,
\]
with $2\le m\le(5D/\varepsilon)^d$. The same bound holds for $\sup_{x\in\cK}\E\sum_t(f_t(X_t)-f_t(x))$ when $f_t=A_t(X_1,B_1,\dots,X_{t-1},B_{t-1})$ for maps $A_t$ into $\cG(\cK)$.
\end{theorem}

Against an adaptive adversary this is a bound on the pseudo-regret, the supremum outside the expectation; the proof gives nothing for $\E\sup_x$, the estimates $\hat\ell_t$ being unbounded.

\subsection{Exponential weights and one round}

\begin{lemma}\label{lem:ew}
Let $\eta>0$, $\hat\ell_1,\dots,\hat\ell_n\in\R^\cC$ arbitrary, $q_1$ uniform, $Z_t=\sum_zq_t(z)e^{-\eta\hat\ell_t(z)}$ and $q_{t+1}(y)=q_t(y)e^{-\eta\hat\ell_t(y)}/Z_t$. Then for every $y\in\cC$, $\sum_t\bigl(\hat\ell_t(y)+\eta^{-1}\ln Z_t\bigr)\ge-\eta^{-1}\ln m$.
\end{lemma}

\begin{proof}
Let $\hat L_t=\hat\ell_1+\dots+\hat\ell_t$ and $W_t=\sum_ze^{-\eta\hat L_t(z)}$, so $W_0=m$. By induction $q_t(z)=e^{-\eta\hat L_{t-1}(z)}/W_{t-1}$, so $Z_t=W_t/W_{t-1}$, and $\sum_t\ln Z_t=\ln W_n-\ln m\ge-\eta\hat L_n(y)-\ln m$.
\end{proof}

\begin{lemma}\label{lem:br}
Let $q\in\Delta^\circ(\cC)$, $(p,g)\in\mathcal S_q$, $f\in[0,1]^\cC$ and $y\in\cC$. Let $X$ have law $p$, $B\in\{0,1\}$ with $\Prob(B=1\mid X)=f(X)$, and $\hat\ell=g(X,B)/p(X)$. Then
\[
\E\Bigl[f(X)-f(y)+\hat\ell(y)+\eta^{-1}\ln\sum_zq(z)e^{-\eta\hat\ell(z)}\Bigr]\le\mathcal R(p,g;f,y).
\]
Moreover, with $c_y(x)=p(x)-\mathbf 1\{x=y\}+g(x,1)(y)-g(x,0)(y)$,
\begin{equation}\label{eq:affine}
\mathcal R(p,g;f,y)=\sum_xg(x,0)(y)+\ip{c_y,f}.
\end{equation}
\end{lemma}

\begin{proof}
With $\pi(1\mid x)=f(x)$ and $\pi(0\mid x)=1-f(x)$, the left side is
\[
\sum_xp(x)f(x)-f(y)+\sum_{x:p(x)>0}\sum_b\pi(b\mid x)\Bigl[g(x,b)(y)+\frac{p(x)}\eta\ln\sum_zq(z)e^{-\eta g(x,b)(z)/p(x)}\Bigr].
\]
For $p(x)>0$ the logarithm is at most $0$ by \eqref{eq:Sq}; for $p(x)=0$ the terms $\pi(b\mid x)g(x,b)(y)$ of $\mathcal R$ are nonnegative, since $g(x,b)\ge0$ there. Dropping the logarithms and adding these terms gives $\mathcal R(p,g;f,y)$. The identity \eqref{eq:affine} is a rearrangement of \eqref{eq:Rround}.
\end{proof}

The coin makes $\mathcal R$ affine in $f$. Had the learner fed its estimator with the observation $f(X)$ itself, $g(X,f(X))$ would enter at the value $f(X)$, the round's payoff would not be affine in $f$, the adversary's mixed strategies would not reduce to their conditional means, and the minimax step below would need compactness in infinite dimension, as in \cite{LG21}.

\subsection{The minimax step of a round}

\begin{lemma}\label{lem:mm1}
Let $E$ and $E'$ be real topological vector spaces, $\mathcal A\subset E$ nonempty, compact and convex, $Y\subset E'$ convex, $c$ real, and $\Phi\colon E'\times E\to\R$ such that $\Phi(y,\cdot)$ is continuous on $E$ and concave on $\mathcal A$ for every $y\in Y$, and $\Phi(\cdot,a)$ is continuous and convex on $Y$ for every $a\in\mathcal A$. If every $a\in\mathcal A$ has some $y\in Y$ with $\Phi(y,a)<c$, then some $y\in Y$ has $\Phi(y,a)<c$ for every $a\in\mathcal A$.
\end{lemma}

\begin{proof}
For $y\in Y$ the set $K_y=\{a\in\mathcal A:\Phi(y,a)\ge c\}$ is closed in $\mathcal A$, and by assumption no point lies in all of them. Since $\mathcal A$ is compact, finitely many, $K_{y_1},\dots,K_{y_k}$, have empty intersection, so every $a\in\mathcal A$ has some $i$ with $\Phi(y_i,a)<c$. The minimax theorem of Sion \cite{Sion58}, in its form for finitely many points of $Y$ \cite{Komiya88}, applied to the function $(a,y)\mapsto-\Phi(y,a)$, which is continuous and convex in $a$ on the compact convex set $\mathcal A$ and continuous and concave in $y$ on the convex set $Y$, gives some $y\in Y$ with $\Phi(y,a)<c$ for every $a\in\mathcal A$.
\end{proof}

Compactness is asked of one side only, the adversary's, which the coin makes finite-dimensional. The adversary's mixed strategies are the pairs $(w,\theta)$ in
\[
\mathcal A=\bigl\{(w,\theta):\ w\in\Delta(\cC),\ \theta=(\theta_y)_{y\in\cC},\ \theta_y\in w_y\mathcal P_L\bigr\}\subset\R^\cC\times(\R^\cC)^\cC,
\]
with the payoff $\Phi(p,g;w,\theta)=\sum_y\bigl[w_y\sum_xg(x,0)(y)+\ip{c_y,\theta_y}\bigr]$. For $w_y>0$ put $f_y=\theta_y/w_y\in\mathcal P_L$; by \eqref{eq:affine}, $\Phi(p,g;w,\theta)=\sum_{y:w_y>0}w_y\mathcal R(p,g;f_y,y)$, the payoff of an adversary that draws the comparator $y$ from $w$ and then a loss vector with conditional mean $f_y$. Put $\underline{\mathcal R}(w,\theta)=\inf_{(p,g)\in\mathcal S_q}\Phi(p,g;w,\theta)$.

\begin{lemma}\label{lem:advside}
The set $\mathcal A$ is convex and compact, and $\mathcal S_q$ is convex. Let $c$ be real, and assume that every $(w,\theta)\in\mathcal A$ has some $(p,g)\in\mathcal S_q$ with $\Phi(p,g;w,\theta)<c$. Then some $(p,g)\in\mathcal S_q$ has
\begin{equation}\label{eq:roundminimax}
\mathcal R(p,g;f,y)<c\qquad\text{for every }y\in\cC\text{ and every }f\in\mathcal P_L .
\end{equation}
\end{lemma}

\begin{proof}
The class $\mathcal P_L$ is convex and compact (Lemma~\ref{lem:class}), so $\mathcal A$ is the image of the compact $\Delta(\cC)\times\mathcal P_L^\cC$ under the continuous map $(w,(f_y))\mapsto(w,(w_yf_y))$. It is convex: for $(w,\theta),(w',\theta')\in\mathcal A$ and $l\in[0,1]$, if $w''_y=lw_y+(1-l)w'_y>0$ then $l\theta_y+(1-l)\theta'_y=w''_yf''_y$ with $f''_y$ the convex combination of $f_y$ and $f'_y$ with the weights $lw_y/w''_y$ and $(1-l)w'_y/w''_y$, and if $w''_y=0$ both terms vanish. The payoff $\Phi$ is linear and continuous in $(w,\theta)$, and in $(p,g)$, $c_y$ is affine and so is $\Phi$, which is therefore continuous. For the convexity of $\mathcal S_q$, consider $E_q=\{(s,u)\in[0,\infty)\times\R^\cC:\ s>0\text{ and }\sum_zq(z)e^{-\eta u(z)/s}\le1,\text{ or }s=0\text{ and }u\ge0\}$. For $s>0$, $(s,u)\mapsto s\ln\sum_zq(z)e^{-\eta u(z)/s}$ is the perspective of a convex function, hence convex, so the part of $E_q$ with $s>0$ is convex. A combination of $(s_1,u_1)$, $s_1>0$, with $(0,u_2)$, $u_2\ge0$, is $(ls_1,lu_1+(1-l)u_2)$, and $\sum_zq(z)e^{-\eta(lu_1+(1-l)u_2)(z)/(ls_1)}\le\sum_zq(z)e^{-\eta u_1(z)/s_1}\le1$; two points with $s=0$ combine to one with $s=0$ and $u\ge0$. So $E_q$ is convex, and $\mathcal S_q$, defined by $(p(x),g(x,b))\in E_q$ for every $(x,b)$, is convex. Lemma~\ref{lem:mm1} with $Y=\mathcal S_q$ gives $(p,g)\in\mathcal S_q$ with $\Phi(p,g;w,\theta)<c$ for every $(w,\theta)\in\mathcal A$. At $w=e_y$, $\theta_y=f$ and $\theta_{y'}=0$ for $y'\ne y$, a point of $\mathcal A$, the formula above gives $\Phi(p,g;w,\theta)=\mathcal R(p,g;f,y)$, which is \eqref{eq:roundminimax}.
\end{proof}

\subsection{The value of a round}

For $a,b\in[0,1]$ let $\mathrm{kl}(a,b)=a\ln(a/b)+(1-a)\ln((1-a)/(1-b))$, with $0\ln0=0$, and for $r\in\Delta(\cC)$ and $q\in\Delta^\circ(\cC)$ let $\KL(r\|q)=\sum_{y:r(y)>0}r(y)\ln(r(y)/q(y))$.

\begin{lemma}\label{lem:gibbs}
For $r\in\Delta(\cC)$, $q\in\Delta^\circ(\cC)$ and $\eta>0$, $\inf\bigl\{\ip{r,e}:\ e\in\R^\cC,\ \sum_zq(z)e^{-\eta e(z)}\le1\bigr\}=-\KL(r\|q)/\eta$.
\end{lemma}

\begin{proof}
Let $R_+=\{r>0\}$. For a feasible $e$, Jensen's inequality for the logarithm with the weights $r(y)$ on $R_+$ gives $0\ge\ln\sum_{y\in R_+}r(y)\frac{q(y)}{r(y)}e^{-\eta e(y)}\ge-\KL(r\|q)-\eta\ip{r,e}$. Conversely, for $M>0$ let $u(y)=-\eta^{-1}\ln(r(y)/q(y))$ and $e_M(y)=\min(u(y),M)$ on $R_+$, $e_M=M$ off $R_+$, and $e=e_M+\eta^{-1}\ln(1+e^{-\eta M})$. On $R_+$, $e^{-\eta e_M}\le r/q+e^{-\eta M}$, so $\sum_zq(z)e^{-\eta e_M(z)}\le1+e^{-\eta M}$ and $e$ is feasible, and $\ip{r,e}\le-\KL(r\|q)/\eta+\eta^{-1}\ln(1+e^{-\eta M})$, which tends to $-\KL(r\|q)/\eta$.
\end{proof}

\begin{lemma}\label{lem:dual}
Let $w\in\Delta(\cC)$, $f_y\in[0,1]^\cC$ and $\theta_y=w_yf_y$ for $y\in\cC$, $\fbar=\sum_yw_yf_y$ and $I_x(w,\theta)=\sum_{y:w_y>0}w_y\,\mathrm{kl}(f_y(x),\fbar(x))$ for $x\in\cC$. For every $x_0\in\cC$ and $\delta>0$ there is $g$ with $(e_{x_0},g)\in\mathcal S_q$ and
\[
\Phi(e_{x_0},g;w,\theta)\le\fbar(x_0)-\frac{I_{x_0}(w,\theta)+\KL(w\|q)}\eta-\sum_yw_yf_y(y)+\delta .
\]
\end{lemma}

Letting $\delta\to0$ and minimising over $x_0$, $\underline{\mathcal R}(w,\theta)\le\min_{x\in\cC}[\fbar(x)-(I_x(w,\theta)+\KL(w\|q))/\eta]-\sum_yw_yf_y(y)$.

\begin{proof}
The kl terms are finite: $\fbar(x)=0$ forces $f_y(x)=0$ for all $y$ with $w_y>0$, and $\fbar(x)=1$ forces $f_y(x)=1$. By \eqref{eq:affine}, $\Phi=\ip{p,\fbar}-\sum_yw_yf_y(y)+\sum_{x,b}\ip{W_{x,b},g(x,b)}$ with $W_{x,1}(y)=w_yf_y(x)$ and $W_{x,0}(y)=w_y(1-f_y(x))$, nonnegative vectors of masses $\fbar(x)$ and $1-\fbar(x)$. For fixed $p$ the constraints \eqref{eq:Sq} bind each $g(x,b)$ separately. If $p(x)>0$, Lemma~\ref{lem:gibbs} with $g(x,b)=p(x)e$ gives $\inf\ip{W_{x,b},g(x,b)}=-p(x)k_{x,b}/\eta$, where $k_{x,b}=\Pi\,\KL(r\|q)$ if $W_{x,b}$ has mass $\Pi>0$ and $r=W_{x,b}/\Pi$, and $k_{x,b}=0$ if $W_{x,b}=0$; if $p(x)=0$ the infimum over $g(x,b)\ge0$ is $0$. Take $p=e_{x_0}$, $g(x,b)=0$ for $x\ne x_0$, and $g(x_0,b)$ within $\delta/2$ of the infimum for $b=0,1$. Then $(e_{x_0},g)\in\mathcal S_q$ and $\Phi(e_{x_0},g;w,\theta)\le\fbar(x_0)-(k_{x_0,0}+k_{x_0,1})/\eta-\sum_yw_yf_y(y)+\delta$. Splitting each logarithm, with sums over $y$ with $w_y>0$,
\[
k_{x,1}+k_{x,0}=\sum_yw_yf_y(x)\ln\frac{w_yf_y(x)}{\fbar(x)q(y)}+\sum_yw_y(1-f_y(x))\ln\frac{w_y(1-f_y(x))}{(1-\fbar(x))q(y)}=I_x(w,\theta)+\KL(w\|q).
\]
With $x=x_0$ this is the claim.
\end{proof}

The bound on $\underline{\mathcal R}(w,\theta)$ after the lemma is an equality, which the formal proof does not contain: for $(p,g)\in\mathcal S_q$ the infima of the proof give $\ip{W_{x,b},g(x,b)}\ge-p(x)k_{x,b}/\eta$, and $\ip{W_{x,b},g(x,b)}\ge0$ where $p(x)=0$, so $\Phi(p,g;w,\theta)\ge\sum_xp(x)[\fbar(x)-(I_x(w,\theta)+\KL(w\|q))/\eta]-\sum_yw_yf_y(y)$, which is at least the minimum over $x$.

The one-round inequality enters through the next lemma. Let $\cF''_L$ be the class of convex $L$-Lipschitz functions $\cJ\to[-L\varepsilon,1+L\varepsilon]$, and let $\Pi\colon\cJ\to\cC$ be the nearest-point map of Section~\ref{sec:ids} for the net $\cC$.

\begin{lemma}\label{lem:conv}
Assume that $\cF(\cJ,L)$ satisfies $\mathrm{(H)}[\alpha,\beta]$ and that every point of $\cJ$ is at distance less than $\varepsilon$ from $\cC\subset\cJ$. Let $(w,\theta)\in\mathcal A$, $f_y=\theta_y/w_y$ for $w_y>0$, and $\fbar=\sum_yw_yf_y$. There is $\xi\in\Delta(\cC)$ with
\[
\sum_x\xi(x)\fbar(x)-\sum_yw_yf_y(y)\le\kappa_\varepsilon\alpha+L\varepsilon\bigl(1+\beta^{1/2}\bigr)+\Bigl(\tfrac12\beta\,I(\xi)\Bigr)^{1/2},\qquad I(\xi)=\sum_x\xi(x)I_x(w,\theta).
\]
\end{lemma}

\begin{proof}
For $w_y>0$ let $h_y\in\cF''_L$ be the extension of $f_y$ of Lemma~\ref{lem:class}(iii), equal to $f_y$ on $\cC$, and $\mu=\sum_yw_y\delta_{h_y}$, with mean $\bar h=\sum_yw_yh_y$, equal to $\fbar$ on $\cC$. The affine map $T(h)=(h+L\varepsilon)/\kappa_\varepsilon$ sends $\cF''_L$ into $\cF(\cJ,L/\kappa_\varepsilon)$, a subclass of $\cF(\cJ,L)$ closed under finite convex combinations, which therefore satisfies $\mathrm{(H)}[\alpha,\beta]$; it commutes with means and minima. Let $\rho$ be the law that $\mathrm{(H)}$ gives for $T_\#\mu$ and $\xi=\Pi_\#\rho$. Lemma~\ref{lem:moving}(b), whose proof uses of the class only the Lipschitz bound and of $\cC$ only that every point of $\cJ$ is within $\varepsilon$ of it, applied to the class $\cF(\cJ,L/\kappa_\varepsilon)$ with the Lipschitz constant $L/\kappa_\varepsilon$, gives \eqref{eq:moved} for $T_\#\mu$; multiplied by $\kappa_\varepsilon$ it reads
\[
\sum_x\xi(x)\bar h(x)-\int h_\star\,d\mu(h)\le\kappa_\varepsilon\alpha+L\varepsilon\bigl(1+\beta^{1/2}\bigr)+\bigl(\beta V_\mu(\xi)\bigr)^{1/2}.
\]
On $\cC$, $\bar h=\fbar$ and $h_y=f_y$, so $V_\mu(\xi)=\sum_x\xi(x)\sum_yw_y(\fbar(x)-f_y(x))^2$. Since $y\in\cC\subset\cJ$, $f_y(y)=h_y(y)\ge\min_\cJ h_y$, so $\sum_yw_yf_y(y)\ge\int h_\star\,d\mu$. Lemma~\ref{lem:pinsker} with the means $f_y(x)$ and $\fbar(x)$ gives $(f_y(x)-\fbar(x))^2\le\frac12\mathrm{kl}(f_y(x),\fbar(x))$, so $V_\mu(\xi)\le I(\xi)/2$.
\end{proof}

\begin{proposition}\label{prop:value}
Under the assumptions of Lemma~\ref{lem:conv}, let $q\in\Delta^\circ(\cC)$ and $\mathcal G=\kappa_\varepsilon\alpha+L\varepsilon(1+\beta^{1/2})+\eta\beta/8$. Every $(p,g)\in\mathcal S_q$ has $\mathcal R(p,g;0,y)\ge0$ for some $y\in\cC$, and for every $\mathrm{tol}>0$ some $(p,g)\in\mathcal S_q$ has $\mathcal R(p,g;f,y)<\mathcal G+\mathrm{tol}$ for every $y\in\cC$ and every $f\in\mathcal P_L$.
\end{proposition}

Since the zero vector lies in $\mathcal P_L$, the infimum of $\overline{\mathcal R}$ over $\mathcal S_q$ is therefore finite and at most $\mathcal G$.

\begin{proof}
Let $(p,g)\in\mathcal S_q$ and $x\in\cC$. If $p(x)>0$, then $1-t\le e^{-t}$ and \eqref{eq:Sq} give $\sum_zq(z)(1-\eta g(x,0)(z)/p(x))\le1$, that is $\sum_zq(z)g(x,0)(z)\ge0$; if $p(x)=0$, then $g(x,0)\ge0$. So the $q$-average of $y\mapsto\mathcal R(p,g;0,y)=\sum_xg(x,0)(y)$ is nonnegative, and some $y$ has $\mathcal R(p,g;0,y)\ge0$.

Fix $\mathrm{tol}>0$ and $(w,\theta)\in\mathcal A$, and let $\xi$ be given by Lemma~\ref{lem:conv}. By Lemma~\ref{lem:dual} and the bound after it, dropping $\KL(w\|q)\ge0$ and bounding the minimum over $x$ by the $\xi$-average,
\[
\underline{\mathcal R}(w,\theta)\le\kappa_\varepsilon\alpha+L\varepsilon\bigl(1+\beta^{1/2}\bigr)+\bigl(\tfrac12\beta I\bigr)^{1/2}-\frac I\eta,\qquad I=I(\xi)\ge0,
\]
and $as-s^2/\eta\le a^2\eta/4$ with $a=(\beta/2)^{1/2}$ and $s=I^{1/2}$ gives $\underline{\mathcal R}(w,\theta)\le\mathcal G$. So some $(p,g)\in\mathcal S_q$ has $\Phi(p,g;w,\theta)<\mathcal G+\mathrm{tol}$, and Lemma~\ref{lem:advside} with $c=\mathcal G+\mathrm{tol}$ gives the claim.
\end{proof}

\subsection{The class on the net and the net}

\begin{lemma}\label{lem:class}
Let $\cC\subset\cJ$ be finite.
\begin{enumerate}[label=(\roman*)]
\item For $F\colon\R^d\to\R$ convex and $L$-Lipschitz and $x\in\R^d$ there is $s\in\R^d$ with $|s|\le L$ and $F(z)\ge F(x)+\ip{s,z-x}$ for every $z$.
\item $\mathcal P_L$ is convex and compact, and contains the restriction to $\cC$ of every $f\in\cF(\cJ,L)$.
\item If every point of $\cJ$ is at distance at most $\varepsilon$ from $\cC$, every $v\in\mathcal P_L$, with witnesses $(s_x)$, is the restriction to $\cC$ of $h_v(z)=\max_{x\in\cC}(v(x)+\ip{s_x,z-x})$, which is convex and $L$-Lipschitz with values in $[-L\varepsilon,1+L\varepsilon]$ on $\cJ$.
\end{enumerate}
\end{lemma}

\begin{proof}
(i) For $r>0$ let $F^r$ be the function of Lemma~\ref{lem:smooth} for $F$ with $\epsilon=r$. It is convex, $L$-Lipschitz and of class $C^2$, so $|\nabla F^r|\le L$, and $|F^r-F|\le r$. The gradient inequality of the convex differentiable function $F^r$ at $x$ gives $F(z)\ge F^r(z)-r\ge F^r(x)+\ip{s_r,z-x}-r\ge F(x)+\ip{s_r,z-x}-2r$ with $s_r=\nabla F^r(x)$. A subsequence of $(s_{1/k})$, which lies in the closed ball of radius $L$, converges to some $s$ with $|s|\le L$, and the inequality passes to the limit for each $z$.

(ii) The defining inequalities are linear in $(v,s)$ and the ball of radius $L$ is convex, so $\mathcal P_L$ is convex. It is bounded, and closed: if $v_k\to v$ with witnesses $s^k$, a subsequence of $s^k$ converges and the inequalities pass to the limit. For $f\in\cF(\cJ,L)$, the extension $\tilde f(z)=\min_{y\in\cJ}(f(y)+L|z-y|)$ is convex, $L$-Lipschitz and equal to $f$ on $\cJ$, as in the proof of Proposition~\ref{prop:body} with $L$ in place of $n$; (i) at each $x\in\cC$ gives the witnesses.

(iii) The function $h_v$ is a maximum of affine functions with gradients of norm at most $L$. At $z\in\cC$ the term $x=z$ gives $h_v(z)\ge v(z)$, and every term is at most $v(z)$ by the inequalities of $\mathcal P_L$. For $z\in\cJ$ and $x\in\cC$ with $|z-x|\le\varepsilon$, $|h_v(z)-v(x)|\le L\varepsilon$.
\end{proof}

\begin{lemma}\label{lem:net}
For $n\ge2$ and $\varepsilon\in(0,1/4]$, the net $\cC$ of Algorithm~\ref{alg:exo} lies in $\cJ$, its distinct points are at distance at least $\varepsilon/2$, every point of $\cJ$ is at distance less than $\varepsilon$ from $\cC$, and $2\le m\le(5D/\varepsilon)^d$.
\end{lemma}

\begin{proof}
By Lemma~\ref{lem:inner}(a), $\cJ$ contains $B(x_0,r_0)$ with $r_0=1-1/n\ge1/2$, and $\cJ\subset\cK\subset B(x_0,D)$. Let $x\in\cJ$, $t=hd^{1/2}/(2r_0)\le1$ and $y=(1-t)x+tx_0$. For $|u|\le tr_0$, $y+u=(1-t)x+t(x_0+u/t)$ with $x_0+u/t\in B(x_0,r_0)$, so $B(y,tr_0)\subset\cJ$. Rounding the coordinates of $y$ gives a point $z\in h\mathbb Z^d$ with $|z-y|\le hd^{1/2}/2=tr_0$, so $z\in B(y,tr_0)\subset\cJ\subset x_0+[-D,D]^d$ is a scanned point, and $|x-z|\le t(|x-x_0|+r_0)\le hd^{1/2}(D/(2r_0)+1/2)\le hd^{1/2}(D+1/2)<\varepsilon/2$. Every scanned point is kept or is within $\varepsilon/2$ of a kept point, so $x$ is within $\varepsilon$ of $\cC$. The set $\cC$ is $(\varepsilon/2)$-separated in $\cJ$, so Lemma~\ref{lem:cover}(a) bounds $m$ by $(1+4\diam\cJ/\varepsilon)^d$, and $\diam\cJ\ge1\ge\varepsilon$ gives $1+4\diam\cJ/\varepsilon\le5D/\varepsilon$. Finally $x_0\pm e/2$, for a unit vector $e$, lie in $\cJ$ at distance $1$, each within $\varepsilon\le1/4$ of $\cC$, so $m\ge2$.
\end{proof}

\subsection{Proof of Theorems~\ref{thm:exo} and~\ref{thm:explicit}}

\begin{proof}[Proof of Theorem~\ref{thm:exo}]
By Lemma~\ref{lem:inner}(b) each $f_t$ restricted to $\cJ$ lies in $\cF(\cJ,n)$, and by Lemma~\ref{lem:class}(ii) its restriction to $\cC$, still written $f_t$, lies in $\mathcal P_L$. By Lemma~\ref{lem:net}, $\cC\subset\cJ$ and every point of $\cJ$ is within $\varepsilon$ of $\cC$. The solver is deterministic, so $(q_t,P_t,G_t)$ is a function of $H_{t-1}=(X_1,B_1,\dots,X_{t-1},B_{t-1})$, which takes finitely many values; $\Prob(H_n=(x_s,b_s)_s)=\prod_tP_t(x_t)\pi_t(b_t\mid x_t)$ with $\pi_t(1\mid x)=f_t(x)$.

Fix $y\in\cC$. Lemma~\ref{lem:ew}, applied on each path, gives $\sum_t(f_t(X_t)-f_t(y))\le\eta^{-1}\ln m+\sum_t\mathrm{Br}_t(y)$ with $\mathrm{Br}_t(y)=f_t(X_t)-f_t(y)+\hat\ell_t(y)+\eta^{-1}\ln\sum_zq_t(z)e^{-\eta\hat\ell_t(z)}$. Given $H_{t-1}$, $X_t$ has law $P_t$, $B_t$ is a coin of mean $f_t(X_t)$ and $(P_t,G_t)\in\mathcal S_{q_t}$, so Lemma~\ref{lem:br}, step~1 of the algorithm and Proposition~\ref{prop:value} give
\[
\E[\mathrm{Br}_t(y)\mid H_{t-1}]\le\mathcal R(P_t,G_t;f_t,y)\le\overline{\mathcal R}(P_t,G_t)\le\mathcal G+\mathrm{tol}.
\]
Hence $\E\sum_tf_t(X_t)-\sum_tf_t(y)\le\eta^{-1}\ln m+n(\mathcal G+\mathrm{tol})$ for every $y\in\cC$. By Lemma~\ref{lem:moving}(a), $\min_\cC\sum_tf_t\le\min_\cJ\sum_tf_t+nL\varepsilon$, and by Lemma~\ref{lem:inner}(c), $\inf_\cJ\sum_tf_t\le\inf_\cK\sum_tf_t+1$. This is the bound, with $m$ bounded by Lemma~\ref{lem:net}.

Against $f_t=A_t(H_{t-1})$, $f_t$ is fixed given $H_{t-1}$, so the conditional bound holds as before and Lemma~\ref{lem:ew} is pathwise: $\E\sum_t(f_t(X_t)-f_t(y))\le\eta^{-1}\ln m+n(\mathcal G+\mathrm{tol})$ for each fixed $y\in\cC$. For $x\in\cK$ let $x'=x_0+(1-1/n)(x-x_0)\in\cJ$ and $y_x=\Pi(x')$; on every path $f_t(y_x)\le f_t(x')+n\varepsilon\le f_t(x)+1/n+n\varepsilon$, as in the proofs of Lemmas~\ref{lem:inner}(c) and~\ref{lem:moving}(a). Taking $y=y_x$ gives the bound for $\E\sum_t(f_t(X_t)-f_t(x))$, for every $x\in\cK$.
\end{proof}

\begin{proof}[Proof of Theorem~\ref{thm:explicit}]
By Lemma~\ref{lem:inner}(a), $\cJ$ is a convex body containing a ball of radius $1/2$, with $\diam\cJ\le D$, so Proposition~\ref{prop:body} gives $\mathrm{(H)}[1/n,\beta]$ for $\cF(\cJ,n)$ with $\beta=C_0d^2\ln(ndD)\ln(2n)$, $C_0=1.2\cdot10^9$ (a larger $\beta$ weakens \eqref{eq:H}, and $\diam\cJ\le D$), and $\beta\ge C_0\ln^24\ge2$. Then $\varepsilon=1/(n^2(2\beta)^{1/2})\le1/4$, and Theorem~\ref{thm:exo} with $\alpha=1/n$, $\eta=(8\ln m/(n\beta))^{1/2}$ and $\mathrm{tol}=1/n$ gives, since $nL\varepsilon=n^2\varepsilon$ and $n\kappa_\varepsilon\alpha=1+2n\varepsilon$,
\[
\frac{\ln m}\eta+\frac{n\eta\beta}8=\Bigl(\tfrac12\beta n\ln m\Bigr)^{1/2},\qquad 1+n^2\varepsilon+(1+2n\varepsilon)+n^2\varepsilon(1+\beta^{1/2})+1=3+2n\varepsilon+n^2\varepsilon(2+\beta^{1/2}),
\]
with $2n\varepsilon=2/(n(2\beta)^{1/2})\le1/2$ and $n^2\varepsilon(2+\beta^{1/2})=2(2\beta)^{-1/2}+2^{-1/2}\le1+2^{-1/2}$, since $n\ge2$ and $(2\beta)^{1/2}\ge2$. So the regret is at most $3+\frac12+1+2^{-1/2}+(\frac12\beta n\ln m)^{1/2}\le6.71+(\frac12\beta n\ln m)^{1/2}$, and $\ln m\le d\ln(5n^2(2\beta)^{1/2}D)$ by Lemma~\ref{lem:net}.

For the order, let $\mathcal L=\ln(nD)\ge\ln4$. If $d\ge n$ the regret is at most $n\le d^{3/2}n^{1/2}\mathcal L^{3/2}$, the losses taking values in $[0,1]$. If $d<n$, then $\ln(ndD)\le2\mathcal L$ and $\ln(2n)\le2\ln n\le2\mathcal L$, so $\beta\le4C_0d^2\mathcal L^2$ and $5n^2(2\beta)^{1/2}D\le5(8C_0)^{1/2}n^2dD\mathcal L$, and since $2\ln n+\ln d+\ln D\le3\mathcal L$, $\ln\mathcal L\le\mathcal L$ and $5(8C_0)^{1/2}\le2^{19}=4^{19/2}$ (checked in Lean), its logarithm is at most $\frac{27}2\mathcal L$. The root is therefore at most $(27C_0)^{1/2}d^{3/2}n^{1/2}\mathcal L^{3/2}=1.8\cdot10^5\,d^{3/2}n^{1/2}\mathcal L^{3/2}$, and $6.71\le10^4\,d^{3/2}n^{1/2}\mathcal L^{3/2}$ since $\mathcal L\ge\ln4\ge1$. This gives the bound with $C_X=1.8\cdot10^5+10^4=1.9\cdot10^5$, and $C_X\ge1$ covers the case $d\ge n$.
\end{proof}

\subsection{The policy}\label{sec:explicit-cost}

The learner keeps its coins in its state, while a policy in the sense of Section~\ref{sec:intro} reads only actions and observations. Marginalising the coins gives one with the same law of play. Write $P_t(\cdot\mid x_{<t},b_{<t})$ for the law the learner computes in round $t$ from the earlier actions and coins and $\mathrm{Bern}(o)$ for the law on $\{0,1\}$ of mean $o$; off the earlier actions and coins of positive probability, $P_t$ is extended arbitrarily, and the formulas below multiply it there by $0$.

\begin{lemma}\label{lem:marg}
For a history $\chi=(x_{<t},o_{<t})\in(\cC\times[0,1])^{t-1}$ put
\begin{gather*}
\phi_t(b_{<t}\mid\chi)=\frac{\prod_{s<t}P_s(x_s\mid x_{<s},b_{<s})\,\mathrm{Bern}(o_s)(b_s)}{\sum_{b'\in\{0,1\}^{t-1}}\prod_{s<t}P_s(x_s\mid x_{<s},b'_{<s})\,\mathrm{Bern}(o_s)(b'_s)},\\
\varpi_t(x\mid\chi)=\sum_{b_{<t}}\phi_t(b_{<t}\mid\chi)P_t(x\mid x_{<t},b_{<t})
\end{gather*}
when the denominator is positive, and $\varpi_t(\cdot\mid\chi)$ uniform on $\cC$ otherwise; for a history in $(\cK\times[0,1])^{t-1}$ with an action outside $\cC$, let $\varpi_t(\cdot\mid\chi)$ be uniform on $\cC$ as well. Then $\varpi$ is a $\cC$-policy, a Borel function of the history, and against every $f\in\cG(\cK)^n$ the actions have the same law under $\varpi$ as under the learner; in particular the two have the same regret.
\end{lemma}

\begin{proof}
For fixed actions $x_{<t}\in\cC^{t-1}$ the denominator is a polynomial in $o_{<t}$ and $\varpi_t$ a ratio of polynomials where it is positive; the histories with given actions $x_{<t}\in\cC^{t-1}$ form a Borel set, there are finitely many such sets, and $\varpi_t$ is uniform off their union; so $\varpi$ is Borel. Fix $f$ and put $o_s=f_s(x_s)$. Under the learner,
\[
\Prob(X_{1:t}=x_{1:t},\,B_{<t}=b_{<t})=\prod_{s<t}\bigl[P_s(x_s\mid x_{<s},b_{<s})\,\mathrm{Bern}(o_s)(b_s)\bigr]\,P_t(x_t\mid x_{<t},b_{<t}),
\]
since $X_s$ is drawn from $P_s$ and then $B_s$ from $\mathrm{Bern}(f_s(X_s))$. Summing over $x_t$ and $b_{<t}$ shows that the denominator of $\phi_t(\cdot\mid x_{<t},o_{<t})$ is $\Prob(X_{<t}=x_{<t})$. When it is positive, summing over $b_{<t}$ and dividing by it gives $\Prob(X_t=x_t\mid X_{<t}=x_{<t})=\varpi_t(x_t\mid x_{<t},o_{<t})$; the chain rule then gives, for every $x_{1:n}$, $\Prob(X_{1:n}=x_{1:n})=\prod_t\varpi_t(x_t\mid x_{<t},o_{<t})$, both sides vanishing from the first $t$ with $\Prob(X_{1:t}=x_{1:t})=0$. This is the law \eqref{eq:play} of $\varpi$ against $f$.
\end{proof}

So the learner of Theorem~\ref{thm:explicit} gives a $\cC$-policy with the same regret bound; its round $t$ is a sum over $2^{t-1}$ strings of coins, and the explicit object is the learner with its coins.

\begin{proof}[Proof of Corollary~\ref{cor:minimax}]
Let $x_0$ be the centre of a unit ball in $\cK$; then $D\ge2$. For $n\ge2$, run the learner of Theorem~\ref{thm:explicit} with this $D$, and let $\varpi$ be the $\cC$-policy of Lemma~\ref{lem:marg}. Then $\Reg_n(\varpi,f)$ is the regret of the learner, at most $C_Xd^{3/2}\sqrt n\ln^{3/2}(nD)$ for every $f$. For $n=1$ the $\{x_0\}$-policy that plays $x_0$ has regret at most $1\le C_Xd^{3/2}\ln^{3/2}2$.
\end{proof}

\section{A learner whose every step takes polynomial time}\label{sec:poly}

This section proves Remark~\ref{rem:upper-prob}, from which Theorem~\ref{thm:upper} follows (Section~\ref{sec:intro}), and Theorem~\ref{thm:polytope} (Section~\ref{sec:poly-polytope}). It defines the machine and the run of a program against a strategy (Section~\ref{sec:poly-machine}), the learner (Section~\ref{sec:poly-learner}) and the positioning it starts with (Section~\ref{sec:poly-pg}); the remark then follows from two propositions, one on the program and one on the regret, by an argument on each run (Section~\ref{sec:poly-proof}). Section~\ref{sec:poly-regret} proves the proposition on the regret. Throughout, $L_\delta=\ln(2dn/\delta)$.

\subsection{The machine}\label{sec:poly-machine}

\begin{definition}\label{def:machine}
A \emph{program} is a finite list of instructions acting on a state made of a counter, natural registers $N_0,N_1,\dots$, real cells $M_0,M_1,\dots$, the number of Gaussian draws made, the number of rounds played and the points played. Its instructions are:
\begin{enumerate}[label=(\roman*)]
\item a rational constant into a real cell; $M_a\circ M_b$ into a cell, for $\circ\in\{+,-,\times,/\}$, with $x/0=0$; the square root, the exponential or the logarithm of a cell, with $\sqrt x=0$ for $x<0$ and $\ln x=\ln|x|$; a jump if $M_a<M_b$;
\item a constant into a natural register; $N_a+N_b$ and the truncated difference $\max(N_a-N_b,0)$; a jump if $N_a<N_b$; the copy of $N_a$ into a real cell; an unconditional jump;
\item a spectral instruction: the symmetric part of a $d\times d$ block of cells is replaced, in a block of cells, by the matrix with the same eigenvectors and with each eigenvalue $x$ replaced by $\max(a,\min(b,x))$, or by $1$ if $x>a$ and $0$ otherwise, where $a,b$ are cells;
\item a membership question: a cell receives $1$ if the point held by $d$ cells lies in $\cK$, and $0$ otherwise;
\item a Gaussian draw: a cell receives the next unused term of the sequence $(g_i)_{i\ge0}$;
\item a play: the point held by $d$ cells is played as $x_t$, where $t$ is the number of rounds played, and a cell receives $o_t=f_t(x_t)+\varepsilon_t$.
\end{enumerate}
The cells are addressed through the natural registers. A run starts with $N_0=d$, $N_1=n$, $M_0=r_0$, $M_1=R_0$, $M_2=\delta$ and every other register and cell at $0$; one \emph{step} is one instruction.
\end{definition}

The sequence $(g_i)$ has the law of independent standard Gaussian variables, and the noise $(\varepsilon_t)$ the law $\nu$, independent of it. A run is determined by $\cK$, the losses, the noise and the Gaussian sequence. Against a strategy $\sigma$ the loss $f_t$ is $\sigma_t(x_1,o_1,\dots,x_{t-1},o_{t-1})$, read off the state when round $t$ is played; the strategy never reads the Gaussian sequence. In Theorem~\ref{thm:upper} and Remark~\ref{rem:upper-prob} the cost clause, ``the first play within $B_0$ steps and each later play within $B$ steps of the one before'', and the clause that every point played lies in $\cK$ hold for every Gaussian sequence and every noise sequence; the regret clause is an inequality for the outer measure of the product of the Gaussian law and $\nu$ in the remark, and for the integral of a measurable majorant in the theorem, so that no measurability of the run is asked.

The machine computes with exact real numbers, as \cite{BEL21} count operations for their Theorem~1; a spectral instruction is one step, as one call of an eigenvalue routine.

\subsection{The learner}\label{sec:poly-learner}

The learner DF2-O first puts $\cK$ in a position (Procedure PG, Section~\ref{sec:poly-pg}), then runs an online Newton method on the positioned body cut by a ball. A \emph{position} is a pair $(\mu_S,A_S)$ with $A_S$ an invertible $d\times d$ matrix; the body in this position is $S(\cK)=\{y:\mu_S+A_Sy\in\cK\}$, so that a membership question for $S(\cK)$ is one for $\cK$. Its parameters, for $d,n\ge1$ and $0<\delta<1$, are
\begin{gather*}
\delta_1=\min\Bigl(\frac\delta3,\frac{e^{-14}}{(dn)^3}\Bigr),\quad \delta_e=\frac{4\delta_1}{5n},\quad L_e=36\ln\frac1{\delta_e},\quad \lambda=\frac1{725\,dL_e^2},\quad \eta=\frac{\sqrt d}{L_e\sqrt n},\quad \sigma^2=\frac1{8L_e},\\
\rho=1450\,d^2L_e^3,\quad \varepsilon=\frac{14500\,d^{3/2}L_e^5}{\sqrt n},\quad \kappa=\alpha=\frac14,\quad \theta=1,\quad c_1=8,\quad \gamma=\frac1{100},\quad \kappa_0=(1-\lambda)^{-2},\\
R=134\,d\,\bigl(\ln(2d+2)\bigr)^{1/2},\quad j=\Bigl\lceil\log_2\frac{nR(R+7)}{\delta_1\varepsilon}\Bigr\rceil,\quad s_b=\frac{\delta_e^2}{R(R+1)}.
\end{gather*}
Let $\cK'=S(\cK)\cap B(0,R)$ and $\cK_\varepsilon=(1-\varepsilon)\cK'$. The gauge $\tilde\pi(y)$ of $\cK'$ at $y\ne0$ is the upper end of the interval left by $j$ steps of bisection for $\inf\{s>0:y/s\in\cK'\}$ started from $[|y|/R,|y|]$, and $\tilde\pi(0)=0$. For a closed convex set $Q\subset\R^d$, a ball $B(c,r)\subset Q$, a positive definite matrix $P$, a point $z$ and an accuracy $e>0$, the subroutine $\mathrm{AP}(Q,c,r,P,z,e)$ is an ellipsoid method with a separation obtained from membership questions (on the line for $d=1$); it fails or returns a point of $Q\cap B(0,R)$, and when $Q\subset B(0,R)$ it returns a point $\mu\in Q$ with $\frac12\ip{P(\mu-z),\mu-z}\le\frac12\ip{P(x-z),x-z}+e$ for every $x\in Q$. Appendix~\ref{app:ap} gives AP and proves this (Proposition~\ref{prop:ap}).

\begin{algorithm}[DF2-O]\label{alg:df2o}
\leavevmode
\begin{center}
\fbox{\begin{minipage}{0.93\textwidth}
\textbf{Input:} $d,n,r_0,R_0,\delta$, membership questions for $\cK$, Gaussian draws.

\textbf{Start.} Run Procedure PG with confidence $\delta_1$ on its own block of the Gaussian sequence. If it returns no position, or if $\varepsilon>\frac12$, play $0$ in every round. Otherwise start an epoch.

\textbf{Start of an epoch:} $\mu=0$, $P=\sigma^{-2}I$, $\Sigma=\sigma^2I$, $A=\sigma I$, $b=0$, $k=0$, $G=0$, $l=0$, the ball $(c,r)=(0,1-\varepsilon)$, and $\Phi(y)=\frac12\ip{Py,y}-\ip{b,y}+k$.

\textbf{Round} $t$, with $\xi\sim N(0,I_d)$ and $\zeta=g/2$, $g\sim N(0,1)$, read from the block of round $t$:
\begin{enumerate}[label=\arabic*.,leftmargin=2em]
\item $w=A\xi$, $X=\mu+(1-\lambda)w$, $p=\max\bigl(1,\tilde\pi(X)/(1-\varepsilon)\bigr)$, $x=\mu_S+A_S(X/p)$. If $x\notin\cK$, play $0$ in this round and every later one. Otherwise play $x$, observe $o$, and put $Y=p(o+\zeta)/2+2(p-1)/\varepsilon$.
\item $u=Pw$, $c_Y=\eta\kappa\lambda\kappa_0Y$, $a=1-c_Y$, $q=a+c_Y|\xi|^2$. If $a\le0$ or $q\le0$, act the guard.
\item $P'=aP+c_Y uu^{\mathsf T}$, $\Sigma'=a^{-1}\bigl(\Sigma-(c_Y/q)ww^{\mathsf T}\bigr)$, $A'=a^{-1/2}\bigl(A-c'w\xi^{\mathsf T}\bigr)$ with $c'=(c_Y/q)/(1+(a/q)^{1/2})$; $b'=b+c_Y\bigl(\ip{u,\mu}u-P\mu\bigr)$, $k'=k+\frac{c_Y}2\bigl(\ip{u,\mu}^2-\ip{P\mu,\mu}\bigr)$, $G'=G+\frac Y{1-\lambda}u$, $l'=l+\frac Y{1-\lambda}\ip{u,\mu}$; $\Phi'(y)=\frac12\ip{P'y,y}-\ip{b',y}+k'$. If $\tr P'>d/\delta_e$ or $\tr\Sigma'>d/\delta_e$, act the guard.
\item $z=\Sigma'(b'-\frac\eta2G')$. If $|z|>R/\delta_e^3$, act the guard.
\item Test. Let $Q^-=\cK_\varepsilon$ in the first round of the epoch, and $Q^-=\{y\in\cK_\varepsilon:\Phi(y)\le(c_1-\gamma)\rho\}$ otherwise; $p^-=\mathrm{AP}(Q^-,c,r,P',z,\min(1,\alpha\rho/2))$; act the guard if it fails. If $\eta(l'-\ip{G',p^-})-2\Phi'(p^-)\le-\theta\rho$, start a new epoch at the next round and skip steps~6 and~7.
\item $c\leftarrow(1-s_b)p^-$, $r\leftarrow s_b(1-\varepsilon)$; $w_\mu=\mu-\frac{\eta Y}{(1-\lambda)q}w$. If $|w_\mu|>R+\eta(16dnL_e/\delta_e)^{1/2}$, act the guard.
\item $\mu\leftarrow\mathrm{AP}\bigl(\{y\in\cK_\varepsilon:\Phi'(y)\le c_1\rho\},c,r,P',w_\mu,1/(16n^2\max(1,4R^2\tr P'))\bigr)$; act the guard if it fails. Replace $(P,\Sigma,A,b,k,G,l,\Phi)$ by the primed values.
\end{enumerate}
After the guard acts, the learner plays $\mu_S$ in every later round, or $0$ if $\mu_S\notin\cK$.
\end{minipage}}
\end{center}
\end{algorithm}

Every point played lies in $\cK$: it has passed a membership question, or it is $0\in B(0,r_0)\subset\cK$. The learner reads the loss of a round only through the observation of that round, so its play in round $t$ depends on the losses of the rounds before $t$ only. The updates of step~3 are the closed forms of $\Sigma'=P'^{-1}$ and $A'A'^{\mathsf T}=\Sigma'$ when $a,q>0$ (Lemma~\ref{lem:r1}).

\subsection{The position}\label{sec:poly-pg}

For a convex body $\cK'$ with $0$ in its interior, let $M(\cK')=\E\,\pi_{\cK'}(g)/\E|g|$, $g\sim N(0,I_d)$, where $\pi_{\cK'}$ is the gauge of $\cK'$. Procedure PG grows balls: from the position $(0,r_0I)$ it runs $m=\max(1,\lceil\log_2(R_0/r_0)\rceil)$ stages, stage $i$ positioning the body $\cK\cap B(0,2^ir_0)$ seen in the current position, with confidence $\delta_1/m$ on its own block of the Gaussian sequence, and composing the position it returns with the current one. A stage runs, a fixed number of times, a projected stochastic descent on pairs $(\mu,T)$, $T$ symmetric with eigenvalues in a fixed interval, for the objective $\E\,\tilde\pi(\mu+Tg)-\frac\kappa d\ln\det T$, with the gauge computed by bisection, followed by a test of the average iterate amplified by a majority vote; it returns the first iterate that passes. Appendix~\ref{app:pg} writes out Procedure PG with its parameters.

\begin{theorem}\label{thm:pg}
Let $d\ge1$, $0<r_0\le R_0$, $0<\delta<1$ and $\cK\subset\R^d$ closed and convex with $B(0,r_0)\subset\cK\subset B(0,R_0)$. With probability at least $1-\delta$ (in outer measure), Procedure PG run with confidence $\delta$ returns a position $(\mu_S,A_S)$ with
\[
B(0,1)\subset S(\cK)\subset B\bigl(0,134\,d\,(\ln(2d+2))^{1/2}\bigr)\quad\text{and}\quad M(S(\cK))\le1.27\,\frac{(d+1)^{1/2}}d.
\]
\end{theorem}

The proof is in Appendix~\ref{app:pg}: Lemma~\ref{lem:l} gives the Gaussian position, Proposition~\ref{prop:cert} certifies it from three moments, Theorem~\ref{thm:pgcore} finds it by a projected stochastic descent and a test, and growing balls remove the dependence on $R_0/r_0$.

\subsection{Proof of Remark~\ref{rem:upper-prob}}\label{sec:poly-proof}

Let $\mathrm{DF2O}_t(\cK,f,\varepsilon,g)$ be the point that Algorithm~\ref{alg:df2o} plays in round $t$ when the losses are a fixed sequence $f$, the noise is $\varepsilon$ and the Gaussian sequence is $g$. Against a strategy $\sigma$ its play $\hat x_t$ is defined by induction: $\hat x_t=\mathrm{DF2O}_t(\cK,f^{\hat x},\varepsilon,g)$, where $f^{x}_s=\sigma_s(x_1,f^x_1(x_1)+\varepsilon_1,\dots,x_{s-1},f^x_{s-1}(x_{s-1})+\varepsilon_{s-1})$ are the losses that $\sigma$ chooses along the points $x$; the induction is well founded since the play of round $t$ reads the losses of the rounds before $t$ only.

\begin{proposition}\label{prop:mach}
There are a program $P$, a real constant $C'$ and an integer $k\ge0$ such that for all $d,n\ge1$, $0<r_0\le R_0$, $0<\delta<1$, every compact convex $\cK$ with $B(0,r_0)\subset\cK\subset B(0,R_0)$, every sequence $f$ of convex functions $\cK\to[0,1]$ and all sequences $\varepsilon$ and $g$: the run of $P$ on $d,n,r_0,R_0,\delta$ makes its first play within $C'(dL_\delta(2+\ln(R_0/r_0)))^k$ steps and each later play of a round $t<n$ within $C'(dL_\delta)^k$ steps of the one before, and its point of round $t<n$ is $\mathrm{DF2O}_t(\cK,f,\varepsilon,g)$; the rounds are numbered from $0$ here, as in the formal statement.
\end{proposition}

\begin{proof}
The program $P$ is written in a structured language over the instructions of Definition~\ref{def:machine} (sequence, the two conditionals and the two loops on the comparisons of (i) and (ii)) and compiled to a list of instructions; the jumps the compiler inserts are counted as steps. It has three parts: the code of Procedure PG, the code of AP, and the code of a round of Algorithm~\ref{alg:df2o}.

\emph{The points played.} The code of PG ends, on every Gaussian sequence, with the position that Procedure PG returns or with the mark of its failure, and reads the Gaussian sequence in the blocks of the procedure; the projection on the domain $D$ of PG-core is one spectral instruction (the clamp of the eigenvalues) and a rescaling of $\mu$, and the test $\hat W\ge\frac58\frac\kappa dI$ is the other one. The code of AP returns, on every input of a call made by Algorithm~\ref{alg:df2o}, the point or the failure of the subroutine of Appendix~\ref{app:ap}. The code of a round makes the updates of steps~1 to~7. Composing the three, the point that $P$ plays in round $t<n$ is $\mathrm{DF2O}_t(\cK,f,\varepsilon,g)$.

\emph{The steps.} The loops of the code of PG have lengths fixed by $d,r_0,R_0$ and $\delta_1$ before any answer is read, so this code takes the same number $k_{\mathrm{PG}}$ of steps on every run. Each parameter of a stage is at most a power of $dL_\delta(2+\ln(R_0/r_0))$, since $\rho_i\le26244\,d^2$ and $n$ enters only through $\delta_1$, inside the logarithm; hence $k_{\mathrm{PG}}\le C(dL_\delta(2+\ln(R_0/r_0)))^{k}$. Outside the calls of AP, the steps before the first play, PG excepted, and the steps of a round are at most $1006\,d^2(j+1)$, with $j\le70L_\delta$. Each call of AP stops within a cap that is a polynomial of degree $5$ in the counts of the ellipsoid method of Appendix~\ref{app:ap}; the parameters and the guards of Algorithm~\ref{alg:df2o} bound every input of the two calls of a round (the radii, $\tr P'$, the norm of the centre, $1/e$) by $10^{18}(2dn/\delta)^{14}$, so these counts are $O(d^2L_\delta)$ and each cap is at most $C(dL_\delta)^k$. Adding up gives the two bounds of the proposition with one constant $C'$ and one exponent $k$.
\end{proof}

\begin{proposition}\label{prop:regret}
There is a constant $C$ such that for all $d,n\ge1$, $0<r_0\le R_0$, $0<\delta<1$, every compact convex $\cK$ with $B(0,r_0)\subset\cK\subset B(0,R_0)$, every strategy $\sigma$ on $\cK$ for $n$ rounds and every law $\nu$ with (N), the set of pairs $(g,\varepsilon)$ on which the plays $\hat x$ of Algorithm~\ref{alg:df2o} against $\sigma$ satisfy
\[
\sum_{t=1}^nf^{\hat x}_t(\hat x_t)-\inf_{x\in\cK}\sum_{t=1}^nf^{\hat x}_t(x)>C\,d^{3/2}\sqrt n\,L_\delta^5
\]
has outer measure at most $\delta$ under the product of the Gaussian law and $\nu$.
\end{proposition}

\begin{proof}[Proof of Remark~\ref{rem:upper-prob}]
Take $P$, $C'$, $k$ from Proposition~\ref{prop:mach} and $C$ from Proposition~\ref{prop:regret}. Fix the data of the theorem and sequences $g$ and $\varepsilon$, and let $x$ be the points that $P$ plays against $\sigma$. The run of $P$ against $\sigma$ reads the loss of round $t$ only at the play of round $t$, after the plays of the rounds before; so it is, step by step, the run of $P$ against the fixed sequence $f^x$ of the losses that $\sigma$ chose on it. These are convex functions $\cK\to[0,1]$, since $\sigma$ is a strategy on $\cK$. Proposition~\ref{prop:mach} applied to $f^x$ gives the cost clause, and $x_t=\mathrm{DF2O}_t(\cK,f^x,\varepsilon,g)$ for $t<n$. The sequence $\hat x$ satisfies the same equations, and they have one solution below round $n$: if $x_s=\hat x_s$ for $s<t$, then $f^x_s=f^{\hat x}_s$ for $s\le t$, and the play of round $t$, which reads these losses only, is the same. So $x_t=\hat x_t$ for $t<n$. Every point $\hat x_t$ lies in $\cK$ (Section~\ref{sec:poly-learner}), and the regret of $P$ against $\sigma$ is the quantity bounded in Proposition~\ref{prop:regret}, on every run.
\end{proof}

For a sequence $f$ fixed in advance, the strategy that ignores the history and returns $f_t$ in round $t$ is a strategy on $\cK$. Since $C'(dL_\delta)^k\le C'(dnL_\delta)^k$, Remark~\ref{rem:upper-prob} gives at once the following form, which we call the (POLY) theorem.

\begin{corollary}[(POLY)]\label{cor:poly}
There are a program $P$, real constants $C,C'$ and an integer $k\ge0$ with the following property. Let $d,n,r_0,R_0,\delta,\cK,\nu$ be as in Remark~\ref{rem:upper-prob} and $f_1,\dots,f_n$ convex functions $\cK\to[0,1]$. Then on every run of $P$ against $f$ every point played lies in $\cK$, the first play comes within $C'(dnL_\delta(2+\ln(R_0/r_0)))^k$ steps and each later play within $C'(dnL_\delta)^k$ steps of the one before; and the runs with $\Reg_n>C\,d^{3/2}\sqrt n\,L_\delta^5$ have outer probability at most $\delta$.
\end{corollary}

\subsection{The regret of the learner}\label{sec:poly-regret}

This section proves Proposition~\ref{prop:regret}. It compares the learner with one that computes the gauge exactly, bounds the regret of the second by four statements, and takes their union; the cut is that of the formal proof. Rounds are numbered $1,\dots,n$, and $\pi=\pi_{\cK'}$ is the gauge of $\cK'$.

\paragraph{Two learners and the case of the analysis.} Let DJ be Algorithm~\ref{alg:df2o}, and let D0 be the same algorithm with the exact gauge $\pi$ in place of $\tilde\pi$ in step~1; D0 serves the proof only. The two read the Gaussian sequence in the same places: the block of Procedure PG, then in round $t$ the $d$ terms of $\xi$ and the term $g$ of the dither $\zeta=g/2$. The \emph{case of the analysis} is the event that Procedure PG returns a position with the two conclusions of Theorem~\ref{thm:pg} and that $\varepsilon\le\frac12$. It is determined by the block of Procedure PG. In this case $S(\cK)\subset B(0,R)$, so $\cK'=S(\cK)$ and $B(0,1)\subset\cK'\subset B(0,R)$.

For a round $t$ let $\cF_t$ be the $\sigma$-field generated by the noise of the rounds before $t$ and by the terms of the Gaussian sequence that the learner reads before round $t$: the block of Procedure PG and the terms of the rounds before $t$. The state of D0 before round $t$ is measurable in $\cF_t$, and the draws $\xi,g$ of round $t$ are independent of $\cF_t$.

\begin{definition}\label{def:lossproc}
A \emph{loss process} on $\cK$ for $n$ rounds is a sequence of random losses $f_t=L_t(\ell_t)$, $t=1,\dots,n$, where for each $t$ the family $(L_t(\ell))_{\ell\in\Lambda}$ consists of convex functions $\cK\to[0,1]$, measurable in $\ell$ at each interior point of $\cK$, and the parameter $\ell_t$ is measurable in $\cF_t$.
\end{definition}

A sequence fixed in advance is a loss process with a constant parameter, and so is the comparison run of Proposition~\ref{prop:gad} below. A run of D0 against a loss process is the run in which the loss of round $t$ is $f_t$; it is well defined because $f_t$ is a function of what the run read before round $t$.

\paragraph{Epochs.} Fix a run of D0 and a round $a\le n$. The \emph{virtual epoch started at $a$} is the sequence of rounds $a,a+1,\dots,n$ of Algorithm~\ref{alg:df2o} run from the state of the start of an epoch, with the draws of these rounds, on the positioned and halved losses $g_s(y)=\frac12f_s(\mu_S+A_Sy)$ and the noise $\frac12(\varepsilon_s+\zeta_s)$. In these terms the observation of step~1 is $Y=p\,(g_s(X/p)+\frac12(\varepsilon_s+\zeta_s))+2(p-1)/\varepsilon$. The first epoch of the run is the virtual epoch started at $1$, and an epoch whose test fires in round $r$ is followed by the virtual epoch started at $r+1$. Let $(\mu_s,A_s)$ be the mean and the factor of the virtual epoch before round $s$, put $p(x)=\max\bigl(1,\pi(x)/(1-\varepsilon)\bigr)$, and
\[
e_s(x)=p(x)\,g_s\bigl(x/p(x)\bigr)+\frac{2(p(x)-1)}\varepsilon,\quad \bar e_s=\E\,e_s\bigl(\mu_s+(1-\lambda)A_s\xi\bigr),\quad \mathrm{eReg}_r(x)=\sum_{s=a}^r\bigl(\bar e_s-e_s(x)\bigr),
\]
the expectation over $\xi\sim N(0,I_d)$ with the state held fixed. The function $e_s$ is the extension of $g_s$ that the learner observes: $e_s=g_s$ on $\cK_\varepsilon$, and $Y=e_s(X)+\frac p2(\varepsilon_s+\zeta_s)$. Say that the virtual epoch started at $a$ \emph{meets its bound} if no guard acts in it and one of the following holds:
\begin{enumerate}[label=(\Alph*)]
\item the test fires in some round $r\le n$, in no round of $a,\dots,r-1$, and $\mathrm{eReg}_r(x)\le0$ for every $x\in\cK_\varepsilon$;
\item the test fires in no round of $a,\dots,n$, and $\mathrm{eReg}_n(x)\le\rho/\eta$ for every $x\in\cK_\varepsilon$.
\end{enumerate}

\paragraph{Four statements.}

\begin{lemma}[one epoch]\label{lem:ve}
Let $d,n,r_0,R_0,\delta,\cK,\nu$ be as in Proposition~\ref{prop:regret}, $f$ a loss process on $\cK$ for $n$ rounds and $1\le a\le n$. For D0 run against $f$, the event that the case of the analysis holds and the virtual epoch started at $a$ does not meet its bound has outer probability at most $\frac{149}{168}\delta_e$.
\end{lemma}

\begin{lemma}[Azuma]\label{lem:az}
Let $d,n,r_0,R_0,\delta,\cK,\nu$ and $f$ be as in Lemma~\ref{lem:ve}, let $x^0_t$ be the point that D0 plays in round $t$ against $f$, and $m_t$ the expectation of $f_t$ at the point that D0 plays in round $t$ from its state before round $t$, over the draw $\xi$ of that round with the state and $f_t$ held fixed. The event that the case of the analysis holds and
\[
\sum_{t=1}^nf_t(x^0_t)>\sum_{t=1}^nm_t+\Bigl(2n\ln\frac7{\delta_e}\Bigr)^{1/2}
\]
has outer probability at most $\delta_e/7$.
\end{lemma}

\begin{lemma}[one run]\label{lem:det}
Let $d,n,r_0,R_0,\delta,\cK$ be as in Proposition~\ref{prop:regret}, and $f_1,\dots,f_n$ convex functions $\cK\to[0,1]$. Consider a run of D0 against $f$ in the case of the analysis, in which every virtual epoch started at a round $a\le n$ meets its bound and the inequality of Lemma~\ref{lem:az} fails. Then
\[
\Reg_n=\sum_{t=1}^nf_t(x^0_t)-\inf_{x\in\cK}\sum_{t=1}^nf_t(x)\le\frac{2\rho}\eta+n\varepsilon+\Bigl(2n\ln\frac7{\delta_e}\Bigr)^{1/2}.
\]
\end{lemma}

Proposition~\ref{prop:gad} compares DJ against a strategy with D0 in a world built from the same draws. Let $\sigma$ be a strategy on $\cK$ for $n$ rounds. The \emph{comparison run} is defined round by round: before round $t$ D0 has a state; $\tilde x_t$ is the point that DJ would play from this state with the draw of round $t$, that is, step~1 with the bisection gauge; the loss of round $t$ is $\tilde f_t=\sigma_t(\tilde x_1,\tilde o_1,\dots,\tilde x_{t-1},\tilde o_{t-1})$, with $\tilde o_s=\tilde f_s(\tilde x_s)+\varepsilon_s$; and D0 plays round $t$ against $\tilde f_t$, at its own point and with its own observation. The strategy is thus fed the points of the bisection gauge and observations without dither, and D0 plays with the exact gauge.

\begin{proposition}\label{prop:gad}
Let $d,n,r_0,R_0,\delta,\cK,\sigma,\nu$ be as in Proposition~\ref{prop:regret}, let $\hat x$ be the plays of DJ against $\sigma$ (Section~\ref{sec:poly-proof}) and $B$ a real number. The outer probability that the case of the analysis holds and the regret of $\hat x$ against $f^{\hat x}$ exceeds $B+1$ is at most the outer probability that the case of the analysis holds and the regret of D0 in the comparison run against $\tilde f$ exceeds $B$, plus $\delta_1$.
\end{proposition}

\paragraph{Proof of Proposition~\ref{prop:regret}.} Take $C=10^{22}$, and let $\hat x$ be the plays of DJ against $\sigma$. The losses $f^{\hat x}_t$ are convex with values in $[0,1]$ and every $\hat x_t$ lies in $\cK$, so $\Reg_n\le n$ on every run.

If $\varepsilon>\frac12$, then $14500\,d^{3/2}L_e^5>\frac12\sqrt n$, and $L_e\le900\,L_\delta$ by the definition of $\delta_1$ and $\delta_e$; so $n<2\cdot14500\,d^{3/2}\sqrt n\,(900L_\delta)^5\le C\,d^{3/2}\sqrt nL_\delta^5$, and the event of the proposition is empty.

Let $\varepsilon\le\frac12$ and $B=2\rho/\eta+n\varepsilon+(2n\ln(7/\delta_e))^{1/2}$. From $2\rho/\eta=2900\,d^{3/2}L_e^4\sqrt n$, $n\varepsilon=14500\,d^{3/2}L_e^5\sqrt n$, $\ln(7/\delta_e)\le6+L_e/36$, $1\le L_e\le900\,L_\delta$ and $L_\delta\ge\ln2$,
\[
B+1\le10^{22}\,d^{3/2}\sqrt n\,L_\delta^5 .
\]
The runs with $\Reg_n>C\,d^{3/2}\sqrt nL_\delta^5$ lie in the union of two sets: the runs outside the case of the analysis, and the runs in the case with $\Reg_n>B+1$. The first set lies in the event that Procedure PG, run with confidence $\delta_1$ on its block, does not return a position with the conclusions of Theorem~\ref{thm:pg}; its outer probability is at most $\delta_1$. By Proposition~\ref{prop:gad}, the second has outer probability at most $\delta_1$ plus that of the case together with a regret of D0 in the comparison run above $B$.

The comparison run is the run of D0 against a loss process: its parameter before round $t$ is the history $(\tilde x_s,\tilde o_s)_{s<t}$, which is a function of what D0 read before round $t$ and of the noise of the rounds before $t$, and $L_t(h)=\sigma_t(h)$ is measurable in $h$ at each point of $\cK$. So Lemmas~\ref{lem:ve} and~\ref{lem:az} apply to it, and Lemma~\ref{lem:det} applies to it at each fixed value of the draws and the noise. In the case of the analysis, a regret of D0 above $B$ forces, by Lemma~\ref{lem:det}, that some virtual epoch started at a round $a\le n$ does not meet its bound or that the inequality of Lemma~\ref{lem:az} holds. The first has probability at most $n\cdot\frac{149}{168}\delta_e$ and the second at most $\delta_e/7$. In total, with $\delta_e=4\delta_1/(5n)$ and $n\ge1$,
\[
2\delta_1+n\cdot\frac{149}{168}\cdot\frac{4\delta_1}{5n}+\frac{4\delta_1}{35n}\le\Bigl(2+\frac{596}{840}+\frac4{35}\Bigr)\delta_1<2.824\,\delta_1\le\delta,
\]
since $\delta_1\le\delta/3$. \qed

\paragraph{Proof of Lemma~\ref{lem:det}.} The case of the analysis gives $\cK'=S(\cK)$ with $B(0,1)\subset\cK'$, and $0\in\cK_\varepsilon$. By hypothesis the virtual epoch started at $1$ meets its bound. If it ends by (A) in a round $r<n$, the run continues with the virtual epoch started at $r+1$, which meets its bound too; so the rounds $1,\dots,n$ split into epochs $[a_k,b_k]$, each the virtual epoch started at $a_k$ up to the round $b_k$, with $\max_{x\in\cK_\varepsilon}\mathrm{eReg}_{b_k}(x)\le0$ for every epoch but the last and $\le\rho/\eta$ for the last. No guard acts, so in round $t$ of the epoch $k$ D0 plays $x^0_t=\mu_S+A_S(X/p(X))$ with $X=\mu_t+(1-\lambda)A_t\xi$ and the state of that epoch.

For every $X$, $p(X)\ge1$ and $g_t\ge0$ give $\frac12f_t(x^0_t)=g_t(X/p(X))\le e_t(X)$. Taking the expectation over $\xi$ with the state fixed, $m_t\le2\bar e_t$. For $y\in\cK_\varepsilon$, $p(y)=1$ and $e_t(y)=g_t(y)=\frac12f_t(\mu_S+A_Sy)$. Summing over the epochs,
\[
\sum_{t=1}^nm_t-\sum_{t=1}^nf_t(\mu_S+A_Sy)\le2\sum_k\mathrm{eReg}_{b_k}(y)\le\frac{2\rho}\eta\qquad(y\in\cK_\varepsilon).
\]
Let $x\in\cK$ and $y=S^{-1}(x)$, so that $y\in\cK'$ and $(1-\varepsilon)y\in\cK_\varepsilon$. Since $\mu_S=\mu_S+A_S0\in\cK$ and $f_t$ is convex with values in $[0,1]$, $f_t(\mu_S+(1-\varepsilon)A_Sy)\le(1-\varepsilon)f_t(x)+\varepsilon f_t(\mu_S)\le f_t(x)+\varepsilon$. With the failure of the inequality of Lemma~\ref{lem:az},
\[
\sum_{t=1}^nf_t(x^0_t)-\sum_{t=1}^nf_t(x)\le\Bigl(2n\ln\frac7{\delta_e}\Bigr)^{1/2}+\frac{2\rho}\eta+n\varepsilon,
\]
and the supremum over $x\in\cK$ is the claim. \qed

\paragraph{Proof of Lemma~\ref{lem:az}.} Let $D_t=f_t(x^0_t)-m_t$. In the case of the analysis $x^0_t$ is a function of the state of D0 before round $t$, measurable in $\cF_t$, and of the draw $\xi$ of round $t$, independent of $\cF_t$; the loss $f_t=L_t(\ell_t)$ has its parameter measurable in $\cF_t$. Every point at which a loss is read lies in the interior of $\cK$, since its gauge in $\cK'$ is at most $1-\varepsilon$, and a family of convex functions measurable in the parameter at each interior point is jointly measurable on the interior (it is continuous in the point there). So the conditional expectation of $f_t(x^0_t)$ given $\cF_t$ is the integral over $\xi$ with $\cF_t$ frozen, which is $m_t$, and $(D_t)$ is a sequence of martingale differences for $(\cF_{t+1})$ with $|D_t|\le1$. The inequality of Azuma and Hoeffding \cite{Azuma67} gives $\Prob\bigl(\sum_tD_t>s\bigr)\le e^{-s^2/(2n)}$, which is $\delta_e/7$ at $s=(2n\ln(7/\delta_e))^{1/2}$. The case of the analysis is measurable and depends on the block of Procedure PG only; outside it nothing is claimed. \qed

\paragraph{Proof of Proposition~\ref{prop:gad}.} Let $\eta_g=R\,2^{-j}$. In the case of the analysis $|y|\le R\,\pi(y)$ on $\cK'\subset B(0,R)$, and $j$ steps of bisection leave an interval of length at most $|y|2^{-j}$; so $\pi\le\tilde\pi\le(1+\eta_g)\pi$, and with $\tilde p=\max(1,\tilde\pi(X)/(1-\varepsilon))$, $p\le\tilde p\le(1+\eta_g)p$. The points $X/p$ and $X/\tilde p$ lie in $\cK_\varepsilon$, at distance $|X/p|(1-p/\tilde p)\le R\eta_g$. A convex function $h:\cK'\to[0,1]$ is $(1/\varepsilon)$-Lipschitz on $\cK_\varepsilon$, since $y+\varepsilon B(0,1)\subset(1-\varepsilon)\cK'+\varepsilon\cK'=\cK'$ for $y\in\cK_\varepsilon$; hence
\[
\bigl|h(X/\tilde p)-h(X/p)\bigr|\le R\eta_g/\varepsilon. \tag{G}
\]
The choice of $j$ gives $n\eta_g(R+7)/\varepsilon\le\delta_1$, and in particular $nR\eta_g/\varepsilon\le1$.

For $0\le k\le n$ the hybrid $H_k$ runs the learner with the bisection gauge in the rounds before $k+1$ and with the exact gauge from round $k+1$ on, on the same draws and noise. In every hybrid the strategy is fed, in round $t$, the point $\tilde x_t$ that the bisection gauge gives from the state of the hybrid and the draw of round $t$, and the observation $\tilde o_t=\tilde f_t(\tilde x_t)+\varepsilon_t$; the learner plays round $t$ against $\tilde f_t$ at the point of its own gauge. Let $\Phi_k$ be the sum over $t$ of $\tilde f_t$ at the point the exact gauge gives from the state of $H_k$ before round $t$, minus $\inf_{x\in\cK}\sum_t\tilde f_t(x)$, and $E_k$ the event that the case of the analysis holds and $\Phi_k>B$. Then $H_0$ is the comparison run, and $E_0$ is the event on the right of the proposition. In $H_n$ the states are those of DJ against $\sigma$, the strategy is fed the plays $\hat x_t$ and the observations of the game, so $\tilde f=f^{\hat x}$, and (G) applied in each round gives $\Reg_n(\hat x)\le\Phi_n+nR\eta_g/\varepsilon\le\Phi_n+1$. So the event on the left lies in $E_n$.

One step: $\Prob(E_{k+1})\le\Prob(E_k)+\eta_g(R+7)/\varepsilon$. The hybrids $H_k$ and $H_{k+1}$ agree before round $k+1$ and differ in round $k+1$ only in the gauge with which the learner forms its observation; the point fed to the strategy, and so the history, is the same in both. The dither $g$ of round $k+1$ is read by nothing but that observation, $Y=a+bg$ with $b=p/4$ and $a=p(h(X/p)+\varepsilon_{k+1})/2+2(p-1)/\varepsilon$, $h(y)=\tilde f_{k+1}(\mu_S+A_Sy)$, where $p$ is the exact factor in $H_k$ and $\tilde p$ in $H_{k+1}$; the rest of the run, and $\Phi$, is the same function of $Y$ and of the other coordinates in both. Fixing the other coordinates, the two sections of the events in $g$ are $\{z:a_1+(1+u)b_0z\in S\}$ and $\{z:a_0+b_0z\in S\}$ for one measurable set $S$, with $b_0=p/4$, $1+u=\tilde p/p$, $0\le u\le\eta_g$. For one standard Gaussian $z$,
\[
\Prob\bigl(a_1+(1+u)b_0z\in S\bigr)\le\Prob\bigl(a_0+b_0z\in S\bigr)+u+\frac{|a_1-a_0|}{2b_0}:
\]
the density of $N(a_1,(1+u)^2b_0^2)$ is at least $(1+u)^{-1}$ times that of $N(a_1,b_0^2)$, which costs $1-(1+u)^{-1}\le u$, and moving the mean by $c$ costs at most $\Prob(|z|\le c/(2b_0))\le c/(2b_0)$. By (G) and $\tilde p-p\le\eta_gp$,
\[
\frac{|a_1-a_0|}{2b_0}\le\frac{\tilde p-p}p\Bigl(1+|\varepsilon_{k+1}|+\frac4\varepsilon\Bigr)+\frac{R\eta_g}\varepsilon\le\eta_g\Bigl(1+|\varepsilon_{k+1}|+\frac{R+4}\varepsilon\Bigr).
\]
Integrating over the other coordinates, with $\E|\varepsilon_{k+1}|\le(\E e^{\varepsilon_{k+1}^2}-1)^{1/2}\le1$ from (N), gives the step with $\eta_g(3+(R+4)/\varepsilon)\le\eta_g(R+7)/\varepsilon$. The events are measurable: in the case of the analysis every point read lies in the interior of $\cK$, where $\sigma_t$ is jointly measurable, and the exact gauge is an infimum over rational points.

Summing the $n$ steps, the event on the left has outer probability at most $\Prob(E_n)\le\Prob(E_0)+n\eta_g(R+7)/\varepsilon\le\Prob(E_0)+\delta_1$. \qed

\paragraph{Proof of Lemma~\ref{lem:ve}.} This is the analysis of one epoch of the online Newton method, and it is where the order $d^{3/2}\sqrt n$ comes from. For one epoch started fresh, with losses fixed in advance, it is the analysis of the method of \cite[Chapter~11]{Lat24}, with four changes: the rank-one updates of steps~2 and~3 in place of a stored covariance; the play with the factor $1-\lambda$ and the extension $e_s$, under which the estimators at the mean are unbiased with no density ratio; a single region $\{\Phi\le c_1\rho\}$ and its lower level $(c_1-\gamma)\rho$ in place of a list of quadratics; and the restart test of step~5, which is a concave quadratic maximised by one projection. On the intersection of the events of the epoch (noise, gauge and Gaussian tails in each round; the concentration of the estimated gradients and Hessians on fixed nets of $B(0,R)$ and of the unit sphere; the second moment of the observations), whose failure probabilities sum to $\frac{149}{168}\delta_e$, the epoch either fires with $\max_{\cK_\varepsilon}\mathrm{eReg}\le0$ or reaches the horizon with $\max_{\cK_\varepsilon}\mathrm{eReg}\le\rho/\eta$, and the guard does not act. The proof is in Appendix~\ref{app:epoch}.
\subsection{Proof of Theorem~\ref{thm:polytope}}\label{sec:poly-polytope}

Let $\cK=\{x\in\R^d:\ip{a_i,x}\le b_i,\ i=1,\dots,m\}$ be as in Theorem~\ref{thm:polytope}: its $m(d+1)$ coefficients are rational numbers whose numerators and denominators in lowest terms are at most $M$ in absolute value, which we call \emph{height at most $M$}. The run of a program on a polytope starts with $N_0=d$, $N_1=n$, $N_2=m$, $M_2=\delta$, the numbers $a_{i1},\dots,a_{id},b_i$ in the cells $3+(d+1)(i-1),\dots,3+(d+1)(i-1)+d$, and every other register and cell at $0$. The hypotheses force $m\ge1$, since with no row $\cK=\R^d$ has unbounded diameter, and then $M\ge1$, since a denominator is at least $1$. The program $P_{\mathrm{pol}}$ is a prefix PRE, which finds a point $a_0$ with $B(a_0,\frac12)\subset\cK$, followed by a simulation $S(P)$ of the program $P$ of Remark~\ref{rem:upper-prob} on $\cK-a_0$, in which each membership question is answered from the $m$ inequalities. Remark~\ref{rem:upper-prob} applies to the body $\cK-a_0$ with $r_0=\frac12$ and $R_0=n$.

\begin{lemma}\label{lem:loc}
Let $d,m,M\ge1$ and let $\cK$ be compact, given by $m$ inequalities of height at most $M$. Then $\cK\subset B(0,R_{\mathrm{loc}})$ with $R_{\mathrm{loc}}=\bigl(dm\,M^{2(d+1)}\bigr)^{d+1}$.
\end{lemma}

\begin{proof}
By the theorem of Krein and Milman, $\cK$ is the closed convex hull of its extreme points, so it suffices to bound the norm of an extreme point $x$. Let $S$ be the set of the rows active at $x$, those with $\ip{a_i,x}=b_i$. No $v\ne0$ is orthogonal to every $a_i$, $i\in S$: otherwise $x\pm sv\in\cK$ for $s>0$ small, and $x$ is not extreme. Multiply the row $i$ by the product of the denominators of its $d+1$ coefficients: this gives integers $\alpha_i\in\mathbb Z^d$ and $\beta_i\in\mathbb Z$ of absolute value at most $M^{d+1}$, with $\ip{\alpha_i,x}=\beta_i$ for $i\in S$. The integer matrix $G=\sum_{i\in S}\alpha_i\alpha_i^{\mathsf T}$ is invertible, since $\ip{Gv,v}=\sum_{i\in S}\ip{\alpha_i,v}^2$, and $Gx=h$ with $h=\sum_{i\in S}\beta_i\alpha_i$. The entries of $G$ and $h$ are at most $X=mM^{2(d+1)}$ in absolute value. By Cramer's rule, with $|\det G|\ge1$ and the bound $d!\,X^d$ on the determinant of a $d\times d$ matrix with entries at most $X$, every coordinate of $x$ is at most $d!\,X^d$ in absolute value, so $|x|\le d^{1/2}d!\,X^d\le(dX)^{d+1}$.
\end{proof}

The radius is far from the best one; only its logarithm, $(d+1)\ln(dm)+2(d+1)^2\ln M$, enters the count of steps.

\paragraph{The centre.} For $\rho>0$, a point $a$ \emph{passes the row $i$} if $\ip{a_i,a}+\rho|a_i|\le b_i$, the maximum of $\ip{a_i,\cdot}$ on $B(a,\rho)$ being the left side; a point that passes every row has $B(a,\rho)\subset\cK$. Procedure CENTRE$(R,\rho,N)$ starts with $E_0=R^2I$ and $a^0=0$, and for $k=0,\dots,N-1$: if $a^k$ passes every row, it returns $a^k$; otherwise let $i$ be the first row that $a^k$ fails; if $a_i=0$ it fails, and else $(E_{k+1},a^{k+1})$ is the central cut of Lemma~\ref{lem:cc} of the ellipsoid $E(E_k,a^k)$ by $\{x:\ip{a_i,x-a^k}\le0\}$, for $d=1$ the centre of that formula and $E_{k+1}=E_k/4$. After $N$ tests it fails.

\begin{lemma}\label{lem:centre}
Let $d\ge1$, $R,\rho>0$ and $N\ge0$. If CENTRE$(R,\rho,N)$ returns a point $a$, then $B(a,\rho)\subset\cK$. If $\cK\subset B(0,R)$, $\cK$ contains a closed ball $B(z,2\rho)$ and $N>2d(d+1)\ln(R/\rho)$, it returns a point.
\end{lemma}

\begin{proof}
The first claim is the remark above. For the second, let $\cK_\rho$ be the set of the points that pass every row. It lies in $\cK\subset B(0,R)=E(E_0,a^0)$, and it contains $B(z,\rho)$: for $|x-z|\le\rho$ the maximum of $\ip{a_i,\cdot}$ on $B(x,\rho)$ is at most its maximum on $B(z,2\rho)\subset\cK$. A row $a_i=0$ is passed, since $\cK\ne\emptyset$ gives $0\le b_i$. If $a^k$ fails the row $i$, every $x\in\cK_\rho$ has $\ip{a_i,x}\le b_i-\rho|a_i|<\ip{a_i,a^k}$. By Lemma~\ref{lem:cc} for $d\ge2$, and by the halving of an interval for $d=1$, induction on $k$ gives $\cK_\rho\subset E(E_k,a^k)$ and $\det E_k\le R^{2d}e^{-k/(d+1)}$ as long as the tests fail, while Lemma~\ref{lem:ball} gives $\det E_k\ge\rho^{2d}$. So the tests that fail number at most $2d(d+1)\ln(R/\rho)<N$.
\end{proof}

The program does not know $M$, so it tries the radii $R_j=2^{2^j}$: stage $j$ runs CENTRE$(R_j,\frac12,N_j)$ with $N_j=\lceil2d(d+1)\ln(2R_j)\rceil+1$, for $j=0,1,\dots$ until a stage returns a point. Under the hypotheses of Theorem~\ref{thm:polytope}, $\cK$ is compact and lies in $B(0,R_{\mathrm{loc}})$ by Lemma~\ref{lem:loc}, and it contains a closed ball of radius $1=2\cdot\frac12$; so the stage $j^*$, the least $j$ with $R_{\mathrm{loc}}\le R_j$, returns a point if no earlier stage does (Lemma~\ref{lem:centre}), and the stages up to $j^*$ make at most $\sum_{j\le j^*}N_j\le40d^2(\ln(2R_{\mathrm{loc}})+1)$ tests.

\paragraph{The simulation.} Let $r$ be larger than every register that $P$ names, and at least $4$. The program $S(P)$ replaces each instruction of $P$ by a block, which uses the registers $r,\dots,r+20$ and the cells above an offset $\mathrm{OFF}$, held in the register $r$. An instruction with real operands adds $\mathrm{OFF}$ to its operand registers and runs on the cells so shifted; an instruction on natural registers is kept, and a jump goes to the block of its target. A membership question for a point $y$ becomes the $m$ tests $\ip{a_i,y}\le b_i'$, with $b_i'=b_i-\ip{a_i,a_0}$, and writes $1$ if all pass and $0$ otherwise; a play of $y$ becomes a play of $y+a_0$. $S(P)$ asks no membership question; a block takes at most
\[
W=7dm+12m+7d+12\le38\,dm
\]
steps, and changes the number of draws, the rounds and the points played at its last step only. A state $s'$ of $S(P)$ is \emph{related} to a state $s$ of $P$ when the registers below $r$ agree, the cell $\mathrm{OFF}+c$ of $s'$ holds the cell $c$ of $s$, the draws and the rounds agree, each point played in $s'$ is the point of $s$ plus $a_0$, the counter of $s'$ is at the block of the instruction of $s$, and $s'$ is framed: the registers $r+1,\dots,r+10$ hold $d$, $m$, addresses and the constants $1$ and $0$, and the cells at these addresses hold the rows $a_i$, the numbers $b_i'$ and the coordinates of $a_0$. For a strategy $\sigma$ on $\cK$ the translated strategy on $\cK-a_0$ is
\[
\sigma^{a_0}_t(x_1,o_1,\dots,x_{t-1},o_{t-1})(y)=\sigma_t(x_1+a_0,o_1,\dots,x_{t-1}+a_0,o_{t-1})(y+a_0).
\]

\begin{lemma}\label{lem:sim}
Let $\sigma$ be a strategy on $\cK$, and let $s_0'$ be related to $s_0$. Fix a noise sequence and a Gaussian sequence. Let $s_j$ be the state of $P$ after $j$ steps from $s_0$ against $\sigma^{a_0}$, with the membership questions answered for $\cK-a_0$, and $s_k'$ the state of $S(P)$ after $k$ steps from $s_0'$ against $\sigma$. There are $0=\kappa_0<\kappa_1<\cdots$ with $\kappa_{j+1}\le\kappa_j+W$ such that $s'_{\kappa_j}$ is related to $s_j$ for every $j$, and $s'_k$ has the rounds of $s_j$ for $\kappa_j\le k<\kappa_{j+1}$.
\end{lemma}

\begin{proof}
By induction on $j$, one block at a time. The operands that a block reads are those of the instruction, by the relation, and it writes the result at the shifted address. The block of a membership question answers $1$ exactly when $\ip{a_i,y}\le b_i-\ip{a_i,a_0}$ for every $i$, that is when $y\in\cK-a_0$. In the block of a play, $S(P)$ plays $y+a_0$ where $P$ plays $y$. The observations of the earlier rounds are the same in both runs, by induction on the rounds, so the loss that $\sigma$ chooses in $S(P)$, at $y+a_0$, is the loss that $\sigma^{a_0}$ chooses in $P$, at $y$, and the two observations are equal. Since a block changes the rounds and the points at its last step only, the loss stays the same during the block. When $P$ has halted, so has $S(P)$, and $\kappa_{j+1}=\kappa_j+1$.
\end{proof}

\paragraph{The prefix.} PRE lays out above the rows the cells of the relation and the offset $\mathrm{OFF}$, computes the $m$ products $\frac12|a_i|$, runs the stages of the centre, writes the numbers $b_i'$ and the coordinates of the point $a_0$ returned, writes $\frac12$, $n$ and $\delta$ in the cells $\mathrm{OFF},\mathrm{OFF}+1,\mathrm{OFF}+2$, resets the register $2$ and passes control to $S(P)$.

\begin{lemma}\label{lem:pre}
For every $r\ge4$ there are a program PRE that asks no membership question, a constant $c_{\mathrm{pre}}$ and an integer $k_{\mathrm{pre}}$ with the following property. Let $d,n,m,M,\delta,\cK$ be as in Theorem~\ref{thm:polytope}, and let PRE stand at the address $0$ of a program $Q$. There are an integer $T\le c_{\mathrm{pre}}(dm\ln(2M))^{k_{\mathrm{pre}}}$ and a point $a_0$ with $B(a_0,\frac12)\subset\cK$ such that on every run, whatever the losses, the noise and the Gaussian sequence, the run of $Q$ on the polytope plays nothing in its first $T$ steps and is then in a state related to the initial state of $P$ on the inputs $d,n,\frac12,n,\delta$, at the centre $a_0$ and with $S(P)$ starting where PRE ends.
\end{lemma}

The proof gives $c_{\mathrm{pre}}=10^9/(\ln2)^{13}$ and $k_{\mathrm{pre}}=13$. The count of steps is that of the stages up to $j^*$, which Lemma~\ref{lem:loc} bounds, with at most $100(d+1)^2(m+1)$ steps for each test of a stage and its cut, and $O(dm)$ steps for the layout and the frame. The integer $T$ and the point $a_0$ depend on the data of the polytope and on nothing that the run reads later.

\begin{proof}[Proof of Theorem~\ref{thm:polytope}]
\textup{(a)} Take $P$, $C$, $C'$ and $k$ from Remark~\ref{rem:upper-prob}, $r$ above every register of $P$ and at least $4$, PRE from Lemma~\ref{lem:pre}, and $P_{\mathrm{pol}}$ the code of PRE followed by $S(P)$. It asks no membership question. Fix the data of the theorem, and let $T$ and $a_0$ be given by Lemma~\ref{lem:pre}.

The body $\cK'=\cK-a_0$ is compact and convex, $B(0,\frac12)\subset\cK'$, and $\cK'\subset B(0,n)$ since $a_0\in\cK$ and the diameter of $\cK$ is at most $n$; also $\frac12\le n$. The strategy $\sigma^{a_0}$ is a strategy on $\cK'$: its losses are those of $\sigma$ composed with $y\mapsto y+a_0$, convex with values in $[0,1]$ on $\cK'$, and at each $y\in\cK'$ it is a Borel function of the history, composed with a translation of the points. So Remark~\ref{rem:upper-prob} applies to $P$ on the inputs $d,n,\frac12,n,\delta$, on $\cK'$ against $\sigma^{a_0}$.

Fix the noise and the Gaussian sequence. By Lemmas~\ref{lem:pre} and~\ref{lem:sim}, after its first $T$ steps the run of $P_{\mathrm{pol}}$ follows block by block the run of $P$ on $\cK'$ against $\sigma^{a_0}$, and its point of each round $t<n$ is $x_t+a_0$, where $x_t\in\cK'$ is the point of $P$; so it lies in $\cK$. By the remark $P$ plays first within $B_0=\lceil C'(dL_\delta(2+\ln(2n)))^k\rceil$ steps, and each later play of a round below $n$ within $B=\lceil C'(dL_\delta)^k\rceil$ steps of the one before. A block takes at most $W$ steps, so $P_{\mathrm{pol}}$ plays first within $T+WB_0$ steps, and each later play within $WB$ steps of the one before. With $W\le38\,dm$, $\ln(2n)\le L_\delta$, and $d$, $m$, $L_\delta/\ln2$ and $\ln(2M)/\ln2$ all at least $1$,
\[
c_{\mathrm{pre}}(dm\ln(2M))^{k_{\mathrm{pre}}}+W(B_0+1)\le C_1\bigl(dmL_\delta\ln(2M)\bigr)^{k_1},\qquad W(B+1)\le C_1(dmL_\delta)^{k_1},
\]
with $k_1=\max(k_{\mathrm{pre}},2k+1)$, which is $\max(13,2k+1)$, and a constant $C_1$ that depends on $C'$, $k$, $c_{\mathrm{pre}}$ and $k_{\mathrm{pre}}$ only.

The observations of the two runs are equal, and the loss of round $t$ of the run of $P_{\mathrm{pol}}$ at $y+a_0$ is the loss of round $t$ of the run of $P$ at $y$. So the sums of the losses at the points played are equal, and the infimum over $\cK$ of the sum of the losses of $P_{\mathrm{pol}}$ is the infimum over $\cK'$ of the sum of those of $P$: the regret of $P_{\mathrm{pol}}$ against $\sigma$ is the regret of $P$ on $\cK'$ against $\sigma^{a_0}$, on every run. The set of the runs on which it exceeds $C\,d^{3/2}\sqrt n\,L_\delta^5$ has outer probability at most $\delta$ by the remark, whose bound contains neither $r_0$ nor $R_0$.

\textup{(b)} Take $P_{\mathrm{pol}}$, $C$, $C_1$ and $k_1$ from \textup{(a)}, and put $C'=\max(C,0)$; \textup{(a)} holds with $C'$ in place of $C$. Run $P_{\mathrm{pol}}$ with $\delta=1/(n+1)$, so that $L_\delta=L$; the points and the steps are those of \textup{(a)}. Every point played lies in $\cK$ and every loss takes values in $[0,1]$ on $\cK$, so $\Reg_n\le n$ on every run. The runs above $C'\,d^{3/2}\sqrt n\,L^5$ lie in a measurable set $A$ of probability at most $1/(n+1)$, and $R=C'\,d^{3/2}\sqrt n\,L^5+n\mathbf 1_A$ is a measurable majorant of the regret with $\E R\le C'\,d^{3/2}\sqrt n\,L^5+n/(n+1)$. So \textup{(b)} holds with $P'_{\mathrm{pol}}=P_{\mathrm{pol}}$, this $C'$, $C'_1=C_1$ and $k'_1=k_1$.
\end{proof}

\section{A lower bound against every learner}\label{sec:lower}

This section proves Theorem~\ref{thm:lower}, cut as its formal proof is cut. Here $B$ is the closed unit ball of $\R^d$, $r_0=\frac12$ and $b=\frac14$. The adversary is first allowed private randomness: a point $x^\circ$, the \emph{tip}, uniform on $r_0B$, and independent standard Gaussian variables $Z_1,Z_2,\dots$, independent of the learner's seed $\omega$. The game is run on the samples $(x^\circ,Z,\omega)$ under the product law. A bound on the expected regret against this randomised adversary gives one value of its randomness with the same bound, and with that value fixed the adversary is a strategy in the sense of Theorem~\ref{thm:lower} (Section~\ref{sec:lower-proof}). In this section $g$, $k$, $\kappa$ and $L$ are the depth, the order, the scale and the Lipschitz constant of the dent, $\sigma$ is the scale of the noise, and $\Pi_t$ is the law the adversary keeps of the tip before round $t$. For a law $\Pi$, a function $X$ and a set $S$, $\E_\Pi[X;S]=\int_SX\,d\Pi$.

\subsection{The dent and the frame}\label{sec:lower-dent}

Let $m\in\R^d$, $M$ a symmetric positive definite matrix, $|v|_M=\ip{Mv,v}^{1/2}$, $\ip{u,v}_M=\ip{Mu,v}$, and $\kappa,k,r_1>0$. Put $\phi(r)=r^2/2+k$ for $0\le r\le r_1$ and $\phi(r)=r_1r-r_1^2/2+k$ for $r>r_1$, and let the \emph{bowl} be $q(x)=\kappa\,\phi(|x-m|_M)$. For a tip $\theta\in\R^d$ the \emph{dent} is the lower convex hull of the bowl and of the point $(\theta,0)$:
\[
D(x;\theta)=\inf_{0<\lambda\le1}\lambda\,q\bigl(\theta+(x-\theta)/\lambda\bigr).
\]
For a play $x$ put $a=|x-m|_M$, $T(\theta)=\ip{x-m,\theta-m}_M/a$ ($T=0$ if $a=0$), $h(x,\theta)=\ip{x-m,\theta-m}_M-a^2/2+k$ and $\rho(\theta)^2=|\theta-m|_M^2+2k$, and let $\hat h(x,\theta)$ be the tangent plane of the bowl at $x$, read at $\theta$ and divided by $\kappa$: $\hat h=h$ for $a\le r_1$ and $\hat h=r_1T-r_1^2/2+k$ for $a>r_1$.

\begin{lemma}[the dent]\label{lem:dent}
For every $\theta\in\R^d$:
\begin{enumerate}[label=(\alph*)]
\item $D(\cdot;\theta)$ is convex, $0\le D(\cdot;\theta)\le q$ and $D(\theta;\theta)=0$;
\item $D(\cdot;\theta)$ is Lipschitz with constant $\kappa r_1\|M^{1/2}\|$;
\item if $|\theta-m|_M+\rho(\theta)\le r_1$, then $D(x;\theta)=q(x)$ where $h(x,\theta)\le0$, and where $h(x,\theta)>0$ one has $a\le r_1$, $h\le\rho^2/2$ and $D(x;\theta)=q(x)-\kappa\,\mathrm{dev}(x;\theta)$ with
\[
\mathrm{dev}=\rho^2-h-\rho(\rho^2-2h)^{1/2},\qquad0\le\mathrm{dev}\le\min\bigl(h,2h^2/\rho^2\bigr);
\]
\item $D(x;\theta)\ge q(x)-\kappa\,\hat h(x,\theta)_+$ at every $x$;
\item if $T(\theta)\le(2k)^{1/2}$, then $\hat h(x,\theta)\le2k$;
\item if $T(\theta)>(2k)^{1/2}$, then $\min\bigl(\hat h(x,\theta)_+,\phi(a)\bigr)^2\le(4T^2+r_1T)\,\phi(a)$.
\end{enumerate}
\end{lemma}

\begin{proof}
(a) The points $(1-\lambda)(\theta,0)+\lambda(y,s)$ with $\lambda\in[0,1]$ and $s\ge q(y)$ form the convex hull of the epigraph of $q$ and of one point, and $D(\cdot;\theta)$ is its lower boundary, hence convex. $D\ge0$ since $q\ge0$, $D\le q$ by $\lambda=1$, and $D(\theta;\theta)=0$ by letting $\lambda\to0$. (b) The bowl has the Lipschitz constant $\kappa r_1\|M^{1/2}\|$ since $\phi'\le r_1$; each function of the infimum has it, and so has the infimum. (c) In the coordinates $z=M^{1/2}(x-m)$, $p=M^{1/2}(\theta-m)$, on the ray $z=p+te$ with $|e|=1$ and $t>0$, $D/\kappa=t\inf_{s\ge t}F(s)$ with $F(s)=\phi(|p+se|)/s$. Then $s^2F'(s)=s\frac{d}{ds}\phi(|p+se|)-\phi(|p+se|)$, which is nondecreasing in $s$ and equals $(s^2-\rho^2)/2$ for $s\le\rho$, where $|p+se|\le r_1$. So $F$ decreases up to $\rho$ and increases after, $D/\kappa=\phi(|z|)$ for $t\ge\rho$ and $D/\kappa=t(\rho+\ip{p,e})$ for $t\le\rho$. Since $|z-p|^2=\rho^2-2h$, the case $t\le\rho$ is $h\ge0$, and there $\phi(|z|)-D/\kappa=\rho^2-h-\rho(\rho^2-2h)^{1/2}=h^2/\bigl(\rho^2-h+\rho(\rho^2-2h)^{1/2}\bigr)$, which gives the bounds. (d) For $\lambda\in(0,1]$ and $y=\theta+(x-\theta)/\lambda$, the convexity of $q$ gives $\lambda q(y)\ge\lambda q(x)+(1-\lambda)\ip{\nabla q(x),x-\theta}\ge q(x)-\kappa\hat h_+$. (e) For $a\le r_1$, $\hat h=aT-a^2/2+k\le2k-(a-(2k)^{1/2})^2/2$, and for $a>r_1$, $\hat h\le2k-(r_1-(2k)^{1/2})^2/2$. (f) Let $a\le r_1$. If $a\le T$, then $\hat h_+^2\le2a^2T^2+2k^2\le4T^2\phi(a)$, using $k<T^2/2$; if $a>T$, then $\hat h<T^2$ and $\phi(a)>T^2/2$. For $a>r_1$, $\hat h<r_1T+T^2/2$ and $\min(\hat h_+,\phi)^2\le\hat h_+\phi$.
\end{proof}

The \emph{frame} of a law on $r_0B$ with mean $m$ and covariance $\Sigma=\sum_i\lambda_ie_ie_i^{\mathsf T}$, at a floor $w>0$, is the pair $(m,M)$ with $M=\sum_ie_ie_i^{\mathsf T}/\max(\lambda_i,w^2)$, a continuous function of $\Sigma$. Let $W=\Sigma+w^2I$, $d_{\mathrm{live}}$ the number of $i$ with $\lambda_i\ge w^2$ and $D_M=\tr(M\Sigma)=\sum_i\min(1,\lambda_i/w^2)$.

\begin{lemma}[the frame]\label{lem:frame}
(F1) $w^2M\le I$, so $\|M^{1/2}\|\le1/w$; (F2) $W^{-1}\le M\le2W^{-1}$ and $W\le2M^{-1}$; (F3) $d_{\mathrm{live}}\le D_M\le d$.
\end{lemma}

\begin{proof}
On each eigenvector, from $w^2\le\max(\lambda,w^2)\le\lambda+w^2\le2\max(\lambda,w^2)$, and from $D_M=\sum_i\min(1,\lambda_i/w^2)$.
\end{proof}

\subsection{The strategy and the genie}\label{sec:lower-strategy}

Fix a depth $g>0$, an order $k\ge240$, $0<L\le\frac14$ and $\sigma>0$ with $d\ge2048k$ and $\frac32L+g\le\frac12$, and put $\kappa=g/k$, $r_1=5d^{1/2}$ and the floor $w_*=\kappa r_1/L$. The frame of a law $\Pi$ is taken at the floor $w_*$, and the dent $D_\Pi$ of $\Pi$ is the dent with the frame of $\Pi$ and the constants $\kappa$, $k$, $r_1$. A law is \emph{live} if $d_{\mathrm{live}}\ge d/2$.

Given $\Pi$ and a play $x$, with $T$ as above and $R(\theta)^2=|\theta-m|_M^2$, the four \emph{cells} of $(\Pi,x)$ are the tail $\{T>(2k)^{1/2}\}$ and, outside the tail, the outer cell $\{R^2>4d\}$, the inner cell $\{R^2<D_M/64\}$ and the regular cell $\{D_M/64\le R^2\le4d\}$. With $\varepsilon_1=k^3/d^3$, $\varepsilon_2=k^2/d^2$ and $\varepsilon_3=k^3/d^2$, the tail is \emph{ripe} if $\Pi(\mathrm{tail})\ge\varepsilon_1$ or $\E_\Pi[T^2;\mathrm{tail}]\ge\varepsilon_3$, the outer or the inner cell is ripe if its mass is at least $\varepsilon_2$, and the regular cell never is. The \emph{rest} is the union of the regular cell and of the cells that are not ripe.

The genie's rule maps a law $\Pi$, a play $x$ and a tip $\theta$ to an \emph{answer} and a \emph{signal}. If $\Pi$ is live, the answer is the cell of $(\Pi,x)$ that contains $\theta$ if that cell is ripe and ``rest'' otherwise, and the signal is $b+D_\Pi(x;\theta)$. If $\Pi$ is not live, the answer is ``rest'' and the signal is $b$. On a sample $(x^\circ,Z,\omega)$, against a learner, $\Pi_1$ is the uniform law on $r_0B$. In round $t$ the learner plays $x_t$ from its seed and its earlier plays and observations, the genie answers $A_t$, the answer at $(\Pi_t,x_t,x^\circ)$, and the learner observes $o_t=s_t(x^\circ)+\sigma Z_t$, with $s_t(\theta)$ the signal at $(\Pi_t,x_t,\theta)$. With $\varphi_\sigma$ the density of $N(0,\sigma^2)$ and $c_t(\theta)$ the answer at $(\Pi_t,x_t,\theta)$,
\[
\frac{d\Pi_{t+1}}{d\Pi_t}(\theta)=\frac{\mathbf 1[c_t(\theta)=A_t]\,\varphi_\sigma\bigl(o_t-s_t(\theta)\bigr)}{\int\mathbf 1[c_t=A_t]\,\varphi_\sigma(o_t-s_t)\,d\Pi_t},
\]
and $\Pi_{t+1}=\Pi_t$ where the denominator is $0$ or infinite. Round $t$ is \emph{live} if $\Pi_t$ is. A round that is not live has a constant answer and a constant signal, so $\Pi_{t+1}=\Pi_t$ and no later round is live. The answers stay with the adversary: the learner sees its observations only. This is the \emph{Gaussian game}. For a tip $x^\circ$ and a sequence $z$ of reals, the strategy $S_{x^\circ,z}$ gives in round $t$, after the plays $x_1,\dots,x_{t-1}$, the loss $f_t(y)=s_t(y;x^\circ)+\sigma z_t$, where $s_t(y;\theta)$ is the signal at $(\Pi_t,y,\theta)$ and $\Pi_t$ is computed by the recursion on these plays, with the observations $s_u(x_u;x^\circ)+\sigma z_u$ and the answers at $x^\circ$. The clamp of $z$ at $\ell=(6\ln n)^{1/2}$ is $z'_t=\max(-\ell,\min(\ell,z_t))$.

\begin{lemma}[measurability]\label{lem:meas}
For every family $(\rho_h)$ of probability densities on $\Pi_1$ with finite values, indexed by a measurable space and jointly measurable in $(h,\theta)$, the answer and the signal at $(\rho_h\Pi_1,x,\theta)$ are measurable in $(h,x,\theta)$. Hence, in the Gaussian game against a learner whose maps are jointly measurable, $\Pi_t$ has a density on $\Pi_1$ that is measurable in the history and the tip, and the history of the first $t\le n$ rounds is a measurable function of the sample; and the loss of $S_{x^\circ,z}$ in round $t$ at each point is a measurable function of the earlier plays.
\end{lemma}

\begin{proof}
For each rational $\lambda\in(0,1]$ the function $\lambda q(\theta+(x-\theta)/\lambda)$ is continuous in $(m,M,x,\theta)$, and the dent is the infimum of these functions. $M$ and $D_M$ are continuous functions of $\Sigma$, and $d_{\mathrm{live}}$, the trace of a spectral projection of $\Sigma$, is the limit of the traces of $f_j(\Sigma)$ with $f_j(s)=\min(1,\max(0,1+j(s-w_*^2)))$. The mean, the covariance, the masses of the cells and the second moment of the tail are integrals of the density, measurable in $h$ by Fubini's theorem. The density of $\Pi_{t+1}$ is that of $\Pi_t$ times measurable factors, divided by their integral, off the measurable set where that integral is $0$ or infinite.
\end{proof}

\subsection{The lemmas}\label{sec:lower-lemmas}

\begin{lemma}[legality]\label{lem:legal}
At every law $\Pi$ on $r_0B$ and every tip $\theta$, the dent $D_\Pi(\cdot;\theta)$ is convex and $L$-Lipschitz, and $0\le D_\Pi(x;\theta)\le\frac32L+g$ for $x\in B$. For every tip $x^\circ$ and every sequence $z$ with $|\sigma z_t|\le\frac18$ for all $t$, $S_{x^\circ,z}$ is a strategy on $B$ for $n$ rounds whose loss of round $t$ depends only on $x_1,\dots,x_{t-1}$ and is $\frac14$-Lipschitz on $B$. If $x^\circ\in B$, then for every learner and every seed $\omega$, the regret against $S_{x^\circ,z}$ in the game without noise is
\[
\Reg_n=\sum_{t=1}^n\bigl(s_t(x_t;x^\circ)-b\bigr),
\]
with $\Pi_t$ and $x_t$ those of the Gaussian game on the sample $(x^\circ,z,\omega)$.
\end{lemma}

\begin{proof}
Convexity is Lemma~\ref{lem:dent}~(a). The Lipschitz constant is $\kappa r_1\|M^{1/2}\|\le\kappa r_1/w_*=L$ by Lemma~\ref{lem:dent}~(b) and (F1). The mean of $\Pi$ lies in $r_0B$, so for $x\in B$, $|x-m|\le\frac32$, and $0\le D\le q\le\kappa(r_1|x-m|/w_*+k)\le\frac32L+g$, using $\phi(a)\le r_1a+k$ and (F1). With $b=\frac14$, $\frac32L+g\le\frac12$ and $|\sigma z_t|\le\frac18$ the losses take values in $[\frac18,\frac78]$ on $B$, and they are measurable by Lemma~\ref{lem:meas}. The losses at $x^\circ$ are $b+\sigma z_t$, their least values, since $D\ge0=D(x^\circ;x^\circ)$; so the infimum in the regret is attained at $x^\circ$, and the run in the game without noise is that of the Gaussian game with the noise $z$.
\end{proof}

\begin{lemma}[the posterior]\label{lem:post}
In the Gaussian game against a learner whose maps are jointly measurable, for every $t\le n+1$ and every measurable $F\ge0$, with $G_{t-1}$ the plays, observations and answers of the rounds before $t$,
\[
\E\,F(x^\circ,Z_t,G_{t-1},\omega)=\E\int\!\!\int F(\theta,z,G_{t-1},\omega)\,N(0,1)(dz)\,\Pi_t(d\theta).
\]
\end{lemma}

\begin{proof}
Induction on $t$. In one round, with the tip of law $\Pi$ and the pair (answer, observation) of law $\delta_{c(\theta)}\otimes N(s(\theta),\sigma^2)$ given the tip, Bayes' formula with the normaliser written out disintegrates the joint law along the posterior; the normaliser is $0$ or infinite on an event of probability $0$, where the convention $\Pi_{t+1}=\Pi_t$ changes nothing. The noise of the next round is independent of everything before it.
\end{proof}

For a law $\Pi$ of the tip, a measurable answer map $c$ with finitely many values, a bounded measurable signal $s$ and the observation $y=s(\theta)+\sigma z$, let $Z(a,y)=\int\mathbf 1[c=a]\varphi_\sigma(y-s)\,d\Pi$ and
\[
J=\E\ln\frac{\varphi_\sigma(\sigma z)}{Z(c(\theta),y)},
\]
the mutual information of the tip and of the pair (answer, observation), with $\theta\sim\Pi$ and $z\sim N(0,1)$ independent. Let $H(A)=\sum_ap_a\ln(1/p_a)$, $p_a=\Pi(c=a)$. For a probability density $\rho$ on a law $\Pi_0$ put $K(\rho)=\int\rho\ln\rho\,d\Pi_0$, the divergence of $\rho\,\Pi_0$ from $\Pi_0$.

\begin{lemma}[the chain rule]\label{lem:chain}
(i) For all reals $\hat c_a$,
\[
J\le H(A)+\sum_a\frac{1}{2\sigma^2}\int_{\{c=a\}}(s-\hat c_a)^2\,d\Pi.
\]
(ii) If $\Pi=\rho\,\Pi_0$ with $\rho$ bounded and measurable, and $\rho'$ is the density after the round, $\rho\,\mathbf 1[c=a]\varphi_\sigma(y-s)$ normalised on $\Pi_0$ at the answer $a$ and the observation $y$, then $\E K(\rho')=K(\rho)+J$.
\end{lemma}

\begin{proof}
(i) Write $\ln(\varphi_\sigma(\sigma z)/Z)=\ln(1/p_a)+\ln(\varphi_\sigma(\sigma z)\,p_a/Z)$ with $a=c(\theta)$. The first term has mean $H(A)$. Given $a$, $Z(a,\cdot)/p_a$ is the density of $y$ under $\Pi(\cdot\mid c=a)$, and the mean of the second term is the mean divergence of $N(s(\theta),\sigma^2)$ from that law, which is at most its mean divergence from $N(\hat c_a,\sigma^2)$ by Gibbs' inequality, that is $\int_{\{c=a\}}(s-\hat c_a)^2\,d\Pi/(2\sigma^2p_a)$. (ii) At the true tip, $\ln\rho'(\theta)=\ln\rho(\theta)+\ln(\varphi_\sigma(\sigma z)/Z(c(\theta),y))$; integrate against the law of the round, using the disintegration of Lemma~\ref{lem:post} for one round.
\end{proof}

The bound (i) takes the place of the bound $\frac12\ln(1+\mathrm{Var}\,s/\sigma^2)$ of the Gaussian channel, so that no maximum of the entropy is needed. With the signal in an interval of length $\frac32L+g$ and $\hat c_a$ its midpoint, the cell $a$ costs at most $p_aB_\sigma$ nats, $B_\sigma=(\frac32L+g)^2/(8\sigma^2)$.

For the next two lemmas fix a law $\Pi$ on $r_0B$ and a play $x$, with the frame of $\Pi$, $q=q_\Pi(x)$, $D=D_\Pi(x;\cdot)$, the masses $p_{\mathrm{rest}}$ and $p_{\mathrm{reg}}$ of the rest and of the regular cell, and $P_R=d^2/(4194372\,k^2)$.

\begin{lemma}[the rest]\label{lem:rest}
Let $k\ge1$ and $d\ge2048k$.
\begin{enumerate}[label=(\roman*)]
\item If $D_M\ge d/2$, then $p_{\mathrm{reg}}\,g/2\le\E_\Pi D$. In all cases $p_{\mathrm{reg}}\ge p_{\mathrm{rest}}-3k^2/d^2$.
\item If $D_M\ge d/2$ and $p_{\mathrm{rest}}\ge\frac12$, then
\[
\frac{\E_\Pi[(q-D)^2;\mathrm{rest}]}{2\sigma^2p_{\mathrm{rest}}}\le\frac{g}{\sigma^2P_R}\cdot\frac{\E_\Pi[D;\mathrm{rest}]}{p_{\mathrm{rest}}}.
\]
\end{enumerate}
\end{lemma}

\begin{proof}
A regular tip satisfies the condition of Lemma~\ref{lem:dent}~(c), since $R\le2d^{1/2}$ and $\rho\le(4d+2k)^{1/2}$. By (c) and (e), $q-D\le\kappa\cdot2h_+^2/\rho^2\le\kappa\cdot8k^2/\rho^2\le\kappa\cdot512k^2/D_M$, as $\rho^2\ge R^2\ge D_M/64$. With $D_M\ge d/2$ and $d\ge2048k$ this is at most $\kappa k/2=g/2\le q/2$, since $q\ge\kappa k=g$; so $D\ge q/2\ge g/2$ on the regular cell, which gives the first claim of (i). The cells of the rest that are not regular are not ripe, of masses below $\varepsilon_1$, $\varepsilon_2$ and $\varepsilon_2$, which gives the second.

(ii) Divide by $\kappa^2$ and write $\phi=\phi(a)\ge k$. On the regular cell $(q-D)^2/\kappa^2\le(1024k^2/d)^2\le2^{20}(k^3/d^2)\phi$. On an inner or outer cell that is not ripe, $T\le(2k)^{1/2}$, and (d), (e) give $q-D\le2\kappa k$; these cells contribute at most $8(k^3/d^2)\phi$. On a tail that is not ripe, (d), (f) and $q-D\le q$ give $(q-D)^2/\kappa^2\le(4T^2+r_1T)\phi$, and $\E_\Pi[T^2;\mathrm{tail}]<\varepsilon_3$, $\E_\Pi[T;\mathrm{tail}]<(\varepsilon_1\varepsilon_3)^{1/2}$ bound its contribution by $9(k^3/d^2)\phi$. So $\E_\Pi[(q-D)^2;\mathrm{rest}]\le(2^{20}+17)\kappa^2(k^3/d^2)\phi$. By (i), $p_{\mathrm{reg}}\ge p_{\mathrm{rest}}/2$, so $\E_\Pi[D;\mathrm{rest}]\ge p_{\mathrm{reg}}q/2\ge p_{\mathrm{rest}}\kappa\phi/4$. With $\kappa k^3=gk^2$, $1/p_{\mathrm{rest}}\le2$ and $4(2^{20}+17)=4194372$ the claim follows.
\end{proof}

Let $\Lambda(\Pi)=\frac12\ln\det W$, the \emph{ledger}. For a cell $c$ of mass $p_c>0$ let $U_c=p_c^{-1}\E_\Pi[(\theta-m)(\theta-m)^{\mathsf T};c]+w_*^2I$ and $\Psi_c=\tr(W^{-1}U_c)-d-\ln\det U_c+\ln\det W$, which is $\sum_i\psi(u_{c,i})$ with $\psi(u)=u-1-\ln u$ and $u_{c,i}$ the eigenvalues of $W^{-1/2}U_cW^{-1/2}$. The \emph{payment} of a cell is $\frac12p_c\Psi_c$, and the payment $P'$ of a round is the sum of the payments of its ripe cells.

\begin{lemma}[the ledger]\label{lem:ledger}
\begin{enumerate}[label=(\roman*)]
\item For every law $\Pi$ carried by a ball, every measurable answer map with finitely many values, its level sets as the cells $c$, and every measurable signal, with $\Pi'$ the posterior after the answer and the observation $s(\theta)+\sigma z$, $\E\Lambda(\Pi')\le\Lambda(\Pi)-\frac12\sum_cp_c\Psi_c$.
\item Let $k\ge6$. A cell inside the tail pays at least $p_ck/4$ and at least $\E_\Pi[T^2;c]/8$. A cell inside $\{R^2>4d\}$ pays at least $p_cd/62$, and so does a cell inside $\{R^2<D_M/64\}$ if $\Pi$ is live.
\item Let $k\ge1$ and $d\ge2048k$. If every ripe cell pays as in (ii), then $H(A)\le\frac{4(3\ln d+1)}k\,P'$ and the mass of the ripe cells is at most $\frac4kP'$.
\item In the Gaussian game, for every learner whose maps are jointly measurable, $\sum_{t\le n}\E P'_t\le R_L=\frac d2\ln\bigl(1+r_0^2/(dw_*^2)\bigr)$, with $P'_t$ the payment of round $t$, which is $0$ on the rounds that are not live.
\end{enumerate}
\end{lemma}

\begin{proof}
Since $m$ is the mean, $\sum_cp_cU_c=W$, so $\sum_cp_c\Psi_c=\ln\det W-\sum_cp_c\ln\det U_c$. (i) For positive definite $X$ and $U$, $\ln\det X\le\ln\det U+\tr(U^{-1}X)-d$. Take $X=\Sigma(\Pi')+w_*^2I$ and $U=U_c$ for the answer $c$. $\Sigma(\Pi')$ is at most the second moment of $\Pi'$ about $m$, and given the answer $c$ the mean of that second moment over the observation is $U_c-w_*^2I$, by the disintegration of Lemma~\ref{lem:post} for one round; so the mean of $\tr(U_c^{-1}X)-d$ is at most $0$. (ii) For a cell in the tail, the largest $u_{c,i}$ is at least the Rayleigh quotient in the direction $W^{1/2}Me$, $e$ the $M$-unit vector along $x-m$, which is at least $\E_\Pi[T^2\mid c]/2>k$ by (F2); and $\psi(u)\ge u/2$ for $u\ge6$. For a cell in $\{R^2>4d\}$, $\sum_iu_{c,i}\ge\E_\Pi[R^2\mid c]/2>2d$ by (F2), and the convexity of $\psi$ gives $\Psi_c\ge d\,\psi(2)\ge0.306\,d$. For a cell in $\{R^2<D_M/64\}$ on a live law, $\sum_iu_{c,i}\le D_M/64+\frac34d\le\frac{49}{64}d$ by (F2) and (F3), and $\Psi_c\ge d\,\psi(\frac{49}{64})\ge0.0326\,d$. (iii) $H(A)\le\sum_{c\text{ ripe}}p_c(\ln(1/p_c)+1)$, since $p_{\mathrm{rest}}\ln(1/p_{\mathrm{rest}})\le1-p_{\mathrm{rest}}$. Each ripe cell is compared with its own payment: a tail of mass at least $\varepsilon_1$ with $p_ck/4$; a tail of mass below $\varepsilon_1$, ripe by its second moment, with $\varepsilon_3/8$, while $p_c(\ln(1/p_c)+1)\le\varepsilon_1(3\ln d+1)$ and $\varepsilon_1=(8/d)(\varepsilon_3/8)$; an inner or outer cell of mass at least $\varepsilon_2$ with $p_cd/62$, while $p_c(\ln(1/p_c)+1)\le p_c(2\ln d+1)$. Since $d\ge2048k$, the largest of the ratios is $4(3\ln d+1)/k$, and $8/d$ and $62/d$ are at most $4/k$. (iv) Sum (i) over the rounds, the cells being the level sets of the answer: a round that is not live has the single answer ``rest'' and pays nothing. Then $\Lambda\ge\frac d2\ln w_*^2$, and $\Lambda(\Pi_1)\le\frac d2\ln(r_0^2/d+w_*^2)$ since $\tr\Sigma(\Pi_1)\le r_0^2$.
\end{proof}

\begin{lemma}[the count]\label{lem:count}
Let $w>0$ with $X=r_0^2/(dw^2)\ge e/2$, and
\[
N_*=\frac d4\ln\frac{r_0^2}{2dw^2}-\frac12\ln(\pi d)-\frac1d.
\]
(i) If $\rho$ is a bounded measurable probability density on $\Pi_1$ and fewer than $d/2$ eigenvalues of the covariance of $\rho\,\Pi_1$ are at least $w^2$, then $K(\rho)\ge N_*$. (ii) In the Gaussian game, for every learner whose maps are jointly measurable, $N_*\,\Prob(\Pi_{n+1}\text{ is not live})\le\sum_{t\le n}\E J_t$, with $J_t$ the information of round $t$ and the floor $w$ in place of $w_*$ in the word live.
\end{lemma}

\begin{proof}
(i) Let $(m,M)$ be the frame of $\rho\,\Pi_1$ at the floor $w$, and $G$ the Gaussian density of mean $m$ and precision $M$ divided by the density of $\Pi_1$ on $r_0B$; its integral on $\Pi_1$ is at most $1$. Gibbs' inequality gives
\[
K(\rho)\ge\int\rho\ln G\,d\Pi_1=\ln\mathrm{vol}(r_0B)-\frac12\sum_i\ln\max(\lambda_i,w^2)-\frac d2\ln(2\pi)-\frac12D_M.
\]
Let $j<d/2$ eigenvalues be at least $w^2$. Their sum is at most $r_0^2$ and the others are below $w^2$, so $\sum_i\ln\max(\lambda_i,w^2)\le j\ln(r_0^2/j)+(d-j)\ln w^2$, which increases in $j$ up to $d/2$ since $X\ge e/2$ and so is at most $\frac d2\ln(2r_0^2w^2/d)$. Further $D_M\le d$, and $\ln\mathrm{vol}(r_0B)=d\ln r_0+\frac d2\ln\pi-\ln\Gamma(\frac d2+1)$ with $\ln\Gamma(\frac d2+1)\le\frac d2\ln\frac d2-\frac d2+\frac12\ln(\pi d)+\frac1d$, from Robbins' bound on $\ln m!$ at the integers and the logarithmic convexity of $\Gamma$ between them. (ii) Lemma~\ref{lem:chain}~(ii) summed over the rounds gives $\sum_{t\le n}\E J_t=\E K(\rho_{n+1})$, $\rho_{n+1}$ the density of $\Pi_{n+1}$; apply (i) where $\Pi_{n+1}$ is not live, and $K\ge0$ elsewhere.
\end{proof}

The term $1/d$ in $N_*$, where the remainder in Binet's formula would give $1/(6d)$, comes from this cruder bound on $\ln\Gamma$.

\subsection{The balance}\label{sec:lower-balance}

Let $N_*$ and $R_L$ be taken at the floor $w_*=\kappa r_1/L$. In the Gaussian game against a learner, for each round $t\le n$, with $S_t$ the event that round $t$ is live, let $J^m_t=\E[J_t;S_t]$, $P^m_t=\E[P'_t;S_t]$, $a^m_t=\E[a_t;S_t]$ and $D^m_t=\E\bigl[\E_{\Pi_t}D_{\Pi_t}(x_t;\cdot);S_t\bigr]$, where $a_t$ is the mass of the ripe cells of round $t$; let $l_t=\Prob(S_t)$, and $p$ the probability that $\Pi_{n+1}$ is not live.

\begin{proposition}[the balance]\label{prop:balance}
Let $k\ge1$, $g,\sigma,L>0$, $B\ge0$, $40(3\ln d+1+2B)\le k$, $2048k\le d$ and $g\le kL/(40d)$. If reals $J^m_t,P^m_t,a^m_t,D^m_t,l_t$ ($t\le n$) and $p$ satisfy, for every $t\le n$,
\begin{gather*}
J^m_t\le\frac{4(3\ln d+1)+8B}k\,P^m_t+\frac g{\sigma^2P_R}\,D^m_t,\qquad\frac g2\Bigl(\frac{999}{1000}\,l_t-a^m_t\Bigr)\le D^m_t,\qquad a^m_t\le\frac4kP^m_t,\\
1-p\le l_t,\qquad\sum_{t\le n}P^m_t\le R_L,\qquad N_*\,p\le\sum_{t\le n}J^m_t,
\end{gather*}
then
\[
\sum_{t\le n}D^m_t\ge\max\Bigl(\frac{\sigma^2P_R}g\,N_*\Bigl(p-\frac3{10}\Bigr),\ \frac g2\Bigl(\frac{999}{1000}\,n(1-p)-\frac{12N_*}k\Bigr)\Bigr).
\]
If $L=\sigma\le\frac14$, $B=B_\sigma$, $k=120\ln d+68\le d/2048$, $n\ge e^4d^5$, $\Gamma=\ln(n/d^5)$ and $g=2\cdot10^{-4}\sigma d^{3/2}\Gamma^{1/2}/(kn^{1/2})$, then the six conditions alone give
\[
\sum_{t\le n}D^m_t\ge2.5\cdot10^{-5}\,\frac{\sigma d^{3/2}(n\Gamma)^{1/2}}k.
\]
\end{proposition}

\begin{proof}
The first condition on $k$ gives $k\ge40$, hence $d\ge81920$, and $g\le kL/(40d)$ gives $X=(kL)^2/(100g^2d^2)\ge16$; then $R_L\le3N_*$. Summing the first condition and using $40(3\ln d+1+2B)\le k$ gives $N_*p\le\frac3{10}N_*+(g/(\sigma^2P_R))\sum_tD^m_t$, the first bound. Summing the second, with $\sum_tl_t\ge n(1-p)$ and $\sum_ta^m_t\le12N_*/k$, gives the second. For the rate, $X=k^4n/(4\cdot10^{-6}d^5\Gamma)$, $g\le\sigma/81920$, $B_\sigma\le0.2813$ and $40(3\ln d+1+2B_\sigma)\le k$, and $0.1248\,d\Gamma\le N_*\le\frac d4(\Gamma+4\ln k+11.8)$. If $p\ge\frac12$, the first bound is at least the rate. If $p<\frac12$, then $12N_*/k\le0.2495\,n$ and the second bound is at least $gn/8$, which is the rate.
\end{proof}

\begin{proposition}[the game]\label{prop:game}
Let $d\ge4\cdot10^6$, $n\ge d^{10}$, $\sigma=L=1/(8(6\ln n)^{1/2})$, $k=120\ln d+68$, $\Gamma=\ln(n/d^5)$ and $g=2\cdot10^{-4}\sigma d^{3/2}\Gamma^{1/2}/(kn^{1/2})$. These parameters satisfy the conditions of Section~\ref{sec:lower-strategy}, and those of Proposition~\ref{prop:balance} with $B=B_\sigma$; and $n\ge e^4d^5$. For every learner whose maps are jointly measurable, the numbers $J^m_t,P^m_t,a^m_t,D^m_t,l_t,p$ of the Gaussian game satisfy the six conditions of Proposition~\ref{prop:balance}, and
\[
\sum_{t\le n}D^m_t=\E\sum_{t=1}^n\bigl(s_t(x_t;x^\circ)-b\bigr).
\]
\end{proposition}

\begin{proof}
On a live round, given the history and the seed, $D_M\ge d_{\mathrm{live}}\ge d/2$ by (F3). Apply Lemma~\ref{lem:chain}~(i) with $\hat c_a$ the midpoint $b+\frac12(\frac32L+g)$ of the range of the signal on the ripe cells and, on the rest, the midpoint if $p_{\mathrm{rest}}<\frac12$ and $b+q$ if $p_{\mathrm{rest}}\ge\frac12$. The ripe cells cost at most $B_\sigma a_t$. The rest costs at most $p_{\mathrm{rest}}B_\sigma\le B_\sigma a_t$ in the first case and, by Lemma~\ref{lem:rest}~(ii) multiplied by $p_{\mathrm{rest}}$, at most $(g/(\sigma^2P_R))\,\E_{\Pi_t}D_{\Pi_t}(x_t;\cdot)$ in the second. With Lemma~\ref{lem:ledger}~(ii) and~(iii) this is the first condition on that round. Lemma~\ref{lem:rest}~(i), with $3k^2/d^2\le\frac1{1000}$ and $p_{\mathrm{rest}}=1-a_t$, gives the second, and the third is Lemma~\ref{lem:ledger}~(iii). These hold on $S_t$ and are integrated over it. A round that is not live leaves the law unchanged, so the events $S_t$ decrease in $t$ and $1-p\le l_t$. The payment is Lemma~\ref{lem:ledger}~(iv) and the count Lemma~\ref{lem:count}~(ii), with $J_t=0$ and $P'_t=0$ on the rounds that are not live. The last identity is Lemma~\ref{lem:post}: the mean of $s_t(x_t;x^\circ)-b$ is the mean of its integral against $\Pi_t$, which is $\E_{\Pi_t}D_{\Pi_t}(x_t;\cdot)$ on $S_t$ and $0$ off it.
\end{proof}

\subsection{Proof of Theorem~\ref{thm:lower}}\label{sec:lower-proof}

\begin{proof}[Proof of Theorem~\ref{thm:lower}]
Let $d\ge4\cdot10^6$, $n\ge d^{10}$, the parameters of Proposition~\ref{prop:game}, and a learner. On the sample $(x^\circ,Z,\omega)$ let $\Reg^G(x^\circ,Z,\omega)=\sum_{t\le n}(s_t(x_t;x^\circ)-b)$, the regret of the Gaussian game; it is measurable, with values in $[0,n/2]$. The runs with the noise $Z$ and with its clamp $Z'$ coincide unless $|Z_t|>\ell$ for some $t\le n$, an event of probability at most $2n\,e^{-\ell^2/2}=2n^{-2}$, so $\E\Reg^G(x^\circ,Z',\omega)\ge\E\Reg^G(x^\circ,Z,\omega)-1/n$. Propositions~\ref{prop:game} and~\ref{prop:balance} give $\E\Reg^G(x^\circ,Z,\omega)\ge2.5\cdot10^{-5}\sigma d^{3/2}(n\Gamma)^{1/2}/k$. With $\Gamma\ge\frac12\ln n$, $\sigma=1/(8(6\ln n)^{1/2})$, $\ln d\ge15$ and $k\le125\ln d$,
\[
\E\,\Reg^G(x^\circ,Z',\omega)\ge9\cdot10^{-7}\,\frac{d^{3/2}n^{1/2}}k-\frac1n\ge7\cdot10^{-9}\,\frac{d^{3/2}n^{1/2}}{\ln d}.
\]
The function $\Reg^G(x^\circ,Z',\omega)$ is bounded and measurable for the product of the law of $(x^\circ,Z)$ and of $\mu$, so there is a value of $(x^\circ,Z)$, with $x^\circ\in r_0B$, at which its integral in $\omega$ is at least its mean. Since $\sigma\ell=\frac18$, Lemma~\ref{lem:legal} applies to $S_{x^\circ,Z'}$: it is a strategy on $B$ whose loss of round $t$ depends only on $x_1,\dots,x_{t-1}$ and is $\frac14$-Lipschitz, and the regret of the learner against it in the game without noise is $\Reg^G(x^\circ,Z',\omega)$, which is integrable.
\end{proof}

\section{What the proofs do not give}\label{sec:limits}

Theorem~\ref{thm:upper} fixes neither $C'$ nor $k$, and its machine computes with exact real numbers; nothing is claimed for a machine that computes with a bounded number of bits. Its bound is on the upper expectation of the regret, with $\ln^5(2dn(n+1))$ where Theorem~\ref{thm:explicit} has $\ln^{3/2}(nD)$, and the constant $C=10^{22}$ that its proof gives is far from the $C_X=1.9\cdot10^5$ of the explicit learner, which pays for it with a time exponential in $d$: the bound $(5n^2(2\beta)^{1/2}D)^d$ on the size $m$ of its net. Theorem~\ref{thm:explicit} holds against an adaptive adversary for the pseudo-regret only, the supremum over the comparator taken outside the expectation, which is also the form of the bound of \cite{XZ25}; it gives nothing for the expected supremum, since the loss estimates are unbounded, and no bound with high probability.

Theorem~\ref{thm:lower} is proved on the unit ball, for $d\ge4\cdot10^6$ and $n\ge d^{10}$, with a factor $1/\ln d$, and its adversary reads the plays. The ball is not a polytope, and no lower bound is proved on the polytopes of Theorem~\ref{thm:polytope}. It bounds the expected regret of the learners whose maps are jointly measurable, and the program of Theorem~\ref{thm:upper} is not shown to be one of them. Against losses fixed in advance the proofs give no lower bound beyond the order $d\sqrt n$ of linear losses \cite{DHK08}: the exponent of $d$ in the minimax regret for losses fixed in advance, in the game without noise, lies between $1$ and $3/2$, and the question is open. For this question the lower bound of \cite{BLS25} quoted in Section~\ref{sec:intro} bounds what the information-ratio route can give, as \cite[Section~8.8]{Lat24} notes, and the monograph poses as an exercise the same lower bound, of order $d^2$ up to logarithms, for the quantity that governs exploration by optimisation; neither is a lower bound on the regret.


\section{The formal proof}\label{sec:formal}

This section says what the Lean formalization of the companion repository \cite{AdvBcoLean} states and proves, and what it trusts.

\subsection{The formal statements}\label{sec:formal-lean}
The formalization is written in Lean~4 \cite{Lean4} with Mathlib \cite{Mathlib20}, in the versions of Section~\ref{sec:formal-soft}. Seventeen theorems state results of the paper, each from definitions written out in a statement file of the repository:
\begin{itemize}
\item \texttt{AdversarialBco.theoremConj1UpperExp}: Theorem~\ref{thm:upper};
\item \texttt{AdversarialBco.theoremConj1UpperExp\_of\_conj1}: Theorem~\ref{thm:upper} from Remark~\ref{rem:upper-prob};
\item \texttt{AdversarialBco.theoremConj1Upper}: Remark~\ref{rem:upper-prob};
\item \texttt{AdversarialBco.theoremPoly32O\_holds}: Corollary~\ref{cor:poly};
\item \texttt{AdversarialBco.theoremPoly32O\_of\_conj1}: Corollary~\ref{cor:poly} from Remark~\ref{rem:upper-prob}, a second proof;
\item \texttt{AdversarialBco.theoremConj1UpperPoly}: Theorem~\ref{thm:polytope}\textup{(a)};
\item \texttt{AdversarialBco.theoremConj1UpperPoly\_of\_conj1}: Theorem~\ref{thm:polytope}\textup{(a)} from Remark~\ref{rem:upper-prob};
\item \texttt{AdversarialBco.theoremConj1UpperPolyExp}: Theorem~\ref{thm:polytope}\textup{(b)};
\item \texttt{AdversarialBco.theoremConj1UpperPolyExp\_of\_poly}: Theorem~\ref{thm:polytope}\textup{(b)} from Theorem~\ref{thm:polytope}\textup{(a)};
\item \texttt{AdversarialBco.theoremConj1Lower}: Theorem~\ref{thm:lower};
\item \texttt{AdversarialBco.Explicit.theoremExplicit}: Theorem~\ref{thm:explicit}, with its constants and its bound against an adaptive adversary;
\item \texttt{AdversarialBco.Explicit.lemmaMarg}: Lemma~\ref{lem:marg};
\item \texttt{AdversarialBco.Explicit.corollaryExplicit}: Corollary~\ref{cor:minimax}, with $C_X=1.9\cdot10^5$ and every $n\ge1$;
\item \texttt{AdversarialBco.BodySpec.dilTheorem}: Theorem~\ref{thm:dil}, for every family;
\item \texttt{AdversarialBco.BodySpec.bodyExplicit}: Proposition~\ref{prop:body}, with $C_0=1.2\cdot10^9$;
\item \texttt{AdversarialBco.BodySpec.bodyHExplicit}: Proposition~\ref{prop:body} for the inner body $\cJ$ of a convex body $\cK$ that contains a ball of radius $1$, with $\diam\cJ$ replaced in $\beta$ by any $D\ge\diam\cK$, the form that the proof of Theorem~\ref{thm:explicit} uses;
\item \texttt{AdversarialBco.corollaryADV}: Corollary~\ref{cor:minimax} for another class of learners, with a constant that is not given (below).
\end{itemize}
Theorem~\ref{thm:upper} is stated in \texttt{StatementConj1Exp.lean}, over \texttt{StatementConj1.lean}, where Remark~\ref{rem:upper-prob} is stated, over the machine, the noise and the regret of \texttt{StatementPoly.lean}, where Corollary~\ref{cor:poly} is stated; Theorem~\ref{thm:polytope}\textup{(a)} is stated in \texttt{StatementConj1Poly.lean} and Theorem~\ref{thm:lower} in \texttt{StatementConj1Lower.lean}, each over \texttt{StatementConj1.lean}, and Theorem~\ref{thm:polytope}\textup{(b)} in \texttt{StatementConj1PolyExp.lean}, over \texttt{StatementConj1Poly.lean}; of the others, the next three are stated in \texttt{StatementExplicit.lean}, the next three in \texttt{StatementBody.lean} and the last in \texttt{Statement.lean}.

In Theorem~\ref{thm:upper}, Remark~\ref{rem:upper-prob} and Corollary~\ref{cor:poly} the program is a list of instructions of Definition~\ref{def:machine}, the rounds are numbered from $0$, a strategy is a function of the round, the earlier points and the earlier observations, asked for $t<n$ to give a function convex on $\cK$ with values in $[0,1]$ on $\cK$ and Borel in the history at each point of $\cK$, and the law $\nu$ of the noise is asked that, given $\varepsilon_1,\dots,\varepsilon_{t-1}$, $\varepsilon_t$ has mean $0$ and $\E e^{\varepsilon_t^2}\le2$, with $e^{\varepsilon_t^2}$ integrable. The probability is the outer measure of the product of the law of independent standard Gaussian variables and $\nu$. In Theorem~\ref{thm:upper} the majorant $R$ is a measurable function with values in $[0,\infty)$, above the regret at every point, and its expectation is its Lebesgue integral.

In Theorem~\ref{thm:polytope} the program, the rounds, the strategies, the noise, the probability and, in part \textup{(b)}, the majorant $R$ are those of Remark~\ref{rem:upper-prob} and Theorem~\ref{thm:upper}, on the polytope. The polytope is given by an $m\times d$ matrix and a vector of rational numbers, and the height is a bound on the absolute values of the numerator and of the denominator in lowest terms of each entry; the closed ball of radius $1$ may have any centre, and the diameter hypothesis is written as $|x-y|\le n$ for all $x,y\in\cK$. The program is asked to contain no membership question, and it is run from the state of Section~\ref{sec:poly-polytope}, which holds $d$, $n$, $m$, $\delta$ and the coefficients and nothing else, with $\delta=1/(n+1)$ in part \textup{(b)}. The constants $C,C_1,k_1$ of part \textup{(a)} and $C',C'_1,k'_1$ of part \textup{(b)} are the $C,C',k$ of the two formal statements.

In Theorem~\ref{thm:lower} the rounds are numbered from $0$ and the game is that of Theorem~\ref{thm:upper} on $B$ with the noise $0$. A learner is a probability space $(\Omega,\mu)$ of seeds and maps from the seed, the earlier points and the earlier observations to $B$, asked for $t<n$ to be jointly measurable. A strategy is one of the class of Theorem~\ref{thm:upper} for $\cK=B$ whose loss, for $t<n$ and every history, does not read the earlier observations and is $\frac14$-Lipschitz on $B$. The regret is that of Theorem~\ref{thm:upper} at the points of the run, against the losses chosen on the run, and the statement asks that it be integrable in the seed and bounds its Lebesgue integral.

In Theorem~\ref{thm:explicit}, Lemma~\ref{lem:marg} and Corollary~\ref{cor:minimax} the body $\cK$ is a compact convex subset of $\R^d$ that contains a closed ball of radius $1$, a loss is a function convex on $\cK$ with values in $[0,1]$ on $\cK$, and the rounds are numbered from $0$. Algorithm~\ref{alg:exo} is written out. Where the algorithm scans the points of the lattice in a fixed order and computes step~1 by a deterministic procedure, the formal statement takes the order and the procedure as arguments: it holds for every order and for every procedure that meets step~1, and it states that both exist. The procedure is asked to meet step~1 at every positive probability vector, which loses nothing, since one that meets it at the weights the run reaches can be changed at the others without changing the law of play. The probability space is the finite set of the histories of actions and coins, with the product law of the proof of Theorem~\ref{thm:exo}. Theorem~\ref{thm:explicit} is stated with, in addition, $m\ge2$. In Theorem~\ref{thm:dil} the conclusion is stated for every entropic Gaussian at the temperature $\lambda_h$, with the statement that one exists, and $s_h$ is the operator norm of $S_h$; by Lemma~\ref{lem:ec} this is the theorem of the text. In Proposition~\ref{prop:body} a finitely supported law on $\cF(\cJ,n)$ is a finite family of weights on functions on $\R^d$ that are convex, $n$-Lipschitz and with values in $[0,1]$ on $\cJ$, and the law $\rho$ charges only $\cJ$ and integrates the functions.

The last theorem states Corollary~\ref{cor:minimax} for learners of another class. A learner there is a probability space $(\Omega,\mu)$ of seeds and, for each round $t$, a jointly measurable map from the seed and the observations $f_1(x_1),\dots,f_{t-1}(x_{t-1})$ to $\cK$; its earlier actions are functions of these, so it has the information of a learner of Section~\ref{sec:intro}. The minimax regret is the infimum, over the learners whose regret is integrable against every sequence of losses fixed in advance, of the supremum over the sequences of the expected regret~\eqref{eq:regret}; the restriction to integrable regret removes learners and so can only raise the infimum. The statement is that there is a constant $C$ such that this minimax regret is at most $C\,d^{3/2}\sqrt n\,\ln^{3/2}(n\diam\cK)$ for all $d\ge1$, $n\ge2$ and every such body $\cK$. Neither class of learners contains the other, since a $\cC$-policy asks no measurability and a measurable learner need not play in a finite set. Its proof runs through Theorem~\ref{thm:exo} against sequences fixed in advance and Proposition~\ref{prop:body} with constants it does not compute.

\subsection{The other statements}\label{sec:formal-proved}
Every other numbered statement is proved in Lean as a step of the proofs of these theorems, in the form in which the paper states it. The repository lists, for each numbered statement, the Lean declarations that prove it. Every inequality between explicit numbers that these proofs use is proved there as a step of the same proofs; the text marks as checked in Lean those it states without derivation. The formal proof departs from the text in the following places.
\begin{itemize}
\item The regular case of Theorem~\ref{thm:dil} is proved for every minimiser of $\mathcal E_{\lambda_h}$ over the invertible $S$, with $s_h$ the operator norm of $S$, and Lemma~\ref{lem:compact} for minimisers over the invertible $S$; by Lemma~\ref{lem:ec} the entropic Gaussian is one of them.
\item Theorem~\ref{thm:exo} is proved for every finite $\cC\subset\cJ$ of covering radius less than $\varepsilon$, the nearest point of Lemma~\ref{lem:moving} being a measurable selection.
\item Theorem~\ref{thm:gauss} is proved for $F(b+SZ)$, with $F$ an $L$-Lipschitz function on $\R^d$, $b$ a point and $S$ a linear map from $\R^k$ to $\R^d$; the theorem of the text is the case $b=0$, $S=\mathrm{Id}$. Its proof is the argument of Maurey and Pisier for continuously differentiable functions (Lemma~\ref{lem:mp}), in the formalization of R.~Degenne \cite{Maiti2026Power}, with two lemmas of the Tau Ceti project \cite{TauCeti}, both under the Apache License~2.0 and ported to the versions of Section~\ref{sec:formal-soft} with no proof changed; the smoothing of a Lipschitz function and the passage to the limit are written for this paper.
\item Section~\ref{sec:poly} and Appendices~\ref{app:epoch} to~\ref{app:pg} follow the cut of the formal proof of Theorem~\ref{thm:upper}, and the results of \cite{Lat24} they use are proved there in the form the text states (Proposition~\ref{prop:sifacts} and the tools before the proof of Lemma~\ref{lem:events}). Two of these proofs take routes of their own. Gaussian concentration with the constant $\frac12$, cited from \cite{TIS76}, is proved by the covariance identity along a rotation of a pair of independent Gaussian vectors, Gaussian integration by parts in the rotated variable and the argument of Herbst, then extended to Lipschitz functions by the smoothing of Theorem~\ref{thm:gauss}. Freedman's inequality \cite[Theorem~B.20]{Lat24} is proved with $4$ in its logarithm, where the source prints $2$.
\item Lemma~\ref{lem:fd} is proved for exact values, $\eta_0=\eta_i=0$; the proof of Lemma~\ref{lem:sep} adds the errors of the values it reads.
\end{itemize}

\subsection{What the formal proof trusts}\label{sec:formal-trust}
The seventeen theorems depend on the axioms \texttt{propext}, \texttt{Classical.choice} and \texttt{Quot.sound} only, and the package has no \texttt{sorry}: a source lint of every file of the package other than the challenges below finds no \texttt{sorry} and no \texttt{native\_decide}, and of the constructs it looks for it finds only 29 tactic macros.
The Lean kernel checks every declaration when the package is built. Comparator \cite{Comparator} checks that the theorems prove, with no other axiom, the statements of challenge modules that import only Mathlib and each other. There is one challenge for each statement file: a copy of the file that differs from it only in its comments, over copies of the statement files it imports, and the theorems to prove, stated with no proof, in the same module or, for \texttt{StatementPoly.lean}, \texttt{StatementConj1.lean} and \texttt{StatementConj1Poly.lean}, whose copies other challenges import, in a module of their own above the copy. The copies of \texttt{StatementExplicit.lean} and \texttt{StatementPoly.lean} carry with them, verbatim, the definitions of \texttt{Statement.lean} they use; the copy of \texttt{StatementConj1.lean} is a second module over that of \texttt{StatementPoly.lean}, those of \texttt{StatementConj1Exp.lean}, \texttt{StatementConj1Lower.lean} and \texttt{StatementConj1Poly.lean} are each a third module over the two, and that of \texttt{StatementConj1PolyExp.lean} is a fourth module over that of \texttt{StatementConj1Poly.lean}. Within Comparator, nanoda \cite{nanoda}, a second implementation of the kernel, checks again every declaration the theorems depend on, Mathlib included. Comparator accepts the seventeen theorems: on the release tree, each of the nine runs, one for each challenge, ends with nanoda accepting the solution.

\subsection{Software and hardware}\label{sec:formal-soft}
The formalization uses Lean~4.34.1 and Mathlib~v4.34.1. The runs below were made on one machine, an AMD Ryzen~9 5900X with 12 cores and 64~GB of memory. With Mathlib from its cache, the build of the whole package from an empty build directory took 25 minutes of wall-clock time on three cores, pinned and under a CPU quota of 300 per cent, and its largest process used 4.6 GiB of memory. Each Comparator run, with nanoda, on four cores and under a cap of 8 GiB, took between 69 and 178 seconds, with a peak memory between 2763 and 5210 MiB.

\appendix

\section{Gaussian concentration for Lipschitz functions}\label{app:gauss}

Gaussian concentration for Lipschitz functions is due to \cite[Theorem~1A and Corollary~1]{TIS76}, with the constant $1/2$ in the exponent. This appendix proves the form used in Section~\ref{sec:body}, with the constant $2/\pi^2$, by the interpolation argument of Maurey and Pisier \cite[Theorem~2.2, p.~176, and p.~177]{Pis86}. For a standard Gaussian vector $W$ of $\R^k$, $\E e^{b|W|}\le\prod_i\E e^{b|W_i|}\le2^ke^{kb^2/2}$, since $\E e^{aW_1}=e^{a^2/2}$; so a function on $\R^k$ bounded by $c+\ell|w|$, and its exponential $e^{s(c+\ell|w|)}$, are integrable against the law of $W$. If $X$ and $Y$ are independent standard Gaussian vectors of $\R^k$ and $\theta$ is real, then $(\sin\theta\,X+\cos\theta\,Y,\ \cos\theta\,X-\sin\theta\,Y)$ is again a pair of independent standard Gaussian vectors, since its law is the image of the standard Gaussian law of $\R^{2k}$ under an orthogonal map.

\begin{theorem}\label{thm:gauss}
Let $Z$ be a standard Gaussian vector of $\R^k$ and $G\colon\R^k\to\R$ an $\ell$-Lipschitz function with $\ell>0$. Then $G(Z)$ is integrable and, for every $t\ge0$,
\[
\Prob\bigl(G(Z)-\E G(Z)\ge t\bigr)\le e^{-2t^2/(\pi^2\ell^2)}.
\]
\end{theorem}

The proof treats continuously differentiable functions first.

\begin{lemma}\label{lem:mp}
Let $Z$ be a standard Gaussian vector of $\R^k$ and $G\colon\R^k\to\R$ of class $C^1$ with $|\nabla G|\le\ell$. Then $G(Z)$ is integrable and $\E\,e^{s(G(Z)-\E G(Z))}\le e^{\pi^2s^2\ell^2/8}$ for every real $s$.
\end{lemma}

\begin{proof}
Since $|G(x)|\le|G(0)|+\ell|x|$, $G(Z)$ and $e^{sG(Z)}$ are integrable. Let $X$ and $Y$ be independent standard Gaussian vectors of $\R^k$. Jensen's inequality in $Y$ gives $\E e^{s(G(X)-\E G(Y))}\le\E e^{s(G(X)-G(Y))}$. For $u,v\in\R^k$ and real $\theta$ put $U_\theta=\sin\theta\,u+\cos\theta\,v$ and $V_\theta=\cos\theta\,u-\sin\theta\,v$; the derivative of $\theta\mapsto G(U_\theta)$ is $\ip{\nabla G(U_\theta),V_\theta}$, continuous in $\theta$, so
\[
G(u)-G(v)=\int_0^{\pi/2}\ip{\nabla G(U_\theta),V_\theta}\,d\theta .
\]
Jensen's inequality for the uniform probability on $[0,\pi/2]$ gives
\[
e^{s(G(u)-G(v))}\le\frac2\pi\int_0^{\pi/2}\exp\Bigl(\frac{\pi s}2\ip{\nabla G(U_\theta),V_\theta}\Bigr)\,d\theta .
\]
Take $u=X$ and $v=Y$, and take expectations (the integrand is nonnegative, so Fubini's theorem applies). For each $\theta$ the pair $(U_\theta,V_\theta)$ is a pair of independent standard Gaussian vectors, so given $U_\theta$ the variable $\ip{\nabla G(U_\theta),V_\theta}$ is a centred Gaussian of variance $|\nabla G(U_\theta)|^2\le\ell^2$, and
\[
\E\exp\Bigl(\frac{\pi s}2\ip{\nabla G(U_\theta),V_\theta}\Bigr)=\E\exp\Bigl(\frac{\pi^2s^2}8|\nabla G(U_\theta)|^2\Bigr)\le e^{\pi^2s^2\ell^2/8}.
\]
The weight $2/\pi$ on an interval of length $\pi/2$ gives $\E e^{s(G(X)-G(Y))}\le e^{\pi^2s^2\ell^2/8}$, and the first inequality of the proof gives the claim.
\end{proof}

\begin{proof}[Proof of Theorem~\ref{thm:gauss}]
Since $|G(x)|\le|G(0)|+\ell|x|$, $G(Z)$ and $e^{sG(Z)}$ are integrable for every real $s$. Fix $\epsilon>0$, let $\psi_\epsilon$ be a nonnegative function of class $C^\infty$ on $\R^k$ with support in the closed ball $B(0,\epsilon)$ and $\int\psi_\epsilon=1$, and $G_\epsilon(x)=\int G(x-y)\psi_\epsilon(y)\,dy$. As in the proof of Lemma~\ref{lem:smooth}, which uses no convexity for these two facts, $G_\epsilon$ is of class $C^\infty$ and $\ell$-Lipschitz, so $|\nabla G_\epsilon|\le\ell$; and $|G_\epsilon-G|\le\ell\epsilon$, since $|G(x-y)-G(x)|\le\ell\epsilon$ on the support of $\psi_\epsilon$. The centred variables $G(Z)-\E G(Z)$ and $G_\epsilon(Z)-\E G_\epsilon(Z)$ then differ by at most $2\ell\epsilon$, and Lemma~\ref{lem:mp} gives, for every real $s$,
\[
\E\,e^{s(G(Z)-\E G(Z))}\le e^{2|s|\ell\epsilon}\,\E\,e^{s(G_\epsilon(Z)-\E G_\epsilon(Z))}\le e^{2|s|\ell\epsilon}e^{\pi^2s^2\ell^2/8}.
\]
Letting $\epsilon\to0$ gives $\E\,e^{s(G(Z)-\E G(Z))}\le e^{\pi^2s^2\ell^2/8}$. For $t\ge0$ and $s=4t/(\pi^2\ell^2)$, Markov's inequality gives
\[
\Prob\bigl(G(Z)-\E G(Z)\ge t\bigr)\le e^{-st+\pi^2s^2\ell^2/8}=e^{-2t^2/(\pi^2\ell^2)}.
\]
\end{proof}

In Section~\ref{sec:body} the theorem is applied to $Z\mapsto r(b_h+S_hZ)$, which is $s_h$-Lipschitz since $r$ is $1$-Lipschitz and $|S_hv|\le s_h|v|$.

The formal proofs of Lemmas~\ref{lem:half} and~\ref{lem:ibp}(b) use Gaussian integration by parts in one coordinate, the next lemma; the proofs of the text do not. Write $\gamma_k$ for the standard Gaussian density of $\R^k$. A function on $\R^k$ has \emph{exponential growth} if it is bounded by $c\,e^{b|w|}$ for some constants $b$, $c$; by the bound on $\E e^{b|W|}$ above, such functions, and their products with polynomials, are integrable against $\gamma_k$.

\begin{lemma}\label{lem:gibp}
Let $W$ be a standard Gaussian vector of $\R^k$ and $H:\R^k\to\R$ differentiable at every point, with $|H(w)|\le Ae^{B|w|}$ and $\|DH(w)\|\le Ae^{B|w|}$ for every $w$, for some constants $A,B$. Then $\E[W_iH(W)]=\E[\partial_iH(W)]$.
\end{lemma}

\begin{proof}
For $k=1$, integrate by parts on $[-r,r]$ with $\gamma_1'(u)=-u\gamma_1(u)$; the boundary terms are at most $2Ae^{Br-r^2/2}\to0$ and both integrals converge absolutely. For general $k$, apply the case $k=1$ in the coordinate $w_i$ for fixed other coordinates $w'$, with the constant $Ae^{B|w'|}$, since $|w|\le|w_i|+|w'|$ and $|\partial_iH|\le\|DH\|$, and integrate in $w'$ (Fubini).
\end{proof}

\section{Pinsker's inequality for two Bernoulli laws}\label{app:info}

The proof of Lemma~\ref{lem:conv} compares two coins through the following inequality, with $\mathrm{kl}$ as in Section~\ref{sec:explicit}. The hypotheses exclude the pairs at which $\mathrm{kl}(a,b)$ is infinite.

\begin{lemma}\label{lem:pinsker}
Let $a,b\in[0,1]$ with $a=0$ if $b=0$ and $a=1$ if $b=1$. Then $(a-b)^2\le\frac12\,\mathrm{kl}(a,b)$.
\end{lemma}

\begin{proof}
If $b\in\{0,1\}$ then $a=b$ and both sides vanish. Let $0<b<1$ and $g(b')=\mathrm{kl}(a,b')-2(a-b')^2$ for $b'\in(0,1)$. Then $g(a)=0$ when $0<a<1$, and
\[
g'(b')=\frac{b'-a}{b'(1-b')}-4(b'-a)=(b'-a)\Bigl(\frac1{b'(1-b')}-4\Bigr),
\]
which has the sign of $b'-a$, since $b'(1-b')\le\frac14$. So $g$ is nonincreasing on $(0,a]$ and nondecreasing on $[a,1)$, and $g(b)\ge g(a)=0$. When $a\in\{0,1\}$, $g$ is monotone on $(0,1)$ with limit $0$ at $a$, and the same conclusion holds.
\end{proof}

\section{From regular families to every family}\label{app:calc}

This appendix proves the uniqueness of the entropic Gaussian and Theorem~\ref{thm:dil} for every family, from the regular case of Section~\ref{sec:core}. The tools are a smoothing, a uniqueness lemma and a compactness lemma; no gradient of a nonsmooth function enters. Symmetrisation is the inequality $\E G(b+w)\ge G(b)$ of the proof of Lemma~\ref{lem:ec}, for $G$ convex and $w$ an integrable random vector whose law is even, since $\E G(b+w)=\E\frac12(G(b+w)+G(b-w))$; here $w$ is a centred Gaussian vector (Lemmas~\ref{lem:unique} and~\ref{lem:compact}).

\begin{lemma}[smoothing]\label{lem:smooth}
Let $F\colon\R^d\to\R$ be convex and $L$-Lipschitz, and $\epsilon>0$. There is a function $F_\epsilon\colon\R^d\to\R$ of class $C^2$, convex, $L$-Lipschitz and with a Lipschitz gradient, such that $|F_\epsilon-F|\le\epsilon$.
\end{lemma}

\begin{proof}
Let $s=\epsilon/(L+1)$, $\psi$ a nonnegative function of class $C^\infty$ with support in the closed ball $B(0,s)$ and $\int\psi=1$ (a multiple of $y\mapsto\exp(-1/(1-|y|^2/s^2))$ on the open ball, extended by $0$, is one), and $F_\epsilon(x)=\int F(x-y)\psi(y)\,dy$. Then $F_\epsilon$ is an average of the convex functions $x\mapsto F(x-y)$, and $|F_\epsilon(x)-F_\epsilon(x')|\le\int|F(x-y)-F(x'-y)|\,\psi(y)\,dy\le L|x-x'|$. It is of class $C^\infty$ with $\nabla F_\epsilon(x)=\int F(x-y)\nabla\psi(y)\,dy$, so $|\nabla F_\epsilon(x)-\nabla F_\epsilon(x')|\le\int|F(x-y)-F(x'-y)|\,|\nabla\psi(y)|\,dy\le L|x-x'|\int|\nabla\psi|$. Finally $|F(x-y)-F(x)|\le L|y|\le Ls<\epsilon$ on the support of $\psi$, so $|F_\epsilon-F|\le\epsilon$.
\end{proof}

\begin{lemma}[uniqueness]\label{lem:unique}
Let $\Fbar\colon\R^d\to\R$ be convex, Lipschitz and coercive, and $\lambda>0$. Then $\mathcal E_\lambda$ has at most one minimiser.
\end{lemma}

\begin{proof}
For each value of $Z$, $(b,S)\mapsto\Fbar(b+SZ)$ is convex, so $\E\Fbar(b+SZ)$ is convex in $(b,S)$. The function $S\mapsto-\ln\det S$ is strictly convex on positive definite matrices: with $S_t=(1-t)S_1+tS_2$ and $\nu_i$ the eigenvalues of $S_1^{-1/2}(S_2-S_1)S_1^{-1/2}$, $\ln\det S_t=\ln\det S_1+\sum_i\ln(1+t\nu_i)$, strictly concave in $t$ unless every $\nu_i$ vanishes. So two minimisers $(b_1,S_1)$ and $(b_2,S_2)$ have $S_1=S_2=S$, since otherwise their midpoint would give a smaller value. Suppose $b_1\ne b_2$ and put $v=b_2-b_1$. Then $b\mapsto\mathcal E_\lambda(b,S)$ is convex and minimal at $b_1$ and $b_2$, so $t\mapsto\E\Fbar(b_1+tv+SZ)$ is constant on $[0,1]$. With $A=S^2$, $\eta_v=A^{-1}v/(v^{\mathsf T}A^{-1}v)$ and $\zeta_v=\ip{\eta_v,SZ}$, write $SZ=\zeta_vv+W'$; $\zeta_v$ and $W'$ are jointly Gaussian with $\E[W'\zeta_v]=A\eta_v-v/(v^{\mathsf T}A^{-1}v)=0$, hence independent, and $\zeta_v$ has a Gaussian law with positive variance. Let $\phi_v(r)=\E\Fbar(b_1+rv+W')$, a convex function with $\phi_v(r)\ge a_1|b_1+rv|-a_2$ by coercivity and symmetrisation, so $\phi_v$ is not affine. By Fubini's theorem $\E\Fbar(b_1+tv+SZ)=\E\phi_v(t+\zeta_v)$, and the value at $t=1/2$ is the mean of the values at $0$ and $1$:
\[
\E\Bigl[\tfrac12\bigl(\phi_v(\zeta_v)+\phi_v(1+\zeta_v)\bigr)-\phi_v\bigl(\tfrac12+\zeta_v\bigr)\Bigr]=0 .
\]
The integrand is continuous in $\zeta_v$ and nonnegative by convexity, and $\zeta_v$ has a positive density, so $\phi_v(r+\frac12)=\frac12(\phi_v(r)+\phi_v(r+1))$ for every real $r$. A convex function equal to its chord at the midpoint of an interval is affine on that interval: if it were strictly below the chord at some point of $(r,r+\frac12)$, convexity on the interval from that point to $r+1$ would put its value at $r+\frac12$ strictly below the chord, and symmetrically on the other side. So $\phi_v$ is affine on every interval $[r,r+1]$, hence on $\R$, a contradiction.
\end{proof}

\begin{lemma}[compactness]\label{lem:compact}
Let $G_k$ $(k\ge1)$ and $G$ be convex functions on $\R^d$, each $G_k$ being $L'$-Lipschitz, with $G_k\to G$ locally uniformly and $G_k(x)\ge a_1|x|-a_2$ for all $k$ and $x$, where $a_1>0$. Let $\lambda_k\to\lambda>0$ and $(b_k,S_k)$ be the entropic Gaussian of $G_k$ at $\lambda_k$. Then a subsequence of $(b_k,S_k)$ converges to the entropic Gaussian $(b,S)$ of $G$ at $\lambda$, and along it $\E H_k(b_k+S_kZ)\to\E H(b+SZ)$ for all continuous $H_k$ converging locally uniformly to $H$ (which is then continuous) with $|H_k(y)|\le c(1+|y|^2)$ for all $k$ and $y$.
\end{lemma}

\begin{proof}
Write $\mathcal E_k$ for the objective \eqref{eq:entropic} of $G_k$ at $\lambda_k$ and $\mathcal E$ for that of $G$ at $\lambda$; $G$ is convex, $L'$-Lipschitz and coercive with the constants $a_1$, $a_2$. Since $G_k(0)\to G(0)$, $|G_k(y)|\le c_0+L'|y|$ for all $k$. So $\mathcal E_k(b_k,S_k)\le\mathcal E_k(0,\mathrm{Id})\le c_0+L'\E|Z|=:c_2$. Let $s_k$ and $s_k'$ be the largest and smallest eigenvalues of $S_k$. As in the proof of Lemma~\ref{lem:ec}, $\E G_k(b_k+S_kZ)\ge a_1\max(|b_k|,(2/\pi)^{1/2}s_k)-a_2$, and $\lambda_k\ln\det S_k\le\lambda_kd\lnp s_k$; so $a_1\max(|b_k|,(2/\pi)^{1/2}s_k)-a_2-\lambda_kd\lnp s_k\le c_2$, which bounds $s_k$ and $|b_k|$. Next $\lambda_k\ln\det S_k=\E G_k(b_k+S_kZ)-\mathcal E_k(b_k,S_k)\ge-a_2-c_2$, so $\ln s_k'\ge\ln\det S_k-(d-1)\lnp s_k$ is bounded below. A subsequence converges to some $(b,S)$ with $S$ positive definite. For each value of $Z$, $y_k=b_k+S_kZ\to y=b+SZ$ and $H_k(y_k)\to H(y)$, with $|H_k(y_k)|\le c'(1+|Z|^2)$; dominated convergence gives $\E H_k(y_k)\to\E H(y)$. With $H_k=G_k$ this gives $\mathcal E_k(b_k,S_k)\to\mathcal E(b,S)$, and for fixed $(b',S')$ the same argument gives $\mathcal E_k(b',S')\to\mathcal E(b',S')$. So $\mathcal E(b,S)\le\mathcal E(b',S')$: $(b,S)$ minimises $\mathcal E$, and it is the entropic Gaussian of $G$ by Lemma~\ref{lem:unique}.
\end{proof}

\begin{lemma}\label{lem:approx}
Theorem~\ref{thm:dil} holds for every family.
\end{lemma}

\begin{proof}
Let $(\mu_c,F_c)$ be a family with $\Delta_\star>0$ and $b^\star$ a minimiser of $\Fbar$. For $k\ge1$ let $\epsilon_k=1/k$, and let $F_c^k$ be the function that Lemma~\ref{lem:smooth} gives for the convex $(L+\epsilon_k)$-Lipschitz function $F_c+\epsilon_k|\cdot-b^\star|$ with $\epsilon=\epsilon_k^2$. It is convex, $(L+\epsilon_k)$-Lipschitz, of class $C^2$ with a Lipschitz gradient, and
\[
\bigl|F_c^k-F_c-\epsilon_k|\cdot-b^\star|\bigr|\le\epsilon_k^2 .
\]
So $F_c^k$ is coercive and attains its minimum, and $(\mu_c,F_c^k)$ is a regular family whose mean $\Fbar_k$ has $|\Fbar_k-\Fbar-\epsilon_k|\cdot-b^\star||\le\epsilon_k^2$, so $\Fbar_k\ge\Fbar-1$ is coercive with the constants $a_1$ and $a_2+1$. As $k\to\infty$, $F_c^k\to F_c$ and $\Fbar_k\to\Fbar$ locally uniformly, and so do $u_c^k=F_c^k-\Fbar_k$ and $V_k=\sum_c\mu_c(u_c^k)^2$; since $u_c^k$ is $2(L+1)$-Lipschitz and $u_c^k(0)$ is bounded, $|V_k(y)|\le c(1+|y|^2)$ for all $k$. The minima converge: $\min F_c-\epsilon_k^2\le\min F_c^k\le F_c^k(x_c)\le\min F_c+\epsilon_k|x_c-b^\star|+\epsilon_k^2$. A minimiser $b^\star_k$ of $\Fbar_k$ satisfies $\Fbar(b^\star)+\epsilon_k|b^\star_k-b^\star|-\epsilon_k^2\le\Fbar_k(b^\star_k)\le\Fbar_k(b^\star)\le\Fbar(b^\star)+\epsilon_k^2$, so $|b^\star_k-b^\star|\le2\epsilon_k\to0$ and $\min\Fbar_k\to\min\Fbar$. Hence the floor $\Delta_\star^k$ of the family $k$ tends to $\Delta_\star$, its gains $g_c^k=F_c^k(b^\star_k)-\min F_c^k$ tend to $g_c$, and $g^k_{\max}\to g_{\max}$.

For $k$ large, $\Delta^k_\star>0$, and the regular case of Theorem~\ref{thm:dil} for the family $k$ and its minimiser $b^\star_k$ gives a scale $h_k=g^k_{c_k}\in[\Delta^k_\star/2,g^k_{\max}]$ and the entropic Gaussian $(b_k,S_k)$ of $\Fbar_k$ at $h_k/(16d)$, with \eqref{eq:P3} for the family $k$. Along a subsequence $c_k$ is a constant $c_0$, so $h_k\to h=g_{c_0}\in[\Delta_\star/2,g_{\max}]$ and $h>0$. By Lemma~\ref{lem:compact} with $G_k=\Fbar_k$ and $G=\Fbar$, a further subsequence has $(b_k,S_k)\to(b_h,S_h)$, the entropic Gaussian of $\Fbar$ at $\lambda_h=h/(16d)$. Along it $V_k(b^\star_k)\to V(b^\star)$, $V_k(b_k)\to V(b_h)$ and, by Lemma~\ref{lem:compact} with $H_k=V_k$, $\E V_k(b_k+S_kZ)\to\E V(b_h+S_hZ)$; the largest eigenvalue of $S_k$ tends to that of $S_h$, so the quantity $T$ of the family $k$, formed with its Lipschitz constant $L+\epsilon_k$, tends to $T_h$; and $\ln(2g^k_{\max}/\Delta^k_\star)\to\ln(2g_{\max}/\Delta_\star)$. Every term of \eqref{eq:P3} converges, so \eqref{eq:P3} holds for $(\mu_c,F_c)$ at the scale $h$, one of the values $g_c$.

(S4) holds for $(b_h,S_h)$ by Lemma~\ref{lem:ec}, and so does (S3) in the limit: for each $k$ and $z$, Lemma~\ref{lem:ec} for the mean $\Fbar_k$, of class $C^1$, gives $\E\Fbar_k(b_k+S_kZ)\le\Fbar_k(z)+\lambda_kd$ with $\lambda_k=h_k/(16d)$; along the subsequence the left side tends to $\E\Fbar(b_h+S_hZ)$, by Lemma~\ref{lem:compact} with $H_k=\Fbar_k$, and the right side to $\Fbar(z)+\lambda_hd$. The other two conclusions follow from \eqref{eq:P3}, (S3) and (S4): Step~2 of the proof of Theorem~\ref{thm:dil} uses only (S3) and (S4) to bound $\Delta(p_h)$, and \eqref{eq:P3} gives $\mathcal I(p_h)\ge\frac q2\cdot h\Delta_\star/(2K_S(1+T_h)^2d^2\ln(2g_{\max}/\Delta_\star))=\Delta_\star^2/(16K_S(1+T_h)^2d^2\ln(2g_{\max}/\Delta_\star))$, from which the last display of the proof gives the bound with $K_D$.
\end{proof}

The perturbation $\epsilon_k|\cdot-b^\star|$ makes the prescribed $b^\star$ the limit of the minimisers $b^\star_k$; without it the limit would be some minimiser of $\Fbar$, which would be enough for Proposition~\ref{prop:body}, since every minimiser of the penalised mean there lies in $\cJ$.

\section{One epoch of the online Newton method}\label{app:epoch}

This appendix proves Lemma~\ref{lem:ve}. Section~\ref{app:epoch-setting} states one epoch on an abstract filtration and reduces Lemma~\ref{lem:ve} to it; Sections~\ref{app:epoch-round} and~\ref{app:epoch-ub} prove the two facts on one round that the learner is built on; Section~\ref{app:epoch-events} lists the events of an epoch with their probabilities.

\subsection{One epoch}\label{app:epoch-setting}

Fix $d,n,\delta$ with $\varepsilon\le\frac12$ and the parameters of Section~\ref{sec:poly-learner}; inside an epoch $L$ stands for $L_e=36\ln(1/\delta_e)$. Let $1\le m\le n$, let $(\Omega,\mu)$ be a probability space with $\sigma$-fields $\mathcal U_0\subset\mathcal U_1\subset\cdots$, and let $\Theta$ be a measurable space carrying, for each $\theta\in\Theta$, a closed convex body $\tilde\cK(\theta)$ with $B(0,1)\subset\tilde\cK(\theta)\subset B(0,R)$ and $M(\tilde\cK(\theta))\le1.27(d+1)^{1/2}/d$, and convex functions $g_s(\theta,\cdot):\tilde\cK(\theta)\to[0,1]$ for $1\le s\le m$. For $1\le s\le m$:
\begin{enumerate}[label=(E\arabic*)]
\item $\xi_s$ is a random vector of $\R^d$, measurable in $\mathcal U_s$, standard Gaussian and independent of $\mathcal U_{s-1}$;
\item $\theta_s:\Omega\to\Theta$ is measurable in $\mathcal U_{s-1}$, and $\tilde\cK(\theta_s)=\tilde\cK(\theta_1)$;
\item $e_s^{\circ}$ is a real random variable, measurable in $\mathcal U_s$, with $\E[e_s^\circ\mid\mathcal U_{s-1},\xi_s]=0$ and $\E[\exp((e_s^\circ)^2)\mid\mathcal U_{s-1},\xi_s]\le2$.
\end{enumerate}
The \emph{epoch} is the run of the rounds of Algorithm~\ref{alg:df2o} on $\tilde\cK=\tilde\cK(\theta_1)$, from the state of the start of an epoch, with the exact gauge $\pi=\pi_{\tilde\cK}$ in step~1, the draw $\xi_s$, the loss $g_s=g_s(\theta_s,\cdot)$, and the observation $Y_s=p(X_s)(g_s(X_s/p(X_s))+e_s^\circ)+2(p(X_s)-1)/\varepsilon$ in round $s$, until its first restart or the end of round $m$. The routine AP may be any routine that meets the specification of Section~\ref{sec:poly-learner}, and the round is asked to be measurable jointly in the parameter, the state, the draw and the observation, which holds for the bodies and losses below and for the routine of the learner. With $e_s$ the extension of $g_s$ of Section~\ref{sec:poly-regret}, $\bar e_s$ the integral of $e_s(\mu_s+(1-\lambda)A_sz)$ over the standard Gaussian law of $z$, and $\mathrm{eReg}_r(x)=\sum_{s\le r}(\bar e_s-e_s(x))$, the epoch \emph{meets its bound at horizon $m$} as in Section~\ref{sec:poly-regret}: no guard acts, and either the test fires in a round $r\le m$ with $\mathrm{eReg}_r\le0$ on $\cK_\varepsilon$, or it fires in no round and $\mathrm{eReg}_m\le\rho/\eta$ on $\cK_\varepsilon$.

\begin{proposition}[one epoch]\label{prop:epoch}
In this setting the epoch fails to meet its bound at horizon $m$ with probability at most $\frac{149}{168}\delta_e$.
\end{proposition}

\begin{proof}[Proof of Lemma~\ref{lem:ve} from Proposition~\ref{prop:epoch}]
Let $f_t=L_t(\ell_t)$ be the loss process and $a\le n$. Put $m=n-a+1$ and, for the $s$-th round $t=a+s-1$ of the epoch, $\mathcal U_s=\cF_{t+1}$ ($\mathcal U_0=\cF_a$), $\xi_s$ the draw of round $t$ and $e^\circ_s=\frac12(\varepsilon_t+\zeta_t)$. The parameter $\theta_s$ is the pair of the position returned by Procedure PG and of $\ell_t$, with $\tilde\cK(\theta)=S(\cK)$ and $g_s(\theta,y)=\frac12L_t(\ell)(\mu_S+A_Sy)$; on the complement of the case of the analysis, an event of $\mathcal U_0$, the position is replaced by a fixed position of the case, which changes nothing on the case. The losses are measurable in the parameter at each interior point, and convex, so jointly measurable on the interior, where every point read lies. (E1) holds because the draws of round $t$ are independent of $\cF_t=\mathcal U_{s-1}$, and (E2) because the position is read before round $1$ and $\ell_t$ is measurable in $\cF_t$. (E3) holds: given $\cF_t$ and $\xi_s$, the noise $\varepsilon_t$ has mean $0$ by (N) and the dither $\zeta_t=g/2$ is an independent centred Gaussian, and $(e^\circ_s)^2\le\frac12(\varepsilon_t^2+\zeta_t^2)$ gives $\E[\exp((e^\circ_s)^2)\mid\mathcal U_{s-1},\xi_s]\le(\E[e^{\varepsilon^2_t}\mid\mathcal U_{s-1},\xi_s])^{1/2}\,\E e^{g^2/8}\le2^{1/2}(4/3)^{1/2}<2$. In the case of the analysis the virtual epoch started at $a$ is the epoch of this instance, so the event of Lemma~\ref{lem:ve} lies in the event of Proposition~\ref{prop:epoch}.
\end{proof}

Inside the proof of Proposition~\ref{prop:epoch} the loss of round $s$ is read only through integrals over $\xi_s$ given $\mathcal U_{s-1}$, and by the freezing lemma (the conditional expectation, given $\mathcal U_{s-1}$, of a jointly measurable function of $\mathcal U_{s-1}$-measurable data and of $\xi_s$ is the integral over $\xi_s$ with the data held fixed) each such integral is the conditional expectation with $g_s$ frozen. So every bound on one round proved for one convex loss holds almost surely for the loss of the round.

\subsection{One round}\label{app:epoch-round}

In round $s$ of an epoch write $P=P_s$, $\Sigma=\Sigma_s$, $A=A_s$, $\xi=\xi_s$, $w=A\xi$, $u=Pw$, $c=\eta\kappa\lambda\kappa_0Y_s$, $a=1-c$, $q=a+c|\xi|^2$ and $P'=aP+cuu^{\mathsf T}$, as in steps~2 and~3 of Algorithm~\ref{alg:df2o}.

\begin{lemma}[rank-one round]\label{lem:r1}
Let $P$ be positive definite and $AA^{\mathsf T}=\Sigma=P^{-1}$.
\begin{enumerate}[label=(\roman*)]
\item $u=A^{-\mathsf T}\xi$, $\Sigma u=w$, $\ip{u,w}=|\xi|^2$, and $\ip{P'x,x}=\ip{(aI+c\xi\xi^{\mathsf T})A^{-1}x,A^{-1}x}\ge\min(a,q)|A^{-1}x|^2$; so $a>0$ and $q>0$ make $P'$ positive definite; and $\det P'=a^{d-1}q\det P$.
\item If $a,q>0$, then $\Sigma'=a^{-1}(\Sigma-(c/q)ww^{\mathsf T})$ is the inverse of $P'$.
\item If $a,q>0$, then $c'=(c/q)/(1+(a/q)^{1/2})$ satisfies $c'^2|\xi|^2-2c'+c/q=0$, and $A'=a^{-1/2}(A-c'w\xi^{\mathsf T})$ is invertible with $A'A'^{\mathsf T}=\Sigma'$.
\item For $\hat g=Y_su/(1-\lambda)$, $\Sigma'\hat g=\frac{Y_s}{(1-\lambda)q}w$ and $\ip{\Sigma'\hat g,\hat g}=\kappa_0Y_s^2|\xi|^2/q$.
\item $\tr P'=a\tr P+c|u|^2$ and $\tr\Sigma'=a^{-1}(\tr\Sigma-(c/q)|w|^2)$.
\end{enumerate}
\end{lemma}

\begin{proof}
(i) $PA=A^{-\mathsf T}$, so $u=PA\xi=A^{-\mathsf T}\xi$, $\Sigma u=A\xi=w$ and $\ip{u,w}=\ip{A^{-\mathsf T}\xi,A\xi}=|\xi|^2$. Then $P=A^{-\mathsf T}A^{-1}$ and $uu^{\mathsf T}=A^{-\mathsf T}\xi\xi^{\mathsf T}A^{-1}$, so $P'=A^{-\mathsf T}(aI+c\xi\xi^{\mathsf T})A^{-1}$; the middle matrix has the eigenvalue $q$ on $\xi$ and $a$ on its orthogonal complement, so its determinant is $a^{d-1}q$ and $\det P'=a^{d-1}q\det P$. (ii) With $\Sigma u=w$, $w^{\mathsf T}P=u^{\mathsf T}$ and $w^{\mathsf T}u=|\xi|^2$,
\[
\Sigma'P'=\frac1a\Bigl(aI+cwu^{\mathsf T}-\frac{ac}qwu^{\mathsf T}-\frac{c^2|\xi|^2}qwu^{\mathsf T}\Bigr)=I+\frac c{aq}\,wu^{\mathsf T}\bigl(q-a-c|\xi|^2\bigr)=I.
\]
(iii) Since $A\xi=w$, $(A-c'w\xi^{\mathsf T})(A-c'w\xi^{\mathsf T})^{\mathsf T}=\Sigma-(2c'-c'^2|\xi|^2)ww^{\mathsf T}$, and $2c'-c'^2|\xi|^2=c/q$ is the quadratic, whose discriminant is $1-c|\xi|^2/q=a/q>0$; $c'$ is its root $(1-(a/q)^{1/2})/|\xi|^2$, written without division by $|\xi|^2$. Then $A-c'w\xi^{\mathsf T}=A(I-c'\xi\xi^{\mathsf T})$ and $1-c'|\xi|^2=(a/q)^{1/2}>0$, so $A'$ is invertible. (iv) $\Sigma'u=a^{-1}(w-(c/q)|\xi|^2w)=w/q$, and $\ip{w,u}=|\xi|^2$, $\kappa_0=(1-\lambda)^{-2}$. (v) Traces of (i) and (ii).
\end{proof}

So a round needs the four operations, two square roots and $O(d^2)$ operations, and the factor $A_s$ is a function of the transcript of the epoch. Given $\mathcal U_s$ the play $X_s=\mu_s+(1-\lambda)A_s\xi_s$ has the law $p'_s=N(\mu_s,(1-\lambda)^2\Sigma_s)$ for every factor $A_s$ of $\Sigma_s$.

\subsection{The surrogate and its estimators}\label{app:epoch-ub}

For the law $p_s=N(\mu_s,\Sigma_s)$ the surrogate of round $s$ is
\[
s_s(z)=\Bigl(1-\frac1\lambda\Bigr)\E\,e_s(X)+\frac1\lambda\,\E\,e_s\bigl((1-\lambda)X+\lambda z\bigr),\qquad X\sim p_s,
\]
the surrogate of \cite[Chapter~11]{Lat24}. It is a function of $(g_s,\mu_s,\Sigma_s)$, measurable in $\mathcal U_s$ at each point, and smooth in $z$. The estimators of round $s$ at the mean are
\[
\hat g_s=\kappa_0Y_sP_s(X_s-\mu_s)=\frac{Y_su_s}{1-\lambda},\quad \hat s''_s=\lambda\kappa_0Y_s\bigl(\kappa_0P_s(X_s-\mu_s)(X_s-\mu_s)^{\mathsf T}P_s-P_s\bigr)=\lambda\kappa_0Y_s(u_su_s^{\mathsf T}-P_s),
\]
so that $\eta\kappa\hat s''_s=c(u_su_s^{\mathsf T}-P_s)$ and $P_{s+1}=P_s+\eta\kappa\hat s''_s$: the learner is the online Newton method of \cite[Chapter~11]{Lat24} with these estimators, its Newton point $\mu_s-\eta\Sigma_{s+1}\hat g_s$ being the point $w_\mu$ of step~6 by Lemma~\ref{lem:r1}~(iv).

\begin{lemma}[the derivatives of the surrogate]\label{lem:ub}
Let $0<\lambda<1$, $\mu\in\R^d$, $A$ an invertible $d\times d$ matrix, $P=(AA^{\mathsf T})^{-1}$, and $e:\R^d\to\R$ measurable with $|e(y)|\le B(1+|y|)$ for some $B$; let $s$ be the surrogate of $e$ at $N(\mu,AA^{\mathsf T})$. Then $s$ is twice continuously differentiable, and for every $z$, with $g$ a standard Gaussian vector, $X_z=(1-\lambda)(\mu+Ag)+\lambda z$ and $m_z=(1-\lambda)\mu+\lambda z$,
\[
s'(z)=\E\bigl[\kappa_0e(X_z)P(X_z-m_z)\bigr],\qquad s''(z)=\lambda\,\E\bigl[\kappa_0e(X_z)\bigl(\kappa_0P(X_z-m_z)(X_z-m_z)^{\mathsf T}P-P\bigr)\bigr].
\]
\end{lemma}

\begin{proof}
The first term of $s$ does not depend on $z$, and $X_z$ has the law $N(m_z,(1-\lambda)^2AA^{\mathsf T})$, whose density $p'_z$ has the inverse covariance $\kappa_0P$; so $s(z)=(1-\frac1\lambda)\E\,e(\mu+Ag)+\frac1\lambda\int e(x)p'_z(x)\,dx$. The mean $m_z$ has derivative $\lambda$ in $z$, so
\[
\nabla_zp'_z(x)=\lambda\kappa_0P(x-m_z)\,p'_z(x),\qquad\nabla_z^2p'_z(x)=\lambda^2\kappa_0\bigl(\kappa_0P(x-m_z)(x-m_z)^{\mathsf T}P-P\bigr)p'_z(x).
\]
For $z$ in a bounded set these are dominated by a polynomial in $x$ times a fixed Gaussian density, and $e$ grows at most linearly, so the derivatives of order at most two pass under the integral and are continuous in $z$. Dividing by $\lambda$ gives the two formulas.
\end{proof}

At $z=\mu_s$, $X_{\mu_s}$ has the law $p'_s$ of the play. The extension $e_s$ is Lipschitz, so the lemma applies to it; with $\E[Y_s\mid\mathcal U_s,\xi_s]=e_s(X_s)$, which is (E3), and the freezing lemma, it gives $\E[\hat g_s\mid\mathcal U_s]=s'_s(\mu_s)$ and $\E[\hat s''_s\mid\mathcal U_s]=s''_s(\mu_s)$: the estimators at the mean are unbiased. The estimators of \cite[Chapter~11]{Lat24} are these integrands under $p_s$ weighted by the density ratio $p'_s/p_s$, with a clipping of the ratio; sampling from $p'_s$ itself removes the ratio, the clipping and the bias it causes.

The rest of the proof reads the extension and the surrogate through seven facts. In \cite{Lat24} they are stated with constants left implicit, as exercises, or with a last step that asks $L$ to be large; the next proposition states them with the constants of the learner, as the formal proof proves them.

\begin{proposition}[the surrogate with constants]\label{prop:sifacts}
Let $d,n\ge1$ and $0<\delta<1$ with $\varepsilon\le\frac12$, with the parameters of Section~\ref{sec:poly-learner} and $L=L_e$. Let $\tilde\cK$ be closed and convex with $B(0,1)\subset\tilde\cK\subset B(0,R)$ and $M(\tilde\cK)\le1.27(d+1)^{1/2}/d$, let $g:\tilde\cK\to[0,1]$ be convex, $p=\max(1,\pi_{\tilde\cK}/(1-\varepsilon))$ and $e=p\,g(\cdot/p)+2(p-1)/\varepsilon$.
\begin{enumerate}[label=(\roman*)]
\item $e$ is convex on $\R^d$, nonnegative and $(8/\varepsilon)$-Lipschitz, and $e=g$ on $(1-\varepsilon)\tilde\cK$.
\end{enumerate}
Let further $\mu\in(1-\varepsilon)\tilde\cK$, $A$ invertible with $AA^{\mathsf T}\le2\sigma^2I$ and $\delta_eI\le P\le\delta_e^{-1}I$, where $P=(AA^{\mathsf T})^{-1}$, and let $s$ be the surrogate of $e$ at $N(\mu,AA^{\mathsf T})$ and $W$ a standard Gaussian vector.
\begin{enumerate}[label=(\roman*),resume]
\item $s$ is twice continuously differentiable and convex, and $s\le e$.
\item $-10/(\lambda\varepsilon)\le s\le1$ on $(1-\varepsilon)\tilde\cK$.
\item $\E\,e(\mu+AW)-s(\mu)\le\frac2\lambda\tr(s''(\mu)AA^{\mathsf T})+\frac{2\delta_ed}\lambda$.
\item For all $z,v\in\R^d$, with $c_H=8\lambda/((1-\lambda)\varepsilon)$, $|\ip{s''(z)v,v}|\le c_H(2d\sigma^2)^{1/2}\ip{Pv,v}$ and $|\ip{s''(z)v,v}|\le c_H(d/\delta_e)^{1/2}|v|^2$.
\item If $\lambda^2\ip{P(x-y),x-y}\le1/L$, then $s''(x)\le2s''(y)+\delta_eP$.
\item If $\lambda^2\ip{P(y-\mu),y-\mu}\le1/L$, then $s(\mu)-s(y)\le\ip{s'(\mu),\mu-y}-\frac14\ip{s''(\mu)(y-\mu),y-\mu}+\delta_e/\lambda^2$.
\end{enumerate}
\end{proposition}

\begin{proof}
Each item is a result of \cite{Lat24}; the changes are these. (i) is \cite[Proposition~3.23]{Lat24}, stated there for the coefficient $(1-\varepsilon)/\varepsilon$ with the Lipschitz constant $2/(\varepsilon(1-\varepsilon)r)$ for a body that contains a ball of radius $r$, at most $4/\varepsilon$ here, where $r=1$ and $\varepsilon\le\frac12$, and resting on facts of the gauge and of the perspective that the source cites (\cite[Lemma~3.3]{Lat24}, parts 1 to 3 and 7); the formal proof proves these facts, convexity for every coefficient at least $(1-\varepsilon)/\varepsilon$, and the constant $8/\varepsilon$ at the coefficient $2/\varepsilon$. (ii) is \cite[Lemma~12.3]{Lat24}, with the smoothness from Lemma~\ref{lem:ub}. (iii) is \cite[Lemma~11.16]{Lat24}, an exercise there: $s\le e=g\le1$ on $(1-\varepsilon)\tilde\cK$, and $s(x)\ge-\frac1\lambda\E\,e(\mu+AW)\ge-\frac1\lambda(3+\frac4\varepsilon)\ge-\frac{10}{\lambda\varepsilon}$, since $p(\mu+w)-1\le2\pi(w)$ and $\E\pi(AW)\le M(\tilde\cK)(d\|AA^{\mathsf T}\|)^{1/2}\le1$ by \cite[Lemma~10.5]{Lat24}. (iv) is \cite[Proposition~12.10]{Lat24}, whose last step asks $L$ large: with $\Prob(|W|^2\ge t)\le2e^{-3t/(8d)}$ the error term is at most $\delta_ed/\lambda$ once $L\ge12.45\ln(1/\delta_e)$, and $L=36\ln(1/\delta_e)$. (v): the second bound is \cite[Proposition~12.6]{Lat24}, part 1, with $\|P\|\le1/\delta_e$; the first is its part 2, an exercise there, proved by Gaussian integration by parts in the direction of $v$, which gives $|\ip{s''(z)v,v}|\le\frac{\lambda}{1-\lambda}\frac8\varepsilon(2\|AA^{\mathsf T}\|/\pi)^{1/2}\ip{Pv,v}$. (vi) is \cite[Proposition~12.5]{Lat24}, whose proof fixes the smoothing radius $e^{-L/200}$ and then needs $L\ge234\ln(1/\delta_e)$; the radius that balances the two error terms needs only $L\ge9.82\ln(1/\delta_e)$. (vii) follows from (vi) as \cite[Proposition~12.9]{Lat24} does.
\end{proof}

\subsection{The events of an epoch}\label{app:epoch-events}

Let $x=\ln(14n/\delta_e)$. The proof of Proposition~\ref{prop:epoch} runs on the intersection of five events, each written for the rounds of the epoch up to the stopping time $\tau$ at which the first of four conditions fails: the minimiser of the sum of the surrogates over $\cK_\varepsilon$ lies in the lower region of the next round; $a_s>0$ and $q_s>0$; $\delta_eI\le\frac12\bar P_{s+1}\le P_{s+1}\le2\bar P_{s+1}\le\delta_e^{-1}I$, where $\bar P_1=P_1$ and $\bar P_{s+1}=\bar P_s+\eta\kappa s''_s(\mu_s)$; and the epoch does not restart at the end of round $s$.
\begin{lemma}[the events of an epoch]\label{lem:events}
In the setting of Section~\ref{app:epoch-setting}, the event (a) fails with probability at most $\frac37\delta_e$, and each of the events (b) to (e) fails while (a) holds with at most the probability written with it.
\begin{enumerate}[label=(\alph*)]
\item Noise, gauge and Gaussian tails: in every round $s\le\tau$, $|e^\circ_s|\le x^{1/2}$, $\pi(X_s)\le L^{1/2}$ and $|\xi_s|^2\le\frac83dx$.
\item The second moment of the observations: $\sum_{s\le\tau}Y_s^2\le22n$; probability $\frac17\delta_e$.
\item The linear part of the estimated surrogates: $\bigl|\sum_{u\le t}\ip{\hat g_u-s'_u(\mu_u),y-\mu_u}\bigr|\le b_1$ for all $t\le\tau$ and all $y\in B(0,R)$ in the ellipsoids of (e) below, with $b_1=1+\lambda^{-1}\bigl(\frac{11}{20}(10dnL)^{1/2}+\frac{13}{100}d^{3/2}L^2/\varepsilon\bigr)$; probability $\frac1{21}\delta_e$.
\item The matrices: $|\ip{(P_{t+1}-\bar P_{t+1})v,v}|\le4\kappa\eta B_n\ip{\bar P_tv,v}$ for all $t\le\tau$ and $v\in\R^d$, with $B_n=\lambda L^2\bigl(1+(10dn)^{1/2}+d^2L/\varepsilon\bigr)$; probability $\frac17\delta_e$.
\item The quadratic part: $\bigl|\sum_{u\le t}\ip{(\hat s''_u-s''_u(\mu_u))(y-\mu_u),y-\mu_u}\bigr|\le B_n/(L\lambda^2)$ for all $t\le\tau$ and all $y\in B(0,R)$ in the ellipsoids $\{\lambda^2\ip{P_u(y-\mu_u),y-\mu_u}\le1/(2L)\}$, $u\le t$; probability $\frac18\delta_e$.
\end{enumerate}
\end{lemma}

By the union bound the epoch leaves the intersection of the five events with probability at most $(\frac37+\frac17+\frac1{21}+\frac17+\frac18)\delta_e=\frac{149}{168}\delta_e$. The proof of the lemma uses four tools, each proved in the formal proof.

\paragraph{Tails.} For a standard Gaussian vector $W$ of $\R^d$ and $t\ge0$, $\Prob(|W|^2\ge t)\le2e^{-3t/(8d)}$, since $\E\exp(3|W|^2/(8d))=(1-3/(4d))^{-d/2}\le2$; a real variable $r$ with $\E e^{r^2}\le2$ has $\Prob(|r|\ge b)\le2e^{-b^2}$ for $b\ge0$; and a $k$-Lipschitz function $F$ of $m+SW$ satisfies $\Prob(F\ge\E F+t)\le\exp(-t^2/(2k^2\|S\|^2))$, which is Gaussian concentration with the constant $\frac12$ \cite[Theorem~1A and Corollary~1]{TIS76} (\cite[Theorem~B.16]{Lat24} states a two-sided form, with $\ln(2/\delta)$ in place of $\ln(1/\delta)$); the formal proof uses this constant here, not the constant $2/\pi^2$ of Theorem~\ref{thm:gauss}.

\paragraph{Lemma TR.} For $Y,Z\ge0$ square integrable and $T\ge0$, $\E[YZ]\le T\,\E Y+(\E Y^2)^{1/2}(\E[Z^21\{Z>T\}])^{1/2}$: split on $\{Z\le T\}$ and apply the inequality of Cauchy and Schwarz on the rest. For $z$ standard Gaussian in $\R^d$, a vector $b$ and $g=\ip{z,b}/|b|$, with $\E[g^{2k}1\{|g|>a\}]\le((4k-1)!!)^{1/2}e^{-a^2/4}$, this gives
\begin{gather*}
\E[Y\ip{z,b}^2]\le|b|^2\bigl(T\,\E Y+(\E Y^2)^{1/2}(105^{1/2}e^{-T/4})^{1/2}\bigr),\\
\E\bigl[Y(\ip{z,b}^2-|b|^2)^2\bigr]\le|b|^4\bigl(T\,\E Y+(\E Y^2)^{1/2}(2027025^{1/2}e^{-T^{1/2}/4})^{1/2}\bigr),
\end{gather*}
the second for $T\ge1$.

\paragraph{Freedman's inequality.} \cite[Theorem~B.20]{Lat24} at a stopping time $\tau'\le N$: if $\E[X_t\mid\cF_{t-1}]=0$ for $t\le\tau'$, then with probability at least $1-\delta'$, $|\sum_{t\le\tau'}X_t|\le3(V\ell)^{1/2}+2B\ell$, with $V$ the predictable variance up to $\tau'$, $B=\max(1,\max_{t\le\tau'}|X_t|)$ and $\ell=\ln(4\max(B,V^{1/2})/\delta')$; the source prints $2$ in place of $4$ in $\ell$, and its proof gives $4$. Its one-sided companion \cite[Theorem~B.22]{Lat24}: if $0\le Z_t\le B$, then with probability at least $1-\delta'$, $\sum_{t\le T}Z_t<2\sum_{t\le T}\E[Z_t\mid\cF_{t-1}]+B\ln(1/\delta')$; here it is proved for every $T\le N$ at once, the union over the horizons being free at the first horizon where the bound fails.

\paragraph{Lemma U.} Let $\mathcal C$ be a finite set, and for $c\in\mathcal C$ and $s\le m$ let
\[
Z_{c,s}=A_{c,s}(\varsigma_s,\xi_s)+B_{c,s}(\varsigma_s,\xi_s)\,e^\circ_s-\E_z A_{c,s}(\varsigma_s,z),
\]
with $\varsigma_s$ measurable in $\mathcal U_{s-1}$, $z$ standard Gaussian, and $A_{c,s},B_{c,s}$ measurable and bounded by a constant times a power of $1+|\xi|$. Let $\kappa'>0$, $0<\delta'\le1$, and let $V$, $R'$ be any functions of the path, of $c$ and of $t$. Then with probability at least $1-\delta'$: for every $c\in\mathcal C$ and $t\le m$ such that $V\ge0$, $\kappa'R'\ge1$, $\sum_{s\le t}\E_z(A_{c,s}^2+B_{c,s}^2)(\varsigma_s,z)\le V$ and $|Z_{c,s}|\le R'$ for $s\le t$,
\[
\Bigl|\sum_{s\le t}Z_{c,s}\Bigr|\le3(V\ell)^{1/2}+2R'\ell,\qquad\ell=\ln\frac{4n|\mathcal C|\max(\kappa'R',\kappa'V^{1/2})}{\delta'}.
\]
Proof: by (E1), (E3) and the freezing lemma, $\E[Z_{c,s}\mid\mathcal U_{s-1}]=0$ and $\E[Z_{c,s}^2\mid\mathcal U_{s-1}]\le\E_z(A_{c,s}^2+B_{c,s}^2)(\varsigma_s,z)$, since $\E[(e^\circ_s)^2\mid\mathcal U_{s-1},\xi_s]\le1$; Freedman's inequality applied to $\kappa'Z_{c,\cdot}$, for each $c$ and each horizon, with the union over the $n|\mathcal C|$ pairs, gives the bound with the predictable variance and the range in place of $V$ and $R'$, and the bound grows with both.

\begin{proof}[Proof of Lemma~\ref{lem:events}]
Write $\ell_e=\ln(1/\delta_e)=L/36$. By the definition of $\delta_1$ and $\delta_e$, $\ell_e\ge14+3\ln(dn)$, so $x\le\frac43\ell_e$; Lemma~\ref{lem:rows} gives also $n\delta_e\le\lambda^2$. Before $\tau$ the state is good (the paragraph before Lemma~\ref{lem:closure}): $\mu_s\in\cK_\varepsilon$, $\Sigma_s\le2\sigma^2I$ and $\delta_eI\le P_s\le\delta_e^{-1}I$.

(a) Fix a round $s\le m$. By (E3) and the tower property $\E\exp((e^\circ_s)^2)\le2$, so $\Prob(|e^\circ_s|>x^{1/2})\le2e^{-x}=\delta_e/(7n)$; and $\Prob(|\xi_s|^2>\frac83dx)\le2e^{-x}=\delta_e/(7n)$. Neither reads $\tau$. The gauge is read at $s\le\tau$ only, where the state is good. On a good state $X_s=\mu_s+(1-\lambda)A_s\xi_s$ with $\pi(\mu_s)\le1$ and $\|(1-\lambda)A_s\|\le(2\sigma^2)^{1/2}=(4L)^{-1/2}$. By \cite[Lemma~10.5]{Lat24}, $\E\pi((1-\lambda)A_s\xi)\le(4L)^{-1/2}M(\tilde\cK)\,d^{1/2}\le1$, using $M(\tilde\cK)\le1.27(d+1)^{1/2}/d$, so $\E\pi(X_s)\le2$; the gauge is $1$-Lipschitz since $B(0,1)\subset\tilde\cK$, and Gaussian concentration gives $\Prob(\pi(X_s)>L^{1/2})\le\exp(-2L(L^{1/2}-2)^2)\le\delta_e/(7n)$. The event that the state before round $s$ is good is measurable in $\mathcal U_{s-1}$, and $\xi_s$ is independent of $\mathcal U_{s-1}$, so by the freezing lemma the bound holds given $\mathcal U_{s-1}$ on that event. The union over the three tails and the rounds $s\le m\le n$ costs $3n\cdot\delta_e/(7n)$.

(b) Let $p_u=p(X_u)$, $v_u=p_u-1$, $B=4L^2/\varepsilon^2$ and $Z_u=1\{\pi(X_u)\le L^{1/2}\}\min(Y_u^2,B)$ for $u\le m$, so that $0\le Z_u\le B$. Since $Y_u=e_u(X_u)+p_ue^\circ_u$ with $g_u\in[0,1]$ and $\E[(e^\circ_u)^2\mid\mathcal U_{u-1},\xi_u]\le1$, $\E[Z_u\mid\mathcal U_{u-1},\xi_u]\le4p_u^2+8v_u^2/\varepsilon^2\le8+10v_u^2/\varepsilon^2$, and by the freezing lemma
\[
\E[Z_u\mid\mathcal U_{u-1}]\le W_u=8+\frac{10}{\varepsilon^2}\,\E_z v\bigl(\mu_u+(1-\lambda)A_uz\bigr)^2,\qquad v=p-1 .
\]
By \cite[Theorem~B.22]{Lat24} uniform in the horizon, with probability at least $1-\delta_e/7$, $\sum_{u\le t}Z_u<2\sum_{u\le t}W_u+B\ln(7/\delta_e)$ for every $t\le m$. On (a), for $u\le\tau$: $\pi(X_u)\le L^{1/2}$, and $|Y_u|\le2L/\varepsilon$ by the tails of (a), so $Z_u=Y_u^2$. The sum of the $W_u$ is bounded on every path: for every $t\le\tau$,
\begin{equation}\label{eq:vsum}
\sum_{s\le t}\E_z v\bigl(\mu_s+(1-\lambda)A_sz\bigr)^2\le\frac{n\varepsilon^2}4 .
\end{equation}
For $s<\tau$: $\E v^2\le2L\,\E v+e^{-L}$, by the truncation of $v$ at $2L$ and the Gaussian tail of the gauge computed in (a); $\E v\le\varepsilon(\bar e_s-e_s(\mu_s))$, since $e_s=h_s+v/\varepsilon$ with $h_s=p\,g_s(\cdot/p)+(p-1)/\varepsilon$ convex, $v(\mu_s)=0$ and Jensen's inequality; and $\bar e_s-e_s(\mu_s)\le\bar e_s-s_s(\mu_s)\le\frac2\lambda\tr(s''_s(\mu_s)\Sigma_s)+\frac{2\delta_ed}\lambda$, as in (W1) below. Lemma~\ref{lem:ld} over the rounds $u<\tau$, with $\bar P_\tau\le I/(2\delta_e)$ from the third condition of the stopping time in round $\tau-1$, bounds $\frac2\lambda\sum_{u<\tau}\tr(s''_u(\mu_u)\Sigma_u)$ by $dL/(\lambda\eta)$, and $2L\varepsilon\,dL/(\lambda\eta)=n\varepsilon^2/10$ by the definitions of $\lambda,\eta,\varepsilon$. For $s=\tau$ the crude bound $\E v^2\le4\|(1-\lambda)A_\tau\|^2d\le d/L$ suffices. The rows of Lemma~\ref{lem:rows} close the sum at $n\varepsilon^2/4$. Hence on (a), $\sum_{u\le\tau}W_u\le8n+\frac{10}{\varepsilon^2}\frac{n\varepsilon^2}4=\frac{21}2n$, and $B\ln(7/\delta_e)\le n$ by Lemma~\ref{lem:rows}; so $\sum_{u\le\tau}Y_u^2>22n$ lies in the exceptional event. The horizon $\tau$ depends on the path, which the uniformity over horizons absorbs. \cite[Lemma~10.13]{Lat24} has $21n$ for the play $p$; here the play is $p'$ and the sum \eqref{eq:vsum} replaces \cite[Lemma~10.12]{Lat24}, whose second part bounds it in the source through \cite[Lemma~10.6]{Lat24}.

(c) Let $\mathcal C$ be a set of points of $B(0,R)$ with every point of $B(0,R)$ within $\delta_e^4$ of one of them and $|\mathcal C|\le(1+2R/\delta_e^4)^d$. For $\bar y\in\mathcal C$ take $Z_{\bar y,s}=\ip{\hat g_s-s'_s(\mu_s),\bar y-\mu_s}$, with $A_{\bar y,s}=e_s(X_s)\ip{P_sA_s\xi_s,\bar y-\mu_s}/(1-\lambda)$ and $B_{\bar y,s}=p(X_s)\ip{P_sA_s\xi_s,\bar y-\mu_s}/(1-\lambda)$, both multiplied by the indicator that the state is good, which is $1$ for $s\le\tau$. The mean of $A_{\bar y,s}$ is $\ip{s'_s(\mu_s),\bar y-\mu_s}$ by Lemma~\ref{lem:ub}. Let $\mathrm y(z)=e_s(X(z))^2+p(X(z))^2$ with $X(z)=\mu_s+(1-\lambda)A_sz$; then $\E_z(A^2+B^2)=\kappa_0\E_z[\mathrm y\,\ip{z,A_s^{\mathsf T}P_s(\bar y-\mu_s)}^2]$ and $|A_s^{\mathsf T}P_sh|^2=\ip{P_sh,h}$. Three facts hold on every path for $s\le t\le\tau$: $\sum_{s\le t}\E_z\mathrm y\le9n$, from $\E_z\mathrm y\le6+\frac{10}{\varepsilon^2}\E_zv^2$ and \eqref{eq:vsum}; $\E_z\mathrm y^2\le72+25600(R^4+3d^2)/\varepsilon^4$; and, for $\bar y$ within $\delta_e^4$ of a point $y$ of the ellipsoids of (e), $\ip{P_s(\bar y-\mu_s),\bar y-\mu_s}\le q_L=\frac{1+\delta_e}{2L\lambda^2}+2\delta_e^6$. Lemma~TR at $T=8\ell_e$ bounds the variance by $V_L=\kappa_0q_L(8\ell_e\cdot9n+n\,t_L)$ with $t_L\le1$ the tail term, and $V_L\le\frac{11}{10}n/\lambda^2$. On (a), $|Y_s|\le L/(2\varepsilon)$, $|\xi_s|^2\le\frac{32}9d\ell_e$ and $|s'_s|\le8/\varepsilon$, so $|Z_{\bar y,s}|\le R_L=(\kappa_0q_L\frac{32}9d\ell_e)^{1/2}\frac L{2\varepsilon}+\frac{16R}\varepsilon$. Lemma~U with $\kappa'=1/R_L$ and $\delta'=\delta_e/21$, whose logarithm is at most $7d\ell_e$, bounds every sum at a point of the net. For $y$ within $\delta_e^4$ of $\bar y$ the two sums differ by at most $nG\delta_e^4$, where $G$ bounds $|\hat g_s-s'_s(\mu_s)|$ on (a). The arithmetic $3(V_L\ell)^{1/2}+2R_L\ell+nG\delta_e^4\le b_1$ closes (c).

(e) The same net and the increments $Z_{\bar y,s}=\ip{(\hat s''_s-s''_s(\mu_s))h,h}$ with $h=\bar y-\mu_s$: by $\hat s''_s=\lambda\kappa_0Y_s(u_su_s^{\mathsf T}-P_s)$ they are of the form of Lemma~U with the direction $\lambda\kappa_0(\ip{P_sA_s\xi_s,h}^2-\ip{P_sh,h})$, whose mean against $e_s$ is $\ip{s''_s(\mu_s)h,h}$ by Lemma~\ref{lem:ub}. The second display of Lemma~TR at $T^{1/2}=16\ell_e$ bounds the variance by $V_D=(\lambda\kappa_0)^2q_L^2((16\ell_e)^2\cdot9n+n\,t_D)\le n/\lambda^2$; on (a) the range is at most $R_D=\frac L{2\varepsilon}\lambda\kappa_0q_L(\frac{32}9d\ell_e+1)+\frac{8\lambda}{(1-\lambda)\varepsilon}(2d\sigma^2)^{1/2}q_L$, the second term by Proposition~\ref{prop:sifacts}~(v); and with $|\ip{(\hat s''_s-s''_s(\mu_s))w,w}|\le K_D|w|^2$ on (a), the sums at $y$ and at $\bar y$ differ by at most $nK_D\delta_e^4\cdot4R$. Lemma~U with $\kappa'=1/R_D$ and $\delta'=\delta_e/8$, and the arithmetic $3(V_D\ell)^{1/2}+2R_D\ell+4nRK_D\delta_e^4\le B_n/(L\lambda^2)$, close (e).

(d) For $t\le\tau$, $P_{t+1}-\bar P_{t+1}=\eta\kappa\mathcal A_t$ with $\mathcal A_t=\sum_{u\le t}(\hat s''_u-s''_u(\mu_u))$. Let $q_t(v)=2\ip{\bar P_tv,v}$; for $u\le t$, $\ip{P_uv,v}\le q_t(v)$ (since $P_u\le2\bar P_u\le2\bar P_t$) and $4\delta_e|v|^2\le q_t(v)\le|v|^2/\delta_e$. On (a), $|\ip{\hat s''_uv,v}|+|\ip{s''_u(\mu_u)v,v}|\le R_c\ip{P_uv,v}$ with $R_c=\lambda\kappa_0\frac L{2\varepsilon}(\frac{32}9d\ell_e+1)+\frac{8\lambda}{(1-\lambda)\varepsilon}(2d\sigma^2)^{1/2}$. Take a net $\mathcal C$ of the unit ball at scale $\delta_e^4$, $|\mathcal C|\le(1+2/\delta_e^4)^d$, the increments $\ip{(\hat s''_s-s''_s(\mu_s))\bar v,\bar v}$ for $\bar v\in\mathcal C$, and Lemma~U with $\kappa'=\delta_e^{-2}$, $\delta'=\delta_e/7$, the variance level $(\lambda\kappa_0)^2q_t(\bar v)^2X_v$ with $X_v=(16\ell_e)^2\cdot9n+n\,t_D$ (Lemma~TR as in (e)) and the range level $R_cq_t(\bar v)+\delta_e^2$; both levels depend on the path, which Lemma~U allows. If (d) fails at $t\le\tau$, it fails at a unit vector $v$: $|\ip{\mathcal A_tv,v}|>2B_nq_t(v)$. The point $\bar v$ of the net next to $v$ has $\ip{\mathcal A_t\bar v,\bar v}$ and $q_t(\bar v)$ within explicit multiples of $\delta_e^4$ of those at $v$, and with $3\lambda\kappa_0(X_v\ell)^{1/2}\le B_n/2$ and $2R_c\ell\le B_n/2$ the sum at $\bar v$ exceeds the level of Lemma~U, an event of probability at most $\delta_e/7$.

In the setting of Section~\ref{app:epoch-setting} the loss of round $s$ is read at a parameter measurable in $\mathcal U_{s-1}$; every step above conditions on $\mathcal U_{s-1}$ and integrates over $\xi_s$ alone, so the freezing lemma applies at each round with the loss of that round frozen. The source's concentration for the estimated surrogates (\cite[Propositions~12.22, 12.23, 12.24, 12.27 and~12.30]{Lat24}) is replaced by Lemma~U on fixed nets and Lemma~TR; the moments of the estimators are read under the play $p'$, with no density ratio (Section~\ref{app:epoch-ub}).
\end{proof}

\subsection{The region}\label{app:epoch-region}

Write $|v|_M^2=\ip{Mv,v}$. In an epoch let $\Phi_0(y)=|y|^2/(2\sigma^2)$ and let $\Phi_t$ be the quadratic $\Phi'$ of round $t$ of Algorithm~\ref{alg:df2o}. The updates of step~3 give $\Phi_t-\Phi_{t-1}=\frac{\eta\kappa}2|\cdot-\mu_t|^2_{\hat s''_t}$, so
\[
\Phi_t=\Phi_0+\hat\Omega_t,\qquad\hat\Omega_t(y)=\frac{\eta\kappa}2\sum_{u\le t}|y-\mu_u|^2_{\hat s''_u},\qquad\Omega_t(y)=\frac{\eta\kappa}2\sum_{u\le t}|y-\mu_u|^2_{s''_u(\mu_u)},
\]
and $\Phi_t$ has the Hessian $P_{t+1}$ and $\Phi_t(\mu_t)=\Phi_{t-1}(\mu_t)$. The regions are $\mathcal R_t=\{y\in\cK_\varepsilon:\Phi_{t-1}(y)\le c_1\rho\}$, on which step~7 of round $t-1$ projects, and $\mathcal R^-_t=\{y\in\cK_\varepsilon:\Phi_{t-1}(y)\le(c_1-\gamma)\rho\}$, the set of the test of round $t$. Let $S_t=\{y:\lambda^2|y-\mu_u|^2_{P_u}\le1/(2L)\text{ for }u\le t\}$, $\mathrm{pri}=R^2/(2\sigma^2)\ge\sup_{\cK_\varepsilon}\Phi_0$ and $\mathrm{dev}'=\eta\kappa B_n/(2L\lambda^2)$. On the event (e), $|\hat\Omega_t-\Omega_t|\le\mathrm{dev}'$ on $S_t\cap B(0,R)$ for $t\le\tau$; the terms of $\Omega_t$ are nonnegative, since $s''_u$ is positive semidefinite, the surrogate being convex.

\begin{lemma}[midpoints]\label{lem:dfi}
For a quadratic $\Phi$ with Hessian $P$, $\Phi(y)+\Phi(y')-2\Phi(\frac{y+y'}2)=\frac14|y-y'|^2_P$. So $\Phi(y)\le c$, $\Phi(y')\le c'$ and $\Phi(\frac{y+y'}2)\ge-e$ give $|y-y'|^2_P\le4(c+c'+2e)$.
\end{lemma}

\begin{proof}
The linear and constant parts cancel, and $\frac12|y|^2_P+\frac12|y'|^2_P-|\frac{y+y'}2|^2_P=\frac14|y-y'|^2_P$.
\end{proof}

\begin{lemma}[monotonicity]\label{lem:mon}
Let the sets $S_t$ decrease in $t$, let $\Omega_u\le\Omega_t$ for $u\le t$, and let $|\hat\Omega_t-\Omega_t|\le\mathrm{dev}'$ on $S_t\cap B(0,R)$ for $t\le\tau$. Then $\Phi_0+\hat\Omega_u\le\Phi_0+\hat\Omega_t+2\,\mathrm{dev}'$ on $S_t\cap B(0,R)$ for $u\le t\le\tau$.
\end{lemma}

\begin{proof}
Let $y\in S_t\cap B(0,R)$; it lies in $S_u$. Then $\hat\Omega_u(y)\le\Omega_u(y)+\mathrm{dev}'\le\Omega_t(y)+\mathrm{dev}'\le\hat\Omega_t(y)+2\,\mathrm{dev}'$.
\end{proof}

\begin{proposition}[the region]\label{prop:inc2}
On the event (e), if $8(c_1\rho+2\,\mathrm{dev}')<1/(2L\lambda^2)$ and $\mathrm{pri}\le\rho$, then for every $u$ with $1\le u\le\tau$: $(a_u)$ $\mu_u\in\mathcal R_u$; $(b_u)$ every $y\in\mathcal R_u$ has $|y-\mu_s|^2_{P_s}\le8(c_1\rho+2\,\mathrm{dev}')$ for every $s\le u$. In particular $\mathcal R_u\subset S_u\cap B(0,R)$.
\end{proposition}

\begin{proof}
By induction on $u$. For $u=1$, $\mu_1=0\in\cK_\varepsilon=\mathcal R_1$, and $|y|^2_{P_1}=|y|^2/\sigma^2\le2\,\mathrm{pri}\le8c_1\rho$ on $\cK_\varepsilon$. Let $2\le u\le\tau$ and assume $(a_s)$, $(b_s)$ for $s<u$. No restart happened before round $u$ and $P_u$ is positive definite, so $\mathcal R_u$ is closed and convex.

$(a_u)$: $\mu_{u-1}\in\cK_\varepsilon$ and $\Phi_{u-1}(\mu_{u-1})=\Phi_{u-2}(\mu_{u-1})\le c_1\rho$ by $(a_{u-1})$, so $\mathcal R_u$ is not empty and the projection $\mu_u$ lies in it.

$(b_u)$ for $s\le u-1$: let $y\in\mathcal R_u$ and $x_0=\mu_{u-1}$, which lies in $\mathcal R_u$ and, by $(b_{u-1})$, in $S_{u-1}$. The segment $[x_0,y]$ lies in $\mathcal R_u\subset\cK_\varepsilon\subset B(0,R)$. Let $y'$ be its last point, from $x_0$, in the closed convex set $S_{u-1}$. For $s=1$, $|y'-\mu_1|^2_{P_1}\le2\,\mathrm{pri}$. For $2\le s\le u-1$, Lemma~\ref{lem:mon} gives $\Phi_{s-1}(y')\le\Phi_{u-1}(y')+2\,\mathrm{dev}'\le c_1\rho+2\,\mathrm{dev}'$, and $\Phi_{s-1}(\mu_s)\le c_1\rho$ by $(a_s)$. The midpoint of $y'$ and $\mu_s$ lies in $S_{s-1}\cap B(0,R)$ ($\mu_s\in\mathcal R_s\subset S_s$ by $(b_s)$), where $\Phi_{s-1}\ge\Phi_0+\Omega_{s-1}-\mathrm{dev}'\ge-\mathrm{dev}'$; Lemma~\ref{lem:dfi} with the Hessian $P_s$ gives $|y'-\mu_s|^2_{P_s}\le8(c_1\rho+2\,\mathrm{dev}')<1/(2L\lambda^2)$. So $y'$ lies in the interior of every ellipsoid of $S_{u-1}$; if $y'\ne y$, the points of the segment just after $y'$ would still lie in $S_{u-1}$, against the choice of $y'$. So $y'=y$.

$(b_u)$ for $s=u$: $\Phi_{u-1}\le c_1\rho$ at $y$ and at $\mu_u$; their midpoint lies in $\mathcal R_u\subset S_{u-1}\cap B(0,R)$, where $\Phi_{u-1}\ge-\mathrm{dev}'$, and Lemma~\ref{lem:dfi} with the Hessian $P_u$ gives $|y-\mu_u|^2_{P_u}\le8(c_1\rho+\mathrm{dev}')$.
\end{proof}

The event (e) quantifies over all points of $B(0,R)$ in the past ellipsoids, so the points $y'$ and the midpoints, which depend on the whole run, need no measurability. The condition of the proposition is the row $16c_1\rho L\lambda^2+4\eta B_n<1$, since $32\,\mathrm{dev}'L\lambda^2=4\eta B_n$ at $\kappa=\frac14$.

\subsection{The test}\label{app:epoch-test}

For $t\le\tau$ and $y\in\R^d$ let $\mathrm{sReg}_t(y)=\sum_{u\le t}(s_u(\mu_u)-s_u(y))$, $\mathrm{Lin}_t(y)=\sum_{u\le t}\ip{s'_u(\mu_u),\mu_u-y}$, $W_t(y)=\sum_{u\le t}|y-\mu_u|^2_{s''_u(\mu_u)}$ and $T_t=\mathrm{Lin}_t-\frac14W_t$; their estimates are $\widehat{\mathrm{Lin}}_t(y)=\sum_{u\le t}\ip{\hat g_u,\mu_u-y}=l_t-\ip{G_t,y}$, with $l_t,G_t$ the values of $l',G'$ in round $t$, and $\hat T_t=\widehat{\mathrm{Lin}}_t-\frac14\sum_{u\le t}|y-\mu_u|^2_{\hat s''_u}$. Put $\nu=n\delta_e/\lambda^2$.

\begin{lemma}[the sandwich]\label{lem:sandwich}
For $t\le\tau$ and $y\in S_t$: (i) $\mathrm{sReg}_t(y)\le T_t(y)+\nu$; (ii) $T_t(y)\le\mathrm{sReg}_t(y)+\frac34W_t(y)+\nu$.
\end{lemma}

\begin{proof}
For $u\le t$, $\lambda^2|y-\mu_u|^2_{P_u}\le1/(2L)$. (i) By Proposition~\ref{prop:sifacts}~(vii), $\lambda^2|y-\mu|^2_P\le1/L$ gives $s(\mu)-s(y)\le\ip{s'(\mu),\mu-y}-\frac14|y-\mu|^2_{s''(\mu)}+\delta_e/\lambda^2$; sum over $u\le t$. (ii) By Taylor's formula there is $\zeta_u\in[\mu_u,y]$ with $s_u(y)=s_u(\mu_u)+\ip{s'_u(\mu_u),y-\mu_u}+\frac12|y-\mu_u|^2_{s''_u(\zeta_u)}$, and $\zeta_u$ is as close to $\mu_u$ as $y$ is, so Proposition~\ref{prop:sifacts}~(vi) gives $s''_u(\zeta_u)\le2s''_u(\mu_u)+\delta_eP_u$. Hence $s_u(\mu_u)-s_u(y)\ge\ip{s'_u(\mu_u),\mu_u-y}-|y-\mu_u|^2_{s''_u(\mu_u)}-\frac{\delta_e}2|y-\mu_u|^2_{P_u}$, and the sum over $u\le t$ is $\mathrm{sReg}_t(y)\ge\mathrm{Lin}_t(y)-W_t(y)-n\delta_e/(4L\lambda^2)\ge T_t(y)-\frac34W_t(y)-\nu$.
\end{proof}

\begin{lemma}[the dichotomy]\label{lem:d2}
Let $Q$ be a set and $\mathrm{sReg},T,\hat T,W,\Xi$ real functions on $Q$. Let $\kappa<a<c_{\mathrm{up}}$, $r(a-\kappa)=c_{\mathrm{up}}-a$, $\eta>0$, $h>0$, $e\ge0$ and $\nu,\mathrm{dev},\rho'',\pi_0$ real, and suppose that on $Q$: (i) $\mathrm{sReg}\le T+\nu$; (ii) $T\le\mathrm{sReg}+\frac{c_{\mathrm{up}}-a}2W+\nu$; (iii) $\eta T\le\rho''-\frac{a-\kappa}2\eta W$; (iv) $|\hat T-T|\le\mathrm{dev}$; (v) $\eta\hat T-\pi_0\le\Xi\le\eta\hat T$. Let $m$ be a real number with $\Xi\le m+e$ on $Q$ and $m\le b$ for every upper bound $b$ of $\Xi$ on $Q$. Then:
\begin{enumerate}[label=(\alph*)]
\item if $m\le-h$, then $\eta\,\mathrm{sReg}\le-(h-e)+\pi_0+\eta\,\mathrm{dev}+\eta\nu$ on $Q$;
\item if $\eta\,\mathrm{sReg}\le-\bigl((1+r)h+r\rho''+(1+r)\eta\,\mathrm{dev}+\eta\nu\bigr)$ on $Q$, then $m\le-h$.
\end{enumerate}
\end{lemma}

\begin{proof}
(a) On $Q$, $\eta\,\mathrm{sReg}\le\eta T+\eta\nu\le\eta\hat T+\eta\,\mathrm{dev}+\eta\nu\le\Xi+\pi_0+\eta\,\mathrm{dev}+\eta\nu\le m+e+\pi_0+\eta\,\mathrm{dev}+\eta\nu$. (b) Let $k_c$ be the bound of the hypothesis, $w=r(\rho''+h+\eta\,\mathrm{dev})$, and $y\in Q$. If $\frac{c_{\mathrm{up}}-a}2\eta W(y)\le w$, then by (ii) $\eta T(y)\le\eta\,\mathrm{sReg}(y)+w+\eta\nu\le-k_c+w+\eta\nu=-h-\eta\,\mathrm{dev}$. Otherwise $\frac{a-\kappa}2\eta W(y)>w/r$, and by (iii) $\eta T(y)<\rho''-w/r=-h-\eta\,\mathrm{dev}$. In both cases $\Xi(y)\le\eta\hat T(y)\le\eta T(y)+\eta\,\mathrm{dev}\le-h$, so $-h$ is an upper bound of $\Xi$ on $Q$ and $m\le-h$.
\end{proof}

The mismatch $\frac{c_{\mathrm{up}}-a}2W$ between $T$ and $\mathrm{sReg}$ is never bounded: where it is large, (iii) puts $T$ below the threshold, and what pays is the curvature $a-\kappa$ that the learner leaves unused. The learner takes $a=\frac12$, $c_{\mathrm{up}}=2$ and $\kappa=\frac14$, so $r=6$.

\paragraph{The test is a projection.} At $\kappa=\frac14$, $2\Phi_t=2\Phi_0+\frac\eta4\sum_{u\le t}|\cdot-\mu_u|^2_{\hat s''_u}$, so the statistic of step~5 is
\[
\Xi_t=\eta\,\widehat{\mathrm{Lin}}_t-2\Phi_t=\eta\hat T_t-2\Phi_0 .
\]
It is a quadratic with Hessian $-2P_{t+1}$ on every path: $\Xi_t(y)=\Xi_t(z_t)-|y-z_t|^2_{P_{t+1}}$ with $z_t=\Sigma_{t+1}(b_t-\frac\eta2G_t)$, the point $z$ of step~4. Its maximum over a closed convex set is attained at the projection of $z_t$ in the norm of $P_{t+1}$, and a point at which $\frac12|\cdot-z_t|^2_{P_{t+1}}$ is within $e$ of its least value has $\Xi_t$ within $2e$ of the maximum. Step~5 calls AP on $\mathcal R^-_t$ with the accuracy $\min(1,\alpha\rho/2)$, and $\mathcal R^-_t\subset\cK_\varepsilon\subset B(0,R)$; so in a round where AP does not fail, the number $m=\Xi_t(p^-)$ is as in Lemma~\ref{lem:d2} with $Q=\mathcal R^-_t$ and $e=\alpha\rho$. On $\cK_\varepsilon$, $0\le\eta\hat T_t-\Xi_t=2\Phi_0\le2\,\mathrm{pri}$: this is (v) with $\pi_0=2\,\mathrm{pri}$. The statistic $\eta\hat T_t$ itself is not concave, the Hessian estimates having no sign; subtracting $2\Phi_0$ costs $2\,\mathrm{pri}$ in (a) and nothing in (b).

\paragraph{The hypotheses in the run.} On the events (c) and (e), for $t\le\tau$ and $Q=\mathcal R^-_t$: $\mathcal R^-_t\subset\mathcal R_t\subset S_t\cap B(0,R)$ by Proposition~\ref{prop:inc2}, so (i) and (ii) hold by Lemma~\ref{lem:sandwich}, (v) by the identity above, and (iv) with $\mathrm{dev}=b_1+B_n/(4L\lambda^2)$, since
\[
(\hat T_t-T_t)(y)=\sum_{u\le t}\ip{\hat g_u-s'_u(\mu_u),\mu_u-y}-\frac14\sum_{u\le t}|y-\mu_u|^2_{\hat s''_u-s''_u(\mu_u)}
\]
and the two sums are bounded by the events (c) and (e). Hypothesis (iii), $\eta T_t\le\rho''-\Omega_t$ with $\rho''=\mathrm{pri}+\mathrm{stab}+\eta\,\mathrm{dev}_\kappa$, is the online Newton inequality of Section~\ref{app:epoch-end}.

\subsection{The end of the epoch}\label{app:epoch-end}

This section proves that the epoch meets its bound on the events (a) to (e) of Section~\ref{app:epoch-events}. Let $S_t=\sum_{u\le t}s_u$, and let $\mathrm{xs}_t$ be the minimiser of least norm of $S_t$ over $\cK_\varepsilon$, so that $\mathrm{sReg}_t(x)\le\mathrm{sReg}_t(\mathrm{xs}_t)$ for $x\in\cK_\varepsilon$; the first condition of the stopping time asks $\mathrm{xs}_t\in\mathcal R^-_{t+1}$. Put
\begin{gather*}
\mathrm{pri}=\frac{R^2}{2\sigma^2},\qquad\mathrm{stab}=2C_g\eta^2dnL,\qquad\rho''=\mathrm{pri}+\mathrm{stab}+1+\eta\,\mathrm{dev}_\kappa,\\
\mathrm{dev}_\kappa=b_1+\frac\kappa2\,\frac{B_n}{L\lambda^2},\qquad\mathrm{dev}=b_1+\frac{B_n}{4L\lambda^2},
\end{gather*}
with $C_g=4$.

\paragraph{Before the stopping time.} For $s\le\tau$ the conditions of the stopping time in the rounds before $s$ give: $P_s$ is positive definite and $\delta_eI\le\frac12\bar P_s\le P_s\le2\bar P_s\le\delta_e^{-1}I$ (at $s=1$, $P_1=\bar P_1=I/\sigma^2$, and $4\delta_e\le\sigma^2\le\frac12$). Each $s_u$ is convex, so $s''_u(\mu_u)\ge0$ and $\bar P$ is nondecreasing; hence $\bar P_s\ge I/\sigma^2$, $\Sigma_s\le2\sigma^2I$, and $\Sigma_t\le4\Sigma_s$ for $s\le t\le\tau$.

\begin{lemma}[the closure]\label{lem:closure}
On the events (a) and (d), the second and the third conditions of the stopping time hold in round $\tau$.
\end{lemma}

\begin{proof}
On the event (a), $|Y_\tau|$ and $|\xi_\tau|$ are small enough that $a_\tau\ge\frac12$ and $q_\tau\ge\frac14$. For the third condition, $\bar P_{\tau+1}=\bar P_\tau+\eta\kappa s''_\tau(\mu_\tau)\ge\bar P_\tau$, and the event (d) at $t=\tau$ with the row $4\kappa\eta B_n\le\frac12$ gives $|\ip{(P_{\tau+1}-\bar P_{\tau+1})v,v}|\le\frac12\ip{\bar P_\tau v,v}\le\frac12\ip{\bar P_{\tau+1}v,v}$, so $\frac12\bar P_{\tau+1}\le P_{\tau+1}\le\frac32\bar P_{\tau+1}$; and $\frac12\bar P_{\tau+1}\ge\frac12\bar P_\tau\ge\delta_eI$. Finally Proposition~\ref{prop:sifacts}~(v) bounds $s''_u(\mu_u)$ by $\lambda(8/\varepsilon)(1-\lambda)^{-1}(d/\delta_e)^{1/2}I$ for $u\le\tau$, so
\[
2\bar P_{\tau+1}\le2\Bigl(\frac1{\sigma^2}+\eta\kappa n\,\frac{8\lambda}{(1-\lambda)\varepsilon}\Bigl(\frac d{\delta_e}\Bigr)^{1/2}\Bigr)I\le\frac1{\delta_e}I
\]
by the conditions $4\delta_e\le\sigma^2$ and $16\eta\lambda n(d\delta_e)^{1/2}\le\varepsilon$ of Lemma~\ref{lem:rows}.
\end{proof}

\begin{lemma}[the log-determinant]\label{lem:ld}
Let $\sigma^2,c_m>0$, $\bar P_1=\sigma^{-2}I$ and $\bar P_{u+1}=\bar P_u+H_u$ for $1\le u\le T$, with $0\le H_u\le\bar P_u$ and $\bar P_{T+1}\le c_mI$. Then $\sum_{u\le T}\tr(\bar P_u^{-1}H_u)\le2d\ln(c_m\sigma^2)$. In the epoch, on the events (a) and (d),
\[
\frac2\lambda\sum_{u\le\tau}\tr\bigl(s''_u(\mu_u)\Sigma_u\bigr)\le\frac{dL}{\lambda\eta}.
\]
\end{lemma}

\begin{proof}
Let $Z_u=\bar P_u^{-1/2}H_u\bar P_u^{-1/2}$; its eigenvalues lie in $[0,1]$, where $z\le2\ln(1+z)$, so $\tr(\bar P_u^{-1}H_u)=\tr Z_u\le2\ln\det(I+Z_u)=2\ln\det\bar P_{u+1}-2\ln\det\bar P_u$. The sum telescopes to $2\ln\det(\sigma^2\bar P_{T+1})\le2d\ln(c_m\sigma^2)$. In the epoch take $H_u=\eta\kappa s''_u(\mu_u)$: Proposition~\ref{prop:sifacts}~(v), $P_u\le2\bar P_u$ and the row $8\eta\lambda(2d\sigma^2)^{1/2}\le\varepsilon$ give $H_u\le\bar P_u$, and Lemma~\ref{lem:closure} gives $\bar P_{\tau+1}\le I/(2\delta_e)$. With $\Sigma_u\le2\bar P_u^{-1}$, $\kappa=\frac14$ and $\sigma^2\le\frac12$, the left side of the display is at most $\frac4{\lambda\eta\kappa}\,2d\ln\frac{\sigma^2}{2\delta_e}\le\frac{32\,d}{\lambda\eta}\ln\frac1{\delta_e}\le\frac{dL}{\lambda\eta}$.
\end{proof}

\begin{lemma}[online Newton step]\label{lem:ons2}
Let $t\ge1$, $\eta$ real, and for $1\le s\le t$ let $P_s,H_s$ be symmetric matrices with $P_{s+1}=P_s+\eta H_s$, $P_{s+1}$ positive definite with inverse $\Sigma_{s+1}$, and $g_s,\mu_s,y\in\R^d$, $e_s\in\R$ with $|\mu_{s+1}-y|^2_{P_{s+1}}\le|\mu_s-\eta\Sigma_{s+1}g_s-y|^2_{P_{s+1}}+2e_s$ for $s<t$. Then
\[
\eta\sum_{s=1}^t\Bigl(\ip{g_s,\mu_s-y}-\frac12|y-\mu_s|^2_{H_s}\Bigr)\le\frac12|\mu_1-y|^2_{P_1}+\sum_{s=1}^{t-1}e_s+\frac{\eta^2}2\sum_{s=1}^t|g_s|^2_{\Sigma_{s+1}} .
\]
\end{lemma}

\begin{proof}
Expanding the square, $|\mu_s-\eta\Sigma_{s+1}g_s-y|^2_{P_{s+1}}=|\mu_s-y|^2_{P_s}+\eta|\mu_s-y|^2_{H_s}-2\eta\ip{g_s,\mu_s-y}+\eta^2|g_s|^2_{\Sigma_{s+1}}$. For $s<t$ the hypothesis gives
$\eta(\ip{g_s,\mu_s-y}-\frac12|y-\mu_s|^2_{H_s})\le\frac12|\mu_s-y|^2_{P_s}-\frac12|\mu_{s+1}-y|^2_{P_{s+1}}+e_s+\frac{\eta^2}2|g_s|^2_{\Sigma_{s+1}}$, and for $s=t$ the same holds without the terms in $\mu_{t+1}$ and $e_t$, since the square is nonnegative. The sum telescopes.
\end{proof}

\paragraph{Projections with defects.} Let $Q\subset B(0,R)$ be closed and convex, $P$ positive definite, $w$ a point, $e>0$, and $\mu\in Q$ with $\frac12|\mu-w|^2_P\le\frac12|x-w|^2_P+e$ for $x\in Q$, as AP returns. Let $\mu^*$ be the projection of $w$ on $Q$ in the norm of $P$. The function $\frac12|\cdot-w|_P^2$ exceeds its value at $\mu^*$ by at least $\frac12|\cdot-\mu^*|^2_P$ on $Q$, so $|\mu-\mu^*|_P\le(2e)^{1/2}$, and for $y\in Q$, $|\mu^*-y|^2_P\le|w-y|^2_P$ gives $|\mu-y|^2_P\le|w-y|^2_P+2\bigl((2e)^{1/2}\,2R\|P\|^{1/2}+e\bigr)$. With the accuracy $e=1/(16n^2\max(1,4R^2\tr P'))$ of step~7 and $\|P'\|\le\tr P'$, the defect $(2e)^{1/2}2R\|P'\|^{1/2}+e$ is at most $\frac1{2\sqrt2\,n}+\frac1{16n^2}<\frac1{2n}$, so the defects of an epoch sum to less than $\frac12$.

\paragraph{The inequality (iii).} Let $t\le\tau$ and $y\in\mathcal R^-_t$. By Proposition~\ref{prop:inc2}, $y\in S_{t}\cap B(0,R)$, and for $s\le t-1$ Lemma~\ref{lem:mon} gives $\Phi_s(y)\le\Phi_{t-1}(y)+2\,\mathrm{dev}'\le(c_1-\gamma)\rho+2\,\mathrm{dev}'\le c_1\rho$ by the margin row $2\,\mathrm{dev}'\le\gamma\rho$, so $y\in\mathcal R_{s+1}$, the set on which step~7 of round $s$ projects the point $w_\mu=\mu_s-\eta\Sigma_{s+1}\hat g_s$ (Lemma~\ref{lem:r1}). Lemma~\ref{lem:ons2} with $H_s=\kappa\hat s''_s$, $g_s=\hat g_s$ and the defects above gives
\[
\eta\sum_{s\le t}\Bigl(\ip{\hat g_s,\mu_s-y}-\frac\kappa2|y-\mu_s|^2_{\hat s''_s}\Bigr)\le\mathrm{pri}+\frac12+\frac{\eta^2}2\sum_{s\le t}|\hat g_s|^2_{\Sigma_{s+1}} .
\]
Here $|\hat g_s|^2_{\Sigma_s}=\kappa_0Y_s^2|\xi_s|^2$, $\Sigma_{s+1}\le4\Sigma_s$ before $\tau$ and at $\tau$ on the events, and the events (a) and (b) bound $\sum_{s\le\tau}\kappa_0Y_s^2|\xi_s|^2$ by $C_gdnL$ under the row R15p, so the last term is at most $\mathrm{stab}$. On the events (c) and (e) the left side differs from its value with $s'_s(\mu_s)$, $s''_s(\mu_s)$ in place of the estimates by at most $\eta\,\mathrm{dev}_\kappa$. Since $\frac\kappa2=\frac18$ and $\frac\eta8W_t=\Omega_t$, this is (iii) of Lemma~\ref{lem:d2}: $\eta T_t(y)\le\rho''-\Omega_t(y)$, with $\frac{a-\kappa}2=\frac18$.

\begin{lemma}[the restart wrapper]\label{lem:w}
Let $\cK_\varepsilon$ be a convex subset of a normed space, $\tau\ge1$, $s_1,\dots,s_\tau$ real functions with $S_\tau$ convex on $\cK_\varepsilon$, $\sigma_1,\dots,\sigma_\tau$ reals, $\mathrm{sReg}_t(x)=\sum_{u\le t}\sigma_u-S_t(x)$, and $\mathrm{xs}_{\tau-1},\mathrm{xs}_\tau\in\cK_\varepsilon$ minimisers of $S_{\tau-1}$ and $S_\tau$ on $\cK_\varepsilon$. Let $K_{\tau+1}\subset K_\tau\subset\cK_\varepsilon$ with $K_\tau$ closed, $\beta\ge0$ on $K_\tau$, $E$ a real number, Fire a truth value and $A_1,A_2,A_5,\kappa_s,\kappa_c,b$ reals with $A_1+A_5\le\kappa_s$ and $A_2+\kappa_c\le b$. Assume:
\begin{enumerate}[label=(W\arabic*),start=0]
\item $\mathrm{xs}_{\tau-1}\in K_\tau$;
\item $E\le\mathrm{sReg}_\tau(\mathrm{xs}_\tau)+A_1$;
\item $\mathrm{sReg}_\tau\le A_2-\beta$ on $K_\tau$;
\item $\beta\ge b$ at every point of $K_\tau$ in the closure of $\cK_\varepsilon\setminus K_\tau$, and every $y\in K_\tau$ with $\beta(y)<b$ lies in $K_{\tau+1}$;
\item if Fire, then $\mathrm{sReg}_\tau\le-\kappa_s$ on $K_\tau$; if $\mathrm{sReg}_\tau\le-\kappa_c$ on $K_\tau$, then Fire;
\item $s_\tau(\mathrm{xs}_{\tau-1})-s_\tau(\mathrm{xs}_\tau)\le A_5$.
\end{enumerate}
Then: (a) Fire implies $E\le0$; (b) if $\mathrm{xs}_\tau\notin K_\tau$, then Fire; (c) if not Fire, then $\mathrm{xs}_\tau\in K_{\tau+1}$ and $E\le A_1+A_2$.
\end{lemma}

\begin{proof}
(a) Since $\mathrm{xs}_{\tau-1}$ minimises $S_{\tau-1}$, $\mathrm{sReg}_\tau(\mathrm{xs}_\tau)=\mathrm{sReg}_{\tau-1}(\mathrm{xs}_\tau)+\sigma_\tau-s_\tau(\mathrm{xs}_\tau)\le\mathrm{sReg}_\tau(\mathrm{xs}_{\tau-1})+s_\tau(\mathrm{xs}_{\tau-1})-s_\tau(\mathrm{xs}_\tau)$, which is at most $-\kappa_s+A_5$ by (W0), (W4) and (W5); with (W1), $E\le A_1+A_5-\kappa_s\le0$. (b) If $K_\tau$ is empty, (W4) gives Fire. Otherwise let $y_0\in K_\tau$, and let $y_1$ be the last point of the segment from $y_0$ to $\mathrm{xs}_\tau$ in the closed set $K_\tau$; it lies in the closure of $\cK_\varepsilon\setminus K_\tau$, and by convexity $S_\tau(y_1)\le\max(S_\tau(y_0),S_\tau(\mathrm{xs}_\tau))=S_\tau(y_0)$. So $\mathrm{sReg}_\tau(y_0)\le\mathrm{sReg}_\tau(y_1)\le A_2-b\le-\kappa_c$ by (W2) and (W3), for every $y_0\in K_\tau$, and (W4) gives Fire. (c) By (b), $\mathrm{xs}_\tau\in K_\tau$, and $\mathrm{sReg}_\tau(\mathrm{xs}_\tau)>-\kappa_c$ by (W4), since $\mathrm{xs}_\tau$ maximises $\mathrm{sReg}_\tau$ on $K_\tau$. By (W2), $\beta(\mathrm{xs}_\tau)<A_2+\kappa_c\le b$, so $\mathrm{xs}_\tau\in K_{\tau+1}$ by (W3), and $E\le A_1+A_2$ by (W1), (W2) and $\beta\ge0$.
\end{proof}

\paragraph{The instance.} Take $K_\tau=\mathcal R^-_\tau$, $K_{\tau+1}=\mathcal R^-_\tau\cap\mathcal R^-_{\tau+1}$, $\sigma_u=s_u(\mu_u)$, $\beta=\Omega_\tau/\eta$, $E=\sup_{\cK_\varepsilon}\mathrm{eReg}_\tau$, Fire the firing of the test in round $\tau$, and
\begin{gather*}
\eta A_1=\eta+\frac{dL}\lambda,\qquad\eta A_2=\eta+\rho'',\qquad A_5=1+\frac{10}{\varepsilon\lambda},\qquad\eta b=(c_1-\gamma)\rho-\mathrm{pri}-\mathrm{dev}',\\
\eta\kappa_s=(\theta-\alpha)\rho-2\,\mathrm{pri}-\eta\,\mathrm{dev}-\eta,\qquad\eta\kappa_c=7\theta\rho+6\rho''+7\eta\,\mathrm{dev}+\eta .
\end{gather*}
On the events:
(W0) is the first condition of the stopping time in round $\tau-1$ (for $\tau=1$, $\mathcal R^-_1=\cK_\varepsilon$).
(W1): the law $p'$ of the play is that of $(1-\lambda)X+\lambda\mu$ for $X$ of law $p=N(\mu,\Sigma)$, so $s(\mu)=(1-\frac1\lambda)\E_pe+\frac1\lambda\E_{p'}e$ and $\E_{p'}e-s(\mu)=(1-\lambda)(\E_pe-s(\mu))\le\frac2\lambda\tr(s''(\mu)\Sigma)+\frac{2\delta_ed}\lambda$ by Proposition~\ref{prop:sifacts}~(iv); with $s_u\le e_u$, the sum over $u\le\tau$ and Lemma~\ref{lem:ld} it gives $E\le\mathrm{sReg}_\tau(\mathrm{xs}_\tau)+A_1$.
(W2): for $y\in\mathcal R^-_\tau$, Lemma~\ref{lem:sandwich}~(i) and the inequality (iii) give $\eta\,\mathrm{sReg}_\tau(y)\le\eta T_\tau(y)+\eta\nu\le\rho''+\eta-\Omega_\tau(y)$, as $\nu\le1$.
(W3): a point $y$ of $\mathcal R^-_\tau$ in the closure of $\cK_\varepsilon\setminus\mathcal R^-_\tau$ has $\Phi_{\tau-1}(y)\ge(c_1-\gamma)\rho$, so $\tau\ge2$ and $\Omega_\tau(y)\ge\Omega_{\tau-1}(y)\ge\hat\Omega_{\tau-1}(y)-\mathrm{dev}'\ge(c_1-\gamma)\rho-\mathrm{pri}-\mathrm{dev}'$ by the event (e); and if $\Omega_\tau(y)<\eta b$, then $\Phi_\tau(y)\le\mathrm{pri}+\Omega_\tau(y)+\mathrm{dev}'<(c_1-\gamma)\rho$, so $y\in\mathcal R^-_{\tau+1}$.
(W4): Lemma~\ref{lem:d2} on $Q=\mathcal R^-_\tau$ with $h=\theta\rho$, $e=\alpha\rho$, $\pi_0=2\,\mathrm{pri}$ and $\eta\nu\le\eta$: part (a) for the first half, part (b) for the second, the test firing when $m=\Xi_\tau(p^-)\le-\theta\rho$.
(W5): $-10/(\lambda\varepsilon)\le s_\tau\le1$ on $\cK_\varepsilon$ by Proposition~\ref{prop:sifacts}~(iii).

\begin{lemma}[the rows]\label{lem:rows}
For all $d,n\ge1$ and $0<\delta<1$ with $\varepsilon\le\frac12$, the parameters of Algorithm~\ref{alg:df2o} satisfy
\begin{align*}
&2\eta+\frac{dL}\lambda+\mathrm{pri}+\mathrm{stab}+\eta\,\mathrm{dev}_\kappa+1\le\rho, \qquad 2\,\mathrm{dev}'\le\gamma\rho,\qquad 16c_1\rho L\lambda^2+4\eta B_n<1,\\
&3\eta+\frac{dL}\lambda+\frac{10\eta}{\varepsilon\lambda}+2\,\mathrm{pri}+\eta\,\mathrm{dev}\le(\theta-\alpha)\rho,\\
&2\eta+8\,\mathrm{pri}+7\,\mathrm{stab}+7\eta\,\mathrm{dev}_\kappa+7\eta\,\mathrm{dev}+\mathrm{dev}'+7\theta\rho+7\le(c_1-\gamma)\rho,\\
&3+2\,\mathrm{pri}+\mathrm{stab}+\eta\,\mathrm{dev}_\kappa+\eta\,\mathrm{dev}+\mathrm{dev}'+\theta\rho\le(c_1-\gamma)\rho,
\end{align*}
together with the other fields of the structure Df2Rows: the positivity of the parameters and fifteen inequalities among them that the probability bounds of Section~\ref{app:epoch-events} read.
\end{lemma}

The row of the second line is $A_1+A_5\le\kappa_s$, the row of the third line is $A_2+\kappa_c\le b$, and the first row of the first line is $A_1+A_2\le\rho/\eta$. Lemma~\ref{lem:w} therefore gives: if the test fires in round $\tau$, then $\mathrm{eReg}_\tau\le0$ on $\cK_\varepsilon$; if it does not, then $\mathrm{xs}_\tau\in\mathcal R^-_{\tau+1}$ and $\mathrm{eReg}_\tau\le\rho/\eta$ on $\cK_\varepsilon$.


\paragraph{The guard.} On the events the guard does not act in the rounds up to $\tau$. In a round $t\le\tau$: $a_t,q_t>0$ and the two traces are at most $d/\delta_e$, by the second and third conditions of the stopping time for $t<\tau$ and by Lemma~\ref{lem:closure} at $t=\tau$; $|z_t|\le R/\delta_e^3$ by Lemma~\ref{lem:z}; the point $w_\mu$ of step~6 stays below its bound, since $|w_\mu-\mu_t|^2=\eta^2|\Sigma_{t+1}\hat g_t|^2\le\eta^2\|\Sigma_{t+1}\|\,\kappa_0Y_t^2|\xi_t|^2/q_t$ by Lemma~\ref{lem:r1}~(iv), with $\|\Sigma_{t+1}\|\le1/\delta_e$, $q_t\ge\frac14$, $Y_t^2\le22n$ on the event (b), $|\xi_t|^2\le\frac83dx$ on the event (a) and $|\mu_t|\le R$, and the product is at most $16dnL/\delta_e$; and AP does not fail, by the next two lemmas: the ball of step~6 lies in the regions of the next round.

\begin{lemma}[a ball under quadratics]\label{lem:ibgeo}
Let $Q_0\subset\R^d$ be convex with $\{|x|<r_0\}\subset Q_0\subset\{|x|<R\}$, $\Phi$ a quadratic with symmetric Hessian $P$, $|\ip{Pv,v}|\le D|v|^2$, and $c$ real, $m,G,D>0$. Let $y\in Q_0$ with $\Phi(y)\le c-m$ and $|\nabla\Phi(y)|\le G$, and $0<s\le1$ with $s(R+r_0)\le\min\bigl(m/(2G),(m/D)^{1/2}\bigr)$. Then the open ball of radius $sr_0$ about $(1-s)y$ lies in $Q_0$ and in $\{\Phi\le c\}$.
\end{lemma}

\begin{proof}
A point of the ball is $(1-s)y+sb$ with $|b|<r_0$, a convex combination of two points of $Q_0$. Its distance to $y$ is $s|b-y|\le s(r_0+R)$, and for $|v|\le s(r_0+R)$, $\Phi(y+v)\le c-m+G|v|+\frac D2|v|^2\le c-m+\frac m2+\frac m2$.
\end{proof}

\begin{lemma}[the inner ball]\label{lem:ib2}
On the events, let $t\le\tau$ be a round whose test does not fire, the calls of step~7 of the rounds before $t$ having met their accuracy. Then $\Phi_t(p^-)\le(c_1-\gamma)\rho-2$, and the ball of radius $s_b(1-\varepsilon)$ about $(1-s_b)p^-$ lies in $\mathcal R^-_{t+1}$, hence in $\mathcal R_{t+1}$.
\end{lemma}

\begin{proof}
No fire means $\Xi_t(p^-)>-\theta\rho$. As $p^-\in\mathcal R^-_t$, the hypotheses of Lemma~\ref{lem:d2} in the run give $\eta T_t(p^-)\ge\eta\hat T_t(p^-)-\eta\,\mathrm{dev}\ge\Xi_t(p^-)-\eta\,\mathrm{dev}$, and the inequality (iii) gives $\eta T_t(p^-)\le\rho''-\Omega_t(p^-)$; so $\Omega_t(p^-)<\rho''+\theta\rho+\eta\,\mathrm{dev}$. The point lies in $S_t\cap B(0,R)$, where the event (e) gives $\Phi_t(p^-)\le\mathrm{pri}+\Omega_t(p^-)+\mathrm{dev}'$, and the last row of Lemma~\ref{lem:rows} bounds the sum by $(c_1-\gamma)\rho-2$. The gradient $\nabla\Phi_t(y)=P_1y+\eta\kappa\sum_{u\le t}\hat s''_u(y-\mu_u)$ is at most $R/\delta_e^2$ on $B(0,R)$: the points have norm at most $R$, the norms of the $\hat s''_u$ sum to at most $\delta_e^{-2}$ on the event (a), and $4\delta_e\le\sigma^2$. The Hessian $P_{t+1}$ has norm at most $1/\delta_e$. Lemma~\ref{lem:ibgeo} with $Q_0=\cK_\varepsilon$, $r_0=1-\varepsilon$, $c=(c_1-\gamma)\rho$, $m=2$, $G=R/\delta_e^2$, $D=1/\delta_e$ and $s=s_b=\delta_e^2/(R(R+1))$, for which $s(R+r_0)\le\delta_e^2/R=\min(m/(2G),(m/D)^{1/2})$, puts the ball in $\cK_\varepsilon$ and below the level $(c_1-\gamma)\rho$ of $\Phi_t$, that is in $\mathcal R^-_{t+1}$.
\end{proof}

\begin{lemma}[the centre of the test]\label{lem:z}
Let $\Sigma$ be a linear map with $|\Sigma v|\le|v|/\delta_e$, and $b,G$ vectors with $|b|\le R/(4\delta_e^2)$ and $\frac\eta2|G|\le\frac34R/\delta_e^2$, $\eta\ge0$. Then $|\Sigma(b-\frac\eta2G)|\le R/\delta_e^3$. On the events (a) and (d), the point $z_t$ of step~4 meets these bounds in every round $t\le\tau$.
\end{lemma}

\begin{proof}
The first part is the triangle inequality. For the second: $\|\Sigma_{t+1}\|\le1/\delta_e$ by the third condition of the stopping time and, at $\tau$, Lemma~\ref{lem:closure}; $b_t=\eta\kappa\sum_{s\le t}\hat s''_s\mu_s$ with $|\mu_s|\le R$ and $\sum_s\|\hat s''_s\|\le\delta_e^{-2}$ on (a), with $\eta\kappa\le\frac14$; and $|\hat g_s|\le|Y_s|\,\|P_s\|^{1/2}|\xi_s|/(1-\lambda)$, which on (a) makes $\frac\eta2|G_t|\le\frac34R/\delta_e^2$ under the condition $\frac14\kappa_0^{1/2}C_0(\frac{32}9)^{1/2}e^{-3x/2}n^{3/2}d^{1/2}x^{3/2}\le\frac34$ for $x\ge14+3\ln(dn)$, which is the statement Df2RowZ.
\end{proof}

\begin{proposition}[the epoch on the events]\label{prop:det}
In the setting of Section~\ref{app:epoch-setting}, on every run of the epoch on which the events (a) to (e) hold, the epoch meets its bound at horizon $m$.
\end{proposition}

\begin{proof}
On the events, Proposition~\ref{prop:inc2}, the hypotheses of Lemma~\ref{lem:d2} and the instance of Lemma~\ref{lem:w} hold in every round up to $\tau$, and the guard does not act. If the test fires in round $\tau$, the epoch ends there with $\mathrm{eReg}_\tau\le0$ on $\cK_\varepsilon$. If not, $\mathrm{xs}_\tau\in\mathcal R^-_{\tau+1}$ and the learner does not restart, so the first and the fourth conditions of the stopping time hold in round $\tau$, and the second and the third hold there by Lemma~\ref{lem:closure}; so $\tau=m$, and $\mathrm{eReg}_m\le\rho/\eta$ on $\cK_\varepsilon$.
\end{proof}

Proposition~\ref{prop:epoch} follows: the epoch fails to meet its bound only off the intersection of the events (a) to (e), which has probability at most $\frac{149}{168}\delta_e$.

\section{The projection routine}\label{app:ap}

This appendix gives the routine AP of Algorithm~\ref{alg:df2o} and proves that it meets the specification of Section~\ref{sec:poly-learner}. For $d\ge2$ it is the central-cut ellipsoid method with objective cuts, with a separating direction computed from values of the gauge by finite differences; for $d=1$ it is a bisection. The reduction from membership to separation is in \cite[Theorems~4.3.2 and~4.4.4]{GLS93}, through the shallow-cut ellipsoid method and the polar, and deterministic minimisation of a convex function from its values with a number of evaluations logarithmic in the accuracy is in \cite{Pro96}. What is written here is a direct finite-difference step, with no polar and no shallow cut, the loops capped so that the routine halts on every input, and the Lipschitz constant $\tr(P)(R+|w|)$, which needs no spectral norm. All norms are Euclidean, $|v|_1=\sum_i|v_i|$, $e_i$ is the $i$-th coordinate vector, $F(x)=\frac12|x-w|^2_P$, and for $A$ positive definite $E(A,a)=\{x:\ip{A^{-1}(x-a),x-a}\le1\}$.

\subsection{A subgradient from values}

\begin{lemma}[finite differences]\label{lem:fd}
Let $h:\R^d\to\R$ be convex, $y\in\R^d$, $\delta_1,\dots,\delta_d>0$, and $v_0,v_i^\pm$ numbers with $|v_0-h(y)|\le\eta_0$ and $|v_i^\pm-h(y\pm\delta_ie_i)|\le\eta_i$. Put $g_i=(v_i^+-v_i^-)/(2\delta_i)$ and $D_i=v_i^++v_i^--2v_0$. Then for every $z$,
\[
h(z)\ge h(y)+\ip{g,z-y}-\sum_i|z_i-y_i|\,\frac{D_i+2\eta_0+4\eta_i}{2\delta_i}.
\]
\end{lemma}

\begin{proof}
With exact values: let $t=z-y$, $s_i$ the sign of $t_i$ and $a=1/(1+\sum_i|t_i|/\delta_i)$. Then $y=az+\sum_ia(|t_i|/\delta_i)(y-s_i\delta_ie_i)$ is a convex combination, so $h(z)\ge h(y)+\sum_i(|t_i|/\delta_i)(h(y)-h(y-s_i\delta_ie_i))$, and $h(y)-h(y-s_i\delta_ie_i)=s_ig_i\delta_i-D_i/2$. With errors, the exact $g_i$ is within $\eta_i/\delta_i$ of the measured one, and the exact $D_i$ is at most the measured one plus $2\eta_i+2\eta_0$.
\end{proof}

\paragraph{Procedure LINE$(\varphi,\ell,\tau)$,} for $\varphi$ convex on $\R$ with values available to a prescribed error, $\ell>0$, $\tau>0$ and a cap $J$. A window is five points $p_0+k\lambda$, $0\le k\le4$, with measured values $v_k$; its measured spread is $(v_4-v_3-v_1+v_0)/\lambda$. Start with $p_0=0$, $\lambda=\ell/4$, the five values asked with error $\tau\lambda/400$. While $v_3+v_1-2v_2>\tau\lambda/2$: if $J$ halvings are done, fail; otherwise cut $[p_1,p_3]$ into eight intervals of length $\lambda/4$, ask the six new values with error $\tau\lambda/1600$, and keep, of the two windows of spacing $\lambda/4$ that start at $p_1$ and at $p_2$, the one with the smaller measured spread (the first on a tie). Return $s=p_2$ and $\delta=\lambda$.

\begin{lemma}[LINE]\label{lem:line}
Let $\varphi$ be convex and $L$-Lipschitz on $\R$, $0<\tau\le L$, $\ell>0$ and $J=\lceil\log_2(8L/\tau)\rceil$, and let every value asked at a point of $[0,\ell]$ with an error $\eta\ge\tau\ell4^{-J}/1600$ be returned within $\eta$. Then LINE with the cap $J$ returns $(s,\delta)$ with $0\le s\le\ell$, $\ell4^{-(J+1)}\le\delta\le\ell/4$ and $\varphi(s+\delta)+\varphi(s-\delta)-2\varphi(s)\le\frac{51}{100}\tau\delta$.
\end{lemma}

\begin{proof}
The chord slopes of a convex function over consecutive intervals are nondecreasing. Let $\sigma_1\le\dots\le\sigma_4$ be those of a window and $\tau_1\le\dots\le\tau_8$ those of the eight small intervals of $[p_1,p_3]$; then $\tau_1\ge\sigma_1$ and $\tau_8\le\sigma_4$, so the exact spreads of the two children add up to $(\tau_8-\tau_1)-(\tau_5-\tau_4)\le\sigma_4-\sigma_1$, and the exact second difference at $p_2$ is $\lambda(\sigma_3-\sigma_2)\le\lambda(\sigma_4-\sigma_1)$. In every window the three inner values have error at most $\tau\lambda/400$ and the two outer ones at most $\tau\lambda/100$, so a measured slope is within $\tau/80$ of the exact one and a measured spread within $\tau/40$. The kept child has measured spread at most half that of its parent plus $3\tau/80$, so after $j$ halvings it is at most $(2L+\tau/100)2^{-j}+3\tau/40$, and the measured second difference over $\lambda$ is at most that plus $7\tau/200$. For $j\ge\log_2(8L/\tau)$ this is below $0.37\,\tau$, so the loop stops within $J$ halvings. The windows are nested in $[0,\ell]$; at the stop the measured second difference is at most $\tau\delta/2$, and the exact one at most $\tau\delta/2+4\tau\delta/400$.
\end{proof}

\paragraph{Procedure SEP$(h,y^0,\tau,\ell_1)$,} for $h$ convex and $L$-Lipschitz on $\R^d$, with the cap $J=\lceil\log_2(8L/\tau)\rceil$ in each LINE. For $i=1,\dots,d$: run LINE on $s\mapsto h(y^{i-1}+se_i)$ with length $\ell_i$, get $(s_i,\delta_i)$, and put $y^i=y^{i-1}+s_ie_i$ and $\ell_{i+1}=\tau\delta_i/(16L)$. At $y=y^d$ ask $v_0$ with error $\tau\min_i\delta_i/400$ and $v_i^\pm$ at $y\pm\delta_ie_i$ with error $\tau\delta_i/400$, and return $y$ and $g$, $g_i=(v_i^+-v_i^-)/(2\delta_i)$.

\begin{lemma}[SEP]\label{lem:sep}
Let $d\ge1$, $h$ convex and $L$-Lipschitz on $\R^d$, $0<\tau\le L$, $\ell_1>0$, and let every value asked at a point within $2\ell_1$ of $y^0$ with an error $\eta\ge\eta_{\min}=(\tau\ell_1/1600)2^{-3d(J+1)}$ be returned within $\eta$. Then SEP returns $(y,g)$ with $|y-y^0|\le\frac{64}{63}\ell_1$ and, for every $z$,
\[
h(z)\ge h(y)+\ip{g,z-y}-\tau|z-y|_1 .
\]
\end{lemma}

\begin{proof}
$\ell_{i+1}\le\ell_i/64$, since $\delta_i\le\ell_i/4$ and $\tau\le L$, so $|y-y^i|\le\frac{64}{63}\ell_{i+1}$, and every point asked is within $2\ell_1$ of $y^0$; $\delta_i\ge\ell_i4^{-(J+1)}$ and $16L/\tau\le2^{J+1}$ give the bound on the errors. At $y^i$ the exact second difference in direction $i$ at step $\delta_i$ is at most $0.51\,\tau\delta_i$ (Lemma~\ref{lem:line}); as a function of the base point it is $4L$-Lipschitz, so at $y$ it is at most $0.51\,\tau\delta_i+\frac{64}{63}\tau\delta_i/4\le0.764\,\tau\delta_i$, and the measured $D_i$ at most $0.774\,\tau\delta_i$. Lemma~\ref{lem:fd} with $\eta_0,\eta_i\le\tau\delta_i/400$ gives an error of at most $(0.774+0.015)\tau/2<\tau$ per unit of $|z_i-y_i|$.
\end{proof}

\subsection{The ellipsoid step}

\begin{lemma}[central cut]\label{lem:cc}
Let $d\ge2$, $A$ positive definite, $c\ne0$, $b=Ac/\ip{Ac,c}^{1/2}$, $a'=a-b/(d+1)$ and $A'=\frac{d^2}{d^2-1}\bigl(A-\frac2{d+1}bb^{\mathsf T}\bigr)$. Then $A'$ is positive definite, $E(A,a)\cap\{x:\ip{c,x-a}\le0\}\subset E(A',a')$, and $\det A'=\det A\,\bigl(\frac{d^2}{d^2-1}\bigr)^{d-1}\bigl(\frac d{d+1}\bigr)^2\le e^{-1/(d+1)}\det A$.
\end{lemma}

\begin{proof}
The map $x=a+A^{1/2}u$ and a rotation reduce to $A=I$, $a=0$, $c=e_1$. Then $A'$ has the eigenvalue $d^2/(d+1)^2$ once and $d^2/(d^2-1)$ otherwise, which gives the positive definiteness and the determinant, and $1+x\le e^x$ gives the bound. For $|u|\le1$ with $u_1\le0$: $\frac{(d+1)^2}{d^2}(u_1+\frac1{d+1})^2+\frac{d^2-1}{d^2}(1-u_1^2)=1+\frac{2(d+1)}{d^2}u_1(u_1+1)\le1$.
\end{proof}

\begin{lemma}[inscribed ball]\label{lem:ball}
If $A$ is positive definite, $\rho\ge0$ and $B(b,\rho)\subset E(A,a)$, then $\det A\ge\rho^{2d}$.
\end{lemma}

\begin{proof}
Let $v$ be a unit eigenvector of $A$ for its least eigenvalue $m$. The mean of the inequalities of $E(A,a)$ at $b\pm\rho v$ is $\ip{A^{-1}(b-a),b-a}+\rho^2/m\le1$, so every eigenvalue is at least $\rho^2$.
\end{proof}

\paragraph{The gauge.} For $Q$ closed and convex with $B(a_0,r)\subset Q\subset B(0,R)$, the gauge $\gamma(x)=\inf\{t>0:a_0+(x-a_0)/t\in Q\}$ is convex and $(1/r)$-Lipschitz, $|x-a_0|/(2R)\le\gamma(x)\le|x-a_0|/r$, and $\gamma(x)\le1$ exactly when $x\in Q$. Its value at $x\ne a_0$ with error $\eta$ is the upper end of $\lceil\log_2(|x-a_0|/(r\eta))\rceil$ steps of bisection on $[|x-a_0|/(2R),|x-a_0|/r]$, each step one membership question; the number of steps is fixed before any answer is read.

\paragraph{Procedure ELL-cap$(Q,a_0,r,P,w,e)$,} for $d\ge2$, a set $Q$ given by membership questions, $r>0$ and $e>0$. If $r>R$, fail. Otherwise put $L_F=\tr(P)(R+|w|)$, $\theta=\min(1,e/(3L_FR))$, $\varepsilon_Q=\theta r/2$, $\tau=\varepsilon_Q/(36R^2d^{1/2})$, $\ell_1=r\varepsilon_Q/(24R)$, $N=\lceil2d(d+1)\ln(2R/(\theta r))\rceil+1$, $J=\lceil\log_2(8/(r\tau))\rceil$, $\eta_{\min}=(\tau\ell_1/1600)2^{-3d(J+1)}$ and $B=\lceil\log_2(3R/(r\eta_{\min}))\rceil$. Start with $A_0=R^2I$ and $a^0=0$, and for $k=0,\dots,N-1$:
\begin{enumerate}[label=\arabic*.,leftmargin=2em]
\item if $|a^k|>R$, put $c=a^k$;
\item else, if $a^k\notin Q$ (one question), run SEP on the gauge from $a^k$ with $\tau$, $\ell_1$, $L=1/r$ and the cap $J$, each value read with its bisection and failing if that needs more than $B$ steps; fail if SEP fails or returns $g=0$; put $c=g$;
\item else, record $a^k$; if $P(a^k-w)=0$, return $a^k$; otherwise put $c=P(a^k-w)$;
\end{enumerate}
and let $(A_{k+1},a^{k+1})=(A',a')$ of Lemma~\ref{lem:cc}. After $N$ steps, return the recorded point with the least value of $F$, or fail if there is none.

\begin{proposition}[ELL-cap]\label{prop:ell}
(a) On every input ELL-cap fails or returns a point of $Q$ of norm at most $R$. (b) If $d\ge2$, $Q$ is closed and convex with $B(a_0,r)\subset Q\subset B(0,R)$, $r>0$, $P$ is positive definite and $e>0$, then ELL-cap returns a point $\mu\in Q$ with $F(\mu)\le F(x)+e$ for every $x\in Q$.
\end{proposition}

\begin{proof}
(a) A returned point is a recorded one: the oracle accepted it and its norm is at most $R$. (b) $r\le R$, so the constants are positive. The gauge is convex and $(1/r)$-Lipschitz, so Lemma~\ref{lem:line} rules out the failure of a search, and the values that SEP asks lie within $3R$ of $a_0$ with errors at least $\eta_{\min}$, so their bisections need at most $B$ steps. Let $p$ minimise $F$ on $Q$ and $p_\theta=(1-\theta)p+\theta a_0$; by convexity every $z\in B'=B(p_\theta,\theta r/2)$ has $B(z,\varepsilon_Q)\subset Q$.

The cut of step~2 keeps $B'$. Here $\gamma(a^k)>1$ and $|y-a^k|\le2\ell_1$, so $\gamma(y)>1-\varepsilon_Q/(12R)$. For $z\in Q$, $|z-y|_1\le3Rd^{1/2}$, and Lemma~\ref{lem:sep} gives $\ip{g,z-y}<\varepsilon_Q/(6R)=:\varepsilon_1\le\frac1{12}$. At $z=a_0$ this gives $\ip{g,y-a_0}>1-\varepsilon_1$ with $|y-a_0|\le2.05R$, so $|g|\ge4/(9R)$, and in particular $g\ne0$. For $z$ with $B(z,\varepsilon_Q)\subset Q$, the bound at $z+\varepsilon_Qg/|g|$ gives $\ip{g,z-a^k}<\varepsilon_1-(\varepsilon_Q-2\ell_1)|g|\le\frac{\varepsilon_Q}{6R}-\frac{11}{12}\cdot\frac49\cdot\frac{\varepsilon_Q}R<0$. The cut of step~1 keeps $B(0,R)\supset Q$, since $\ip{a^k,z-a^k}\le|a^k|(R-|a^k|)<0$ for $|z|\le R$.

So, by Lemma~\ref{lem:cc} and induction, $B'\subset E(A_k,a^k)$ up to the first step~3 whose cut removes a point of $B'$. If there were none, Lemma~\ref{lem:ball} would give $(\theta r/2)^{2d}\le\det A_N\le R^{2d}e^{-N/(d+1)}$, against the choice of $N$. So some recorded $a^j$ either has $P(a^j-w)=0$ and then minimises $F$ on $\R^d$, or has $\ip{P(a^j-w),z-a^j}>0$ for some $z\in B'$; then $F(a^j)<F(z)\le F(p)+L_F|z-p|\le F(p)+L_F\theta(2R+r/2)\le F(p)+e$, since $F$ is $L_F$-Lipschitz on $B(0,R)$ ($\|P\|\le\tr P$). The returned point is recorded, so it lies in $Q$, and its value is at most $F(a^j)$.
\end{proof}

By its caps, ELL-cap makes at most $N\bigl(1+(d(5+6J)+2d+1)B\bigr)$ membership questions on every input with $r\le R$: one for the centre of each of the $N$ steps and, in step~2, the values of SEP, $d$ runs of LINE of at most $5+6J$ values and $2d+1$ more values, each value at most $B$ questions. The Lean statements above do not state this count; the machine layer counts the steps of the program through the same caps.

\paragraph{Procedure ELL-1$(Q,a_0,r,p,w,e)$,} for $d=1$. If $r>R$ or $|a_0|>R$, or if $a_0\notin Q$, fail. Let $w'$ be the point of $[-R,R]$ nearest to $w$; if $w'\in Q$, return $w'$. Otherwise put $\xi=e/(p(R+|w|))$, $\mathrm{lo}=a_0$, $\mathrm{hi}=w'$, and repeat $m=\lceil\log_2(|w'-a_0|/\xi)\rceil$ times: the midpoint of $\mathrm{lo}$ and $\mathrm{hi}$ replaces $\mathrm{lo}$ if $Q$ contains it and $\mathrm{hi}$ otherwise. Return $\mathrm{lo}$. The routine AP is ELL-1 on the coordinate line when $d=1$ and ELL-cap when $d\ge2$.

\begin{proposition}[the routine AP]\label{prop:ap}
On every input AP fails or returns a point of $Q$ of norm at most $R$. If $d\ge1$, $Q$ is closed and convex with $B(a_0,r)\subset Q\subset B(0,R)$, $r>0$, $P$ is positive definite and $e>0$, then AP returns a point $\mu\in Q$ with $F(\mu)\le F(x)+e$ for every $x\in Q$. So AP meets the specification of Section~\ref{sec:poly-learner}.
\end{proposition}

\begin{proof}
For $d\ge2$ this is Proposition~\ref{prop:ell}. For $d=1$: a returned point is $w'$, which $Q$ contains, or $\mathrm{lo}$, which starts at the accepted $a_0$ and moves only to accepted midpoints; both lie between $a_0$ and $w'$, so in $[-R,R]$. Under the hypotheses, $Q$ is an interval and the minimiser $q$ of $F$ on $Q$ is the point of $Q$ nearest to $w$, which is also the point of $Q$ nearest to $w'$. If $w'\in Q$, then $q=w'$. Otherwise $q$ is the end of $Q$ between $a_0$ and $w'$; by induction $\mathrm{lo}\in Q$ and $\mathrm{hi}\notin Q$, so $q$ lies between them, and at the end $|\mathrm{lo}-q|\le|w'-a_0|2^{-m}\le\xi$. $F$ is $p(R+|w|)$-Lipschitz on $[-R,R]$, so $F(\mathrm{lo})\le F(q)+e$.
\end{proof}

\section{The position}\label{app:pg}

This appendix writes out Procedure PG and proves Theorem~\ref{thm:pg}. Throughout, $\kappa=\frac18$, $g$ is a standard Gaussian vector of $\R^d$, $\xi$ a standard Gaussian variable, $\pi$ the gauge of the body $\cK$ about $0$, and for $\mu\in\R^d$ and $A$ invertible the body $\cK$ in the position $(\mu,A)$ is $\cK'=\{y:\mu+Ay\in\cK\}$, with gauge $\pi_{\cK'}$. The hypothesis on the body in Lemma~\ref{lem:l} and Proposition~\ref{prop:cert} is: $\cK\subset\R^d$ is compact and convex, contains the open ball of radius $r_0>0$ about $0$, and $-\cK\subset\alpha\cK$ for some $\alpha\ge0$.

\subsection{The Gaussian position}

For a position $(\mu,A)$ let $F(y)=\pi(\mu+Ay)$ and
\[
k=\E F(g),\qquad b=\E[F(g)g],\qquad W=\E[(gg^{\mathsf T}-I)F(g)],\qquad h(\mu,A)=k-\frac\kappa d\ln|\det A| .
\]
For $\alpha\ge0$ and $t\in\R$ let $C_\alpha(t)=(1+\alpha)\varphi(t)+t\bigl(1-(1+\alpha)Q(t)\bigr)$. The bathtub principle (among the functions $w$ with values in $[0,1]$ and $\E w(\xi)=p$, $\E[-\xi w(\xi)]$ is largest for the indicator of the lower $p$-quantile) gives: a function $h_1$ with values in $[-\alpha,1]$ satisfies $|\E[\ip{g,u}h_1(g)]|\le C_\alpha(t)+|t|\,|\E h_1(g)|$ for every unit vector $u$ and every $t$. At the $t\ge0$ with $Q(t)=1/(1+\alpha)$, $t\le C_\alpha(t)=(1+\alpha)\varphi(t)\le(2\ln(1+\alpha))^{1/2}$.

\begin{lemma}[the Gaussian position]\label{lem:l}
Let $d\ge1$ and $\cK$ be as above. Then $h$ has a minimiser over all $\mu\in\R^d$ and all invertible $A$, and one with $A$ symmetric positive definite. At every minimiser $k=\kappa$, $b=0$ and $W=\frac\kappa dI$. A position with these three values satisfies $\cK'\subset B(0,d(1+\kappa)C_\alpha(t)/\kappa)$ for every $t$, $B(0,a)\subset\cK'$ for every $a$ with $\kappa\le Q(a)$, and $\E\pi_{\cK'}(g)\le1/(1-\kappa)$.
\end{lemma}

\begin{proof}
The gauge $\pi$ is convex, positively homogeneous and $(1/r_0)$-Lipschitz, hence differentiable almost everywhere; where it is differentiable, $\ip{\nabla\pi(x),x}=\pi(x)$ and $\ip{\nabla\pi(x),z}\le\pi(z)$ for every $z$.

\emph{Existence.} Let $\cK\subset B(0,R_0)$. For $T$ symmetric positive definite with extreme eigenvalues $t_1\ge t_d$, Jensen's inequality and $\pi(x)\ge|x|/R_0$ give $h(\mu,T)\ge\max\bigl(|\mu|,t_1\E|\xi|\bigr)/R_0-\frac\kappa d\bigl((d-1)\ln t_1+\ln t_d\bigr)$, so the sublevel sets of $h$ on these pairs are compact and $h$ has a minimiser there. Since $Ag$ and $(AA^{\mathsf T})^{1/2}g$ have the same law and $|\det A|=\det(AA^{\mathsf T})^{1/2}$, $h(\mu,A)=h(\mu,(AA^{\mathsf T})^{1/2})$, and the minimiser is one over all invertible $A$.

\emph{Stationarity.} Let $(\mu,A)$ be a minimiser and $X=\mu+Ag$, which has a density. By dominated convergence $(\mu,A)\mapsto\E\pi(\mu+Ag)$ has at $(\mu,A)$ the directional derivative $\ip{\E\nabla\pi(X),\nu}+\tr(S^{\mathsf T}G)$, $G=\E[\nabla\pi(X)g^{\mathsf T}]$, in every direction $(\nu,S)$, $S$ any matrix, and $\ln|\det A|$ has the derivative $\tr(A^{-1}S)$. The derivative of $h$ vanishes in every direction, so $\E\nabla\pi(X)=0$ and $G=\frac\kappa dA^{-\mathsf T}$. Gaussian integration by parts for the Lipschitz function $F$, with $\nabla F(g)=A^{\mathsf T}\nabla\pi(X)$, turns these into $b=\E\nabla F(g)=0$ and $W=\E[g\nabla F(g)^{\mathsf T}]=G^{\mathsf T}A=\frac\kappa dI$, and $k=\E\ip{\nabla\pi(X),X}=\ip{\E\nabla\pi(X),\mu}+\tr(A^{\mathsf T}G)=\kappa$.

\emph{Radii.} The last two claims are Proposition~\ref{prop:cert}(b) and~(c) with $k=\kappa$. The first is the inequality $q|y|\le(1+k)C_\alpha(t)+|t|\,|b|\bigl(|y-y_0|+|y_0|\bigr)$ of the proof of Proposition~\ref{prop:cert}(a), with $b=0$ and $q=\kappa/d$.
\end{proof}

\begin{proposition}[the certificate]\label{prop:cert}
Let $\cK$ be as above, $\mu\in\R^d$ and $A$ invertible.
\begin{enumerate}[label=(\alph*)]
\item If $W\ge qI$ with $q>0$ and $6|t|\,|b|\le q$, then $\cK'\subset B\bigl(0,2(1+k)C_\alpha(t)/q\bigr)$.
\item If $k\le Q(a)$, then $B(0,a)\subset\cK'$.
\item If $k<1$, then $\E\pi_{\cK'}(g)\le1/(1-k)$.
\end{enumerate}
\end{proposition}

\begin{proof}
(a) By integration by parts $b=\E\nabla F(g)$ and $W=\E[g\nabla F(g)^{\mathsf T}]$, $\nabla F(g)=A^{\mathsf T}\nabla\pi(X)$. Let $y\in\cK'$, $z=\mu+Ay\in\cK$, and $y_0=-A^{-1}\mu$, which lies in $\cK'$ since $0\in\cK$. Then $Wy=\E[gh_1]-\E[gh_2]$ with $h_1=\ip{\nabla\pi(X),z}\in[-\alpha,1]$ and $h_2=\ip{\nabla\pi(X),\mu}\in[-\alpha k,k]$, as $\pi(\mu)\le k$ by Jensen's inequality; their means are $\E h_1=\ip{b,y-y_0}$ and $\E h_2=-\ip{b,y_0}$. The bathtub bound, applied to $h_1$ and to $h_2/k$, gives
\[
q|y|\le|Wy|\le(1+k)C_\alpha(t)+|t|\,|b|\bigl(|y-y_0|+|y_0|\bigr).
\]
At a point $y$ of largest norm $R'$ in the compact set $\cK'$ the last term is at most $3|t|\,|b|R'\le qR'/2$.

(b) $\cK'=\{y:F(y)\le1\}$, and $\Prob(F(g)>1)\le k$ by Markov's inequality, so $\Prob(g\in\cK')\ge1-k$. A point $y\notin\cK'$ is separated from $\cK'$ by a half-space $\{\ip{u,\cdot}\le s\}$ containing $\cK'$, $u$ a unit vector, with $s<|y|$; then $1-k\le\Prob(\xi\le s)<1-Q(|y|)$, that is $Q(|y|)<k\le Q(a)$, and $|y|>a$.

(c) For $x=\mu+v$ and $t\ge1$, $\pi(\mu+v/t)\le(1-\frac1t)\pi(\mu)+\frac1t\pi(x)$, so $\pi_{\cK'}(y)\le1+F(y)/(1-\pi(\mu))\le1+F(y)/(1-k)$, and $\E\pi_{\cK'}(g)\le1+k/(1-k)$.
\end{proof}

\subsection{Procedure PG-core}

Let $d\ge1$, $\rho\ge1$, $0<\delta<1$, and let $\cK$ be closed and convex with $B(0,1)\subset\cK\subset B(0,\rho)$, so that $|x|/\rho\le\pi(x)\le|x|$ and $-\cK\subset\rho\cK$. Put $\lambda_\rho=(2\ln(1+\rho))^{1/2}$ and $P_B=1.254$, a rational number above $1/\E|\xi|$. The parameters are
\begin{gather*}
L_D=t_{\max}=\kappa P_B\rho,\qquad t_{\min}=\frac1{9d\lambda_\rho},\qquad\Delta=\kappa\rho\,(4+P_B^2d)^{1/2},\\
G^2=\frac{5L_D^2d(d+2)(d+5)+2\kappa^2/d}{t_{\min}^2},\qquad e_h=\frac{\kappa^2}{922\,d^2\lambda_\rho^2L_D},\qquad\eta=\frac{e_ht_{\min}}{7L_D(1+d^{1/2})\Delta},\\
\tau_b=\frac\kappa{12d\lambda_\rho},\qquad s=\frac\Delta{GN^{1/2}},\qquad j=\lceil\log_2(\rho/\eta)\rceil,\qquad N=\lceil16\Delta^2G^2/e_h^2\rceil,\\
J=\lceil\log_2(4/\delta)\rceil,\qquad m_T=\lceil8\ln(2J/\delta)\rceil,\\
n_T=\Bigl\lceil\frac54\,\frac{L_D^2d^3}{\kappa^2}\max\bigl(6144(d+1)(d+5),\,221184(d+3)\lambda_\rho^2\bigr)\Bigr\rceil.
\end{gather*}

\emph{The gauge by bisection.} For $x\ne0$, $\tilde\pi(x)$ is the upper end of the interval left by $j$ steps of bisection on $[|x|/\rho,|x|]$ with the question ``$s>0$ and $x/s\in\cK$''; $\tilde\pi(0)=0$. The interval always contains $\pi(x)$ and its length is at most $2^{-j}|x|\le2^{-j}\rho\,\pi(x)$, so $\pi\le\tilde\pi\le(1+\eta)\pi$.

\emph{The domain.} $D$ is the set of pairs $(\mu,T)$, $T$ symmetric with $t_{\min}I\le T\le t_{\max}I$ and $|\mu|\le\kappa\rho$, in the space $\R^d\times\mathrm{Sym}_d$ with the norm $(|\mu|^2+\|T\|_F^2)^{1/2}$; it is compact, convex, of diameter at most $\Delta$, and $\mathrm{proj}_D$ is its nearest-point map. On $D$, $0\le F(y)\le L_D(1+|y|)$ and $F$ is $L_D$-Lipschitz.

\emph{The estimate.} At $x=(\mu,T)\in D$ and a Gaussian vector $g$, with $\tilde F(g)=\tilde\pi(\mu+Tg)$,
\[
v(x,g)=\Bigl(\tilde F(g)\,T^{-1}g,\ \frac{\tilde F(g)}2\bigl(T^{-1}(gg^{\mathsf T}-I)+(gg^{\mathsf T}-I)T^{-1}\bigr)-\frac\kappa dT^{-1}\Bigr).
\]

\emph{The test.} The basic verdict at $x$ on $n_T$ Gaussian vectors $g_1,\dots,g_{n_T}$ forms the averages $\hat k,\hat b,\hat W$ of $\tilde F(g_i)$, $\tilde F(g_i)g_i$ and $\tilde F(g_i)(g_ig_i^{\mathsf T}-I)$, and accepts when $\hat k\le1.075\kappa$, $|\hat b|\le\frac34\tau_b$ and $\hat W\ge\frac58\frac\kappa dI$. The test accepts when more than $m_T/2$ of $m_T$ basic verdicts, on independent blocks of vectors, accept.

\paragraph{Procedure PG-core$(\cK,\rho,\delta)$.} For $r=1,\dots,J$, on a fresh block of Gaussian vectors: let $x_1=(0,t_{\min}I)$ and $x_{i+1}=\mathrm{proj}_D\bigl(x_i-s\,v(x_i,g_i)\bigr)$ for $i<N$; let $\bar x=(\mu,T)$ be the average of $x_1,\dots,x_N$; run the test at $\bar x$, and if it accepts, return the position $(\mu,T)$. If no repeat accepts, fail.

\begin{theorem}[PG-core]\label{thm:pgcore}
Let $d\ge1$, $\rho\ge1$, $0<\delta<1$, and $\cK\subset\R^d$ closed and convex with $B(0,1)\subset\cK\subset B(0,\rho)$. The set of streams on which PG-core does not return a position $(\mu,A)$ with $\det A\ne0$,
\[
B(0,1)\subset\cK'\subset B(0,36.4\,d\lambda_\rho)\quad\text{and}\quad M(\cK')\le1.27\,\frac{(d+1)^{1/2}}d
\]
has outer measure at most $\delta$ under the law of independent standard Gaussian vectors.
\end{theorem}

\begin{proof}
\emph{Soundness of the test.} Call estimates accurate when $|\hat k-k|\le\kappa/40$, $|\hat b-b|\le\tau_b/4$ and $\|\hat W-W\|\le\kappa/(8d)$. If accurate estimates are accepted, then $k\le1.1\kappa$, $|b|\le\tau_b$ and $W\ge\frac\kappa{2d}I$: call this the sound condition. Under it, Proposition~\ref{prop:cert} applies with $\alpha=\rho$, $q=\kappa/(2d)$ and the $t\ge0$ with $Q(t)=1/(1+\rho)$, since $6|t|\,|b|\le6\lambda_\rho\tau_b=q$. Part (a) gives $\cK'\subset B(0,4(1+1.1\kappa)d\lambda_\rho/\kappa)=B(0,36.4\,d\lambda_\rho)$; part (b) gives $B(0,1)\subset\cK'$, as $Q(1)=0.15866\ge0.1375\ge k$; part (c) gives $\E\pi_{\cK'}(g)\le1/0.8625=1.1594\ldots$, and with $\E|g|\ge d/(d+1)^{1/2}$ this is $M(\cK')\le1.27(d+1)^{1/2}/d$.

\emph{The test.} At a point of $D$, $0\le\tilde F-F\le\eta F$, so the biases of the three averages are at most $u$, $u$ and $2^{1/2}u$ with $u=\eta L_D(1+d^{1/2})$, below $\kappa/80$, $\tau_b/8$ and $\kappa/(16d)$. By Chebyshev's inequality, with the second moments of $\tilde F(g)$, $\tilde F(g)g$ and $\tilde F(g)(gg^{\mathsf T}-I)$ computed from $\E|g|^2=d$, $\E|g|^4=d(d+2)$, $\E|g|^6=d(d+2)(d+4)$, each average is within the same margins of its mean except with probability $\frac1{12}$ at $n_T$ samples. A basic verdict is therefore accurate with probability at least $\frac34$. At a point of $D$ that violates the sound condition, a basic verdict accepts only if it is not accurate, and by Hoeffding's inequality the test accepts with probability at most $e^{-m_T/8}\le\delta/(2J)$. At a point of $D$ with $k\le1.05\kappa$, $|b|\le\tau_b/2$ and $W\ge\frac34\frac\kappa dI$, accurate estimates are accepted, and the test rejects with probability at most $\delta/(2J)$.

\emph{The descent.} The function $h$ is convex on $D$, and for $(\mu,T)\in D$ the pair $\bigl(T^{-1}b,\frac12(T^{-1}W+WT^{-1})-\frac\kappa dT^{-1}\bigr)$ is a subgradient of $h$ at $(\mu,T)$. The symmetric minimiser $(\mu^*,T^*)$ of Lemma~\ref{lem:l} lies in $D$: $|\mu^*|\le\rho\,\pi(\mu^*)\le\rho k=\kappa\rho$; $\E|\xi|\,|T^*u|=\E|\ip{u,T^*g}|\le\rho\,\E\pi(\mu^*+T^*g)=\kappa\rho$ for every unit $u$; and $T^*\ge t_{\min}I$, because $[-u,u]\subset\cK$ lies in $\mu^*+T^*B(0,9d\lambda_\rho)$ by the outer radius of Lemma~\ref{lem:l} with $\alpha=\rho$. Given the past, $v(x_i,g_i)$ has a mean within $\beta=3^{1/2}u/t_{\min}$ of that subgradient and a second moment at most $G^2$. Projected stochastic subgradient descent then gives $\E h(\bar x)-\min h\le\Delta G/N^{1/2}+\beta\Delta\le\frac{e_h}4+\frac{3^{1/2}e_h}7\le\frac{e_h}2$, so $h(\bar x)\le\min h+e_h$ with probability at least $\frac12$ by Markov's inequality. Such a point has $k\le1.05\kappa$ (compare with $(s\mu,sT)$, $s=\kappa/k$; $e_h<0.0012\kappa<(1.05-1-\ln1.05)\kappa$), $|b|^2\le2\E|\xi|\,L_De_h\le(\tau_b/2)^2$ (compare with a translation) and $W\ge\frac34\frac\kappa dI$ (compare with $T(I+\sigma uu^{\mathsf T})$ for a unit $u$, using the curvature bound $8\varphi(2^{1/2})L_D\le1.1742L_D$). So each repeat accepts with probability at least $\frac12(1-\delta/(2J))$, whatever the repeats before it did.

\emph{Failure.} Some repeat accepts a point that violates the sound condition with probability at most $J\cdot\delta/(2J)=\delta/2$, and no repeat accepts with probability at most $\bigl(\frac12(1+\delta/(2J))\bigr)^J\le\frac\delta4e^{1/2}$; the sum is at most $0.92\delta$. Outside these two events PG-core returns a point of $D$ under the sound condition, and the first paragraph concludes.
\end{proof}

\subsection{Growing balls}

\paragraph{Procedure PG$(\cK,r_0,R_0,\delta)$.} Let $m=\max(1,\lceil\log_2(R_0/r_0)\rceil)$, $\cK_i=\cK\cap B(0,2^ir_0)$ for $i=1,\dots,m$ (so $\cK_m=\cK$), $\rho_1=2$ and $\rho_{i+1}=3\cdot54\,d\,(\ln(1+\rho_i))^{1/2}$. Start from the position $(0,r_0I)$. For $i=1,\dots,m$, on its own block of Gaussian vectors, run PG-core$(\cdot,\rho_i,\delta/m)$ on the body $\cK_i$ in the current position $(\mu,A)$, whose membership question at $y$ is the question $\mu+Ay\in\cK$ and the comparison $|\mu+Ay|\le2^ir_0$, and replace the current position by its composition $(\mu+A\mu',AA')$ with the position $(\mu',A')$ returned. If a stage fails, PG fails. It returns the last position.

\begin{proof}[Proof of Theorem~\ref{thm:pg}]
For a position $P$ write $\cK_i^P$ for $\cK_i$ in the position $P$; let $P_0=(0,r_0I)$ and $P_i$ be the position after stage $i$. Call stage $i$ good if it returns a position and $B(0,1)\subset\cK_i^{P_i}\subset B(0,36.4\,d\lambda_{\rho_i})$; since $36.4\cdot2^{1/2}=51.48\le54$, then $\cK_i^{P_i}\subset B(0,R_i)$ with $R_i=54\,d\,(\ln(1+\rho_i))^{1/2}$.

(i) Stage $1$ runs on $\cK_1^{P_0}=\cK_1/r_0$, which lies between $B(0,1)$ and $B(0,2)$. Let stage $i-1$ be good. As $\cK_{i-1}\subset\cK_i$, $B(0,1)\subset\cK_i^{P_{i-1}}$. As $0\in\cK$, $\cK_i\subset2\cK_{i-1}$, so a point $y$ of $\cK_i^{P_{i-1}}$ is $2z-y_0$ with $z$ and $y_0$ in $\cK_{i-1}^{P_{i-1}}$ ($y_0$ the image of $0$), and $|y|\le3R_{i-1}=\rho_i$. So stage $i$ runs under the hypothesis of Theorem~\ref{thm:pgcore}.

(ii) Let $\rho^\star=400\,d\,(\ln(2d+2))^{1/2}$. Then $\rho_i\le\rho^\star$ for every $i$: the map $\rho\mapsto162\,d\,(\ln(1+\rho))^{1/2}$ is increasing, $\rho_1=2\le\rho^\star$, and it maps $\rho^\star$ below $\rho^\star$, because $1+\rho^\star\le200.5\,(2d+2)^{3/2}\le(2d+2)^{5.5}$ (as $200.5\le4^4$ and $2d+2\ge4$) and $162\cdot5.5^{1/2}=379.92\le400$.

(iii) By (i), given that stages $1$ to $i-1$ are good, stage $i$ runs under the hypothesis of Theorem~\ref{thm:pgcore} on vectors not read before, and is good except on a set of outer measure at most $\delta/m$. All $m$ stages are good outside a set of outer measure at most $\delta$. Then $S(\cK)=\cK_m^{P_m}$ satisfies $B(0,1)\subset S(\cK)\subset B(0,R_m)$ and the bound on $M$ of the last stage, and $R_m\le54\,d\,(\ln(1+\rho^\star))^{1/2}\le54\cdot5.5^{1/2}\,d\,(\ln(2d+2))^{1/2}$ with $54\cdot5.5^{1/2}=126.64\le134$.
\end{proof}

Since the parameters of PG-core are nondecreasing in $\rho$, PG asks at most $mJ(N+m_Tn_T)j$ membership questions, the parameters taken at $\rho^\star$ and $\delta/m$; this is polynomial in $d$, $\ln(R_0/r_0)$ and $\ln(1/\delta)$. Theorem~\ref{thm:pg} does not state this count; the step count of the program is Proposition~\ref{prop:mach}.
\section*{Acknowledgements}

We thank Hedi Hadiji for suggesting the Bayesian approach of Lattimore through minimax duality, and for the remark that the analysis of exploration by optimisation should apply to the one-round inequality, which led to the explicit learner.

\section*{Contact}
Questions, corrections and comments are welcome as issues on the repository: \url{\repo/issues/new/choose}.

\bibliographystyle{amsplain}
\bibliography{refs}

\end{document}
